% !TeX program = lualatex
% Fichier LaTEI généré depuis XML-TEI Métopes / Commons-Publishing.
% Zone technique (préambule, macros, mappings) : régénérable depuis le XML source.
% NE PAS MODIFIER la zone technique — elle sera écrasée à la prochaine génération.
% Zone éditoriale réversible : l'environnement lateiDocument dans ce fichier.
% Corriger uniquement dans la zone réversible.

\documentclass[11pt,twoside,openany]{book}

\usepackage[
  paperwidth=155mm,
  paperheight=230mm,
  top=30mm,
  bottom=19mm,
  inner=20mm,
  outer=30mm,
  headheight=14pt,
  headsep=8mm,
  footskip=10mm
]{geometry}
% Référentiel PURH v0.6 §6.2 : « le bandeau généré est situé environ 3,1 mm
% plus bas que dans le PDF imprimeur » — écart de rendu, pas une remise en
% cause de la marge du haut mesurée sur le maître InDesign (profile.
% margin_top_mm, qui reste la source de vérité pour la position du corps de
% texte). \topmargin remonte le bandeau de titre courant de 3,1 mm ;
% \headsep compense d'autant pour que le corps de texte, lui, démarre
% exactement à la même position qu'avant ce correctif.
\addtolength{\topmargin}{-3.1mm}
\addtolength{\headsep}{3.1mm}
\raggedbottom

\usepackage{fontspec}
\usepackage[french]{babel}
\usepackage{csquotes}
\usepackage{amsmath}
\usepackage{microtype}
\usepackage{indentfirst}
\usepackage{emptypage}
\usepackage[normalem]{ulem}
% [table] : nécessaire pour \rowcolor sur les lignes d'entête de tableau
% (référentiel §11.3, "fond foncé noir 30 %") — charge colortbl en plus des
% commandes \color habituelles (compatible avec l'usage \color[gray]{}
% déjà en place pour le titre courant).
\usepackage[table]{xcolor}

% Référentiel PURH v0.6 §12.1 : « le texte courant du PDF imprimeur
% correspond à un noir process à 90 % [K]. La sortie actuelle utilise du
% noir plein. » CMYK explicite (0,0,0,0.9), pas une valeur de gris RVB :
% "K" désigne le noir process de la quadrichromie, un concept distinct du
% gris \color[gray]{} utilisé ailleurs (titre courant) qui n'a pas de sens
% en CMJN. Appliqué globalement dès le début du document.
\definecolor{PURHBodyBlack}{cmyk}{0,0,0,0.9}
\AtBeginDocument{\color{PURHBodyBlack}}

\IfFontExistsTF{Chaparral Pro}
  {\setmainfont{Chaparral Pro}}
  {\setmainfont{TeX Gyre Pagella}}

\IfFontExistsTF{Josefin Sans}
  {\newfontfamily\PURHTitleFont{Josefin Sans}}
  {\newfontfamily\PURHTitleFont{TeX Gyre Heros}}

% Titraille PURH (référentiel §2.3-§2.5) : partie, article et section
% observés en Josefin Sans Thin, capitales, distinct du \PURHTitleFont
% "normal weight" ci-dessus utilisé ailleurs. Chargée comme famille à part
% (et non via \bfseries sur \PURHTitleFont) car "Thin" n'est pas une série
% NFSS standard que fontspec puisse sélectionner automatiquement — "Josefin
% Sans Thin" existe en revanche comme nom de famille indépendant portant
% elle-même ses propres graisses Thin (romain) et Thin Italic.
\IfFontExistsTF{Josefin Sans Thin}
  {\newfontfamily\PURHTitreFont{Josefin Sans Thin}}
  {\newfontfamily\PURHTitreFont{TeX Gyre Heros}}

\IfFontExistsTF{Latin Modern Mono}
  {\setmonofont{Latin Modern Mono}}
  {\setmonofont{TeX Gyre Cursor}}

% Romain, pas italique (référentiel PURH §2.3 : titres courants "Josefin
% Sans Thin/Light, 10 pt, romain" — l'italique systématique précédente était
% un défaut confirmé, pas un choix).
%
% Cinq vérifications humaines successives (2026-08-04/06) ont porté sur ce
% même réglage : la première jugeait le gris trop clair (corrigé en passant
% de \PURHTitreFont, famille Thin, à \PURHTitleFont, famille standard, sans
% \bfseries) ; la seconde a trouvé ce résultat trop noir et visuellement
% gras — le PDF imprimeur, lui, n'a « pas de graisse » — d'où un retour à la
% famille Thin avec un gris explicite (\color[gray]{0.25}, approximatif) ;
% la troisième a jugé ce gris encore trop clair, d'où un premier passage au
% même système que le fond d'entête de tableau et le texte courant
% (§11.3/§12.1) — une teinte CMJN noir X % plutôt qu'un gris RVB — à 50 %
% noir ; la quatrième a de nouveau jugé ce résultat trop clair au regard du
% PDF imprimeur, où le titre courant est « presque noir » : remonté à 85 %
% noir. La cinquième vérification, sur une page réelle complète d'un livre
% du corpus (pas seulement les liminaires), a de nouveau jugé le résultat
% trop clair à 85 % — poussé au maximum possible, 100 % noir (K plein),
% sans changer de famille (\PURHTitreFont, Josefin Sans Thin) : le résultat
% restait visiblement plus clair que le PDF imprimeur, confirmant qu'à ce
% stade le levier couleur seul était épuisé — un trait Thin, même à pleine
% saturation d'encre, couvre trop peu de surface pour lire aussi sombre
% qu'un trait plus épais à la même teinte.
%
% Sixième vérification (2026-08-06), comparaison directe côte à côte avec
% le PDF imprimeur (recadrage du bandeau à haute résolution) : la graisse
% du PDF imprimeur n'est PAS aussi fine qu'un trait Thin — plutôt un trait
% "Regular" (ni fin ni gras). Changement de levier plutôt que nouvel
% incrément de couleur : \PURHTitleFont (famille Josefin Sans complète,
% graisse Regular par défaut sans \bfseries — PAS \PURHTitreFont/Thin) à
% 75 % noir, un compromis délibéré entre les deux réglages précédents (ni
% le gris à 50 % jugé trop clair, ni l'aspect "gras" que la graisse
% Regular pourrait donner à 100 %) — comparé visuellement à un jeu
% d'échantillons Thin/Light/Regular/Medium × 75 %/100 % avant d'être
% retenu comme la combinaison la plus proche du PDF imprimeur.
%
% Septième vérification (2026-08-06) : le corps (Regular) validé comme
% correct par l'utilisateur — seule la teinte, encore un peu trop noire à
% 75 %, redescendue légèrement à 65 % (un ajustement fin, pas un nouveau
% changement de levier comme aux passes précédentes).
\newcommand{\PURHHeaderFont}{\PURHTitleFont\small\color[cmyk]{0,0,0,0.65}}

% Le corps et son pas de ligne sont fixés explicitement au lieu de dépendre
% de la table de tailles du \documentclass{book} choisi : elle donne un
% pas voisin mais pas identique au pas de grille cible (référentiel PURH
% §2.4, §5.3 : corps 11 pt sur une grille de 13,5 pt).
\renewcommand{\normalsize}{\fontsize{11pt}{13.5pt}\selectfont}
\normalsize

\newcommand{\PURHBookTitle}{Dissimuler pour mieux régner}
\newcommand{\PURHBookSubtitle}{<em>Locus politicus, locus secretus</em> en littérature (XVIIe-XIXe siècles)}
\newcommand{\PURHBookAuthor}{Floriane Daguisé ; Florence Fix}
\newcommand{\PURHPublisher}{PURH}
\newcommand{\PURHYear}{}
\newcommand{\PURHDOI}{}
\newcommand{\PURHISBN}{}

\title{\PURHBookTitle}
\author{\PURHBookAuthor}
\date{\PURHYear}

\setlength{\parindent}{5mm}
\setlength{\parskip}{0pt}
\linespread{1.0}
\pretolerance=100
\tolerance=500
\hyphenpenalty=500
\exhyphenpenalty=500
\emergencystretch=3em
\clubpenalty=10000
\widowpenalty=10000
\displaywidowpenalty=10000

\usepackage{enumitem}

\setlist[itemize,1]{
  label=\textendash,
  leftmargin=1.5em,
  itemsep=0.2\baselineskip,
  topsep=0.4\baselineskip
}

\setlist[enumerate,1]{
  leftmargin=1.8em,
  itemsep=0.2\baselineskip,
  topsep=0.4\baselineskip
}

\usepackage[nobottomtitles*]{titlesec}
\usepackage{titletoc}

\setcounter{secnumdepth}{0}
% Référentiel PURH v0.6 §9 ("Table des matières") : la cible exclut les
% sections internes (intertitres) de la TDM — seuls le niveau des parties
% (\part, niveau -1) et celui des ouvertures de contribution/front matter
% (\addcontentsline{toc}{chapter}, niveau 0) doivent y figurer. tocdepth=2
% incluait à tort les <div type="section1">/<div type="section2"> (niveaux
% \section=1, \subsection=2), gonflant la TDM à trois pages au lieu de deux.
\setcounter{tocdepth}{0}

% Josefin Sans Bold, 16 pt, capitales — le référentiel indiquait Josefin
% Sans Thin (§2.5, §5.3, §4.3), contredit par vérification humaine directe
% du PDF généré face au PDF imprimeur : les titres de partie y sont noirs et
% gras, alors que le Thin rendait un texte maigre et grisâtre, peu lisible.
% Observation directe suivie ici, comme pour le séparateur de note de bas de
% page défini plus bas dans ce même préambule. \PURHTitleFont (pas
% \PURHTitreFont, qui ne charge que la graisse Thin et ne peut donc pas
% produire de \bfseries réel) charge la famille Josefin Sans complète, dont
% fontspec sélectionne automatiquement la graisse Bold via NFSS.
\titleformat{\chapter}[display]
  {\PURHTitleFont\bfseries\fontsize{16pt}{19pt}\selectfont\raggedright}
  {\chaptertitlename~\thechapter}
  {10pt}
  {\MakeUppercase}

% Titre de partie : même correctif Bold que \chapter ci-dessus. Toujours
% \part* (pas de numéro affiché) : le label reste vide plutôt que d'exposer
% un numéro de partie non requis.
\titleformat{\part}[display]
  {\PURHTitleFont\bfseries\fontsize{16pt}{19pt}\selectfont\centering}
  {}
  {0pt}
  {\MakeUppercase}

% Titre de section (intertitre), 12 pt, capitales. Le référentiel §2.5
% indiquait Josefin Sans Thin ; corrigé en Bold (même correctif que la
% titraille partie/chapitre/contribution ci-dessus) après vérification
% humaine directe du PDF généré face au PDF imprimeur : les intertitres y
% apparaissent noirs et gras, pas maigres (chantier de parité v0.6,
% 2026-08-04). Alignement et espacement non chiffrés par le référentiel pour
% ce niveau : conservés tels quels (raggedright, séparations existantes).
\titleformat{\section}[block]
  {\PURHTitleFont\bfseries\fontsize{12pt}{14pt}\selectfont\raggedright}
  {}
  {0pt}
  {\MakeUppercase}

% Sous-section (référentiel v0.6 §4.3 : seul le niveau "section" — c'est-à-
% dire section1 — est chiffré ; aucune valeur propre à section2/section3
% n'est fournie). Un correctif antérieur avait délibérément retiré le gras
% à ce niveau (confirmé alors sur le livre réel — <div type="section2"> y
% porte de vrais sous-titres phrastiques, jamais des libellés courts) ; la
% vérification humaine directe du 2026-08-04 demande au contraire le même
% traitement noir et gras que les autres intertitres, suivie ici. Capitales
% non demandées à ce niveau, contrairement à \section : inchangé. Tailles
% conservées telles quelles (\large/\normalsize) faute de mesure référentiel.
\titleformat{\subsection}[block]
  {\PURHTitleFont\bfseries\large\raggedright}
  {}
  {0pt}
  {}

\titleformat{\subsubsection}[block]
  {\PURHTitleFont\bfseries\normalsize\raggedright}
  {}
  {0pt}
  {}

\titlespacing*{\part}{0pt}{0pt}{30pt}
\titlespacing*{\chapter}{0pt}{20pt}{18pt}
\titlespacing*{\section}{0pt}{18pt}{10pt}
\titlespacing*{\subsection}{0pt}{14pt}{8pt}
\titlespacing*{\subsubsection}{0pt}{12pt}{6pt}

\addto\captionsfrench{
  \renewcommand{\contentsname}{Table des matières}
}

% Entrées de contribution/front matter (référentiel PURH v0.6 §9.1,
% vérification humaine directe du 2026-08-04) : « titre de la communication
% sans graisse, bas de casse, Chaparral, calé à gauche, série de points
% puis numéro de ligne ». Chaparral Pro reste la fonte ambiante (aucune
% famille à sélectionner) : \PURHTitleFont\bfseries retiré, \hfill remplacé
% par le même filet pointillé que \section ci-dessous. \addvspace{8pt}
% inconditionnel retiré : « pas de saut de ligne entre les références sauf
% changement de section » — l'espacement avant une nouvelle partie vient
% désormais du bloc \part ci-dessous, pas d'ici.
%
% Filet pointillé + numéro de page : laissés dans ce 4e argument dédié (pas
% déplacés dans le texte de l'entrée transmis à \addcontentsline). Une
% première tentative avait déplacé \titlerule*/\contentspage dans ce texte
% pour les faire tomber sur la ligne du titre plutôt que sur celle de
% l'auteur (vérification humaine directe du 2026-08-04 : « les points de
% suite doivent être au niveau du titre, pas de l'auteur ») — abandonnée :
% le paquet bookmark, qui construit automatiquement les signets PDF depuis
% CE MÊME texte d'entrée, ne tolère pas \contentspage/\\ dans cet argument
% ("Token not allowed in a PDF string", puis désynchronisation de
% titlesec — bug réel constaté par compilation). Solution retenue : le nom
% d'auteur n'est plus concaténé DANS le texte de l'entrée de chapitre du
% tout — voir \latei_finish_contribution_toc_entry: (latei_macros.tex), qui
% l'écrit désormais comme une ligne de TDM séparée, via \addtocontents
% (jamais capturée par le mécanisme de signets de bookmark, qui n'observe
% que \addcontentsline/\contentsline).
\titlecontents{chapter}
  [0pt]
  {}
  {}
  {}
  {\titlerule*[0.5pc]{.}\contentspage}

% Entrées de partie (le référentiel dit « titres de section », mais
% désigne bien ici le niveau \part de ce document — les véritables
% intertitres/sections sont exclus de la TDM depuis §9/tocdepth=0) : Josefin
% Sans Bold, capitales (vérification humaine directe du 2026-08-04 : à la
% différence du nom d'auteur sous chaque entrée de contribution, qui doit
% rester bas de casse — voir \lateiTocAuthorLine dans latei_macros.tex —
% les titres de ce niveau doivent au contraire être en petites capitales).
% \scshape ici n'a AUCUN effet sur \PURHTitleFont (Josefin Sans, qui n'a pas
% de véritables petites capitales OpenType — confirmé par compilation
% isolée, "Font shape .../b/sc undefined", substitution silencieuse) : la
% mise en majuscules se fait donc en amont, à la source du texte, via
% \MakeUppercase{#1} dans \lateiRenderHead (latei_macros.tex) — seul le
% corps réduit ici (12 pt contre 16 pt sur la page de la partie elle-même)
% distingue ce niveau de véritables grandes capitales.
\titlecontents{part}
  [0pt]
  {\addvspace{1\baselineskip}\PURHTitleFont\bfseries\fontsize{12pt}{14pt}\selectfont\centering}
  {}
  {}
  {}
  [\addvspace{1\baselineskip}]

\titlecontents{section}
  [1.5em]
  {}
  {}
  {}
  {\titlerule*[0.5pc]{.}\contentspage}

\titlecontents{subsection}
  [3em]
  {\small}
  {}
  {}
  {\titlerule*[0.5pc]{.}\contentspage}

\usepackage{fancyhdr}

\pagestyle{fancy}
\fancyhf{}
% Folios à l'extérieur (référentiel PURH §2.3) : LE (verso) et RO (recto).
% Les titres courants intérieurs (RE, LO) et \chaptermark sont pris en
% charge par latei_macros.tex, qui distingue verso (livre/partie) et recto
% (contribution en cours) — les définir ici aussi serait dupliqué et, pire,
% écrasé silencieusement puisque ce fichier est \input après ce préambule.
\fancyhead[LE]{\PURHHeaderFont\thepage}
\fancyhead[RO]{\PURHHeaderFont\thepage}
\renewcommand{\headrulewidth}{0pt}
\renewcommand{\footrulewidth}{0pt}

\fancypagestyle{plain}{%
  \fancyhf{}%
  \renewcommand{\headrulewidth}{0pt}%
  \renewcommand{\footrulewidth}{0pt}%
}

\usepackage[flushmargin]{footmisc}

\setlength{\footnotesep}{0.6\baselineskip}
% Espace avant les notes (référentiel §5.1 : 3 mm) : \skip\footins régit
% l'espace entre la fin du corps de texte et le début de la zone de notes
% (filet inclus) — 1,2\baselineskip (~5,7 mm à 11/13,5 pt) en donnait trop.
\setlength{\skip\footins}{3mm}
% Filet de notes (référentiel §5.1 : 0,25 pt de large sur 72 pt = 25,4 mm de
% long). \footnoterule par défaut de book.cls fait 0,4 pt sur
% 0,4\columnwidth (~42 mm sur l'empagement de ce profil) — les deux
% constats "environ 0,40 pt" / "environ 42 mm" du §5.2 correspondent très
% exactement à cette valeur par défaut, jamais personnalisée jusqu'ici.
\renewcommand{\footnoterule}{%
  \kern-3pt
  \hrule width 72pt height 0.25pt
  \kern 2.6pt
}
% \footnotelayout (footmisc) n'a plus d'effet : la redéfinition plus bas du
% texte de note pour le retrait négatif de première ligne applique
% directement \fontsize{8.5pt}{10.2pt} et court-circuite
% le mécanisme normal de footmisc qui l'invoque.

\usepackage{etoolbox}

% Citations observées : 9/11 pt, retrait gauche 10 mm (pas de retrait
% droit), ~4 mm avant/après (référentiel PURH §5.3 ; état antérieur :
% 11/14 pt, retraits gauche ET droit ≈1,5em, 8pt avant/après).
\renewenvironment{quote}
  {%
    \par\begingroup
    \fontsize{9pt}{11pt}\selectfont
    \list{}{\leftmargin=10mm\rightmargin=0pt}%
    \item\relax
  }
  {%
    \endlist
    \endgroup
  }

\AtBeginEnvironment{quote}{\vspace*{4mm plus 1pt minus 1pt}}
\AtEndEnvironment{quote}{\vspace*{4mm plus 1pt minus 1pt}}

\usepackage{graphicx}
\usepackage{caption}

\graphicspath{{media/}{images/}{assets/images/}}

\captionsetup{
  font=small,
  labelfont=bf,
  labelsep=period,
  skip=11pt
}

% Référentiel PURH v0.6 §11.3 : titre de tableau 9/11 pt, centré, 10 mm
% avant, 3,5 mm après — réglage propre aux tableaux (\captionsetup[table]),
% distinct du réglage générique ci-dessus qui reste seul à s'appliquer aux
% figures (aucune cible équivalente donnée pour elles dans le référentiel).
\DeclareCaptionFont{PURHTableCaptionFont}{\fontsize{9pt}{11pt}\selectfont}
\captionsetup[table]{
  font=PURHTableCaptionFont,
  labelfont=bf,
  labelsep=period,
  justification=centering,
  aboveskip=10mm,
  belowskip=3.5mm
}

\addto\captionsfrench{
  \renewcommand{\figurename}{Figure}
  \renewcommand{\tablename}{Tableau}
}

\usepackage{array}
\usepackage{longtable}
\usepackage{tabularx}
\usepackage{booktabs}

\usepackage{verse}
\usepackage{ragged2e}
\usepackage{xurl}
\urlstyle{same}

\usepackage[
  unicode=true,
  hidelinks,
  pdfusetitle
]{hyperref}

\usepackage{bookmark}

\hypersetup{
  pdftitle={\PURHBookTitle},
  pdfauthor={\PURHBookAuthor},
  pdfsubject={\PURHPublisher},
  pdfcreator={Impressions},
  pdfproducer={LuaLaTeX}
}

% Numéro calé à gauche avec retrait négatif de première ligne (observé
% directement sur le PDF imprimeur — le référentiel indiquait un point
% après le numéro, contredit par cette observation directe, suivie ici) :
% \leftskip aligne toutes les lignes de la note sur la même marge gauche ;
% \parindent négatif ramène uniquement la première ligne (celle du numéro)
% en deçà de cette marge, jusqu'au bord. \@thefnmark seul, sans point.
% Surtout PAS de \noindent ici : \noindent annule précisément l'effet du
% \parindent négatif qu'il est censé appliquer à la première ligne — bug
% réel constaté après vérification humaine du PDF généré (la première ligne
% restait alignée sur \leftskip comme les suivantes, aucun retrait négatif
% visible) ; laisser LaTeX indenter naturellement la première ligne de
% \parindent (donc la ramener à 0, à la marge) est ce qui produit le retrait
% négatif recherché.
% \AtBeginDocument, pas une simple redéfinition de préambule : hyperref/
% bookmark redéfinissent eux-mêmes \@makefntext pour y ajouter leurs
% ancres, mais le font via leur propre \AtBeginDocument — un
% \renewcommand direct ici, même placé après leur \usepackage, se faisait
% donc encore écraser au début du document (bug réel rencontré et vérifié :
% le correctif n'avait aucun effet tant qu'il n'était pas, lui aussi,
% différé). Étant chargés plus haut, leur crochet s'exécute avant celui-ci.
% \fontsize{8.5pt}{10.2pt} explicite ici, pas laissé au
% \footnotelayout de footmisc (référentiel §5.1 : Chaparral Pro 8,5/10,2 pt) :
% cette redéfinition complète de \@makefntext remplace celle de footmisc et
% n'appelle donc plus \footnotelayout — la taille du corps du texte de note
% retombait silencieusement sur celle ambiante (~9/11 pt observés, cf. §5.2)
% tant que ce défaut n'était pas identifié.
\makeatletter
\AtBeginDocument{%
  \renewcommand{\@makefntext}[1]{%
    \fontsize{8.5pt}{10.2pt}\selectfont
    \setlength{\leftskip}{1.2em}%
    \setlength{\parindent}{-1.2em}%
    \@thefnmark\enskip#1%
  }%
}
\makeatother

% -----------------------------------------------------------------
% Macros utilitaires PURH
% -----------------------------------------------------------------
% Visibilité auteur/affiliation sur l'ouverture de contribution (référentiel
% PURH v0.6 §7.2/§17 P1 item 3) : piloté par le profil de mise en page, pas
% par une valeur fixe — « le profil doit distinguer conservation des
% métadonnées et visibilité sur la page ». Doit être déclaré ici, avant que
% les macros LaTEI ne soient chargées : elles lisent ce drapeau dans
% \lateiContributionAuthor / \lateiContributionAffiliation mais ne le
% déclarent pas elles-mêmes.
\newif\iflateiShowContributionAuthor
\lateiShowContributionAuthorfalse

\newcommand{\PURHSeparator}{%
  \par\addvspace{1.5\baselineskip}%
  \noindent\rule{5cm}{0.4pt}%
  \par\addvspace{1.5\baselineskip}%
}

% Titre principal de la page de titre (référentiel PURH v0.6 §8.1,
% vérification humaine directe du 2026-08-04) : même hauteur de départ que
% le faux-titre (\PURHTitlePage partage désormais \vspace*{0.25\textheight}
% avec \PURHFalseTitle), mais corps plus grand — Josefin Sans Bold
% capitales, comme le faux-titre, juste agrandi. Aucune mesure
% millimétrique donnée par l'utilisateur pour ce corps précis : 22/26 pt
% choisi pour rester nettement plus grand que les 12/14 pt du faux-titre.
\newcommand{\PurhTitleMain}[1]{%
  {\PURHTitleFont\bfseries\fontsize{22pt}{26pt}\selectfont\centering\MakeUppercase{#1}\par}
}

% Sous-titre : Josefin Sans Bold bas de casse (pas de \MakeUppercase, à la
% différence du titre), un peu plus petit, centré sur deux lignes — largeur
% de boîte pour forcer le retour à la ligne, même principe que
% \PURHContributionTitleWidth pour les titres d'ouverture de contribution
% (référentiel §7.3 : reproduire la coupure plutôt que l'exiger en dur).
\newcommand{\PurhSubtitle}[1]{%
  \par\vspace{0.6\baselineskip}%
  \begin{center}
  \parbox{88mm}{\PURHTitleFont\bfseries\fontsize{15pt}{18pt}\selectfont\centering #1}
  \end{center}
  \vspace{0.4\baselineskip}%
}

% Responsabilité éditoriale : Chaparral Pro (fonte principale, aucun
% changement de famille) gras bas de casse, plus petit que le sous-titre,
% sur deux lignes explicites — « sous la direction de » toujours seul sur
% sa propre ligne, puis les noms sur la suivante (référentiel §8.1,
% vérification humaine directe du 2026-08-04). #1 est déjà le texte complet
% des deux lignes, séparées par \\ côté appelant (latei_driver.py), pas
% reconstruites ici : cette macro ne fait que les mettre en forme.
\newcommand{\PurhContributors}[1]{%
  \par\vspace{0.5\baselineskip}%
  {\bfseries\fontsize{11pt}{13pt}\selectfont\centering #1\par}%
  \vspace{0.6\baselineskip}%
}

% Mention finale de la page de titre (référentiel v0.7, vérification
% humaine directe du 2026-08-04, complétée le 2026-08-05) : nom complet de
% l'éditeur, majuscules grasses, sur UNE SEULE ligne, Chaparral (fonte
% principale, aucun changement de famille) — pas \PURHTitleFont/Josefin,
% contrairement à toute la titraille (§2.4/§3 du référentiel v0.7) : cette
% mention n'est pas un niveau de titraille, c'est une mention
% institutionnelle calée en bas de page.
%
% \resizebox{0.95\linewidth}{!}{...} plutôt qu'une simple
% \fontsize{14pt}{16pt} fixe (première version, abandonnée le
% 2026-08-05) : à 14 pt, « PRESSES UNIVERSITAIRES DE ROUEN ET DU HAVRE »
% (44 caractères) ne tenait pas sur une seule ligne à la largeur
% d'empagement de ce profil (~105 mm) et retombait sur deux lignes — bug
% réel constaté par vérification humaine directe du PDF généré.
% \resizebox emballe le texte dans une boîte horizontale non coupable
% (empêchant tout retour à la ligne, contrairement à un simple changement
% de \fontsize) puis la met à l'échelle pour occuper exactement 95 % de la
% largeur de la page — garantit à la fois la ligne unique et le remplissage
% « une bonne partie de la largeur de la page » quel que soit le nom
% affiché ou le profil, sans avoir à calculer un corps de police à la main.
\newcommand{\PurhPublisherMention}[1]{%
  \noindent\resizebox{0.95\linewidth}{!}{\bfseries\MakeUppercase{#1}}\par
}

\newenvironment{PurhBlockQuote}
  {%
    \begin{quote}
  }
  {%
    \end{quote}
  }

\newenvironment{PurhBibliography}
  {%
    \par
    \setlength{\parindent}{0pt}%
    \setlength{\leftskip}{0pt}%
  }
  {%
    \par
  }

% -----------------------------------------------------------------
% Liminaires PURH (référentiel v0.6 §8.1, P0 — "construire la séquence
% complète des liminaires") : pages blanches, faux-titre, crédits, page de
% titre. Construites depuis les métadonnées du livre par latei_driver.py,
% jamais depuis le corps LaTEI réversible. Toutes en pagestyle empty (ni
% folio ni titre courant), mais comptées dans la pagination — jamais
% \begin{titlepage} (dont le mode de compatibilité LaTeX 2.09 remet parfois
% \c@page à 1, un risque inutile ici où le compte doit au contraire
% continuer sans interruption depuis la toute première page physique).
% -----------------------------------------------------------------
\newcommand{\PURHBlankPage}{%
  \clearpage
  \thispagestyle{empty}%
  \mbox{}%
  \clearpage
}

% Vérification humaine directe du 2026-08-04 (faux-titre) :
% remonté sur la page (0.35\textheight -> 0.25\textheight, approximatif —
% « un peu plus haut », aucune mesure millimétrique donnée) et repassé en
% Josefin Sans Bold capitales, même corps que les titres de section
% (\titleformat{\section}, 12/14 pt) — le référentiel ne donnait aucune
% cible pour ce niveau spécifique, seule cette vérification fait foi ici,
% comme pour les autres corrections de graisse de ce chantier.
\newcommand{\PURHFalseTitle}[1]{%
  \clearpage
  \thispagestyle{empty}%
  \begin{center}
  \vspace*{0.25\textheight}
  {\PURHTitleFont\bfseries\fontsize{12pt}{14pt}\selectfont\MakeUppercase{#1}\par}
  \end{center}
  \clearpage
}

% Chaparral Pro (fonte principale du document, aucun changement de famille
% nécessaire) 10 pt — vérification humaine directe du 2026-08-04 : \small
% ne correspond pas nécessairement à 10 pt exactement (dépend de l'échelle
% de tailles du \documentclass), remplacé par une taille explicite.
% \vspace*{\fill} (pas \vspace*{0.3\textheight}, ni un \vfill non étoilé)
% : deuxième vérification humaine directe — le colophon doit être calé en
% bas de page, pas centré verticalement. Un \vfill nu en tout début de page
% est silencieusement absorbé par l'algorithme de coupure de page de TeX
% (jamais visible, confirmé par reproduction : le contenu restait centré
% près du haut) — surtout avec \raggedbottom actif dans ce document, qui
% autorise justement les pages à ne pas s'étirer jusqu'à \textheight.
% \vspace* (étoilé) protège la glue même en tête de page.
\newcommand{\PURHCreditsPage}[1]{%
  \clearpage
  \thispagestyle{empty}%
  \begin{center}
  \vspace*{\fill}
  \fontsize{10pt}{11.5pt}\selectfont
  #1
  \end{center}
  \clearpage
}

% Vérification humaine directe du 2026-08-04 : même hauteur de départ que
% le faux-titre (0.25\textheight, pas 0.15) — seul le corps du titre change
% entre les deux pages, pas sa position verticale.
\newcommand{\PURHTitlePage}[1]{%
  \clearpage
  \thispagestyle{empty}%
  \begin{center}
  \vspace*{0.25\textheight}
  #1
  \end{center}
  \clearpage
}

% ============================================================
% Macros LaTEI — régénérables, ne pas modifier
% ============================================================

% Experimental LaTEI macros for the reversible TEI/PURH convergence path.
% These definitions are typographic only; the reversible source remains the
% controlled LaTEI body produced by purh_site.reversible.latex_writer.

\usepackage{xparse}
\usepackage{xstring}
\usepackage{newunicodechar}

% Unicode spaces remain documentary characters in the reversible body; these
% typographic mappings only prevent missing glyphs in direct LaTEI compilation.
\newunicodechar{ }{~}
\newunicodechar{ }{\nobreak\,}
\newunicodechar{ }{\,}
\newunicodechar{‑}{\nobreak\hbox{-}}
\newunicodechar{″}{\ensuremath{\prime\prime}}
% Titres courants recto/verso distincts (référentiel PURH v0.5, "Titres
% courants") : verso (pages paires, RE) porte le titre du livre ou, à
% l'intérieur d'une partie, celui de la partie ; recto (pages impaires, LO)
% porte le titre court de la contribution en cours. Les deux commandes sont
% ré-exécutées à chaque page par fancyhdr — elles doivent donc rester des
% noms de macro, jamais leur valeur expansée au moment de cet \input.
\newcommand{\lateiCurrentRunningTitle}{}
\newcommand{\lateiVersoRunningTitle}{}
\fancyhead[RE]{\PURHHeaderFont\nouppercase{\lateiVersoRunningTitle}}
\fancyhead[LO]{\PURHHeaderFont\nouppercase{\lateiCurrentRunningTitle}}
\AtBeginDocument{%
  \fancyhead[RE]{\PURHHeaderFont\nouppercase{\lateiVersoRunningTitle}}%
  \fancyhead[LO]{\PURHHeaderFont\nouppercase{\lateiCurrentRunningTitle}}%
}

% Ouverture de contribution (référentiel PURH v0.5, "Ouvertures de
% contribution") : titre, sous-titre, auteur et traducteur viennent des
% <p rend="..."> du <div type="titlePage">, seule source de vérité visible —
% ne pas les dupliquer ailleurs. Titre en Josefin Sans Bold 16 pt capitales
% centré (vérification humaine directe du PDF généré face au PDF imprimeur :
% noir et gras là où le Thin — indiqué par le référentiel §2.5/§5.3, mais
% contredit par cette observation — rendait un texte maigre et grisâtre,
% peu lisible ; même correctif que \titleformat{\part}/{\chapter} dans le
% préambule). Sous-titre d'abord laissé en Thin Italic 12 pt (référentiel
% §5.3, non signalé comme défaut lors de la première vérification), puis
% re-corrigé en Bold Italic bas de casse le 2026-08-04 : nouvelle
% vérification humaine directe montrant explicitement gras + italique sur
% le PDF imprimeur pour ce niveau — pas de majuscules forcées, déjà correct
% (aucun \MakeUppercase ici, contrairement au titre). Positions verticales
% exactes (mm) non calibrées ici — hors périmètre de cette passe.
%
% Coupures de ligne du titre (référentiel PURH v0.5/v0.6 §7.3, "coupures
% éditoriales" — P1 item 4, exécuté le 2026-08-04 sur demande explicite,
% objectif : reproduire la règle implicite du PDF imprimeur plutôt qu'exiger
% une annotation manuelle par titre) : le référentiel avertit qu'il ne faut
% pas reconstruire ces coupures « par une largeur de boîte arbitraire ».
% Vérification empirique sur 7 titres réels du PDF imprimeur de *Dissimuler
% pour mieux régner* (7, 9, 15, 27, 26, 51, 191 caractères environ) montre
% cependant que la largeur de boîte N'EST PAS arbitraire : elle correspond
% de très près (104 mm calculé contre 103,4 mm mesuré indépendamment par
% analyse de pixels sur le rendu réel) à la largeur d'empagement du profil
% 155×230 (105 mm), et un simple \parbox à cette largeur reproduit déjà
% correctement, via le même algorithme Knuth-Plass que tout paragraphe
% LaTeX centré, le nombre de lignes et la coupure exacte de plusieurs titres
% tests (ex. "Les lieux de la conjuration : société secrète et hétérotopie
% dans la littérature romantique", 4 lignes identiques au PDF imprimeur).
% Sur d'autres titres, la coupure obtenue diffère de celle du PDF imprimeur
% d'un mot (ex. "Les espaces du secret à Clarens" → "LES ESPACES DU SECRET
% À" / "CLARENS" ici, contre "LES ESPACES DU SECRET" / "À CLARENS" dans le
% PDF imprimeur) — écart probablement dû à une règle InDesign de
% conservation des mots courts (« à », « de »…) avec la ligne suivante, que
% l'algorithme de LaTeX ne reproduit pas. Valeur retenue : 104mm, meilleur
% compromis observé plutôt qu'une largeur arbitraire non vérifiée.
\newcommand{\PURHContributionTitleWidth}{104mm}
\newcommand{\lateiContributionTitle}[1]{%
  \par\vspace{0.5\baselineskip}%
  \begin{center}
  \parbox{\PURHContributionTitleWidth}{\PURHTitleFont\bfseries\fontsize{16pt}{19pt}\selectfont\centering\MakeUppercase{#1}}%
  \end{center}
  \vspace{0.3\baselineskip}%
}
\newcommand{\lateiContributionSubtitle}[1]{%
  {\PURHTitleFont\bfseries\fontsize{12pt}{14pt}\selectfont\itshape\centering #1\par}%
  \vspace{0.4\baselineskip}%
}
% \iflateiShowContributionAuthor (défini dans le préambule, valeur pilotée
% par le profil de mise en page — référentiel PURH v0.6 §7.2/§17 P1 item 3 :
% « auteur et affiliation non imprimés » sur l'ouverture de contribution,
% mais métadonnées conservées) : le contenu réversible <p rend="author-aut">
% / <p rend="authority_affiliation"> reste inchangé dans le corps LaTEI, seul
% son affichage PDF est basculé par ce drapeau, jamais son existence.
% Signature de fin d'article (référentiel PURH v0.6 §8/§17, vérification
% humaine directe du 2026-08-04) : « les articles sont signés à la fin »,
% calés à droite — un emplacement distinct de l'ouverture de contribution
% (où auteur/affiliation restent masqués, cf. \iflateiShowContributionAuthor
% ci-dessus). \lateiContributionAuthor/\lateiContributionAffiliation
% capturent donc leur contenu dans ces macros globales, INDÉPENDAMMENT du
% drapeau d'affichage en tête de page, pour que \lateiRenderContributionSignature
% (appelée après le corps de la contribution, voir \lateiRenderFrontGroup/
% \lateiRenderChapterGroup) puisse la réémettre à la fin. Réinitialisées à
% vide au début de chaque contribution (mêmes macros) pour ne jamais laisser
% fuiter la signature d'une contribution vers la suivante qui n'en a pas.
% \@empty (le sentinelle vide du noyau LaTeX) exige la catégorie de code 11
% pour "@", seulement active entre \makeatletter et \makeatother — ce
% fichier n'en a pas besoin ailleurs et n'ouvre pas ce bloc ici. \@empty y
% serait tokenisé comme "\@" (caractère isolé) suivi du texte ordinaire
% "empty", qui s'imprimerait littéralement (bug réel constaté par
% compilation : "empty empty" visible sur la page). \lateiSignatureEmpty
% propre, sans "@", sert de sentinelle à la place.
\newcommand{\lateiSignatureEmpty}{}
\newcommand{\lateiSignatureAuthor}{}
\newcommand{\lateiSignatureAffiliation}{}
\newcommand{\lateiResetContributionSignature}{%
  \global\let\lateiSignatureAuthor\lateiSignatureEmpty
  \global\let\lateiSignatureAffiliation\lateiSignatureEmpty
}
% Nom d'auteur sous une entrée de TDM (référentiel §9.1) : écrit comme une
% ligne de TDM SÉPARÉE de l'entrée de chapitre elle-même, via
% \addtocontents (voir \latei_finish_contribution_toc_entry: plus bas) —
% jamais concaténé dans le texte transmis à \addcontentsline{{toc}}{{chapter}}
% : le paquet bookmark, qui construit les signets PDF automatiquement
% depuis CE texte, ne tolère pas les macros qu'on y ajouterait pour
% positionner le filet pointillé/numéro de page au niveau du titre plutôt
% que de l'auteur (bug réel constaté par compilation, voir le commentaire
% sur \titlecontents{{chapter}} dans latei_preamble.py). \par au lieu de la
% précédente rupture "\\" à l'intérieur d'un même paragraphe : ici, il
% s'agit bien d'un paragraphe séparé, plus de retrait suspendu de
% \titlecontents à préserver.
%
% \renewcommand{\textsc}[1]{#1} ICI, dans le corps de \lateiTocAuthorLine
% (vérification humaine directe du 2026-08-05, référentiel v0.7) : le nom
% de l'auteur doit apparaître bas de casse dans la TDM, alors que
% \lateiSignatureAuthor — passé tel quel, sans capture intermédiaire — porte
% le nom de famille en petites capitales (<hi rend="small-caps">Nom</hi>,
% routé vers \teiHi[rend={small-caps}]{Nom} puis \textsc{Nom}).
%
% Une première tentative capturait une copie "texte brut" à l'avance, au
% moment de \lateiContributionAuthor, via \protected@xdef + un \textsc local
% — abandonnée : \teiHi (comme toute commande définie par \NewDocumentCommand
% de xparse) est \protected au sens eTeX, donc jamais développée par un
% \edef/\xdef quel que soit l'état de \textsc au même moment ; le fichier
% .toc obtenu contenait encore \teiHi[rend={small-caps}]{Nom} tel quel (bug
% réel constaté : petites capitales toujours visibles dans le PDF malgré la
% capture). \renewcommand fonctionne en revanche PARFAITEMENT ici : ce bloc
% s'exécute au moment où \tableofcontents \input le fichier .toc et exécute
% RÉELLEMENT \lateiTocAuthorLine{...} — une exécution normale de macro, pas
% un \edef — où \teiHi peut s'exécuter (pas seulement se développer) et
% appeler ce \textsc local, correctement redéfini à ce moment précis.
\newcommand{\lateiTocAuthorLine}[1]{%
  \par\noindent\hspace*{1em}%
  {\bfseries\renewcommand{\textsc}[1]{##1}#1\par}%
}
\newcommand{\lateiContributionAuthor}[1]{%
  \global\def\lateiSignatureAuthor{#1}%
  \iflateiShowContributionAuthor
    {\normalsize\bfseries\centering #1\par}%
    \vspace{0.3\baselineskip}%
  \fi
}
\newcommand{\lateiContributionAffiliation}[1]{%
  \global\def\lateiSignatureAffiliation{#1}%
  \iflateiShowContributionAuthor
    {\small\centering #1\par}%
    \vspace{0.3\baselineskip}%
  \fi
}
% Chaparral Pro (fonte principale, aucun changement de famille), sans
% graisse, corps un peu plus petit que le corps du texte (11 pt -> 10 pt),
% calé à droite. Le prénom en bas de casse et le nom en petites capitales
% viennent déjà tels quels du balisage source (<p rend="author-aut">Prénom
% <hi rend="small-caps">Nom</hi></p>, déjà routé vers \teiHi[rend={small-caps}]
% ailleurs) : #1 n'a besoin d'aucune reformulation ici, seulement d'un
% nouvel habillage.
\newcommand{\lateiRenderContributionSignature}{%
  \ifx\lateiSignatureAuthor\lateiSignatureEmpty\else
    \par\vspace{0.8\baselineskip}%
    {\fontsize{10pt}{12pt}\selectfont\raggedleft\lateiSignatureAuthor\par}%
    \ifx\lateiSignatureAffiliation\lateiSignatureEmpty\else
      {\fontsize{10pt}{12pt}\selectfont\raggedleft\lateiSignatureAffiliation\par}%
    \fi
  \fi
}
% Espace vertical entre le bloc de titre d'ouverture (titre, sous-titre,
% auteur/affiliation) et le début du corps : absent du généré alors que le
% PDF imprimeur montre un grand blanc, estimé à 5-6 lignes de titre (16 pt /
% 19 pt d'interligne) par vérification humaine directe (chantier de parité
% v0.6, 2026-08-04) — 5,5 × 19 pt ≈ 105 pt. Valeur volontairement fixe, pas
% pilotée par le profil (comme les autres corps de la titraille), faute de
% mesure millimétrique du référentiel à ce niveau précis (le §7.2 documente
% une position absolue du corps depuis le haut de page, pas cet écart local).
\newcommand{\PURHContributionBodyGap}{\vspace{105pt}}
\newcommand{\lateiContributionTranslator}[1]{%
  {\small\itshape #1\par}%
  \vspace{0.6\baselineskip}%
}
\NewDocumentCommand{\lateiRenderParagraph}{O{} +m}{%
  \iflateiinfootnote
    #2\unskip\space
  \else
    \IfStrEq{\lateiHeadContext}{titlePage}{%
      \IfSubStr{#1}{rend={title-main}}{\lateiContributionTitle{#2}}{%
        \IfSubStr{#1}{rend={title-sub}}{\lateiContributionSubtitle{#2}}{%
          \IfSubStr{#1}{rend={author-aut}}{\lateiContributionAuthor{#2}}{%
            \IfSubStr{#1}{rend={authority\_affiliation}}{\lateiContributionAffiliation{#2}}{%
              \IfSubStr{#1}{rend={editor-trl}}{\lateiContributionTranslator{#2}}{\par #2\par}%
            }%
          }%
        }%
      }%
    }{%
      \IfStrEq{\lateiHeadContext}{figure}{%
        \IfSubStr{#1}{rend={caption}}{\par\small #2\par}{%
          \IfSubStr{#1}{rend={credits}}{\par\footnotesize #2\par}{\par #2\par}%
        }%
      }{\par #2\par}%
    }%
  \fi
}
\NewDocumentCommand{\teiP}{O{} +m}{\lateiRenderParagraph[#1]{#2}}

% Nested footnotes are not valid LaTeX. If a teiNote appears inside another
% note, render a visible symbolic marker and inline parenthesized content.
\newif\iflateiinfootnote
\lateiinfootnotefalse
\NewDocumentCommand{\teiNote}{O{} +m}{%
  \iflateiinfootnote
    \textsuperscript{*}%
  \else
    \begingroup
      \lateiinfootnotetrue
      \footnote{#2}%
    \endgroup
  \fi
}
\NewDocumentCommand{\teiRef}{O{} +m}{\lateiRenderRef[#1]{#2}}
\NewDocumentCommand{\teiTitle}{O{} +m}{%
  \IfSubStr{#1}{level={m}}{\textit{#2}}{%
    \IfSubStr{#1}{level={j}}{\textit{#2}}{%
      \IfSubStr{#1}{level={a}}{\enquote{#2}}{#2}%
    }%
  }%
}

% <formula notation="latex"> : le contenu documentaire porte déjà ses propres
% délimiteurs de mode mathématique (\(...\) ou \[...\]) — passthrough tel
% quel, sans échappement, pour que LaTeX le compose réellement au lieu
% d'afficher le code source échappé comme texte visible.
\NewDocumentCommand{\teiFormula}{O{} +m}{#2}

% Matter switches keep book.cls's structural side effects (cleardoublepage,
% numbered vs unnumbered chapters via \@mainmatterfalse/true) but not its
% page-numbering side effect. PURH pagination is arabe continue from the
% first page to the last (référentiel PURH v0.5 §5.5/§5.6, P0 — état de
% pagination) : liminaries are never roman-numbered, and the switch to the
% body never restarts the counter at 1. \pagenumbering resets the page
% counter every time it is called, so it must be called at most once for
% the whole document, on whichever of front/main/back matter starts first.
% Every "started" flag is set with \global: each group/chapter is rendered
% from inside its own TeX group (teiElement's xparse "+b" body argument,
% IfSubStr's branch braces...), and a *local* \xxxtrue would revert the
% instant that group closes — silently re-arming the guard and making
% \lateiEnsureMainMatter re-run \cleardoublepage + reset the page counter
% on every single chapter instead of only the first one.
\newif\iflateifrontmatterstarted
\newif\iflateimainmatterstarted
\newif\iflateibackmatterstarted
\newif\iflateiinbibliography
\newif\iflateipagenumberinginitialized
\lateifrontmatterstartedfalse
\lateimainmatterstartedfalse
\lateibackmatterstartedfalse
\lateiinbibliographyfalse
\lateipagenumberinginitializedfalse
\NewDocumentCommand{\lateiEnsureContinuousArabicPagination}{}{%
  \iflateipagenumberinginitialized\else
    \pagenumbering{arabic}%
    \global\lateipagenumberinginitializedtrue
  \fi
}
\makeatletter
\NewDocumentCommand{\lateiEnsureFrontMatter}{}{%
  \iflateifrontmatterstarted\else
    \cleardoublepage
    \global\@mainmatterfalse
    \lateiEnsureContinuousArabicPagination
    \global\lateifrontmatterstartedtrue
  \fi
}
\NewDocumentCommand{\lateiEnsureMainMatter}{}{%
  \iflateimainmatterstarted\else
    \cleardoublepage
    \global\@mainmattertrue
    \lateiEnsureContinuousArabicPagination
    \global\lateimainmatterstartedtrue
  \fi
}
\NewDocumentCommand{\lateiEnsureBackMatter}{}{%
  \iflateibackmatterstarted\else
    \if@openright\cleardoublepage\else\clearpage\fi
    \global\@mainmatterfalse
    \lateiEnsureContinuousArabicPagination
    \global\lateibackmatterstartedtrue
  \fi
}
\makeatother

\ExplSyntaxOn
\tl_new:N \l_latei_option_type_tl
\tl_new:N \l_latei_option_page_title_tl
\tl_new:N \l_latei_option_target_tl
\tl_new:N \l_latei_option_url_tl
\tl_new:N \l_latei_option_n_tl
\tl_new:N \l_latei_option_rend_tl
\tl_new:N \l_latei_option_xmlid_tl
\tl_new:N \l_latei_option_internaltarget_tl
\tl_new:N \l_latei_option_numcols_tl
\tl_new:N \l_latei_graphic_path_tl
\tl_new:N \l_latei_graphic_local_path_tl
\tl_new:N \l_latei_running_title_tl
\tl_new:N \g_latei_current_running_title_tl
\tl_new:N \g_latei_verso_running_title_tl
\prop_new:N \g_latei_graphics_map_prop
\prop_new:N \g_latei_running_titles_map_prop
\cs_generate_variant:Nn \prop_get:NnNTF { NVNTF }
% Le verso par défaut est le titre du livre ; \latei_markboth_verso:n le
% remplace par le titre de la partie en cours tant qu'aucune autre partie
% ne l'écrase (référentiel : "verso : titre du livre ou de la partie").
% \PURHBookTitle est déjà défini à ce stade (préambule inséré avant ce
% fichier), mais on stocke le nom de macro, pas sa valeur expansée : la
% même robustesse que pour \lateiCurrentRunningTitle plus haut.
\tl_gset:Nn \g_latei_verso_running_title_tl { \PURHBookTitle }
\RenewDocumentCommand{\lateiCurrentRunningTitle}{}
  {
    \tl_use:N \g_latei_current_running_title_tl
  }
\RenewDocumentCommand{\lateiVersoRunningTitle}{}
  {
    \tl_use:N \g_latei_verso_running_title_tl
  }
\NewDocumentCommand{\lateiDeclareGraphic}{m m}
  {
    \prop_gput:Nnn \g_latei_graphics_map_prop { #1 } { #2 }
  }
\NewDocumentCommand{\lateiDeclareRunningTitle}{m m}
  {
    \prop_gput:Nnn \g_latei_running_titles_map_prop { #1 } { #2 }
  }
% Recto (titre court de la contribution) et verso (titre du livre ou de la
% partie) sont mis à jour indépendamment — jamais la même valeur des deux
% côtés (référentiel : "passe deux fois la même valeur à \markboth").
\cs_new_protected:Npn \latei_markboth_recto:n #1
  {
    \prop_get:NnNTF \g_latei_running_titles_map_prop { #1 } \l_latei_running_title_tl
      { \tl_gset_eq:NN \g_latei_current_running_title_tl \l_latei_running_title_tl }
      { \tl_gset:Nn \g_latei_current_running_title_tl { #1 } }
    \markboth{\tl_use:N \g_latei_verso_running_title_tl}{\tl_use:N \g_latei_current_running_title_tl}
  }
\cs_new_protected:Npn \latei_markboth_verso:n #1
  {
    \prop_get:NnNTF \g_latei_running_titles_map_prop { #1 } \l_latei_running_title_tl
      { \tl_gset_eq:NN \g_latei_verso_running_title_tl \l_latei_running_title_tl }
      { \tl_gset:Nn \g_latei_verso_running_title_tl { #1 } }
    \markboth{\tl_use:N \g_latei_verso_running_title_tl}{\tl_use:N \g_latei_current_running_title_tl}
  }
\NewDocumentCommand{\lateiMarkBoth}{m}
  {
    \latei_markboth_recto:n { #1 }
  }
\NewDocumentCommand{\lateiMarkBothVerso}{m}
  {
    \latei_markboth_verso:n { #1 }
  }
\RenewDocumentCommand{\chaptermark}{m}{\lateiMarkBoth{#1}}
\cs_new_protected:Npn \latei_chapter:n #1
  {
    \prop_get:NnNTF \g_latei_running_titles_map_prop { #1 } \l_latei_running_title_tl
      {
        \tl_gset_eq:NN \g_latei_current_running_title_tl \l_latei_running_title_tl
        \chapter{#1}
        \lateiMarkBoth{#1}
      }
      {
        \tl_gset:Nn \g_latei_current_running_title_tl { #1 }
        \chapter{#1}
        \lateiMarkBoth{#1}
      }
  }
\NewDocumentCommand{\lateiChapter}{m}
  {
    \latei_chapter:n { #1 }
  }
\keys_define:nn { latei / options }
  {
    type .tl_set:N = \l_latei_option_type_tl,
    data-page-title .tl_set:N = \l_latei_option_page_title_tl,
    target .tl_set:N = \l_latei_option_target_tl,
    url .tl_set:N = \l_latei_option_url_tl,
    n .tl_set:N = \l_latei_option_n_tl,
    rend .tl_set:N = \l_latei_option_rend_tl,
    xmlid .tl_set:N = \l_latei_option_xmlid_tl,
    internaltarget .tl_set:N = \l_latei_option_internaltarget_tl,
    numcols .tl_set:N = \l_latei_option_numcols_tl,
    unknown .code:n = {},
  }
\cs_new_protected:Npn \latei_extract_options:n #1
  {
    \tl_clear:N \l_latei_option_type_tl
    \tl_clear:N \l_latei_option_page_title_tl
    \tl_clear:N \l_latei_option_target_tl
    \tl_clear:N \l_latei_option_url_tl
    \tl_clear:N \l_latei_option_n_tl
    \tl_clear:N \l_latei_option_rend_tl
    \tl_clear:N \l_latei_option_xmlid_tl
    \tl_clear:N \l_latei_option_internaltarget_tl
    \tl_clear:N \l_latei_option_numcols_tl
    \keys_set:nn { latei / options } { #1 }
  }
% Pose une cible PDF nommée d'après le xmlid TEI, s'il est présent dans les
% options déjà extraites par \latei_extract_options:n. Utilisé pour que les
% ancres/citations restent atteignables depuis \lateiRenderRef (cf. cibles
% "#id" internes) même quand l'élément lui-même n'a pas de macro de
% sectionnement propre (qui créerait sinon sa propre destination).
\cs_new_protected:Npn \latei_hypertarget_if_xmlid:
  {
    \tl_if_empty:NF \l_latei_option_xmlid_tl
      { \hypertarget{\tl_use:N \l_latei_option_xmlid_tl}{} }
  }
% Ponts sans "_"/":" : \ExplSyntaxOff plus bas rend ces caractères inaptes
% à composer un nom de macro hors de cette zone ; \teiHead et \teiCit, définis
% après \ExplSyntaxOff, doivent donc passer par ces noms classiques.
\NewDocumentCommand{\lateiExtractOptions}{m}
  {
    \latei_extract_options:n { #1 }
  }
\NewDocumentCommand{\lateiHypertargetIfXmlid}{}
  {
    \latei_hypertarget_if_xmlid:
  }
% Nombre de colonnes du tableau en cours (calculé côté Python, cf.
% _table_column_count dans latex_writer.py — LaTeX ne peut pas le déduire
% à la volée sans avoir déjà lu toutes les lignes).
\NewDocumentCommand{\lateiTableNumCols}{}
  {
    \tl_use:N \l_latei_option_numcols_tl
  }
% Rupture structurante d'ouverture de contribution : saut de page, entrée de
% TDM et titre courant à partir de data-page-title — mais aucun titre visible
% ni numérotation de chapitre ici (référentiel PURH v0.5, "Ouvertures de
% contribution" : « sans libellé ni numérotation de chapitre »). Le titre
% réellement affiché vient du <div type="titlePage"> rendu par ailleurs
% (\lateiContributionTitle) ; le dupliquer ici romprait cette règle.
% \cleardoublepage (jamais \clearpage) : « parties et articles sur recto,
% blancs techniques si nécessaire » (référentiel §5.2) — indépendant de
% l'option de classe openany/openright, qui ne régit que le \chapter brut
% du chemin <div type="chapter">, distinct de celui-ci.
% \setcounter{footnote}{0} : référentiel PURH v0.6 §5/§17, « notes remise à 1
% par contribution » — chaque ouverture (front matter, chapitre/article, back
% matter) redémarre sa propre numérotation de notes plutôt que de poursuivre
% celle de la contribution précédente sur tout le livre.
% \thispagestyle{empty} : référentiel PURH v0.6 §7.2/§17 P1 item 2, « aucun
% en-tête ni folio » sur l'ouverture de contribution — seule cette première
% page bascule en pagestyle vide ; les pages suivantes de la contribution
% reprennent le titre courant normalement (référentiel §6).
% \g_latei_pending_toc_title_tl : l'entrée de TDM elle-même est différée
% jusqu'à \latei_finish_contribution_toc_entry: (appelée après le corps de
% la contribution, comme \lateiRenderContributionSignature) — au moment où
% cette rupture s'exécute, \lateiSignatureAuthor n'est pas encore connu (la
% page de titre, qui le capture, n'a pas encore été rendue : elle fait
% partie de #2, qui s'exécute après). Référentiel PURH v0.6 §9.1,
% vérification humaine directe du 2026-08-04 : « prénom et nom de l'auteur,
% en gras » sous chaque titre de contribution dans la TDM.
\tl_new:N \g_latei_pending_toc_title_tl
\cs_new_protected:Npn \latei_add_contribution_opening_break:
  {
    \tl_if_empty:NF \l_latei_option_page_title_tl
      {
        \cleardoublepage
        \thispagestyle{empty}
        \phantomsection
        \tl_gset_eq:NN \g_latei_pending_toc_title_tl \l_latei_option_page_title_tl
        \exp_args:NV \latei_markboth_recto:n \l_latei_option_page_title_tl
        \setcounter{footnote}{0}
      }
  }
\cs_new_protected:Npn \latei_finish_contribution_toc_entry:
  {
    \tl_if_empty:NF \g_latei_pending_toc_title_tl
      {
        % Titre seul ici : filet pointillé + numéro de page restent gérés
        % par le 4e argument de \titlecontents{chapter} (préambule), comme
        % avant — l'auteur, lui, est écrit séparément ci-dessous, jamais
        % concaténé dans CE texte (voir le commentaire sur
        % \lateiTocAuthorLine, plus haut, pour la raison : compatibilité
        % avec les signets PDF automatiques du paquet bookmark).
        \addcontentsline{toc}{chapter}{\tl_use:N \g_latei_pending_toc_title_tl}
        \ifx\lateiSignatureAuthor\lateiSignatureEmpty\else
          \addtocontents{toc}{\protect\lateiTocAuthorLine{\lateiSignatureAuthor}}
        \fi
        \tl_gclear:N \g_latei_pending_toc_title_tl
      }
  }
\NewDocumentCommand{\lateiRenderFrontGroup}{O{} +m}
  {
    \lateiEnsureFrontMatter
    \latei_extract_options:n { #1 }
    \lateiResetContributionSignature
    \latei_add_contribution_opening_break:
    #2
    \latei_finish_contribution_toc_entry:
    \lateiRenderContributionSignature
  }
\NewDocumentCommand{\lateiRenderChapterGroup}{O{} +m}
  {
    \lateiEnsureMainMatter
    \latei_extract_options:n { #1 }
    \lateiResetContributionSignature
    \latei_add_contribution_opening_break:
    #2
    \latei_finish_contribution_toc_entry:
    \lateiRenderContributionSignature
  }
\NewDocumentCommand{\lateiRenderPartGroup}{O{} +m}
  {
    \lateiEnsureMainMatter
    % Toujours forcé ici, pas seulement via \lateiEnsureMainMatter (dont le
    % \cleardoublepage ne joue qu'à la toute première bascule vers le corps) :
    % chaque partie, même la 2e ou la 3e du livre, doit ouvrir en recto.
    \cleardoublepage
    % référentiel PURH v0.6 §7.1 : « page de style vide » sur l'ouverture de
    % partie — aucun en-tête ni folio, comme pour l'ouverture de contribution.
    \thispagestyle{empty}
    \begingroup
      \lateiSetHeadContext{part}
      #2
    \endgroup
  }
\NewDocumentCommand{\lateiRenderBackGroup}{O{} +m}
  {
    \lateiEnsureBackMatter
    \latei_extract_options:n { #1 }
    \latei_add_contribution_opening_break:
    #2
  }
\NewDocumentCommand{\lateiRenderRef}{O{} +m}
  {
    \latei_extract_options:n { #1 }
    \tl_if_empty:NTF \l_latei_option_internaltarget_tl
      {
        \tl_if_empty:NTF \l_latei_option_target_tl
          { #2 }
          { \href{\tl_use:N \l_latei_option_target_tl}{#2} }
      }
      { \hyperlink{\tl_use:N \l_latei_option_internaltarget_tl}{#2} }
  }
\cs_new_protected:Npn \latei_figure_fallback:
  {
    \fbox{\parbox{0.8\linewidth}{\centering\footnotesize Image absente ou non fournie}}
  }
\cs_new_protected:Npn \latei_include_graphic:n #1
  {
    \tl_if_blank:nTF { #1 }
      { \latei_figure_fallback: }
      {
        \IfFileExists{\detokenize{#1}}
          { \includegraphics[width=0.95\linewidth,keepaspectratio]{\detokenize{#1}} }
          { \latei_figure_fallback: }
      }
  }
\NewDocumentCommand{\lateiRenderGraphic}{O{}}
  {
    \latei_extract_options:n { #1 }
    \tl_clear:N \l_latei_graphic_path_tl
    \tl_if_empty:NF \l_latei_option_url_tl
      { \tl_set_eq:NN \l_latei_graphic_path_tl \l_latei_option_url_tl }
    \tl_if_empty:NT \l_latei_graphic_path_tl
      {
        \tl_if_empty:NF \l_latei_option_target_tl
          { \tl_set_eq:NN \l_latei_graphic_path_tl \l_latei_option_target_tl }
      }
    \tl_if_empty:NT \l_latei_graphic_path_tl
      {
        \tl_if_empty:NF \l_latei_option_n_tl
          { \tl_set_eq:NN \l_latei_graphic_path_tl \l_latei_option_n_tl }
      }
    \prop_get:NVNTF \g_latei_graphics_map_prop \l_latei_graphic_path_tl \l_latei_graphic_local_path_tl
      { \exp_args:NV \latei_include_graphic:n \l_latei_graphic_local_path_tl }
      { \exp_args:NV \latei_include_graphic:n \l_latei_graphic_path_tl }
  }
\ExplSyntaxOff

% Référentiel PURH v0.6 §11.2 : Chaparral Pro 10 pt, retrait suspendu 5 mm.
% \hangindent en em dépendait de la taille de fonte ambiante (jamais fixée
% par cette macro) ; passé en mm, valeur absolue conforme à la cible.
\NewDocumentCommand{\lateiBibliographyEntry}{+m}{%
  \par\noindent\fontsize{10pt}{12pt}\selectfont\hangindent=5mm\hangafter=1 #1\par
}
\NewDocumentCommand{\lateiBibliographyBlock}{+m}{%
  \par
  \begingroup
    \lateiinbibliographytrue
    \begin{PurhBibliography}
    #1%
    \end{PurhBibliography}
  \endgroup
  \par
}

% Heads are contextual in TEI: the same <head> can name a chapter, section,
% figure, table, or list. The surrounding LaTEI environment sets this context.
\newcommand{\lateiHeadContext}{fallback}
\newcommand{\lateiSetHeadContext}[1]{\def\lateiHeadContext{#1}}
\NewDocumentCommand{\lateiRenderHead}{m}{%
  % \part* ajoute déjà lui-même son entrée de TDM sous le shape [display] de
  % \titleformat{\part} (vérifié empiriquement) : un \addcontentsline{toc}
  % {part}{...} manuel ici produirait un doublon, pas un renfort.
  %
  % \MakeUppercase{#1} ici (et non \scshape/\textsc dans
  % \titlecontents{part}, préambule) : vérification humaine directe du
  % 2026-08-05 a montré que \scshape n'a AUCUN effet visible sur
  % \PURHTitleFont (Josefin Sans) — confirmé par compilation isolée, le
  % journal LaTeX rapporte explicitement « Font shape .../b/sc undefined »,
  % silencieusement remplacée par la forme normale (bug réel, jamais une
  % erreur bloquante, donc invisible sans vérifier le journal). Cette fonte
  % n'a pas de véritables petites capitales OpenType. Palliatif : capitales
  % pleines (\MakeUppercase), à la taille réduite (12 pt) déjà fixée par
  % \titlecontents{part} — plus petites que le titre de partie affiché sur
  % sa propre page (16 pt), donc lisibles comme un niveau distinct dans la
  % TDM, sans être de véritables petites capitales typographiques.
  % Appliqué ici (à la source du texte de la partie, avant \part*) plutôt
  % que dans \titlecontents{part} : \part* réutilise ce même #1, déjà
  % majuscule, pour son propre \addcontentsline{{toc}}{{part}}{{...}} interne —
  % la mise en majuscules du titre affiché sur la page (via le \MakeUppercase
  % déjà présent dans \titleformat{{\part}}) n'est pas affectée puisqu'elle
  % s'applique une seconde fois, sans effet, sur un texte déjà majuscule.
  \IfStrEq{\lateiHeadContext}{part}{\part*{\MakeUppercase{#1}}\lateiMarkBothVerso{#1}}{%
  \IfStrEq{\lateiHeadContext}{chapter}{\lateiChapter{#1}}{%
    \IfStrEq{\lateiHeadContext}{unnumberedchapter}{\chapter*{#1}\lateiMarkBoth{#1}\addcontentsline{toc}{chapter}{#1}}{%
    \IfStrEq{\lateiHeadContext}{section}{\section{#1}}{%
      \IfStrEq{\lateiHeadContext}{subsection}{\subsection{#1}}{%
        \IfStrEq{\lateiHeadContext}{subsubsection}{\subsubsection{#1}}{%
          \IfStrEq{\lateiHeadContext}{figure}{\par\smallskip\noindent\textbf{#1}\par\smallskip}{%
            \IfStrEq{\lateiHeadContext}{table}{\par\smallskip\noindent\textbf{#1}\par\smallskip}{%
              \IfStrEq{\lateiHeadContext}{list}{\item[]\textbf{#1}}{%
                \par\smallskip\noindent\textbf{#1}\par\smallskip% fallback head
              }%
            }%
          }%
        }%
      }%
    }%
  }%
  }%
  }%
}
\NewDocumentCommand{\teiHead}{O{} +m}{%
  \lateiExtractOptions{#1}%
  \lateiHypertargetIfXmlid
  \lateiRenderHead{#2}%
}

% Apply rend values predictably, including combined values in any order. This
% mirrors the stable renderer order: small caps / position / weight / italic.
\NewDocumentCommand{\lateiIfRendSmallCaps}{m m m}{%
  \IfSubStr{\detokenize{#1}}{\detokenize{small-caps}}{#2}{%
    \IfSubStr{\detokenize{#1}}{\detokenize{small_caps}}{#2}{%
      \IfSubStr{\detokenize{#1}}{\detokenize{small caps}}{#2}{#3}%
    }%
  }%
}
\NewDocumentCommand{\lateiIfRendBold}{m m m}{%
  \IfSubStr{\detokenize{#1}}{\detokenize{bold}}{#2}{\IfSubStr{\detokenize{#1}}{\detokenize{gras}}{#2}{#3}}%
}
\NewDocumentCommand{\lateiIfRendSup}{m m m}{%
  \IfSubStr{\detokenize{#1}}{\detokenize{sup}}{#2}{\IfSubStr{\detokenize{#1}}{\detokenize{exposant}}{#2}{#3}}%
}
\NewDocumentCommand{\lateiIfRendSub}{m m m}{%
  \IfSubStr{\detokenize{#1}}{\detokenize{sub}}{#2}{\IfSubStr{\detokenize{#1}}{\detokenize{indice}}{#2}{#3}}%
}
\NewDocumentCommand{\lateiIfRendUnderline}{m m m}{%
  \IfSubStr{\detokenize{#1}}{\detokenize{underline}}{#2}{\IfSubStr{\detokenize{#1}}{\detokenize{souligne}}{#2}{#3}}%
}
\NewDocumentCommand{\lateiIfRendStrikethrough}{m m m}{%
  \IfSubStr{\detokenize{#1}}{\detokenize{strikethrough}}{#2}{\IfSubStr{\detokenize{#1}}{\detokenize{barre}}{#2}{#3}}%
}
\NewDocumentCommand{\lateiIfRendUppercase}{m m m}{%
  \IfSubStr{\detokenize{#1}}{\detokenize{uppercase}}{#2}{\IfSubStr{\detokenize{#1}}{\detokenize{majuscule}}{#2}{#3}}%
}
\NewDocumentCommand{\lateiHiSmallCaps}{m +m}{%
  \lateiIfRendSmallCaps{#1}{\textsc{#2}}{#2}%
}
\NewDocumentCommand{\lateiHiPosition}{m +m}{%
  \lateiIfRendSup{#1}{\textsuperscript{#2}}{%
    \lateiIfRendSub{#1}{$_{#2}$}{#2}%
  }%
}
\NewDocumentCommand{\lateiHiBold}{m +m}{%
  \lateiIfRendBold{#1}{\textbf{#2}}{#2}%
}
\NewDocumentCommand{\lateiHiItalic}{m +m}{%
  \IfSubStr{\detokenize{#1}}{\detokenize{italic}}{\textit{#2}}{#2}%
}
\NewDocumentCommand{\lateiHiUnderline}{m +m}{%
  \lateiIfRendUnderline{#1}{\underline{#2}}{#2}%
}
\NewDocumentCommand{\lateiHiStrikethrough}{m +m}{%
  \lateiIfRendStrikethrough{#1}{\sout{#2}}{#2}%
}
\NewDocumentCommand{\lateiHiUppercase}{m +m}{%
  \lateiIfRendUppercase{#1}{\MakeUppercase{#2}}{#2}%
}
\NewDocumentCommand{\teiHi}{O{} +m}{%
  \lateiHiItalic{#1}{\lateiHiBold{#1}{\lateiHiPosition{#1}{\lateiHiSmallCaps{#1}{\lateiHiUnderline{#1}{\lateiHiStrikethrough{#1}{\lateiHiUppercase{#1}{#2}}}}}}}%
}

\NewDocumentCommand{\teiForeign}{O{} +m}{\textit{#2}}
\NewDocumentCommand{\teiTerm}{O{} +m}{#2}
\NewDocumentCommand{\teiName}{O{} +m}{#2}
\NewDocumentCommand{\teiPersName}{O{} +m}{#2}
\NewDocumentCommand{\teiPlaceName}{O{} +m}{#2}
\NewDocumentCommand{\teiOrgName}{O{} +m}{#2}
\NewDocumentCommand{\teiDate}{O{} +m}{#2}
\NewDocumentCommand{\teiNum}{O{} +m}{#2}
\NewDocumentCommand{\teiLabel}{O{} +m}{#2}
\NewDocumentCommand{\teiQ}{O{} +m}{\enquote{#2}}
\NewDocumentCommand{\teiSaid}{O{} +m}{\enquote{#2}}
\NewDocumentCommand{\teiAuthor}{O{} +m}{#2}
\NewDocumentCommand{\teiEditor}{O{} +m}{#2}
\NewDocumentCommand{\teiPublisher}{O{} +m}{#2}
\NewDocumentCommand{\teiBiblScope}{O{} +m}{#2}
\NewDocumentCommand{\teiIdno}{O{} +m}{#2}

\NewDocumentCommand{\teiGraphic}{O{}}{\lateiRenderGraphic[#1]}
\NewDocumentCommand{\teiPtr}{O{}}{}
\NewDocumentCommand{\teiLb}{O{}}{\linebreak{}}
\NewDocumentCommand{\teiPb}{O{}}{\ignorespaces}
% <anchor xml:id="..."/> : marqueur de position TEI, sans contenu visible —
% seule sa cible PDF (pour les <ref target="#..."> qui pointent dessus) compte.
\NewDocumentCommand{\teiAnchor}{O{}}{%
  \lateiExtractOptions{#1}%
  \lateiHypertargetIfXmlid
}

\NewDocumentEnvironment{teiDiv}{O{} +b}{%
  \IfSubStr{#1}{type={titlePage}}{%
    \begingroup
    \lateiSetHeadContext{titlePage}%
    #2%
    \PURHContributionBodyGap
    \endgroup
  }{%
    \begingroup
    \lateiSetHeadContext{fallback}%
    \IfSubStr{#1}{type={chapter}}{\lateiSetHeadContext{chapter}}{}%
    \IfSubStr{#1}{type={section}}{\lateiSetHeadContext{section}}{}%
    \IfSubStr{#1}{type={section1}}{\lateiSetHeadContext{section}}{}%
    \IfSubStr{#1}{type={section2}}{\lateiSetHeadContext{subsection}}{}%
    \IfSubStr{#1}{type={subsection}}{\lateiSetHeadContext{subsection}}{}%
    \IfSubStr{#1}{type={section3}}{\lateiSetHeadContext{subsubsection}}{}%
    #2%
    \endgroup
  }%
}{}
\NewDocumentEnvironment{teiList}{O{} +b}{%
  \begingroup
  \lateiSetHeadContext{list}%
  \IfSubStr{#1}{type={ordered}}{\begin{enumerate}}{\begin{itemize}}%
  #2%
  \IfSubStr{#1}{type={ordered}}{\end{enumerate}}{\end{itemize}}%
  \endgroup
}{}
\NewDocumentCommand{\teiItem}{O{} +m}{\item #2}
\NewDocumentEnvironment{teiQuote}{O{} +b}{\begin{PurhBlockQuote}#2\end{PurhBlockQuote}}{}
% <quote> utilisée en cours de phrase (dans <p>/<item>, ou n'importe où
% dans <note> — cf. _write_quote côté Python, qui décide du choix entre
% teiQuote et teiQuoteInline) : \teiQuote/PurhBlockQuote force un
% \begin{quote}, qui coupe le paragraphe même en plein milieu d'une phrase.
% \teiQuoteInline reste une macro simple (aucun \begingroup ni saut de
% paragraphe) et ajoute les guillemets via \enquote (csquotes), comme
% \teiQ déjà.
\NewDocumentCommand{\teiQuoteInline}{O{} +m}{\enquote{#2}}
% Poésie (référentiel PURH v0.5, "Poésie") : de vrais <lg>/<l> TEI passent
% par l'environnement `verse` (aligné à gauche, une ligne par \\, jamais
% justifié) — pas le bloc de citation, qui produisait des blancs excessifs
% en forçant des \linebreak dans un paragraphe justifié. <lg> imbriqués
% (strophes) s'emboîtent naturellement, `verse` gère leur espacement.
\NewDocumentEnvironment{teiLg}{O{} +b}{%
  \begin{verse}
  #2%
  \end{verse}
}{}
\NewDocumentCommand{\teiL}{O{} +m}{#2\\}
\NewDocumentEnvironment{teiFigure}{O{} +b}{%
  \begingroup
  \lateiSetHeadContext{figure}%
  \begin{center}#2\end{center}%
  \endgroup
}{}
\NewDocumentEnvironment{teiCit}{O{} +b}{%
  \lateiExtractOptions{#1}%
  \lateiHypertargetIfXmlid
  #2%
}{}
\NewDocumentEnvironment{teiBibl}{O{} +b}{\lateiBibliographyEntry{#2}}{}

% Tableau réel (longtable/booktabs), pas une simple suite de paragraphes
% centrés : le nombre de colonnes (numcols, calculé côté Python — cf.
% _table_column_count) fixe le gabarit de colonnes. Les teiCell d'une ligne
% sont déjà séparées par un "&" littéral dans la source (écrit côté Python,
% cf. _write_row), et teiRow/teiCell sont des macros (pas des environnements) :
% \multicolumn doit être le tout premier jeton lu par le scanner d'alignement
% de TeX pour une cellule fusionnée (@cols>1) — même non protégée, l'envelopper
% dans une macro (et a fortiori un environnement, qui insère un \begingroup
% avant le corps) le rend invisible à la détection de \omit ("Misplaced
% \omit", confirmé par reproduction isolée). \multicolumn est donc écrit tel
% quel côté Python (cf. _write_cell) pour les cellules fusionnées ; teiCell ne
% gère plus que la mise en forme (gras d'entête), jamais \multicolumn. Une
% ligne d'entête (role=label/header sur la ligne, cf. TEI) reçoit
% \toprule/\midrule ; \bottomrule ferme le tableau.
% Référentiel PURH v0.6 §11.3 : texte 8/9,5 pt, marges internes 2 mm,
% filets 0,25 pt, fond d'entête noir 30 % (un gris clair : 30 % de couverture
% d'encre, pas un fond sombre). Chaparral Pro déjà la fonte principale du
% document (\setmainfont) : aucun changement de famille nécessaire, seule la
% taille change ici. \heavyrulewidth/\lightrulewidth/\cmidrulewidth
% (booktabs) valent par défaut 0,08em/0,05em/0,05em — proportionnels à la
% fonte ambiante, donc jamais 0,25 pt de façon fiable sans réglage explicite.
\newif\iflateirowisheader
\lateirowisheaderfalse
\NewDocumentEnvironment{teiTable}{O{} +b}{%
  \lateiExtractOptions{#1}%
  \begingroup
  \lateiSetHeadContext{table}%
  \fontsize{8pt}{9.5pt}\selectfont
  \setlength{\tabcolsep}{2mm}
  \setlength{\heavyrulewidth}{0.25pt}
  \setlength{\lightrulewidth}{0.25pt}
  \setlength{\cmidrulewidth}{0.25pt}
  \par\begin{center}%
  \begin{longtable}{*{\lateiTableNumCols}{l}}
  \toprule
  #2%
  \bottomrule
  \end{longtable}%
  \end{center}\par%
  \endgroup
}{}
% \rowcolor pour les lignes d'entête est écrit directement par
% purh_site/reversible/latex_writer.py (_write_row), avant même ce \teiRow,
% jamais ici : même contrainte que \multicolumn sur une cellule fusionnée
% (cf. teiCell) — un \rowcolor conditionnel à l'intérieur de cette macro
% n'est plus le premier jeton vu par le scanner d'alignement de TeX et
% produit "Misplaced \noalign" (confirmé par reproduction isolée).
\NewDocumentCommand{\teiRow}{O{} +m}{%
  \IfSubStr{#1}{role={label}}{\lateirowisheadertrue}{%
    \IfSubStr{#1}{role={header}}{\lateirowisheadertrue}{\lateirowisheaderfalse}%
  }%
  #2%
  \\%
  \iflateirowisheader\midrule\lateirowisheaderfalse\fi
}
\NewDocumentCommand{\teiCell}{O{} +m}{%
  \lateiExtractOptions{#1}%
  \iflateirowisheader\bfseries\fi
  \IfSubStr{#1}{role={label}}{\bfseries}{\IfSubStr{#1}{role={header}}{\bfseries}{}}%
  #2%
}
% teiHeader is metadata, not running text. Document wrappers pass through;
% unknown non-header elements keep their content visible rather than vanish.
\NewDocumentEnvironment{teiElement}{O{} +b}{%
  \IfSubStr{#1}{name={teiHeader}}{}{%
    \IfSubStr{#1}{name={group}}{%
      \IfSubStr{#1}{type={section1}}{\lateiRenderPartGroup[#1]{#2}}{%
      \IfSubStr{#1}{type={article}}{\lateiRenderChapterGroup[#1]{#2}}{%
      \IfSubStr{#1}{type={chapter}}{\lateiRenderChapterGroup[#1]{#2}}{%
      \IfSubStr{#1}{type={acknowledgments}}{\lateiRenderFrontGroup[#1]{#2}}{%
      \IfSubStr{#1}{type={abbreviations}}{\lateiRenderFrontGroup[#1]{#2}}{%
      \IfSubStr{#1}{type={introduction}}{\lateiRenderFrontGroup[#1]{#2}}{%
      \IfSubStr{#1}{type={foreword}}{\lateiRenderFrontGroup[#1]{#2}}{%
      \IfSubStr{#1}{type={preface}}{\lateiRenderFrontGroup[#1]{#2}}{%
      \IfSubStr{#1}{type={dedication}}{\lateiRenderFrontGroup[#1]{#2}}{%
      \IfSubStr{#1}{type={conclusion}}{\lateiRenderBackGroup[#1]{#2}}{%
      \IfSubStr{#1}{type={bibliography}}{\lateiRenderBackGroup[#1]{#2}}{#2}%
      }}}}}}}}}}%
    }{%
      \IfSubStr{#1}{name={back}}{\lateiEnsureBackMatter #2}{%
      \IfSubStr{#1}{name={listBibl}}{\lateiBibliographyBlock{#2}}{%
      \IfSubStr{#1}{name={biblStruct}}{\lateiBibliographyEntry{#2}}{#2}%
      }}%
    }%
  }%
}{}

% ---------------------------------------------------------------------------
% Commandes \latei... — corrections de mise en page PDF uniquement.
% Ces commandes ne sont pas restaurées vers XML lors du retour Métopes.
% Le parseur LaTEI parcourt leur contenu sans créer d'élément TEI.
% ---------------------------------------------------------------------------

% Macro interne — traduit les tailles normalisées en dimension PDF.
% Valeurs acceptées : small (~0,5 ligne), medium (~1 ligne), large (~2 lignes).
\NewDocumentCommand{\lateiApplyVSpace}{m}{%
  \IfStrEqCase{#1}{%
    {small}{\addvspace{0.5\baselineskip}}%
    {medium}{\addvspace{\baselineskip}}%
    {large}{\addvspace{2\baselineskip}}%
  }[\PackageError{latei}{%
    Valeur inconnue : #1. Valeurs acceptees : small, medium, large.%
  }{}]%
}

% Indentation — alinéas
% --------------------------------------------------------------------------

% \lateiNoIndent{...} — supprime le retrait de première ligne du contenu.
\NewDocumentCommand{\lateiNoIndent}{+m}{{\parindent=0pt #1}}

% \lateiIndent{...} — force le retrait de première ligne même si LaTeX le supprime.
\newlength{\lateiParindent}
\AtBeginDocument{\setlength{\lateiParindent}{\parindent}}
\NewDocumentCommand{\lateiIndent}{+m}{{\setlength{\parindent}{\lateiParindent}#1}}

% Espacement vertical autonome
% --------------------------------------------------------------------------

% \lateiVSpace{small|medium|large} — insère un blanc vertical entre deux éléments.
\NewDocumentCommand{\lateiVSpace}{m}{\par\lateiApplyVSpace{#1}}

% Espacement vertical enveloppant
% --------------------------------------------------------------------------

% \lateiSpaceBefore{small|medium|large}{...} — ajoute un blanc avant le contenu.
\NewDocumentCommand{\lateiSpaceBefore}{m +m}{\par\lateiApplyVSpace{#1}#2}

% \lateiSpaceAfter{small|medium|large}{...} — ajoute un blanc après le contenu.
\NewDocumentCommand{\lateiSpaceAfter}{m +m}{#2\par\lateiApplyVSpace{#1}}

% Sauts de page autonomes
% --------------------------------------------------------------------------

% \lateiPageBreak — force un saut de page.
\NewDocumentCommand{\lateiPageBreak}{}{\newpage}

% \lateiClearPage — force une nouvelle page en vidant les flottants.
\NewDocumentCommand{\lateiClearPage}{}{\clearpage}

% Sauts de page enveloppants
% --------------------------------------------------------------------------

% \lateiPageBreakBefore{...} — force un saut de page avant le contenu.
\NewDocumentCommand{\lateiPageBreakBefore}{+m}{\newpage#1}

% \lateiPageBreakAfter{...} — force un saut de page après le contenu.
\NewDocumentCommand{\lateiPageBreakAfter}{+m}{#1\newpage}

% \lateiClearPageBefore{...} — vide les flottants et force une nouvelle page avant.
\NewDocumentCommand{\lateiClearPageBefore}{+m}{\clearpage#1}

% \lateiClearPageAfter{...} — force une nouvelle page avec vidage des flottants après.
\NewDocumentCommand{\lateiClearPageAfter}{+m}{#1\clearpage}

% Contrôle des coupures de page
% --------------------------------------------------------------------------

% \lateiKeepWithNext{...} — interdit une coupure de page immédiatement après.
\NewDocumentCommand{\lateiKeepWithNext}{+m}{#1\nopagebreak[4]}

% \lateiKeepTogether{...} — maintient le contenu entier sur une même page.
\NewDocumentCommand{\lateiKeepTogether}{+m}{\begin{samepage}#1\end{samepage}}

% \lateiNoPageBreakBefore{...} — interdit une coupure de page immédiatement avant.
\NewDocumentCommand{\lateiNoPageBreakBefore}{+m}{\nopagebreak[4]#1}

% \lateiNoPageBreakAfter{...} — interdit une coupure de page immédiatement après.
\NewDocumentCommand{\lateiNoPageBreakAfter}{+m}{#1\nopagebreak[4]}

% lateiDocument marks the reversible editorial zone in a monofile.
% Completely transparent at compilation — only serves as a delimiter for the parser.
\newenvironment{lateiDocument}{}{}

% ============================================================
% Mappings graphiques — régénérables, ne pas modifier
% ============================================================

% Experimental LaTEI graphic mapping.
% This file is a compilation artifact, not a reversible source.

% ============================================================
% Mappings titres courants — régénérables, ne pas modifier
% ============================================================

% Experimental LaTEI running-title mapping.
% This file is a compilation artifact, not a reversible source.
\lateiDeclareRunningTitle{Le pouvoir et le cabinet : Wolmar}{Le pouvoir et le cabinet : Wolmar}
\lateiDeclareRunningTitle{Les coulisses et le jardin : Julie}{Les coulisses et le jardin : Julie}
\lateiDeclareRunningTitle{« Malheur à qui n’a plus rien à désirer »}{« Malheur à qui n’a plus rien à désirer »}
\lateiDeclareRunningTitle{Les lieux de la conjuration : société secrète et hétérotopie dans la littérature romantique}{Les lieux de la conjuration : société secrète…}
\lateiDeclareRunningTitle{Fonction politique du locus secretus  : une hétérotopie de compensation}{Fonction politique du locus secretus : une hétérotopie…}
\lateiDeclareRunningTitle{Un caveau en forêt : espace et récit du secret}{Un caveau en forêt : espace et récit du secret}
\lateiDeclareRunningTitle{Le théâtre des opérations : de la politique du coup d’éclat à la politique du secret}{Le théâtre des opérations : de la politique du coup…}
\lateiDeclareRunningTitle{Du refuge au piège : le caveau, miroir du devenir historique}{Du refuge au piège : le caveau, miroir du devenir…}
\lateiDeclareRunningTitle{Le cabinet de toilette dans Son Excellence Eugène Rougon  : incursion dans « le déshabillé de l\lbrack{}a\rbrack{} politique »}{Le cabinet de toilette dans Son Excellence Eugène Rougon…}
\lateiDeclareRunningTitle{Portrait de l’empereur au cabinet de toilette dans La Débâcle}{Portrait de l’empereur au cabinet de toilette…}
\lateiDeclareRunningTitle{Énonciation référentielle et révélation des secrets du pouvoir chez Pontis, Sirot et Lenet}{Énonciation référentielle et révélation des secrets…}
\lateiDeclareRunningTitle{Le je et le pouvoir : une représentation des coulisses}{Le je et le pouvoir : une représentation des coulisses}
\lateiDeclareRunningTitle{Montage narratif : un je entre inclusion et exclusion}{Montage narratif : un je entre inclusion et exclusion}
\lateiDeclareRunningTitle{Visages discursifs du secret : un je acteur ou victime}{Visages discursifs du secret : un je acteur ou victime}
\lateiDeclareRunningTitle{Je narré et je narrant : le problème de la connaissance des secrets}{Je narré et je narrant : le problème de la connaissance…}
\lateiDeclareRunningTitle{Je narré et je narrant : représentation et révélation des secrets}{Je narré et je narrant : représentation et révélation…}
\lateiDeclareRunningTitle{L’immixtion (in)vraisemblable du je dans les secrets d’autrui}{L’immixtion (in)vraisemblable du je dans les secrets…}
\lateiDeclareRunningTitle{L’échange épistolaire, lieu du secret politique dans La Fortune des Rougon}{L’échange épistolaire, lieu du secret politique…}
\lateiDeclareRunningTitle{Le « synoptisme » du pouvoir}{Le « synoptisme » du pouvoir}
\lateiDeclareRunningTitle{Le spectacle oxymoronique du pouvoir oriental : un anti-miroir des princes}{Le spectacle oxymoronique du pouvoir oriental…}
\lateiDeclareRunningTitle{Observation, exploration et espionnage dans la Relation du voyage de la mer du Sud d’Amédée Frézier (1716)}{Observation, exploration et espionnage dans la Relation…}
\lateiDeclareRunningTitle{Être un explorateur : arpenter et affronter le territoire}{Être un explorateur : arpenter et affronter le territoire}
\lateiDeclareRunningTitle{L’observation : voir et savoir}{L’observation : voir et savoir}
\lateiDeclareRunningTitle{L’adresse à Richelieu et à ses descendants : cadrage incitatif des lieux oculaires}{L’adresse à Richelieu et à ses descendants : cadrage…}
\lateiDeclareRunningTitle{Nekuia et convocation : Le Portrait du Grand Cardinal et L’Ombre du grand Armand}{Nekuia et convocation : Le Portrait du Grand Cardinal…}
\lateiDeclareRunningTitle{Épiphanies du visible : Les Muse s et Le Grand Exemple}{Épiphanies du visible : Les Muse s et Le Grand Exemple}
\lateiDeclareRunningTitle{Inscrire et éterniser : Le Temple}{Inscrire et éterniser : Le Temple}
\lateiDeclareRunningTitle{Secrets de théâtre et héros « en déshabillé », de la comédie des Lumières au drame romantique}{Secrets de théâtre et héros « en déshabillé »…}
\lateiDeclareRunningTitle{Deux modèles : Collé et Beaumarchais}{Deux modèles : Collé et Beaumarchais}

\begin{document}

\lateiEnsureContinuousArabicPagination
\PURHBlankPage
\PURHBlankPage
\PURHFalseTitle{Dissimuler pour mieux régner}
\PURHCreditsPage{Suivi éditorial : Anaïs Lebreton\par\vspace{0.5\baselineskip}© Presses universitaires de Rouen et du Havre.\par\vspace{0.1\baselineskip}2 place Émile Blondel – 76821 Mont-Saint-Aignan Cedex\par\vspace{0.1\baselineskip}\url{http://purh.univ-rouen.fr}\par}
\PURHTitlePage{\PurhTitleMain{Dissimuler pour mieux régner}
\PurhSubtitle{<em>Locus politicus, locus secretus</em> en littérature (\textsc{XVII}e-\textsc{XIX}e siècles)}
\vspace{2\baselineskip}
\PurhContributors{sous la direction de\\Floriane Daguisé et Florence Fix}
\vspace*{\fill}
\PurhPublisherMention{Presses universitaires de Rouen et du Havre}}
\PURHBlankPage

% ============================================================
% Zone éditoriale réversible — corrections autorisées ici
% ============================================================

\begin{lateiDocument}
\begin{teiElement}[name={TEI},change={commons\_edition}]

  \begin{teiElement}[name={teiHeader}]

    \begin{teiElement}[name={fileDesc}]

      \begin{teiElement}[name={titleStmt}]

        \teiTitle[type={main}]{Dissimuler pour mieux régner}
        \teiTitle[type={sub}]{<em>Locus politicus, locus secretus</em> en littérature (XVIIe-XIXe siècles)}
        \teiAuthor[role={pbd}]{
          \teiPersName{
            \begin{teiElement}[name={forename}]
Floriane
\end{teiElement}
            \begin{teiElement}[name={surname}]
Daguisé
\end{teiElement}
          }
          \begin{teiElement}[name={affiliation}]

            \teiOrgName{Université de Rouen Normandie, CÉRÉDI, Rouen, France}
          
\end{teiElement}
        }
        \teiAuthor[role={pbd}]{
          \teiPersName{
            \begin{teiElement}[name={forename}]
Florence
\end{teiElement}
            \begin{teiElement}[name={surname}]
Fix
\end{teiElement}
          }
          \begin{teiElement}[name={affiliation}]

            \teiOrgName{Université de Rouen Normandie, CÉRÉDI, Rouen, France}
          
\end{teiElement}
        }
      
\end{teiElement}
      \begin{teiElement}[name={editionStmt}]

        \begin{teiElement}[name={edition}]

          \teiDate{}
        
\end{teiElement}
        \begin{teiElement}[name={respStmt}]

          \begin{teiElement}[name={resp}]

\end{teiElement}
          \teiName{Anaïs Lebreton}
        
\end{teiElement}
        \begin{teiElement}[name={sponsor}]

\end{teiElement}
      
\end{teiElement}
      \begin{teiElement}[name={publicationStmt}]

        \begin{teiElement}[name={ab},type={expression}]

          \begin{teiBibl}[xmlid={bibl-001}]

            \teiPublisher{PURH}
            \begin{teiElement}[name={address}]

              \begin{teiElement}[name={addrLine}]

\end{teiElement}
            
\end{teiElement}
            \begin{teiElement}[name={pubPlace}]

\end{teiElement}
            \begin{teiElement}[name={availability}]

              \begin{teiElement}[name={licence},target={-}]

\end{teiElement}
              \teiP[xmlid={p-001}]{}
            
\end{teiElement}
          
\end{teiBibl}
        
\end{teiElement}
        \begin{teiElement}[name={ab},type={book}]

          \begin{teiBibl}[xmlid={bibl-002}]

            \teiIdno[type={ISBN-13}]{}
            \teiIdno[type={ISSN}]{979-10-240-1887-4}
            \teiDate[type={publishing}]{}
            \teiDate[type={legal-deposit}]{}
            \teiDate[type={imprint}]{}
            \teiBiblScope[unit={page}]{}
            \begin{teiElement}[name={availability}]

              \begin{teiElement}[name={licence},target={-}]

\end{teiElement}
              \teiP[xmlid={p-002}]{}
            
\end{teiElement}
            \begin{teiElement}[name={distributor}]

\end{teiElement}
            \teiName[type={sales-agent}]{}
            \teiRef[type={site},target={-}]{}
          
\end{teiBibl}
          \teiName[type={printer}]{
            \teiName{}
            \teiNum{}
          }
          \begin{teiElement}[name={address}]

            \begin{teiElement}[name={addrLine}]

\end{teiElement}
          
\end{teiElement}
          \begin{teiElement}[name={dim},type={printingvolume}]

\end{teiElement}
          \begin{teiElement}[name={dimensions}]

            \begin{teiElement}[name={width},unit={mm}]

\end{teiElement}
            \begin{teiElement}[name={height},unit={mm}]

\end{teiElement}
            \begin{teiElement}[name={depth},unit={mm}]

\end{teiElement}
          
\end{teiElement}
          \begin{teiElement}[name={dim},type={weight},unit={gr}]

\end{teiElement}
          \begin{teiElement}[name={extent}]

\end{teiElement}
          \begin{teiElement}[name={measure},type={price},unit={EUR}]

\end{teiElement}
        
\end{teiElement}
        \begin{teiElement}[name={ab},subtype={PDF},type={digital\_download}]

          \begin{teiBibl}[xmlid={bibl-003}]

            \teiIdno[type={ISBN}]{}
            \teiIdno[subtype={\#\#},type={DOI}]{}
            \teiRef[type={DOI}]{}
            \teiRef[type={site},target={-}]{}
            \begin{teiElement}[name={distributor}]

\end{teiElement}
            \teiDate[type={publishing}]{}
            \begin{teiElement}[name={availability}]

              \begin{teiElement}[name={licence}]

\end{teiElement}
              \teiP[xmlid={p-003}]{}
            
\end{teiElement}
            \teiName[type={sales-agent}]{}
          
\end{teiBibl}
          \begin{teiElement}[name={dim},type={weight},unit={Mo}]

\end{teiElement}
          \begin{teiElement}[name={measure},type={price},unit={EUR}]

\end{teiElement}
        
\end{teiElement}
        \begin{teiElement}[name={ab},subtype={EPUB},type={digital\_download}]

          \begin{teiBibl}[xmlid={bibl-004}]

            \teiIdno[type={ISBN-13}]{}
            \teiIdno[type={DOI}]{}
            \teiRef[type={DOI}]{}
            \teiRef[type={site},target={-}]{}
            \begin{teiElement}[name={distributor}]

\end{teiElement}
            \teiDate[type={publishing}]{}
            \begin{teiElement}[name={availability}]

              \begin{teiElement}[name={licence}]

\end{teiElement}
              \teiP[xmlid={p-004}]{}
            
\end{teiElement}
            \teiName[type={sales-agent}]{}
          
\end{teiBibl}
          \begin{teiElement}[name={dim},extent={},type={weight},unit={Mo}]

\end{teiElement}
          \begin{teiElement}[name={measure},type={price},unit={EUR}]

\end{teiElement}
          \begin{teiElement}[name={desc},type={accessMode}]
textual
\end{teiElement}
          \begin{teiElement}[name={desc},type={accessibilityFeature}]
ARIA
\end{teiElement}
          \begin{teiElement}[name={desc},type={accessibilityFeature}]
highContrastDisplay
\end{teiElement}
          \begin{teiElement}[name={desc},type={accessibilityFeature}]
readingOrder
\end{teiElement}
          \begin{teiElement}[name={desc},type={accessibilityFeature}]
structuralNavigation
\end{teiElement}
          \begin{teiElement}[name={desc},type={accessibilityFeature}]
tableOfContents
\end{teiElement}
          \begin{teiElement}[name={desc},type={accessibilityFeature}]
unlocked
\end{teiElement}
          \begin{teiElement}[name={desc},subtype={soundHazard},type={accessibilityHazard}]
noSoundHazard
\end{teiElement}
          \begin{teiElement}[name={desc},subtype={flashingHazard},type={accessibilityHazard}]
noFlashingHazard
\end{teiElement}
          \begin{teiElement}[name={desc},subtype={motionSimulationHazard},type={accessibilityHazard}]
noMotionSimulationHazard
\end{teiElement}
          \begin{teiElement}[name={desc},type={accessModeSufficient}]
textual
\end{teiElement}
          \begin{teiElement}[name={desc},type={accessibilitySummary}]

\end{teiElement}
        
\end{teiElement}
        \begin{teiElement}[name={ab},subtype={HTML},type={digital\_online}]

          \begin{teiBibl}[xmlid={bibl-005}]

            \begin{teiElement}[name={distributor}]

\end{teiElement}
            \teiName[type={CMS}]{}
            \teiIdno[type={ISBN}]{}
            \teiIdno[type={DOI}]{}
            \teiRef[type={DOI}]{}
            \teiRef[type={site},internaltarget={}]{}
            \teiDate[type={publishing}]{}
            \teiDate[type={embargoend}]{}
            \begin{teiElement}[name={availability}]

              \begin{teiElement}[name={licence}]

\end{teiElement}
              \teiP[xmlid={p-005}]{}
            
\end{teiElement}
          
\end{teiBibl}
        
\end{teiElement}
      
\end{teiElement}
      \begin{teiElement}[name={sourceDesc}]

        \begin{teiBibl}[xmlid={bibl-006}]

\end{teiBibl}
        \begin{teiBibl}[type={edition1},xmlid={bibl-007}]

          \teiDate{}
          \begin{teiElement}[name={extent},rend={format}]

\end{teiElement}
          \teiPublisher{}
          \teiIdno[type={ISBN}]{}
          \begin{teiElement}[name={edition},rend={type}]

\end{teiElement}
        
\end{teiBibl}
      
\end{teiElement}
    
\end{teiElement}
    \begin{teiElement}[name={profileDesc}]

      \begin{teiElement}[name={langUsage}]

        \begin{teiElement}[name={language},ident={fr-FR}]

\end{teiElement}
      
\end{teiElement}
      \begin{teiElement}[name={textClass}]

        \begin{teiElement}[name={keywords},scheme={keywords},xmllang={fr}]

          \begin{teiList}

            \teiItem{}
            \teiItem{}
            \teiItem{}
          
\end{teiList}
        
\end{teiElement}
      
\end{teiElement}
      \begin{teiElement}[name={abstract},rend={4e-couv}]

        \teiP[xmlid={p-006}]{}
      
\end{teiElement}
      \begin{teiElement}[name={abstract},rend={resume}]

        \teiP[xmlid={p-007}]{}
      
\end{teiElement}
      \begin{teiElement}[name={abstract},rend={abstract},xmllang={\#\#}]

        \teiP[xmlid={p-008}]{}
      
\end{teiElement}
    
\end{teiElement}
    \begin{teiElement}[name={revisionDesc}]

      \begin{teiElement}[name={change},when={2025},who={2025}]

\end{teiElement}
    
\end{teiElement}
  
\end{teiElement}
  \begin{teiElement}[name={text},type={book},xmlid={book}]

    \begin{teiElement}[name={front}]

\end{teiElement}
    \begin{teiElement}[name={group},type={book},xmlid={book-001}]

      \begin{teiElement}[name={group},type={introduction},data-page-title={Introduction},data-page-authors={Floriane Daguisé||Florence Fix@@Université de Rouen Normandie, CÉRÉDI, Rouen, France},data-include-href={../XML/chap\_1\_Introduction\_Locus\_politicus.xml},xmlid={introduction-001}]

        \begin{teiElement}[name={front}]

          \begin{teiDiv}[type={titlePage},xmlid={titlepage-001}]

            \teiP[rend={title-main},xmlid={p-009}]{Introduction}
            \teiP[rend={author-aut},xmlid={p-010}]{Floriane \teiHi[rend={small-caps}]{Daguisé}}
            \teiP[rend={author-aut},xmlid={p-011}]{Florence \teiHi[rend={small-caps}]{Fix}}
            \teiP[rend={authority\_affiliation},xmlid={p-012}]{Université de Rouen Normandie, CÉRÉDI, Rouen, France}
          
\end{teiDiv}
        
\end{teiElement}
        \begin{teiElement}[name={body},xmlid={body-001}]

          \teiP[xmlid={p1}]{« \lbrack{}Q\rbrack{}uiconque ne sait pas se taire est indigne de gouverner\teiNote[xmlid={ftn1},n={1},place={foot}]{\teiP[xmlid={p-013}]{Fénelon, \teiHi[rend={italic}]{Les Aventures de Télémaque}, Jeanne-Lydie Goré (éd.), Paris, Classiques Garnier, 2009, p. 158.}}. » Telles sont les paroles d’Ulysse à Télémaque avant de partir vers Troie. Le troisième livre du \teiHi[rend={italic}]{Télémaque} de Fénelon insiste sur l’importance politique de la maîtrise du secret, laquelle est qualité d’un « homme raisonnable et sûr » et condition de la « confiance\teiNote[xmlid={ftn33},n={2},place={foot}]{\teiP[xmlid={p-014}]{\teiHi[rend={italic}]{Ibid}., p. 158 et 159. Télémaque rapporte ici à Calypso son échange avec Narbal, commandant de la flotte tyrienne. Sur ce thème dans cette œuvre et dans le contexte politique, voir Henk Hillenaar, \teiHi[rend={italic}]{Le Secret de Télémaque}, Paris, PUF, 1994, en particulier le chapitre du même nom, p. 30-37.}} » de ses interlocuteurs, puissants ou sujets. Cette conception s’inscrit dans la tradition machiavélienne d’un prince « \teiHi[rend={italic}]{gran simulatore e dissimulatore}\teiNote[xmlid={ftn34},n={3},place={foot}]{\teiP[xmlid={p-015}]{Machiavel, \teiHi[rend={italic}]{Le Prince}, Christian Bec (éd. et trad.), \teiHi[rend={italic}]{Œuvres complètes}, Paris, Garnier, 1987, t. I, chap. \teiHi[rend={small-caps}]{xviii}, p. 378. Voir Michel Senellart, « Simuler et dissimuler : l’art machiavélien d’être secret à la Renaissance », François Laroque (dir.), \teiHi[rend={italic}]{Histoire et secret à la Renaissance}, Paris, Presses Sorbonne Nouvelle, 1997, p. 99-106.}} » et d’un pouvoir politique moderne ne pouvant être crédible, respecté, qu’à la condition d’une pratique du secret\teiNote[xmlid={ftn35},n={4},place={foot}]{\teiP[xmlid={p-016}]{« Dès lors la culture du secret s’apparente en milieu politique à une culture de la dissimulation qui tend à associer secret et souveraineté », dans Laure Depretto, Christian Renoux, Christophe Speroni \teiHi[rend={italic}]{et al.} (dir.), \teiHi[rend={italic}]{Cultures du secret à l’époque moderne. Raisons, espaces, paradoxes, fabriques}, Paris, Classiques Garnier, « Polen », 2023, p. 12.}}. Cette habileté implique de savoir connaître et conserver les secrets de tous, mais aussi de savoir les confier à quelques-uns quand c’est nécessaire : à la garantie du secret, se joint, sinon lui préexiste, l’accès au secret. La maîtrise du secret est double, tendue entre la possession et la préservation. Afin de manifester son autorité, d’être \teiHi[rend={italic}]{visiblement} souverain, le monarque doit en effet connaître tous les secrets d’autrui et également garder secrètes ses intentions. Louis \teiHi[rend={small-caps}]{XIV} présente ainsi au Dauphin le rôle du Roi :}
          \begin{teiCit}[xmlid={cit01}]

            \begin{teiQuote}
Tout ce qui est le plus nécessaire à ce travail est en même temps agréable ; car, c’est en un mot, mon fils, avoir les yeux ouverts sur toute la terre ; apprendre à toute heure les nouvelles de toutes les provinces et de toutes les nations, le secret de toutes les cours, l’humeur et le faible de tous les princes et de tous les ministres étrangers ; être informé d’un nombre infini de choses qu’on croit que nous ignorons ; pénétrer parmi nos sujets ce qu’ils nous cachent avec le plus de soin ; découvrir les vues les plus éloignées de nos propres courtisans, leurs intérêts les plus obscurs qui viennent à nous par des intérêts contraires. Et je ne sais enfin quel autre plaisir nous ne quitterions point pour celui-là, si la seule curiosité nous le donnait\teiNote[xmlid={ftn36},n={5},place={foot}]{\teiP[xmlid={p-017}]{Louis XIV, \teiHi[rend={italic}]{Mémoires pour l’instruction du Dauphin}, Pierre Goubert (éd.), Paris, Imprimerie nationale, 1992, p. 52. Sur la maîtrise de l’information par Louis XIV comme source de pouvoir et de plaisir, voir l’article de Lucien Bély : « Louis XIV et le plaisir de l’information », \teiHi[rend={italic}]{Dix-septième siècle}, n\teiHi[rend={sup}]{o} 269, 2015, p. 671-684.}}.
\end{teiQuote}
          
\end{teiCit}
          \teiP[xmlid={p2}]{Cette représentation doit beaucoup à un exercice du pouvoir d’Ancien Régime, réputé tout savoir et prévoir, à l’image de Dieu, qui « voit ce qui se passe dans le secret\teiNote[xmlid={ftn37},n={6},place={foot}]{\teiP[xmlid={p-018}]{\teiHi[rend={italic}]{Évangile selon saint Matthieu} (VI, 6), selon la traduction de Lemaître de Sacy, Paris, Robert Laffont, « Bouquins », 2008, p. 1273. Cette présence divine au cœur du secret contredit l’ostentation de la foi et motive le repli, le retrait : « Mais vous, lorsque vous voudrez prier, entrez dans votre chambre, et la porte en étant fermée, priez votre Père dans le secret ; et votre Père, qui voit ce qui se passe dans le secret, vous en rendra récompense. », \teiHi[rend={italic}]{idem}. Sur le caractère divin auquel aucun secret ne résiste, voir le propos de Raphaël Picon, « Dieu et le secret », en particulier la première partie « Dieu, ou l’œil qui voit tout », dans Jacques Guin (dir.),\teiHi[rend={italic}]{ Le Secret}, Paris, Van Dieren, 2007, p. 9-12.}} ». Il ne doit pas y avoir de secret \teiHi[rend={italic}]{pour} le politique et partant aucun lieu qui ne lui soit inconnu. La philosophie politique de Machiavel à Hobbes présente les rapports entre monarques et sujets à l’image d’un corps organique : les membres renvoient au peuple ; la tête, les yeux, le cerveau, au roi. Qu’il soit visible ou non de ses sujets, le roi voit tout et sait tout d’eux. L’hégémonie politique repose sur une hégémonie cognitive, qui identifie ses cibles selon un rétrécissement progressif vers le cœur du pouvoir et les manœuvres courtisanes. Il est « raisonnable » que le souverain sache tout, et pourtant cette surpuissance de l’esprit et de la mémoire est magique, opaque, irrationnelle. Sans doute tient-elle à une part de mensonge ou du moins de faire-accroire : le pouvoir sait tout des secrets d’autrui, c’est l’une des forces de son autorité, fût-ce seulement en laissant croire qu’il sait tout. « Le secret place la personne dans une situation d’exception, il agit comme un charme dont la détermination est purement sociale, indépendant dans son principe du contenu qu’il protège\teiNote[xmlid={ftn38},n={7},place={foot}]{\teiP[xmlid={p-019}]{Georg Simmel, \teiHi[rend={italic}]{Secret et sociétés secrètes} \lbrack{}1908\rbrack{}, Sibylle Muller (trad.), Paris, Circé, « Poche », 1996, p. 43.}} » ; inversement, dans un pouvoir monarchique autoritaire, comme dans toute dictature, les sujets du roi ne savent pas tout de l’exercice politique, tant s’en faut.}
          \teiP[xmlid={p3}]{Ce volume propose d’aborder ce déséquilibre des savoirs sous l’angle des espaces qu’il implique, plus précisément des représentations d’un \teiHi[rend={italic}]{locus politicus }comme\teiHi[rend={italic}]{ locus secretus}, caché, discret ou invisible, connu d’un petit nombre et inintelligible au plus grand nombre. Au regard du souverain ou de l’homme de pouvoir nul lieu – physique comme spirituel – n’est secret, à l’exemple du Cardinal de Richelieu :}
          \begin{teiCit}[xmlid={cit02}]

            \begin{teiQuote}
Cet esprit toujours agissant\teiLb{}Pénétrait d’un regard perçant\teiLb{}Jusques au fond de leur pensée\teiNote[xmlid={ftn39},n={8},place={foot}]{\teiP[xmlid={p-020}]{Georges de Scudéry, « Le Portrait du Grand Cardinal », \teiHi[rend={italic}]{Le Cabinet de Monsieur de Scudéry}, Paris, Augustin Courbé, 1646, v. 131-150, p. 133-134.}}.
\end{teiQuote}
          
\end{teiCit}
          \teiP[xmlid={p4}]{Gouverner la \teiHi[rend={italic}]{chose publique} se conçoit, sous l’Ancien Régime, comme un accès superlatif au privé, un accès néanmoins asymétrique. Cette conception connaît des variations, lesquelles relèvent également des partitions public-privé contestées ou réactivées après la Révolution. Les questions relatives au secret et à l’exercice du pouvoir politique ont fait l’objet de nombreux travaux : ce volume s’intéresse aux lieux du secret politique, conformément à l’étymologie du \teiHi[rend={italic}]{secret}, lequel renvoie à une mise à part, une séparation. En regard de la pénétration absolue, il convient ainsi d’interroger la dynamique de visibilité spatialisée du politique. La mise en spectacle du pouvoir, du \teiHi[rend={small-caps}]{xvii}\teiHi[rend={sup}]{e} au \teiHi[rend={small-caps}]{xix}\teiHi[rend={sup}]{e} siècle, est inhérente à ses pratiques, codes et rituels ; elle s’accompagne tout autant de manœuvres clandestines, parallèles, qui ne seront pas ici celles de l’opposition mais bien celles situées au cœur de l’exercice du pouvoir. La manipulation, l’usage par le pouvoir même des mystères de l’État, pour en imposer, pour faire reconnaître sa puissance sont les enjeux que les contributeurs de ce volume interrogent. Gérard Leclerc en effet distingue, dans l’exposition du pouvoir en majesté, les « spectacles à usage interne, où le pouvoir se présente et se représente lui-même devant un public restreint\teiNote[xmlid={ftn40},n={9},place={foot}]{\teiP[xmlid={p-021}]{Gérard Leclerc, \teiHi[rend={italic}]{Le Regard et le pouvoir}, Paris, PUF, « Sociologie d’aujourd’hui », 2006, p. 115.}} » (celui de la cour notamment), et les « spectacles à usage externe, destinés aux foules, au peuple, où le pouvoir est présent métaphoriquement\teiNote[xmlid={ftn41},n={10},place={foot}]{\teiP[xmlid={p-022}]{\teiHi[rend={italic}]{Ibid}., p. 116.}} », comme les supplices, « violence ostentatoire du pouvoir\teiNote[xmlid={ftn42},n={11},place={foot}]{\teiP[xmlid={p-023}]{\teiHi[rend={italic}]{Ibid}., p. 119.}} », sans oublier, bien sûr, les espaces de représentation que constituent les palais, théâtres de cette « mise en scène du pouvoir ». Mais ces manifestations extrêmement visibles ne sont pas tout : le roi dispose de cabinets secrets, d’espaces fermés, ainsi que le rappelle le dictionnaire de Furetière, qui dans sa définition de « privé* », en tant que « particulier, secret », donne l’exemple du conseil privé du Roi. Le privé considéré serait le lieu d’un secret \teiHi[rend={italic}]{particulier}, spécifiquement politique, voulu, construit. Il en résulte l’exploitation possible d’espaces privés comme des espaces politiques particuliers, en ce sens qu’ils seraient destinés à une pratique du politique par la marge, le retrait, la cachette. Là se tiendraient les arcanes d’un pouvoir dont les vues – étendues – échappent au commun : « Le Roi a tenu un Conseil \teiHi[rend={italic}]{secret} pour quelque grand dessein » indique le \teiHi[rend={italic}]{Dictionnaire de Trévoux} dans l’article « secret » (édition lorraine de 1738-1742). Le lieu du secret politique n’est pas toujours un lieu spécifiquement prévu à cet effet (comme un cabinet de travail ou la réunion d’un petit comité), il le devient toutefois par la présence du pouvoir, qui en retour voit son autorité affermie par le secret. Même lorsque la sur-représentation des dirigeants survient, indique Georges Balandier, dans une surenchère de médiatisation et de spectacularisation, demeure une conviction qu’il existe encore des espaces dérobés aux regards, à l’exposition, au commentaire, et le soupçon que ce que l’on a vu n’était pas « tout » :}
          \begin{teiCit}[xmlid={cit03}]

            \begin{teiQuote}
L’âge des \teiHi[rend={italic}]{media} impose le pouvoir permanent des images, et donc la contrainte de fonder sur elles leurs pouvoirs ; mais le spectaculaire continu banalise, il efface la distance et la séparation sans lesquelles le politique n’a plus d’espace propre, il remplace le secret (l’une des forces du gouvernant) par le bruit\teiNote[xmlid={ftn43},n={12},place={foot}]{\teiP[xmlid={p-024}]{Georges Balandier, \teiHi[rend={italic}]{Le Détour. Pouvoir et modernité}, Paris, Fayard, 1985, p. 11.}}.
\end{teiQuote}
          
\end{teiCit}
          \teiP[xmlid={p5}]{Le bruit, l’apparence, l’apparat du pouvoir en représentation (manifestations publiques, rencontres officielles, alliances et traités, mariages et enterrements, supplices et condamnations à mort en place publique, divertissements officiels, etc.) ont pour corollaire le silence des décisions importantes (secrets d’État, rencontres officieuses, tractations cachées, espionnage), double espace qui emprunte au théâtre ses métaphores de la scène et de la coulisse\teiNote[xmlid={ftn44},n={13},place={foot}]{\teiP[xmlid={p-025}]{Le parallèle est éclairant : Jeremy Bentham, fondateur du panoptique à la fin du \teiHi[rend={small-caps}]{xviii}\teiHi[rend={sup}]{e} siècle dont le modèle de surveillance détermine la réflexion foucaldienne, se serait lui-même inspiré de salles de spectacles, en particulier la salle du théâtre de Besançon, que l’architecte Ledoux fait figurer allégoriquement dans la pupille d’un œil. Michel Foucault n’envisageait pas le lien entre spectacle et panoptique sous cette forme, il considérait plutôt une rupture entre ces deux incarnations du regard (« Notre société n’est pas celle du spectacle, mais de la surveillance \lbrack{}…\rbrack{}. Nous ne sommes ni sur les gradins ni sur la scène, mais dans la machine panoptique », Michel Foucault,\teiHi[rend={italic}]{ Surveiller et punir}, Paris, Gallimard, 1975, p. 252-253). Sur ces liens entre regard spectatorial et regard politique, voir Anthony Vidler, \teiHi[rend={italic}]{L’Espace des Lumières. Architecture et philosophie de Ledoux à Fourier}, Paris, Picard, « Villes et sociétés », 1995, p. 172-175 en particulier.}}, ouvrant aux interprétations sur le dédoublement, quand ce n’est pas la duplicité, l’hypocrisie de l’exercice. La puissance s’exhibe par les codes du spectaculaire, lequel ne suffit pas : chacun sait que les espaces de commémoration, de réception ou de fêtes que le souverain offre aux regards ne sont qu’une partie de son activité, ce qu’il consent à offrir. De ces premiers éléments s’impose la tension entre une visibilité démonstrative et une invisibilité maîtrisée, proclamée sur fond d’exhibition :}
          \begin{teiCit}[xmlid={cit04}]

            \begin{teiQuote}
Le pouvoir semble avoir besoin de visibilité : il souhaite être vu, contemplé, observé, par des regards respectueux, effrayés, timides, éblouis… Rares sont les pouvoirs qui aspirent à l’invisibilité totale, qui cherchent à demeurer dans l’ombre. Ce n’est pas à dire qu’ils ne pratiquent pas le secret, qu’ils ne tendent pas parfois à rendre invisibles certaines de leurs composantes et caractéristiques. Mais l’autre dimension du pouvoir, son autre face de Janus, c’est la recherche ostentatoire de la visibilité, créatrice d’effroi sacré\teiNote[xmlid={ftn45},n={14},place={foot}]{\teiP[xmlid={p-026}]{Gérard Leclerc, \teiHi[rend={italic}]{Le Regard et le pouvoir}, \teiHi[rend={italic}]{op. cit}., p. 4.}}.
\end{teiQuote}
          
\end{teiCit}
          \teiP[xmlid={p6}]{La surexposition du souverain, tout comme celle du prêtre en chaire, du maître à penser ou de la star, supposent, en creux, un autre mode d’existence, celui du lieu écarté, secret : tout laisse croire à une autre vie.}
          \teiP[xmlid={p7}]{L’architecture témoigne, pendant la période considérée, d’une reconfiguration progressive de la partition entre domaine public et domaine privé enracinée dans des dynamiques spatiales. Le dédoublement des appartements du Roi à Versailles tend ainsi à distinguer deux types de publicité du pouvoir : « Les grands appartements sont ouverts au grand nombre, les petits sont réservés à quelques familiers\teiNote[xmlid={ftn46},n={15},place={foot}]{\teiP[xmlid={p-027}]{Michel Delon\teiHi[rend={small-caps}]{,} \teiHi[rend={italic}]{Le Savoir-vivre libertin}, Paris, Hachette, 2000, p. 123.}}. » Plus spécifiquement encore, les intérieurs se morcellent en espaces au retrait gradué, dessinent des refuges plus ou moins isolés\teiNote[xmlid={ftn47},n={16},place={foot}]{\teiP[xmlid={p-028}]{« Qui veut préserver son intimité doit nécessairement cloisonner jusqu’à l’espace intérieur de son logement » explique Youri Carbonnier, dans \teiHi[rend={italic}]{Maisons parisiennes des Lumières}, Paris, PUPS, 2006, p. 319.}}. C’est en particulier le cabinet, à l’indication définitionnelle révélatrice\teiNote[xmlid={ftn48},n={17},place={foot}]{\teiP[xmlid={p-029}]{Voir l’article « Cabinet » de l’\teiHi[rend={italic}]{Encyclopédie} rédigé par l’architecte Jacques-François Blondel, \teiHi[rend={italic}]{Encyclopédie ou Dictionnaire raisonné des sciences, des arts et des métiers, par une société de gens de lettres}, Paris, chez Briasson, David, Le Breton et Durand, 1751-1772, t. II, p. 488-489.}}, qui connaît un investissement architectural sensible aux \teiHi[rend={small-caps}]{xvii}\teiHi[rend={sup}]{e} et \teiHi[rend={small-caps}]{xviii}\teiHi[rend={sup}]{e} siècle, à l’exemple du cabinet du Dauphin élaboré en 1699. Le cabinet de toilette en est l’une des spécialisations les plus éclatantes, en ce qu’il exhibe autant que préserve l’espace d’une intimité qui est lieu de construction d’un être social. Ces pièces-coquilles favorisent également la promotion d’une vie intérieure qui s’épanouit dans le privé ; l’analyse habermassienne implique cette conjonction entre évolution architecturale et histoire des mentalités, spécifiquement celle de la subjectivité\teiNote[xmlid={ftn49},n={18},place={foot}]{\teiP[xmlid={p-030}]{Voir Jürgen Habermas, \teiHi[rend={italic}]{L’Espace public. Archéologie de la publicité comme dimension constitutive de la société bourgeoise} \lbrack{}1962\rbrack{}, Marc Buhot de Launay (trad.), Paris, Payot, 1978, en particulier p. 55-59.}}. Les travaux d’Antoine Lilti approfondissent et réorientent de telles perspectives. Comme il le rappelle, la valorisation d’un corps privé est un puissant matériau littéraire qui se transforme au \teiHi[rend={small-caps}]{xviii}\teiHi[rend={sup}]{e} siècle : « Cette nouvelle écriture biographique est indissociable de l’émergence du roman moderne, qui repose sur le postulat que toute vie, même la plus humble socialement, est digne d’être racontée\teiNote[xmlid={ftn50},n={19},place={foot}]{\teiP[xmlid={p-031}]{Antoine Lilti, \teiHi[rend={italic}]{L’invention de la célébrité. 1750-1850 }\lbrack{}2014\rbrack{}, Paris, Fayard, « Pluriel », 2022, p. 107.}}. » Toute existence étant digne d’être racontée, l’accès à la subjectivité, à l’expérience du corps intime est nécessaire à la singularisation de son récit. Dès lors, les corps des puissants sont interrogés selon le même régime de visibilité. Que ce regard soit porté par l’admiration ou par la contestation, le corps du gouvernant, lieu comme un autre, ses lieux de vie fastueux ou non, ses cabinets de travail sont soumis à l’exigence post-révolutionnaire de la transparence. Le corps du roi comme « lieu de représentation\teiNote[xmlid={ftn51},n={20},place={foot}]{\teiP[xmlid={p-032}]{Georges Balandier, \teiHi[rend={italic}]{Le Pouvoir sur scènes}, Paris, Balland, 1992, p. 40.}} » et le cérémonial très codifié de ses apparitions officielles tirent bien leur pouvoir de fascination de leurs revers et envers supposés, de lieux redoutés, rêvés, imaginés dans lesquels le pouvoir politique agit lorsqu’il sort de l’espace public, matrices de fictions nombreuses dont la portée subversive s’exhibe par des destructions ou entrées en force. Lors de l’intrusion révolutionnaire dans le château de Versailles en 1789, lors de l’incendie des Tuileries par la Commune de Paris en 1870, la destruction spectaculaire s’accompagne d’articles et de pamphlets soulignant qu’on n’en a pourtant pas fini avec le pouvoir : tant que le corps du roi demeure, il peut se déplacer, organiser un autre lieu de pouvoir, se dérober, se reconstruire. Réponse et contre-stratégie à l’exercice dissimulé du pouvoir, la destruction de lieux de spectacle ne va pas sans frustration, tant il est clair que d’autres lieux subsistent. Les cabinets secrets, voire cabinets noirs, les lieux qui n’ont ni nom ni existence officielle sont ceux dans lesquels s’exercerait vraiment le pouvoir, ceux qui rassemblent des initiés : les autres lieux, exhibés au plus grand nombre, n’en seraient que la preuve, démontrant précisément que le pouvoir maîtrise parfaitement ce qu’il montre et ce qu’il ne montre pas, à l’instar du corps-même du roi, qui ne présente aux regards qu’un portrait en majesté. La fascination pour la coulisse, la chambre, le cabinet de toilette ou de conseil, pour les espaces restreints à l’accès limité encouragent le récit de ceux qui les ont vus et s’offrent pour des initiés. Les contre-pouvoirs (sociétés secrètes, conspirations, oppositions) de leur côté ont aussi leur \teiHi[rend={italic}]{locus secretus }; tous en parlent : « Les curieux imaginent à partir du secret, ils créent du mythe, ils mythifient par peur sans doute d’être mystifiés\teiNote[xmlid={ftn52},n={21},place={foot}]{\teiP[xmlid={p-033}]{Pierre Brunel, \teiHi[rend={italic}]{L’Imaginaire du secret}, Grenoble, Ellug, 1998, p. 9.}}. » Mais réside, toujours, une ultime mystification possible : l’évidence – plutôt que l’existence – du secret dissimulé\teiNote[xmlid={ftn53},n={22},place={foot}]{\teiP[xmlid={p-034}]{« Écrire sur le secret à l’Âge classique \lbrack{}…\rbrack{} c’est aussi maintenir la clôture entre l’espace de ceux qui savent et de ceux qui ne savent pas, quitte à ce que savoir consiste surtout à savoir qu’il n’y a rien à savoir, que le secret d’État est un leurre, une mystification » rappelle l’introduction du collectif \teiHi[rend={italic}]{Cultures du secret à l’époque moderne. Raisons, espaces, paradoxes, fabriques}, \teiHi[rend={italic}]{op. cit.}, p. 12.}}. Le pouvoir politique a partie liée avec le spectacle, mais encore faut-il reconnaître que ce spectacle « attire en même temps qu’il dissimule\teiNote[xmlid={ftn54},n={23},place={foot}]{\teiP[xmlid={p-035}]{Jean-Marie Apostolidès, \teiHi[rend={italic}]{Le Roi-machine. Spectacle et politique au temps de Louis XIV}, Paris, Éditions de Minuit, 1981, p. 156. Citation complète : « Le spectacle au \teiHi[rend={small-caps}]{xvii}\teiHi[rend={sup}]{e} siècle possède une fonction d’éblouissement : il attire en même temps qu’il dissimule. »}} ». Si « les représentations solennelles mettent en évidence la capacité surnaturelle du pouvoir » et « suscitent ainsi un sentiment d’appartenance », il convient symétriquement d’envisager la place et les effets de cette part dissimulée, et donc imaginée, du pouvoir politique. La face cachée du spectacle, aussi sinon davantage attractive, objet de toutes les projections, croyances et craintes, relève ainsi naturellement d’une fiction ; c’est en ce cas le domaine fictionnel qui se révèlera peut-être le plus à même d’en traduire à la fois le fonctionnement comme la fascination.}
          \teiP[xmlid={p8}]{Pénétrer le lieu, le détruire ou le décrire, l’épier et le documenter sont dès lors des ressources dramatiques et des actions symboliques placées au cœur des fictions du secret, que le pouvoir en ait encore l’usage et soit saisi sur le fait (espionner une conspiration, être admis à une audience privée et le relater en détails dans un récit de voyage) ou qu’il en soit absent (parvenir à s’immiscer dans le lieu avant que ne commence une cérémonie ou réunion secrète, craindre d’être pris sur le fait font les inlassables ressorts des mélodrames et drames romantiques, par exemple). Parce que la curiosité envers ces lieux du secret et envers les intrigues de ceux qui détiennent le pouvoir en fragilise pour partie l’exercice, elle l’incite aussi à modifier ses pratiques, à imaginer des dispositifs architecturaux – porte dissimulée, escalier dérobé, trappes, fausses cloisons… –, dans la continuité des machineries à destination d’intrigues amoureuses, telle la cheminée tournante du duc de Richelieu réinvestie par la fiction romanesque. Dans ces lieux peuvent exister des objets gigognes qui insistent encore sur le secret : cachettes, meubles à double fond, paravents, etc. Conjointement à ces aménagements, le pouvoir politique est conduit à changer souvent de lieu ou à le rendre extrêmement gardé, ainsi des variations fictionnelles de la cité interdite, stimulées par les récits des voyageurs bredouilles. Le \teiHi[rend={italic}]{locus secretus} est une périphérie du centre visible du pouvoir, un cabinet caché derrière une tenture dans une salle d’apparat, un bosquet dans le parc d’un château, le souterrain ou le caveau étroit et mal éclairé d’un palais grandiose. Sa petitesse et son humilité ne sont que voiles de sa puissance réelle : c’est là que s’exercerait le « vrai » pouvoir, que se jouraient les intrigues déterminantes. En cela, certains lieux prennent des contours inquiétants ; le gouffre, le caveau ou le tombeau, qui font disparaître des adversaires, sont investis ou du fantasme de la dissimulation toute-puissante ou de la tentation des conspirations et des contre-pouvoirs, dans une forme d’isolement tactique et d’hermétisme symbolique. Y être invité (par mot de passe, par lettre cryptée, par un signe, un mot ou un objet à montrer au garde), c’est jouir d’un pouvoir d’autant plus exaltant qu’il ne peut être dit publiquement.}
          \teiP[xmlid={p9}]{La première étape de ce volume envisage naturellement une cartographie du secret. Les différentes contributions mettent en évidence des espaces secrets et du secret parasités par des élans de visibilité et d’opacité. Les espaces secrets et du secret s’animent des tensions ou contradictions de ses habitants, comme en témoigne la transparence en clair-obscur du \teiHi[rend={italic}]{locus secretus} rousseauiste (Henri Portal). Ces qualités ambivalentes s’illustrent plus spécifiquement encore dans les lieux à contre-emploi, tel le cabinet de toilette zolien (Émilie Bauduin), dont la dimension cosmétique est aussi simulacre stratégique. Le secret investit parallèlement des espaces disponibles, au-delà, en-deçà ou en-dessous de la surface visible du pouvoir, à l’image du caveau balzacien que mettent à profit les différents acteurs de soubresauts politiques (Kathia Huynh). Ces espaces traduisent ou trahissent des stratégies politiques en marge du pouvoir ou contre le pouvoir, comme autant d’hétérotopies conjuratrices ou conjuratoires dont la fiction du \teiHi[rend={small-caps}]{xix}\teiHi[rend={sup}]{e} siècle post-révolutionnaire rend compte avec force (Guillaume Cousin). Le caveau ou le cabinet de toilette redeviennent décors anodins lorsqu’en sortent ceux qui s’y sont réunis en secret : aussitôt réversibles, ces lieux sont des dispositifs qui ne fonctionnent qu’avec leurs acteurs.}
          \teiP[xmlid={p10}]{Les publications à prétention documentaire et informationnelle relèvent aussi de la fiction : elles se donnent pour objectif de combler les silences et les secrets du récit historique par d’autres récits. Le titre\teiHi[rend={italic}]{ La Clef du cabinet des Princes de l’Europe ou Recueil historique sur les matières du temps}\teiNote[xmlid={ftn55},n={24},place={foot}]{\teiP[xmlid={p-036}]{ La notice du \teiHi[rend={italic}]{Dictionnaire des journaux} présente ainsi la nature de ce périodique : « Journal d’information générale, \teiHi[rend={italic}]{La Clef }n’a pas de plan tout à fait défini, même si en général, elle donne les nouvelles dans l’ordre des grandes nations (France, Espagne, Italie, Allemagne et Hongrie, Pologne et États du Nord, Grande-Bretagne, Hollande et Pays-Bas, etc.) », Jean-Pierre Kunnert, Jean Sgard (dir.), édition électronique du \teiHi[rend={italic}]{Dictionnaire des Journaux (1600-1789)}, notice 0214. Après 1773, il devient le \teiHi[rend={italic}]{Journal historique et littéraire}.}}, publié de 1704 à 1773, suggère la découverte des tractations diplomatiques offertes à la curiosité et à l’information des lecteurs. \teiHi[rend={italic}]{Le Cabinet secret de l’Histoire entr’ouvert par un médecin}, d’Augustin Cabanès, connaît un grand succès au \teiHi[rend={small-caps}]{xix}\teiHi[rend={sup}]{e} siècle. Ponctuant ses chapitres de phrases comme « La plupart des détails que l’on va lire sont inédits ou peu connus ; ils sont en tout cas d’une indiscutable authenticité\teiNote[xmlid={ftn56},n={25},place={foot}]{\teiP[xmlid={p-037}]{Augustin Cabanès, \teiHi[rend={italic}]{Le Cabinet secret de l’histoire entr’ouvert par un médecin}, Paris, A. Charles libraire, 1895, p. 202.}} », le médecin prétend commenter les grossesses de Marie-Antoinette et de Marie-Louise, les maux de Talleyrand, de Gambetta et de Marat. Il compile plutôt des textes existants et leur donne la caution de la médecine. Certes, reconnaît-il, « il peut sembler téméraire, à un siècle de distance, et sans avoir le document sous les yeux, c’est-à-dire le sujet lui-même, de discuter la nature de son mal\teiNote[xmlid={ftn57},n={26},place={foot}]{\teiP[xmlid={p-038}]{\teiHi[rend={italic}]{Ibid}., p. 133.}} », et pourtant : la lecture et l’observation lui permettent à bon droit de répondre à la question « Quelle était la maladie de Marat ? » Regarder, observer, prétendre avoir vu ou être renseigné ; Montesquieu soulignait déjà l’attrait pour les secrets politiques, symétriquement à l’hégémonie observatrice du pouvoir :}
          \begin{teiCit}[xmlid={cit05}]

            \begin{teiQuote}
À mesure que les princes ont trouvé des arts pour devenir maîtres de nos secrets, par l’art d’ouvrir les lettres sans qu’on s’en aperçoive, nous avons trouvé l’art de publier les leurs par des façons plus secrètes d’imprimer\teiNote[xmlid={ftn58},n={27},place={foot}]{\teiP[xmlid={p-039}]{Montesquieu,\teiHi[rend={italic}]{ Pensées}, Louis Desgraves (éd.), Paris, Robert Laffont, 1991, fragment 791, p. 349. Carole Dornier en examine les sources et les enjeux dans l’article « Secret, espionnage et politique des Modernes chez Montesquieu », Françoise Gevrey, Alexis Lévrier et Bernard Teyssandier (dir.), \teiHi[rend={italic}]{Éthique, poétique et esthétique du secret de l’Ancien Régime à l’époque contemporaine}, Louvain, Peeters, 2015, p. 55-64.}}.
\end{teiQuote}
          
\end{teiCit}
          \teiP[xmlid={p11}]{C’est pourquoi une deuxième étape de ce volume est consacrée à l’énonciation, à la formulation des lieux de secrets : comment se révèlent-ils pour ce qu’ils sont ? Par la lettre, par l’argent, par un objet qui matérialise ce qui n’est autrement que supposition ? De même qu’un secret tient son pouvoir de la crainte des uns et de l’espoir des autres qu’il soit mis au jour, le lieu secret est en réalité marqué par un certain nombre de signes qui le désignent à la curiosité et à l’anxiété. « Secrets de polichinelle » (Marie-Agathe Tilliette), l’argent, la jalousie amoureuse, le désir de vengeance sont bavards. Ils se disent par des lettres secrètes (Étienne Poirier) qui ne sont pas destinées à le rester toujours : le lieu du secret ne saurait se comprendre sans un temps du secret. Aussitôt dit, il n’est plus secret, c’est une évidence, mais il a rempli son office : les conjurés, les initiés, les informés ont su à temps se concerter, agir, obtenir ce qu’ils convoitaient. Révélé, il en remplit un autre : celui du commentaire, souvent à charge, de la contestation, voire du règlement de comptes de la part d’anciens proches du pouvoir écartés des lieux de décision (Cécile Leduc).}
          \teiP[xmlid={p12}]{Comme le rappelle Louis Marin, « l’effet » du secret est bien plus élevé que sa valeur intrinsèque ; il le met en lien avec le simulacre de puissance associé au secret, dont le paradigme fondateur et structurant serait Œdipe\teiNote[xmlid={ftn59},n={28},place={foot}]{\teiP[xmlid={p-040}]{Voir Louis Marin, « Logiques du secret », \teiHi[rend={italic}]{Traverses}, n\teiHi[rend={sup}]{o} 30-31, 1984, p. 60-69.}}. Convaincu que le souverain lui cache quelque chose, le sujet en tire fascination, adhésion ou à l’inverse défiance et motif à contestation. Les lieux du secret sont poreux et vulnérables, peuvent être (ou plutôt le seront immanquablement) dénoncés, dévastés, tout simplement éventés, regardés, décrits. Ils n’opèrent que durant le temps court où le pouvoir s’y loge ou s’en sert ; toutefois leur puissance sur les imaginaires, les traces qu’ils laissent sont infiniment plus durables – ainsi de la fascination pour les tombes des pharaons, pour les cabinets secrets de palais devenus musées, ou pour les mémoires d’intimes ou prétendus tels du pouvoir. Le propos de Louis Marin est donc plus large et, en un sens, suffisamment abstrait pour envisager le secret sous un angle plus vaste : il décrit le « tourniquet métaphysique du secret » avant d’en mesurer la force, soit la part de fictionnalité. Les voyageurs et les observateurs dans les cours royales ou impériales relatent leur étonnement, leur admiration, ils se targuent d’avoir « vu » le pouvoir à l’œuvre, c’est-à-dire ses effets, notamment lors d’exécutions publiques, de bannissements, de sanctions. C’est lors de ces moments d’exposition du pouvoir que celui-ci révèle ce qu’il a de longue date préparé dans l’ombre : la scène de supplice ou de disgrâce exhibe des lieux qui en cachent beaucoup d’autres.}
          \teiP[xmlid={p13}]{C’est pourquoi tout lieu peut être lieu d’un secret à connaître et à utiliser, ainsi que le savent bien les personnages d’espions sollicités par la fiction : on espionne dans les rues, sur les marchés, dans les relais de poste, selon le principe de l’omniscience nécessaire à l’exercice de l’autorité. On espionne la population civile, dans ses activités les plus anodines ; l’espion peut être homme ou femme, jeune ou vieux, étranger ou natif, mais toujours prompt au jeu sur les identités. Une troisième et dernière étape de ce volume s’intéresse ainsi à celles et ceux qui « incarnent » le secret, lui donnent chair : ainsi des tyrans dont les lieux de vie et de décision sont invisibles et bien gardés. Ainsi de ceux qui tentent tout de même de les voir. Les personnages d’espions sont d’autant plus présents dans la fiction que certains auteurs complètent leurs revenus avec cette pratique ou fréquentent des cercles d’espionnage : Courtilz de Sandras aurait été espion, François de Chevrier en est soupçonné, Pierre de Morand également en tant qu’« agent littéraire » du roi de Prusse. Jean-Noël Tardy, à propos de l’âge romantique en France, estime que « le secret caractérise l’homme de lettres et le conspirateur. La séparation du vulgaire qu’implique sa détention rejoint les aspirations des poètes romantiques à se placer au-dessus des foules\teiNote[xmlid={ftn60},n={29},place={foot}]{\teiP[xmlid={p-041}]{Jean-Noël Tardy, \teiHi[rend={italic}]{L’Âge des ombres. Complots, conspirations et sociétés secrètes au }\teiHi[rend={ small-caps italic}]{xix}\teiHi[rend={italic sup}]{e}\teiHi[rend={italic}]{ siècle}, Paris, Les Belles Lettres, 2015, partie « Portrait de l’écrivain en conspirateur (1827-1835) », p. 340-345, ici p. 342.}} », rappelant que Byron, qui se fait initier \teiHi[rend={italic}]{carbonaro} en Italie, ou Hugo en défendant ses œuvres (constitution d’un groupe autour de la bataille d’\teiHi[rend={italic}]{Hernani}), comme ses opinions esthétiques et politiques, ont des pratiques de conspiration et se plaisent à endosser le rôle alors flatteur du conspirateur. Il convient de terminer ce volume sur ces personnages dont le secret est le métier et qui pour l’exercer sont mobiles, qu’ils soient voyageurs et enclins à la comparaison avec leur pays d’origine (Matthias Soubise), peut-être un peu espions, mais toujours aussi fascinés par les tyrans et désireux de faire œuvre avec leur récit (Mathilde Mougin). C’est d’ailleurs rarement le souverain qui incarne le mieux la pleine maîtrise du lieu : l’espace désacralise, sinon dégrade l’image royale au point de la renfermer dans l’espace privé, il révèle en robe de chambre ou en déshabillé des individus vulnérables (Thibaut Julian). À leurs côtés et dans l’ombre, des ministres comme Richelieu (Maxime Cartron) ou Talleyrand, des espions ou espionnes, œuvrent en secret, parce qu’ils connaissent les secrets de l’État et les fragilités de son chef. Ces personnages détiennent un savoir que d’autres n’ont pas, c’est ce qui de fait définit le secret :}
          \begin{teiCit}[xmlid={cit06}]

            \begin{teiQuote}
\lbrack{}…\rbrack{} le secret est un fait social et n’existe qu’en raison de valeurs ou d’intérêts concurrents qui déterminent la ligne de partage entre ceux qui savent (qui sont « dans le secret ») et ceux qui pourraient avoir un intérêt à savoir\teiNote[xmlid={ftn61},n={30},place={foot}]{\teiP[xmlid={p-042}]{« Introduction générale », Laure Depretto, Christian Renoux, Christophe Speroni \teiHi[rend={italic}]{et al.} (dir.), \teiHi[rend={italic}]{Cultures du secret à l’époque moderne…}, \teiHi[rend={italic}]{op. cit}., p. 10.}}.
\end{teiQuote}
          
\end{teiCit}
          \teiP[xmlid={p14}]{Le \teiHi[rend={italic}]{locus secretus} l’est car certains personnages y sont et d’autres pas. Parce qu’il est éminemment craint et désiré, parce qu’on veut ardemment y être, y avoir accès, son aspect importe peu : sa médiocrité, son absence de faste ou d’intérêt apparent assurent même de sa valeur ceux qui y entrent enfin. C’est précisément parce qu’il n’est pas au centre (de la cité, des regards, du palais, etc.) qu’il convainc de son importance. En connaître même l’existence est déjà une marque de proximité avec le pouvoir. Le laisser croire est toute l’activité du courtisan, le côtoyer vraiment toute celle de l’espion. Espace de connivence, le \teiHi[rend={italic}]{locus secretus} est un espace de circulation : on y entre et sort, on y parle et échange des informations et des objets. « C’est pourquoi la possibilité et la tentation de trahir tournent autour du secret, et au risque extérieur d’être découvert vient se mêler le risque intérieur de se découvrir, qui ressemble à la fascination du vide. Le secret met une barrière entre les hommes, mais il éveille en même temps la tentation de la briser par le bavardage ou l’aveu\teiNote[xmlid={ftn62},n={31},place={foot}]{\teiP[xmlid={p-043}]{Georg Simmel,\teiHi[rend={italic}]{ Secret et sociétés secrètes}, \teiHi[rend={italic}]{op. cit.}, p. 44-45.}} » ; sa puissance a pour effet une forme de vulnérabilité et un caractère éphémère : découvert, il a vocation à être remplacé par un autre. Ses objets (lettres, documents, bijoux\teiNote[xmlid={ftn63},n={32},place={foot}]{\teiP[xmlid={p-044}]{Voir Christophe Speroni, « Espaces du secret », \teiHi[rend={italic}]{op. cit.}, p. 119 : « Mais le cabinet c’est aussi le lieu où l’on entrepose des documents confidentiels, journaux intimes et lettres qu’il ne faut pas mettre entre toutes les mains et qui paradoxalement seront souvent les instruments de la révélation des secrets les mieux gardés. Les documents privés sont entreposés dans des malles, dans des cabinets, dans des secrétaires, au fond des tiroirs, parfois fermés à clé, comme si le confinement dans des espaces de plus en plus étroits assurait l’inviolabilité des secrets qu’ils renferment. »}}), volés, sont voués à être niés, accusés d’inauthenticité ou également renouvelés. Les fictions que ce volume analyse ne cessent d’exposer, commenter, échanger et faire circuler les descriptions, les usages et les fonctions de ces lieux dits secrets, qui ne le sont plus dès lors qu’ils sont dits, sans pour autant perdre leur pouvoir de fascination. Probablement le doivent-ils à la part de fiction qu’ils ont dès leur origine et aux effets de publications qui en réactivent le plaisir.}
        
\end{teiElement}
      
\end{teiElement}
      \begin{teiElement}[name={group},type={section1},xmlid={section1-001}]

        \teiHead[xmlid={cartographie-situer-le-secret-politique-001}]{Cartographie : situer le secret politique}
        \begin{teiElement}[name={group},type={chapter},data-page-title={Les espaces du secret à Clarens},data-page-subtitle={Julie ou la Nouvelle Héloïse},data-page-authors={Henri Portal@@Sorbonne Université, CELLF 16-18, UMR 8599, Paris, France},data-include-href={../XML/chap\_2\_Les\_espaces\_du\_secret\_Locus\_politicus.xml},xmlid={chapter-001}]

          \begin{teiElement}[name={front}]

            \begin{teiDiv}[type={titlePage},xmlid={titlepage-002}]

              \teiP[rend={title-main},xmlid={p-045}]{Les espaces du secret à Clarens}
              \teiP[rend={title-sub},xmlid={p-046}]{
                \teiHi[rend={italic}]{Julie ou la Nouvelle Héloïse}
              }
              \teiP[rend={author-aut},xmlid={p-047}]{Henri \teiHi[rend={small-caps}]{Portal}}
              \teiP[rend={authority\_affiliation},xmlid={p-048}]{Sorbonne Université, CELLF 16-18, UMR 8599, Paris, France}
            
\end{teiDiv}
          
\end{teiElement}
          \begin{teiElement}[name={body},xmlid={body-002}]

            \teiP[xmlid={p1-2}]{Au moment de son séjour parisien, au Livre II de \teiHi[rend={italic}]{La Nouvelle Héloïse}, Saint-Preux explique à Julie qu’il est « initié à des mystères plus secrets ». Décrivant ces réunions accessibles uniquement à quelques privilégiés, il ajoute qu’il s’agit « de soupés \lbrack{}\teiHi[rend={italic}]{sic}\rbrack{} priés où la porte est fermée à tout survenant\teiNote[xmlid={ftn1-2},n={1},place={foot}]{\teiP[xmlid={p-049}]{Notre édition de référence sera la suivante : Rousseau, \teiHi[rend={italic}]{La Nouvelle Héloïse}, Henri Coulet (éd.), Paris, Gallimard, « Folio », 1993, 2 tomes.}} (II, XVII, p. 306) ». On le sait, « la pratique du secret et du silence \lbrack{}c\rbrack{}aractéris\lbrack{}e\rbrack{} le comportement de ceux qui se rapprochent des cercles du pouvoir\teiNote[xmlid={ftn35-2},n={2},place={foot}]{\teiP[xmlid={p-050}]{Lucien Faggion, « L’éloquence muette ou le langage du corps », \teiHi[rend={italic}]{Arcana imperii. Gouverner par le secret à l’époque moderne. France, Espagne, Italie}, Paris, les Indes savantes, 2018, p. 124.}} ». En transposant le motif essentiel de l’univers politique d’Ancien Régime qu’est le \teiHi[rend={italic}]{locus secretus }dans un contexte mondain, Rousseau, suivant en cela une tradition ancienne, suggère ainsi que les cercles les plus secrets de la capitale ne sont guère différents des sociétés curiales où le pouvoir tend toujours à s’exercer dans le silence des cabinets et des antichambres. L’espace secret prend certes la forme du lieu privé ; mais il ne perd ni sa capacité à structurer les hiérarchies, ni son rôle déterminant dans les luttes d’influence qui se déploient dans la vie sociale.}
            \teiP[xmlid={p2-2}]{Dans ce passage de l’œuvre, on sent nettement, comme ailleurs dans l’œuvre de Rousseau, une certaine hantise de l’opacité – hantise qui fera toutefois l’objet d’une tentative de conjuration dans la seconde partie du roman, associée par Jean Starobinski à un idéal de transparence et de « limpidité\teiNote[xmlid={ftn36-2},n={3},place={foot}]{\teiP[xmlid={p-051}]{Jean Starobinski, \teiHi[rend={italic}]{Jean-Jacques Rousseau. La Transparence et l’Obstacle}, Paris, Gallimard, « Tel », 1971, p. 105.}} ». Après avoir montré la force des apparences trompeuses qui règnent à Paris lors du séjour de Saint-Preux, le romancier présente en effet au lecteur un monde différent à Clarens, le domaine dirigé par Wolmar, le sage époux de Julie. Si toute trace de duplicité semble avoir disparu de cet espace, reconnaissons qu’« il n’est pas simple de vivre à Clarens\teiNote[xmlid={ftn37-2},n={4},place={foot}]{\teiP[xmlid={p-052}]{Michèle Duchet, « Clarens, “le lac d’amour où l’on se noie” », \teiHi[rend={italic}]{Littérature}, n\teiHi[rend={sup}]{o} 21, 1976, p. 80.}} », et que Rousseau, comme la critique l’a établi\teiNote[xmlid={ftn38-2},n={5},place={foot}]{\teiP[xmlid={p-053}]{Christophe Martin, \teiHi[rend={italic}]{La Philosophie des amants. Essai sur }Julie ou la Nouvelle Héloïse, Paris, Sorbonne Université Presses, 2021, p. 195-214.}}, nous invite à poser un regard distancié sur l’idéal de transparence qui s’y manifeste. Dès lors, il semblerait que le motif du \teiHi[rend={italic}]{locus secretus }n’ait pas été complètement évacué de la seconde partie de \teiHi[rend={italic}]{La Nouvelle Héloïse}. À y regarder attentivement, on n’aura pas de difficulté à repérer des cachettes, des coulisses ou des zones à l’abri du regard qui, non contents de jouer un rôle structurel dans l’organisation du domaine, remplissent également une fonction décisive dans les rapports unissant Wolmar et son épouse. Conditions même de l’exercice du gouvernement politique et de la structuration de l’univers social, ces lieux secrets tendent à recevoir une tonalité inquiétante lorsqu’ils permettent à l’action de Wolmar de prendre une forme intrusive, particulièrement à l’égard de Julie qu’il vise à guérir de son ancienne passion pour Saint-Preux. Cependant, Rousseau semble avoir aussi ménagé les conditions d’une possible réhabilitation du \teiHi[rend={italic}]{locus secretus}, qui, loin de n’être que le synonyme de l’opacité politique, peut s’apparenter à un refuge pour l’individu face aux excès d’un pouvoir politique guidé par des rêves d’omniscience. Récemment, Antoine Lilti, discutant les travaux d’Habermas sur la distinction entre le privé et le public, a mis en évidence que ces espaces ne s’opposaient pas, mais se constituaient plutôt l’un par rapport à l’autre\teiNote[xmlid={ftn39-2},n={6},place={foot}]{\teiP[xmlid={p-054}]{Antoine Lilti, \teiHi[rend={italic}]{L’Héritage des Lumières. Ambivalences de la modernité}, Paris, Seuil-Gallimard-EHESS, « Hautes Études », 2019, p. 192-196.}}. Le roman de Rousseau, fondé sur un entrelacement de la transparence et de l’obstacle, semble ainsi particulièrement adapté pour confirmer cette analyse, \teiHi[rend={italic}]{a fortiori }dans le cadre de Clarens où les exigences de transparence « rend\lbrack{}ent\rbrack{} d’autant plus pressant l’appel à l’intimité et à la solitude\teiNote[xmlid={ftn40-2},n={7},place={foot}]{\teiP[xmlid={p-055}]{\teiHi[rend={italic}]{Ibid.}, p. 193.}} » que Rousseau incarne de façon paradoxalement spectaculaire dans la seconde moitié du siècle.}
            \begin{teiDiv}[type={section1},xmlid={div01}]

              \teiHead[xmlid={le-pouvoir-et-le-cabinet-wolmar-001}]{Le pouvoir et le cabinet : Wolmar}
              \teiP[xmlid={p3-2}]{De prime abord, il semblerait que le \teiHi[rend={italic}]{locus secretus} ne soit pas du tout du goût de Wolmar, qui affectionne plutôt les lieux découverts, les promenades au grand air ou les sains travaux des champs. Privilégiant un rapport direct avec ses ouvriers agricoles, « M. de Wolmar les visite lui-même presque tous les jours, souvent plusieurs fois le jour » (IV, X, p. 57), et, « dans ses tournées », on apprend qu’il aime à discuter avec « quelque bon Vieillard dont le sens et la raison le frappent, et qu’il se plaît à faire causer » (V, II, p. 179). Loin de se confiner dans un cabinet, ou dans le silence d’une antichambre, le maître de Clarens est fréquemment amené à voyager dans le récit. Immédiatement après avoir accueilli Saint-Preux, Wolmar, prétextant des affaires à régler, s’en va à Étanges, non sans susciter les réticences de Julie\teiNote[xmlid={ftn41-2},n={8},place={foot}]{\teiP[xmlid={p-056}]{« Sans dire tout ce que je pensais d’un départ aussi déplacé, j’ai représenté qu’il ne paraît pas assez indispensable pour obliger M. de Wolmar à quitter un hôte qu’il avait lui-même dans sa maison » (IV, XII, p. 116).}}. Ce départ n’est pas anodin ; et les personnages s’étonnent de cette dangereuse absence qui amène à laisser les deux anciens amants en tête-à-tête\teiNote[xmlid={ftn42-2},n={9},place={foot}]{\teiP[xmlid={p-057}]{« Votre conduite n’en est pas plus convenable, et vous jouissez durement de la vertu de votre femme » (IV, XVI, p. 135).}}. Une telle mobilité est pourtant cohérente avec le personnage de Wolmar, qui manifeste quelques lignes plus tôt un curieux désir d’immatérialité : « Si je pouvais changer la nature de mon être et devenir un œil vivant, je ferais volontiers cet échange » (IV, XII, p. 109). En quittant le domaine de Clarens, Wolmar ne renonce pas en effet à exercer son droit de regard sur ses résidents, sachant parfaitement combien Julie a intériorisé cette réalité : « Tu ne veux pas me croire Cousine, mais je te jure qu’il a quelque don surnaturel pour lire au fond des cœurs » (\teiHi[rend={italic}]{ibid.}, p. 116). C’est également le cas en ce qui concerne les rapports entre les domestiques qui sont encouragés à « se servir mutuellement en secret, sans ostentation, sans se faire valoir ». Cette disposition n’est possible que dans la mesure où « ils savent fort bien que le maître, témoin de cette discrétion, les en estime davantage » (IV, X, p. 79). Wolmar ne se rend donc pas invisible en s’enfermant dans le secret d’un cabinet ; il sait se rendre aussi omniprésent qu’impalpable, en alternant présence et absence et en forçant les habitants du domaine à se sentir en permanence sous son regard. Même s’il n’intervient que rarement dans l’économie générale de Clarens, les différents personnages ne s’y trompent pas, et rattachent les évènements ponctuant la vie du domaine à celui qui en est le maître. La soirée concluant la journée de vendanges le montre parfaitement : « La pièce est éclairée de trois lampes, auxquelles M. de Wolmar a seulement fait ajouter des capuchons de fer-blanc pour intercepter la fumée et réfléchir la lumière » (V, VII, p. 239). Le philosophe se présente ainsi sous les aspects d’un dramaturge occupé à élaborer une authentique scène bucolique. L’usage de la périphrase factitive permet au romancier d’indiquer que Wolmar demeure néanmoins à distance. Il apparaît comme le véritable responsable de l’atmosphère champêtre dominant dans cette veillée qui semble se déployer sous l’effet de sa seule volonté.}
              \teiP[xmlid={p4-2}]{Wolmar\teiNote[xmlid={ftn43-2},n={10},place={foot}]{\teiP[xmlid={p-058}]{Rappelons en effet que le roman de Rousseau écarte délibérément toute forme de méchanceté : « pas une méchante action ; pas un méchant homme qui fasse craindre pour les bons » (« Entretien sur les romans », \teiHi[rend={italic}]{Julie et La Nouvelle Héloïse}, \teiHi[rend={italic}]{op. cit.}, t. II, p. 13).}}, on le voit, ne saurait se résumer à la figure rassurante d’un philosophe sage et généreux. Il comporte aussi un aspect plus inquiétant que le récit met en évidence, non sans faire écho aux structures de surveillance qui se développent au \teiHi[rend={small-caps}]{xviii}\teiHi[rend={sup}]{e} siècle et permettent de « voir sans être vu\teiNote[xmlid={ftn44-2},n={11},place={foot}]{\teiP[xmlid={p-059}]{Michel Foucault, \teiHi[rend={italic}]{Surveiller et Punir. Naissance de la prison}, Paris, Gallimard, 1975, p. 173. Voir aussi, sur un registre cette fois spécifiquement littéraire, l’article canonique de Henri Lafon, « “Voir sans être vu”, un cliché, un fantasme », \teiHi[rend={italic}]{Poétique}, n\teiHi[rend={sup}]{o} 29, 1977, p. 50-59.}} ». La dureté du personnage apparaît tout particulièrement au moment où Rousseau, renouant avec l’imaginaire du \teiHi[rend={italic}]{locus secretus}, raconte comment Wolmar révèle à Julie et Saint-Preux qu’il a connaissance de leur correspondance :}
              \begin{teiCit}[xmlid={cit01-2}]

                \begin{teiQuote}
Ensuite il nous a menés dans son cabinet, où j’ai failli tomber de mon haut en lui voyant sortir d’un tiroir, avec les copies de quelques relations de notre ami que je lui avais données, les originaux mêmes de toutes les Lettres que je croyais avoir vu brûler autrefois par Babi dans la chambre de ma mère. Voilà, m’a-t-il dit en nous les montrant les fondements de ma sécurité ; s’ils me trompaient, ce serait une folie de compter sur rien de ce que respectent les hommes. Je remets ma femme et mon honneur en dépôt à celle qui, fille et séduite, préférait un acte de bienfaisance à un rendez-vous unique et sûr. Je confie Julie épouse et mère à celui qui maître de contenter ses désirs sut respecter Julie amante et fille. Que celui de vous deux qui se méprise assez pour penser que j’ai tort le dise, et je me rétracte à l’instant. Cousine, crois-tu qu’il fût aisé d’oser répondre à ce langage ? (IV, XII, p. 117)
\end{teiQuote}
              
\end{teiCit}
              \teiP[xmlid={p5-2}]{Le cabinet, où se nouent les liens entre savoir et pouvoir\teiNote[xmlid={ftn45-2},n={12},place={foot}]{\teiP[xmlid={p-060}]{Voir en particulier la célèbre scène d’intrusion où Wolmar se plaît à faire effraction dans le cabinet de Julie et à surprendre cette dernière en prière. Le maître de Clarens se félicite de savoir parfaitement ce que fait son épouse lorsqu’elle se retire : « Devinez où elle est, me dit son mari voyant que je la cherchais des yeux. » La violence de l’intrusion (« Nous arrivâmes à la porte du cabinet ; elle était fermée. Il l’ouvrit brusquement ») permet ainsi de souligner la façon dont le philosophe met sa science de l’observation au service d’une entreprise démiurgique dans laquelle il s’agit bien d’exercer un pouvoir sur Julie (V, V, p. 225-226).}}, permet à Wolmar de manifester son emprise sur les deux anciens amants. En les entraînant dans le secret du cabinet, il théâtralise le moment de la révélation qui, comme le suggère le recours à un discours hautement dramatisé, intimide fortement ses auditeurs. L’usage du présentatif, du parallélisme de construction, ou du martèlement anaphorique de la première personne, vient de cette manière conférer à son propos une emphase qui laisse Julie sans voix. Le romancier semble en effet nous indiquer combien l’alternative que propose Wolmar n’en est pas une pour les deux héros, tout particulièrement pour Julie qui est blessée par l’exhibition des lettres passées. L’espace secret du cabinet semble occuper un rôle essentiel dans l’univers de Clarens ; s’il représente une manière de « m\lbrack{}ettre\rbrack{} en scène \lbrack{}…\rbrack{} la hiérarchie\teiNote[xmlid={ftn46-2},n={13},place={foot}]{\teiP[xmlid={p-061}]{Georges Balandier, \teiHi[rend={italic}]{Le Pouvoir sur scènes}, Paris, Fayard, 2006, p. 40.}} », le lieu à l’écart permet au romancier de suggérer le caractère intrusif et violent du pouvoir qu’exerce Wolmar sur son épouse.}
            
\end{teiDiv}
            \begin{teiDiv}[type={section1},xmlid={div02}]

              \teiHead[xmlid={les-coulisses-et-le-jardin-julie-001}]{Les coulisses et le jardin : Julie}
              \teiP[xmlid={p6-2}]{L’usage du \teiHi[rend={italic}]{locus secretus }n’est pas réservé au seul Wolmar. Au cours de la présentation de Clarens, Saint-Preux ne manque pas de remarquer que Julie en fait également un usage abondant. Alors que le maître de Clarens, sans toujours en prendre la pleine mesure, tire les espaces secrets du côté de la manipulation, Julie met le lieu privé de la chambre au service du bonheur collectif. C’est par exemple le cas lors des veillées où un vieil homme est convié à la table des maîtres. Julie organise alors une émouvante cérémonie d’échange de dons avec les enfants :}
              \begin{teiCit}[xmlid={cit02-2}]

                \begin{teiQuote}
Elle passe après le dîné \lbrack{}\teiHi[rend={italic}]{sic}\rbrack{} dans sa chambre, et en rapporte un petit présent de quelque nippe convenable à la femme ou aux filles du vieux bonhomme. Elle le lui fait offrir par les enfants, et réciproquement il rend aux enfants quelque don simple et de leur goût dont elle l’a secrètement chargé pour eux. Ainsi se forme de bonne heure l’étroite et douce bienveuillance qui fait la liaison des états divers (V, II, p. 180).
\end{teiQuote}
              
\end{teiCit}
              \teiP[xmlid={p7-2}]{Les fonctions dramaturgiques associées à Wolmar sont déléguées au personnage de Julie. L’atmosphère en est néanmoins différente. Le \teiHi[rend={italic}]{locus secretus }n’est plus un discret cabinet, mais prend dans cette scène la forme de la chambre conjugale, espace familial et maternel que la mère de famille utilise comme les coulisses d’un spectacle théâtral. Ces subtiles manœuvres ne sont de fait nullement inconnues de la collectivité. Le bon vieillard et les enfants savent pertinemment que l’échange des dons n’est pas le fruit d’une intention spontanée, mais qu’il s’inscrit en réalité dans une stratégie subtile visant à garantir l’union entre les générations et les univers sociaux. On répondra sans doute que ce n’est là qu’une scène privée. Encore faut-il ne pas sous-estimer l’écho qu’elle semble susciter auprès de l’ensemble de la communauté de Clarens : « les paysans, voyant leurs vieux pères fêtés dans une maison respectable et admis à la table des maîtres, ne se tiennent point offensés d’en être exclus » (V, II, p. 180). Ainsi le recours à l’espace du secret a-t-il pour fonction, non pas de permettre au pouvoir politique de s’exercer à l’insu de tous, mais plutôt de venir unifier la communauté autour d’une seule et unique joie qui se diffuse par-delà les générations. On songe à la \teiHi[rend={italic}]{Lettre sur les spectacles }(1758) où c’est le père qui introduit le jeune Jean-Jacques au bonheur de la fête\teiNote[xmlid={ftn47-2},n={14},place={foot}]{\teiP[xmlid={p-062}]{« Lettre à d’Alembert », Bernard Gagnebin et Jean Rousset (éd.), \teiHi[rend={italic}]{Œuvres complètes} (t. V), Paris, Gallimard, « Bibliothèque de la Pléiade », 1995, p. 124.}}. Mais ce qui relie ici tous les habitants de la communauté de Clarens, c’est finalement la connaissance d’un même secret (le caractère faussement spontané de la cérémonie) autour duquel se réunit la collectivité. Loin d’être sue de la seule Julie, la gracieuse supercherie de l’échange des dons intègre en réalité les différents membres de la petite société. Le détour par l’espace du secret n’est pas le corolaire d’un pouvoir qui cherche à se soustraire à l’attention de ceux sur lesquels il s’exerce. Bien au contraire, la chambre familiale où Julie s’absente brièvement permet à cette dernière de susciter au sein de la veillée une émotion qui réunit la communauté dans une même joie unificatrice. On retrouve ainsi, comme dans l’épisode des vendanges (V, VII, p. 232-241), « l’ivresse de la joie publique », dans laquelle « chacun est à la fois \teiHi[rend={italic}]{acteur }et \teiHi[rend={italic}]{spectateur}\teiNote[xmlid={ftn48-2},n={15},place={foot}]{\teiP[xmlid={p-063}]{Jean Starobinski, \teiHi[rend={italic}]{Jean-Jacques Rousseau. La Transparence et l’Obstacle}, \teiHi[rend={italic}]{op. cit.}, p. 121.}} » et que Rousseau a sans doute connue lors des fêtes genevoises qui ont bercé son enfance\teiNote[xmlid={ftn49-2},n={16},place={foot}]{\teiP[xmlid={p-064}]{« Bientôt les fenêtres furent pleines de spectatrices qui donnaient un nouveau zèle d’esclaves ; elles ne purent longtemps tenir à leurs fenêtres, elles descendirent » (« Lettre à d’Alembert », \teiHi[rend={italic}]{op. cit.}, p. 123).}}. Les espaces privé et public ne se contentent donc pas de se définir l’un par rapport à l’autre\teiNote[xmlid={ftn50-2},n={17},place={foot}]{\teiP[xmlid={p-065}]{Antoine Lilti, \teiHi[rend={italic}]{L’Héritage des Lumières}, \teiHi[rend={italic}]{op. cit.}, p. 190.}} ; ils s’entrelacent et se rendent mutuellement possibles, sans pour autant se confondre, permettant la constitution d’un univers social où individu et communauté entretiennent des rapports équilibrés.}
              \teiP[xmlid={p8-2}]{Tel n’est pas l’unique exemple de l’usage que fait Julie de l’espace secret ou inaccessible. Tout au long de la seconde partie, Rousseau ne manque pas d’éveiller l’attention de son lecteur sur le fait que Clarens contient de nombreux lieux cachés, qui, à chaque fois, possèdent une fonction précise dans l’organisation de l’existence quotidienne, tout particulièrement dans ses rapports avec le plaisir. C’est tout d’abord le cas du « cabinet » dans lequel Julie, « prenant \lbrack{}Saint-Preux\rbrack{} par la main », l’emmène afin de lui exposer ses grands principes d’éducation « sans être entendus des enfants » (V, III, p. 187). Mais c’est également celui du salon d’Apollon, « petite salle à manger différente de celle où l’on mange ordinairement ». Son accès est restreint. « Les simples hôtes n’y sont point admis » ; et, si Milord Edouard y est convié dès son arrivée, Saint-Preux n’a nullement droit à un tel privilège : « Ce ne fut qu’à mon retour de chez M\teiHi[rend={sup}]{me} d’Orbe que je fus traité dans le salon d’Apollon » (V, II, p. 168). Rousseau fait ici écho à l’admission de Saint-Preux dans les cercles les plus choisis de la capitale dont cet espace de convivialité forme le reflet inverse sur un plan éthique ; mais, plus essentiellement, le romancier s’emploie à suggérer comment cet « asile inviolable de la confiance, de l’amitié, de la liberté » (\teiHi[rend={italic}]{ibid.}) permet d’indiquer les subtiles relations hiérarchiques qui structurent la vie dans le domaine. Il existe toutefois des lieux auxquels les maîtres du logis eux-mêmes n’ont pas un accès permanent. Quoique maîtresse du salon d’Apollon, Julie veille par exemple à ne pas en faire un usage excessif qui empêcherait de goûter le plaisir de l’exceptionnel. Comme Saint-Preux l’interroge à ce sujet, Julie lui fait comprendre que « l’art d’assaisonner les plaisirs n’est que celui d’en être avare » (\teiHi[rend={italic}]{ibid.}, p. 169). On retrouve ainsi l’idée d’un « art de jouir » dans lequel les privations permettent « en réalité d’accroître l’intensité et la qualité des plaisirs\teiNote[xmlid={ftn51-2},n={18},place={foot}]{\teiP[xmlid={p-066}]{Céline Spector, « Rousseau : éthique et économie », \teiHi[rend={italic}]{Cahiers d’économie politique}, n\teiHi[rend={sup}]{o} 53, p. 42.}} ». Remarquable est le traitement qui est réservé à l’espace du gynécée, qui permet également de signifier à Saint-Preux la place qui lui revient dans l’économie générale de Clarens. Contrairement à Paris où la promiscuité qui règne entre les sexes est présentée comme délétère, « les femmes \lbrack{}à Clarens\rbrack{} vivent très séparées des hommes ». Elles ont l’habitude de se retrouver dans « la chambre des enfants », où elles s’occupent des plus jeunes et profitent d’une collation « composée de quelques laitages, de gaufres, d’échaudés, de merveilles, ou d’autres mets du goût des enfants et des femmes » (IV, X, p. 65-66). Si elle s’apparente à un plaisir innocent, la profusion de friandises peut se lire ici comme une forme de compensation érotique à l’absence d’homme, permettant de neutraliser les effets du désir qui menacent de perturber l’ordre social de Clarens. On sait en effet que Wolmar s’emploie à compenser l’absence d’un plaisir dangereux en lui substituant un plaisir dépourvu de toute potentialité subversive, par exemple lorsqu’il remplace les jeux des cabarets par des divertissements où les habitants du domaine ne peuvent pas se corrompre (\teiHi[rend={italic}]{ibid.}, p. 68-69). Quant à Saint-Preux, il semble porter un regard pour le moins hautain sur ces divertissements féminins, se plaisant à mettre sur le même plan les femmes et les enfants. Mais la logique de l’énonciation à la première personne invite à se montrer attentif aux éventuels aveuglements de la voix narratoriale. Le jeune homme oublie que son admission dans le secret du gynécée, espace réservé aux femmes et aux enfants, est en réalité lourde de sens. Elle correspond aussi à une façon de lui signifier que sa renonciation définitive à Julie et plus généralement au désir est bien la condition \teiHi[rend={italic}]{sine qua non }de son admission dans le domaine de Clarens. À ce titre, la critique a pu noter que Saint-Preux trouvait également dans les laitages qu’on lui sert « une jouissance substitutive du corps de Julie\teiNote[xmlid={ftn52-2},n={19},place={foot}]{\teiP[xmlid={p-067}]{Christophe Martin, \teiHi[rend={italic}]{Espaces du féminin}, Oxford, Voltaire Foundation, 2004, p. 414.}} ». En témoigne la dimension fantasmatique du texte au moment où Saint-Preux nous invite à rêver à « ce que doivent être \lbrack{}les laitages\rbrack{} d’une laiterie où Julie préside » (IV, X, p. 66), ce qui permet de voir dans cet innocent goûter une manière de neutraliser le désir de Saint-Preux en même temps que de le révéler, peut-être davantage \teiHi[rend={italic}]{in fine }que celui des femmes dont la voix est en partie tue.}
              \teiP[xmlid={p9-2}]{L’espace secret comporte toutefois une seconde fonction : il permet à Julie d’échapper au regard clairvoyant de son époux, qui – notons-le – n’est pas autorisé à pénétrer dans l’espace du gynécée. On pourrait bien sûr penser que c’est là la spécificité première de ce lieu réservé aux femmes, et que Julie ne fait jamais que se conformer au « contrôle invisible mais sévère\teiNote[xmlid={ftn53-2},n={20},place={foot}]{\teiP[xmlid={p-068}]{\teiHi[rend={italic}]{Ibid.}, p. 255.}} » du philosophe qui semble se plaire à enfermer l’héroïne dans la sphère des plaisirs féminins. Mais le soin que prend Julie à exclure son époux et surtout à inclure Saint-Preux invite à adopter une autre lecture : les allusions de ce dernier à leur ancienne passion\teiNote[xmlid={ftn54-2},n={21},place={foot}]{\teiP[xmlid={p-069}]{« \lbrack{}…\rbrack{} la raison peut s’égarer dans un chalet tout aussi bien que dans un cellier » (IV, X, p. 67).}} semblent faire du gynécée un lieu paradoxal où le désir coupable peut vivre à l’insu même de la vigilance du maître de Clarens. On peut à cet égard y voir un écho à ce célèbre aveu de Wolmar qui déplore l’opacité entourant le cœur de Julie : « on n’en peut parler que par conjecture : un voile de sagesse et d’honnêteté fait tant de replis autour de son cœur, qu’il n’est plus possible à l’œil humain d’y pénétrer » (IV, XIV, p. 129). Le cœur de Julie se fait ici opaque, et semble mettre en échec la souveraineté de la transparence que revendique Wolmar. Les « conjectures » se substituent à l’immédiateté de la connaissance qui semble pourtant dominante à Clarens. Le motif du voile, déjà très présent dans la première partie du roman\teiNote[xmlid={ftn55-2},n={22},place={foot}]{\teiP[xmlid={p-070}]{Voir les analyses de Jean Starobinski, dans \teiHi[rend={italic}]{La Transparence et l’Obstacle}, \teiHi[rend={italic}]{op. cit}., p. 142-148.}}, trouve une manifestation également exemplaire lorsque Saint-Preux évoque le mystérieux jardin de Julie :}
              \begin{teiCit}[xmlid={cit03-2}]

                \begin{teiQuote}
Ce lieu, quoique tout proche de la maison, est tellement caché par l’allée couverte qui l’en sépare, qu’on ne l’aperçoit de nulle part. L’épais feuillage qui l’environne ne permet point à l’œil d’y pénétrer, et il est toujours soigneusement fermé à clef (IV, XI, p. 87).
\end{teiQuote}
              
\end{teiCit}
              \teiP[xmlid={p10-2}]{Cet espace clos est aussi inaccessible à l’extérieur que fermé sur lui-même. En atteste la présence de la négation grammaticale qui vient refléter le refus que Julie semble, en partie inconsciemment, opposer au règne de la transparence voulu par Wolmar. Il est difficile à cet égard de ne pas identifier le jardin à sa créatrice, dont le cœur ne cesse de se soustraire aux investigations. Même si le jardin n’est nullement interdit au maître de Clarens qui a l’habitude d’y travailler comme « garçon jardinier » (\teiHi[rend={italic}]{ibid.}, p. 89), il est clair que la véritable fonction du jardin échappe à l’attention de Wolmar. La critique a depuis longtemps souligné l’importance des connotations érotiques de l’Élysée\teiNote[xmlid={ftn56-2},n={23},place={foot}]{\teiP[xmlid={p-071}]{André Blanc, « Le Jardin de Julie », \teiHi[rend={italic}]{Dix-huitième siècle}, n\teiHi[rend={sup}]{o} 14, 1982, p. 360.}}. Il ne s’agit pas uniquement de l’espace dévolu à « \lbrack{}l\rbrack{}a promenade favorite » de « la plus respectable des mères de famille » (IV, XI, p. 87). Le jardin de Julie, objet d’un fort investissement affectif, offre à cette dernière « un lieu de \teiHi[rend={italic}]{suppléance}, où la jouissance féminine échappe, au moins partiellement, à l’emprise de la Loi\teiNote[xmlid={ftn57-2},n={24},place={foot}]{\teiP[xmlid={p-072}]{Christophe Martin, \teiHi[rend={italic}]{La Philosophie des amants}, \teiHi[rend={italic}]{op. cit.}, p. 212.}} ». On sait que Wolmar, s’appuyant sur une forme d’« \teiHi[rend={italic}]{Ars oblivationis}\teiNote[xmlid={ftn58-2},n={25},place={foot}]{\teiP[xmlid={p-073}]{\teiHi[rend={italic}]{Ibid.}, p. 189.}} », poursuit le dessein de « guérir » son épouse de son ancienne passion pour Saint-Preux, et d’éteindre une fois pour toutes les agitations de sa première jeunesse. De ce point de vue, l’espace du secret n’est pas simplement le lieu de récréation infantilisant d’une vertueuse mère au foyer, il permet à Rousseau d’annoncer le retournement final du roman, et de trahir le fait qu’en dépit de l’entreprise thérapeutique de Wolmar, la passion de Julie pour Saint-Preux demeure le centre rayonnant de l’ensemble du roman.}
            
\end{teiDiv}
            \begin{teiDiv}[type={section1},xmlid={div03}]

              \teiHead[xmlid={malheur-a-qui-na-plus-rien-a-desirer-001}]{« Malheur à qui n’a plus rien à désirer »}
              \teiP[xmlid={p11-2}]{La vocation de cet espace est peut-être, moins de combler un vide, que d’en maintenir la possibilité. Les « chagrins secrets » (V, IV, p. 214) de Julie indiquent que l’existence de l’Élysée ne suffit pas à compenser la tristesse qu’elle éprouve face à l’athéisme de son époux. On connaît à cet égard la formule de Pascal, selon laquelle « un portrait porte absence et présence, plaisir et déplaisir\teiNote[xmlid={ftn59-2},n={26},place={foot}]{\teiP[xmlid={p-074}]{Pascal, \teiHi[rend={italic}]{Pensées}, Philippe Sellier (éd.), Paris, Le Livre de poche, 2000, p. 197, pensée 291.}} ». L’indéniable sensualité du jardin, quoiqu’elle rappelle assurément l’épisode du bosquet, n’est pas en mesure de rivaliser avec celle des « âcres \lbrack{}baisers\rbrack{} » (I, XVI, p. 110) que les amants ont échangés dans la première partie de l’œuvre. Là n’est sans doute pas, cependant, sa véritable destinée, qui est d’entretenir le désir de Julie pour son ancien amant, et non uniquement d’en compenser la perte. Alors que le bosquet, situé significativement à l’opposé de l’Élysée, a pour effet principal de ramener les amants vers un passé ineffaçable, le jardin, vierge de tout souvenir, possède plutôt pour spécificité d’ouvrir les perspectives et d’amener l’héroïne à désirer un avenir qui se situe au-delà des bornes de cette vie. Le nom même de l’Élysée, lourd de connotations mortuaires, semble l’illustrer exemplairement. La jeune femme explique ainsi le plaisir qu’elle prend à imaginer ses enfants travailler en « petits jardiniers » (IV, XI, p. 103). Ce fantasme, qui ne verra pas sa réalisation\teiNote[xmlid={ftn60-2},n={27},place={foot}]{\teiP[xmlid={p-075}]{« \lbrack{}…\rbrack{} dites à Marcellin qu’il ne m’en coutera pas de mourir pour lui » (VI, XII, p. 388).}}, semble en mesure de s’élargir à l’idée d’une vie après la mort, qui pourrait être également le moment des retrouvailles avec Saint-Preux\teiNote[xmlid={ftn61-2},n={28},place={foot}]{\teiP[xmlid={p-076}]{« Non, je ne te quitte pas, je vais t’attendre » (\teiHi[rend={italic}]{ibid.}).}} : « En vérité, mon ami, me dit-elle d’une voix émue, des jours ainsi passés tiennent du bonheur de l’autre vie ; et ce n’est pas sans raison qu’en y pensant j’ai donné d’avance à ce lieu le nom d’Élysée » (IV, XI, p. 104). Le jardin, s’il possède d’évidentes vertus compensatoires, ne permet pas seulement l’apaisement du désir ; il tend aussi à le prolonger, en même temps qu’à le protéger, à l’insu même de Wolmar. Il nous semble donc possible de poser un regard distancié sur le souhait qu’a Julie de ne rien cacher à son époux. Certes, l’héroïne se félicite que son époux dispose de « quelque don surnaturel pour lire au fond des cœurs » (IV, XII, p. 116) ; mais l’ensemble de ses actions indique bien que cela ne l’empêche pas d’opposer une forme de résistance à l’action démiurgique de son époux, sans pour autant qu’un tel dessein ne parvienne à franchir de façon aussi claire le seuil de la verbalisation.}
              \teiP[xmlid={p12-2}]{L’analyse du jardin nous permet en ce sens de mieux comprendre la dernière lettre de Julie. Ce testament spirituel, à première vue, vient confirmer ce que le lecteur pressent depuis longtemps, à savoir la persistante mélancolie de Julie : « Et toutefois j\lbrack{}e\rbrack{} vis inquiète. Mon cœur ignore ce qui lui manque ; il désire, sans savoir quoi » (VI, VIII, p. 334). On sait, au moins depuis les travaux essentiels de Jean Deprun\teiNote[xmlid={ftn62-2},n={29},place={foot}]{\teiP[xmlid={p-077}]{Jean Deprun, \teiHi[rend={italic}]{La Philosophie de l’inquiétude}, Paris, Vrin, 1979, p. 179.}}, combien l’agitation de Julie s’enracine dans une longue tradition théologique qui, dans la lignée de Saint Augustin, fait de l’inquiétude l’indice le plus clair du manque que suscite en nous l’absence de Dieu. Mais, tout en reprenant à son compte cette analyse, Rousseau semble l’infléchir. En effet, le cas emblématique de l’Élysée indique assez que le vide qui habite le cœur de Julie est plutôt le résultat d’un travail en partie souterrain\teiNote[xmlid={ftn63-2},n={30},place={foot}]{\teiP[xmlid={p-078}]{Nous rejoignons donc les analyses récentes de Gabrielle Radica qui a pu suggérer que « le rapport au passé \lbrack{}chez Julie\rbrack{} n’es\lbrack{}t\rbrack{} pas principalement porté par la conscience et la connaissance, mais par un effort et une ressaisie de soi » (« Qu’est devenue Julie ? Sur le perfectionnisme dans \teiHi[rend={italic}]{Julie ou la Nouvelle Héloïse }», \teiHi[rend={italic}]{Annales de la société Jean-Jacques Rousseau}, t. 54, 2021, p. 361).}} visant à maintenir la sensation du manque causé par Saint-Preux malgré la profusion que cherche à faire régner Wolmar. C’est du moins ce que suggère la réponse qu’elle apporte à son ancien amant, lorsqu’il lui demande pourquoi le salon d’Apollon n’est pas plus utilisé : « cela serait trop agréable, et \lbrack{}…\rbrack{} \teiHi[rend={italic}]{l’ennui d’être toujours à son aise} est enfin le pire de tous » (V, II, p. 168). La privation volontaire n’a donc pas pour seule vocation d’intensifier les plaisirs. Elle a également pour fonction de déjouer la logique même de Clarens qui repose sur la satiété de ceux qui y résident, « je ne vois partout que sujets de contentement » (VI, VIII, p. 334). Nulle raison de s’étonner si Julie explique à Saint-Preux : « Malheur à qui n’a plus rien à désirer ! Il perd pour ainsi dire tout ce qu’il possède » (\teiHi[rend={italic}]{ibid.}, p. 333). Dieu, on le voit, possède une tout autre fonction que celle qu’il a chez les moralistes chrétiens. Il ne vient pas étancher la soif spirituelle de Julie, si bien que les temps de prière qu’elle prend permettent surtout de relancer son désir et d’entretenir en elle l’espérance d’un ailleurs où elle pourra goûter une plénitude qu’elle ne trouve nulle part : « dégagée un moment de ses entraves, \lbrack{}mon âme\rbrack{} se console d’y rentrer, par cet essai d’un état plus sublime, qu’elle espère être un jour le sien » (\teiHi[rend={italic}]{ibid.}, p. 334-335). L’élan amoureux et l’élan spirituel se rejoignent ici en une seule force désirante qui refuse de se limiter aux limites de l’univers sensible ; car, comme Julie le confie à son amant dans la dernière lettre, « la vertu qui \lbrack{}les\rbrack{} sépara sur la terre, les unira dans le séjour éternel » (VI, XII, p. 388). Sans naturellement se réduire à ce seul aspect, l’aspiration à une vie après la mort, on le voit, s’unit également à la perspective d’une union parfaite avec Saint-Preux, ce qui invite à inscrire la vocation spirituelle de Julie dans la continuité de sa recherche de bonheur sur la terre. Or l’espérance d’un monde meilleur, que Wolmar associe à « un opium pour l’âme » (VI, VIII, p. 337), ne saurait s’apparenter à un divertissement égocentré ; il semble profiter à la collectivité, s’imposant ainsi comme un élément déterminant dans l’espace social. « En sortant de \lbrack{}s\rbrack{}on cabinet ainsi disposée », Julie n’en est que plus heureuse et plus déterminée à servir son prochain : « \lbrack{}la dévotion\rbrack{} égaie, anime et soutient quand on en prend peu » (\teiHi[rend={italic}]{ibid.}, p. 337). Le \teiHi[rend={italic}]{locus secretus}, qui prend ici la forme du lieu de la prière ou du recueillement, peut se lire comme un cadre favorable à l’émergence d’un « état d’effervescence affective\teiNote[xmlid={ftn64},n={31},place={foot}]{\teiP[xmlid={p-079}]{Henri GouhierGouhier, \teiHi[rend={italic}]{Les Méditations métaphysiques de Jean-Jacques Rousseau}, Paris, Vrin, 2005, p. 129.}} » qui devient la source de l’esprit qui anime et vivifie la construction rationnelle du sage Wolmar : « j’en aime encore mieux ceux que j’aime » (\teiHi[rend={italic}]{ibid.}). La sensation de vide semble pour cette raison indispensable à l’univers de Clarens. Au début de la seconde partie du roman, Julie, soucieuse de tirer Claire de sa solitude après la mort de son époux, lui explique que la « tristesse », lorsqu’elle se diffuse sous la forme de « la communication des cœurs », prend « je ne sais quoi de doux et de touchant que n’a pas le contentement » (IV, I, p. 13). C’est pourquoi Julie est en mesure d’apparaître comme l’âme du domaine, par exemple lorsque Saint-Preux insiste sur le dévouement des habitants pour M\teiHi[rend={sup}]{me} de Wolmar : « son seul regard anime leur zèle, en sa présence ils sont contents, en son absence ils parlent d’elle et s’animent à la servir » (IV, X, p. 58). Lors de la mort de Julie, la seule image de cette dernière couverte d’un voile parvient même à calmer l’émotion populaire que le rationalisme de Wolmar était incapable de maîtriser (VI, XI, p. 381-382), invitant à se souvenir de cette affirmation de Claire : « Ma Julie, tu es faite pour régner » (IV, II, p. 19). Cette scène amène à comprendre que Clarens ne se réduit pas à la seule transparence, mais qu’il repose d’abord sur la présence de ces lieux cachés permettant à Julie de préserver en elle l’élan du désir et de rester cette âme sensible qui invite « tout ce qui l’approche \lbrack{}à\rbrack{} lui ressembler\teiNote[xmlid={ftn65},n={32},place={foot}]{\teiP[xmlid={p-080}]{« Entretien sur les romans », \teiHi[rend={italic}]{op. cit.}, p. 413.}} ».}
              \teiP[xmlid={p13-2}]{Ainsi le \teiHi[rend={italic}]{locus secretus }ne perd-il nullement sa signification politique. S’il diffère fortement du lieu des intrigues et des tractations qu’il était sous l’Ancien Régime, il devient la source d’une « secrète résistance\teiNote[xmlid={ftn66},n={33},place={foot}]{\teiP[xmlid={p-081}]{Christophe Martin, \teiHi[rend={italic}]{La Philosophie des amants}, \teiHi[rend={italic}]{op. cit.}, p. 212.}} » à l’action d’un pouvoir intrusif, déterminé, sous couvert de bonnes intentions, à manipuler les consciences et les mémoires. Là n’est sans doute pas l’unique leçon du roman ; car l’espace secret, en permettant de ménager au sein de la profusion qui règne à Clarens des zones de vide et de frustration, donne également à Julie les moyens d’animer l’ensemble du domaine qui, dépourvu de « la douce influence de cette âme expansive » (V, III, p. 185), resterait une horloge sans vie. On peut ainsi voir dans la lecture, tout particulièrement romanesque, une dernière réactualisation du \teiHi[rend={italic}]{locus secretus}. En invitant le lecteur à se retirer dans sa chambre et à compatir aux malheurs de Julie et de Saint-Preux, la fiction épistolaire n’en joue pas pour autant un rôle secondaire dans l’espace public : elle constitue, un peu à la manière de Julie dans Clarens, « un espace public sensible\teiNote[xmlid={ftn67},n={34},place={foot}]{\teiP[xmlid={p-082}]{Antoine Lilti, \teiHi[rend={italic}]{L’Héritage des Lumières}, \teiHi[rend={italic}]{op. cit.}, p. 190.}} » où les individualités, unies dans une même chaîne d’émotion, sont appelées à vivre une communion qui donne souffle à tout l’univers social.}
            
\end{teiDiv}
          
\end{teiElement}
        
\end{teiElement}
        \begin{teiElement}[name={group},type={chapter},data-page-title={Les lieux de la conjuration : société secrète et hétérotopie dans la littérature romantique},data-page-authors={Guillaume Cousin@@Université d’Artois, Textes et Cultures, Arras, France},data-include-href={../XML/chap\_3\_Les\_lieux\_de\_la\_conjuration\_Locus\_politicus.xml},xmlid={chapter-002}]

          \begin{teiElement}[name={front}]

            \begin{teiDiv}[type={titlePage},xmlid={titlepage-003}]

              \teiP[rend={title-main},xmlid={p-083}]{Les lieux de la conjuration : société secrète et hétérotopie dans la littérature romantique}
              \teiP[rend={author-aut},xmlid={p-084}]{Guillaume \teiHi[rend={small-caps}]{Cousin}}
              \teiP[rend={authority\_affiliation},xmlid={p-085}]{Université d’Artois, Textes et Cultures, Arras, France}
            
\end{teiDiv}
          
\end{teiElement}
          \begin{teiElement}[name={body},xmlid={body-003}]

            \teiP[xmlid={p1-3}]{Le \teiHi[rend={small-caps}]{xix}\teiHi[rend={sup}]{e} siècle français et européen est, pour reprendre l’expression de Jean-Noël Tardy, « l’âge des ombres\teiNote[xmlid={ftn1-3},n={1},place={foot}]{\teiP[xmlid={p-086}]{Jean-Noël Tardy, \teiHi[rend={italic}]{L’Âge des ombres. Complots, conspirations et sociétés secrètes au }\teiHi[rend={ small-caps italic}]{xix}\teiHi[rend={italic sup}]{e}\teiHi[rend={italic}]{ siècle}, Paris, Les Belles Lettres, 2015.}} » : tout au long de la période, les différents régimes politiques sont confrontés à des puissances occultes, qui peuvent se manifester à travers des complots ou des conspirations qui visent le plus souvent à renverser l’ordre établi\teiNote[xmlid={ftn53-3},n={2},place={foot}]{\teiP[xmlid={p-087}]{Nous ne traiterons pas, dans les pages suivantes, les cabinets noirs ou gouvernements parallèles, qui impliquent l’action des gouvernants eux-mêmes.}}. Parmi les forces occultes, les sociétés secrètes occupent une place centrale : leur développement à la fin de l’Ancien Régime amène Barruel à considérer que la Franc-maçonnerie et les Illuminés de Bavière ont une part de responsabilité dans le surgissement de la Révolution\teiNote[xmlid={ftn54-3},n={3},place={foot}]{\teiP[xmlid={p-088}]{Dans ses \teiHi[rend={italic}]{Mémoires pour servir à l’histoire du jacobinisme}, Hambourg, chez P. Fauché, 1798-1799.}}. Héritier de ce culte du secret, le \teiHi[rend={small-caps}]{xix}\teiHi[rend={sup}]{e} siècle voit se multiplier les sociétés secrètes à caractère politique : les Chevaliers de la Foi, la Congrégation, la Charbonnerie… Qu’elles soient ultraroyalistes ou libérales, monarchistes ou républicaines, nationalistes ou internationalistes, elles suscitent fascination et méfiance et charrient avec elles un « imaginaire tentaculaire\teiNote[xmlid={ftn55-3},n={4},place={foot}]{\teiP[xmlid={p-089}]{Nicolas Aude et Marie-Agathe Tilliette, « Introduction. Un imaginaire tentaculaire », \teiHi[rend={italic}]{L’Imaginaire des sociétés secrètes dans la littérature du }\teiHi[rend={ small-caps italic}]{xix}\teiHi[rend={italic sup}]{e}\teiHi[rend={italic}]{ siècle}, Société des études romantiques et dix-neuviémistes, 2021, en ligne : \teiRef[target={https://serd.hypotheses.org/files/2021/05/1.-Introduction-Aude-Tilliette.pdf}]{\teiHi[rend={underline}]{https://serd.hypotheses.org/files/2021/05/1.-Introduction-Aude-Tilliette.pdf}} (consulté le 22 août 2023).}} » qui imprègne la littérature du temps. \teiHi[rend={italic}]{A priori} distinctes, conspiration et littérature partagent en réalité un même régime discursif, la fiction :}
            \begin{teiCit}[xmlid={cit01-3}]

              \begin{teiQuote}
La conspiration possède de multiples dimensions et on peut l’appréhender tantôt comme une pratique de contestation politique, tantôt comme un crime politique ou comme une vision du monde et de l’histoire. L’élément commun à toutes ces manifestations est la présence d’une fiction, ou plutôt de fictions, aux fonctions différentes\teiNote[xmlid={ftn56-3},n={5},place={foot}]{\teiP[xmlid={p-090}]{Jean-Noël Tardy, \teiHi[rend={italic}]{L’Âge des ombres}, \teiHi[rend={italic}]{op. cit.}, p. 18.}}.
\end{teiQuote}
            
\end{teiCit}
            \teiP[xmlid={p2-3}]{L’hyperconscience historique des écrivains romantiques se nourrit de la mythologie de la société secrète, dans laquelle ils perçoivent à la fois des forces agissant au-delà des apparences et des sources d’explication du monde et de l’histoire. Il existe, selon le Vautrin des \teiHi[rend={italic}]{Illusions perdues} (où il se fait appeler Carlos Herrera), deux historiographies possibles : « Il y a deux Histoires : l’Histoire officielle, menteuse qu’on enseigne, l’Histoire \teiHi[rend={italic}]{ad usum delphini }; puis l’Histoire secrète, où sont les véritables causes des événements, une Histoire honteuse\teiNote[xmlid={ftn57-3},n={6},place={foot}]{\teiP[xmlid={p-091}]{Honoré de Balzac, \teiHi[rend={italic}]{Illusions perdues} \lbrack{}1843\rbrack{}, Roland Chollet (éd.), dans Pierre-Georges Castex (dir.), \teiHi[rend={italic}]{La Comédie humaine}, Paris, Gallimard, « Bibliothèque de la Pléiade », 1977, t. V, p. 695.}}. » À cette assertion font écho les trois épisodes de l’\teiHi[rend={italic}]{Histoire des Treize}, symboliquement placés après \teiHi[rend={italic}]{Illusions perdues} dans \teiHi[rend={italic}]{La Comédie humaine }: la société secrète y est l’exemple même de cette autre Histoire.}
            \teiP[xmlid={p3-3}]{Les fictions romantiques se font les relais de ces fictions socio-politiques et exploitent le potentiel romanesque et dramatique qu’offrent les sociétés secrètes. Dans l’introduction de l’ouvrage qu’ils ont dirigé, Nicolas Aude et Marie-Agathe Tilliette dressent la liste des enjeux de la représentation littéraire de ces sociétés secrètes : politiques, historiques, historiographiques, sociaux… mais aussi, pouvons-nous ajouter, topographiques. Reposant sur le culte du secret, la société secrète doit néanmoins posséder un ou des lieux de réunion où les secrets se disent et s’échangent. Nicolas Aude et Marie-Agathe Tilliette écrivent que « des métaphores diverses s’appliquent en effet aux sociétés secrètes, tantôt vues comme des réseaux, comme des toiles ou comme des miroirs concentrés du social », mais ajoutent que « ces collectifs sont toujours fondés sur une cohérence – de personnes ou d’idées – qui manque à la société globale, menacée de dissociations diverses\teiNote[xmlid={ftn58-3},n={7},place={foot}]{\teiP[xmlid={p-092}]{Nicolas Aude et Marie-Agathe Tilliette, « Introduction. Un imaginaire tentaculaire », art. cité, non paginé.}} » ; par la présence de la société secrète dans le lieu où elle tient ses réunions, celui-ci deviendrait alors un \teiHi[rend={italic}]{locus secretus} à la fois intérieur et extérieur à la société, dont l’entrée serait réservée aux membres du groupe et qui serait à l’origine d’une action sur le champ social et politique. En cela, le \teiHi[rend={italic}]{locus secretus} de la société secrète se rapprocherait de l’hétérotopie telle qu’elle est définie par Michel Foucault :}
            \begin{teiCit}[xmlid={cit02-3}]

              \begin{teiQuote}
Il y a également, et ceci probablement dans toute culture, dans toute civilisation, des lieux réels, des lieux effectifs, des lieux qui sont dessinés dans l’institution même de la société, et qui sont des sortes de contre-emplacements, sortes d’utopies effectivement réalisées dans lesquelles les emplacements réels, tous les autres emplacements réels que l’on peut trouver à l’intérieur de la culture sont à la fois représentés, contestés et inversés, des sortes de lieux qui sont hors de tous les lieux, bien que pourtant ils soient effectivement localisables. Ces lieux, parce qu’ils sont absolument autres que tous les emplacements qu’ils reflètent et dont ils parlent, je les appellerai, par opposition aux utopies, les hétérotopies\teiNote[xmlid={ftn59-3},n={8},place={foot}]{\teiP[xmlid={p-093}]{Michel Foucault, « Des espaces autres » \lbrack{}conf. 1966, publ. 1984\rbrack{}, dans Daniel Defert et François Ewald (éd.), \teiHi[rend={italic}]{Dits et écrits, 1954-1988}, Paris, Gallimard, « Quarto », 2017, t. II, p. 1574-1575.}}.
\end{teiQuote}
            
\end{teiCit}
            \teiP[xmlid={p4-3}]{La notion foucaldienne d’hétérotopie permet de penser la poétique spatiale qui participe de la représentation des sociétés secrètes dans la littérature romantique. Pour le montrer, nous nous appuierons sur le corpus dramatique et romanesque suivant, qui va des années 1830 aux années 1850 : \teiHi[rend={italic}]{Hernani} (1830), drame de Hugo ; \teiHi[rend={italic}]{Rome souterraine} (1833), roman de Charles Didier ; \teiHi[rend={italic}]{Léo Burckart} (1839), drame de Nerval et Dumas ; \teiHi[rend={italic}]{Le Capitaine Richard} (1854), roman de Dumas.}
            \begin{teiDiv}[type={section1},xmlid={div01-2}]

              \teiHead[xmlid={le-locus-secretus-une-construction-heterotopique-001}]{Le \teiHi[rend={italic}]{locus secretus}, une construction hétérotopique}
              \teiP[xmlid={p5-3}]{Les quatre œuvres qui nous intéressent mettent toutes en scène une conjuration, mais l’importance qu’elle occupe dans l’économie de l’œuvre varie d’un texte à l’autre. Dans \teiHi[rend={italic}]{Hernani}, la société secrète apparaît à l’acte IV : dès la scène 1, Don Carlos utilise le terme de « ligue\teiNote[xmlid={ftn60-3},n={9},place={foot}]{\teiP[xmlid={p-094}]{Victor Hugo, \teiHi[rend={italic}]{Hernani} \lbrack{}1830\rbrack{}, Anne Ubersfeld (éd.), dans Jacques Seebacher et Guy Rosa (dir.), \teiHi[rend={italic}]{Œuvres complètes. Théâtre I}, Paris, Robert Laffont, « Bouquins », 1985, p. 619.}} » pour qualifier le groupe formé pour empêcher son accession au trône impérial, puis Hugo, à la scène 3, désigne ses personnages comme des « conjurés », à l’exception de trois d’entre eux – Hernani, Don Ruy Gomez et le duc de Gotha. Ces conjurés se retrouvent tous ensemble dans le tombeau de Charlemagne à Aix-la-Chapelle ; or, comme le signale Foucault, le cimetière est une « curieuse hétérotopie », « un lieu autre par rapport aux espaces culturels ordinaires\teiNote[xmlid={ftn61-3},n={10},place={foot}]{\teiP[xmlid={p-095}]{Foucault, « Des espaces autres », art. cité, p. 1576.}} ». Bien sûr, il faut distinguer le cimetière, espace collectif, du tombeau, espace singulier réservé à un mort et à sa famille ; en apparence, le tombeau n’a pas d’enjeu collectif, mais Don Carlos n’affirme-t-il pas, à la scène précédente : « Le pape et l’empereur sont tout. Rien n’est sur terre / Que pour eux et par eux\teiNote[xmlid={ftn62-3},n={11},place={foot}]{\teiP[xmlid={p-096}]{Hugo, \teiHi[rend={italic}]{Hernani}, IV, 2, p. 626.}} » ? La nature omnipotente de l’empereur confère ainsi au tombeau de Charlemagne une dimension sacrée et universelle, qui est aussi celle du cimetière, « où chaque famille possède sa noire demeure\teiNote[xmlid={ftn63-3},n={12},place={foot}]{\teiP[xmlid={p-097}]{Foucault, « Des espaces autres », art. cité, p. 1577.}} ». Le lieu souterrain, que Hugo emprunte à la littérature gothique et au mélodrame, se retrouve dans \teiHi[rend={italic}]{Le Capitaine Richard}, où Dumas décrit une réunion de l’Union de Vertu, l’historique \teiHi[rend={italic}]{Tugendbund}\teiNote[xmlid={ftn64-2},n={13},place={foot}]{\teiP[xmlid={p-098}]{Fondé à Königsberg en 1808, le \teiHi[rend={italic}]{Tugendbund} (ou Ligue de Vertu) est une association antinapoléonienne de lutte pour la nation allemande. Interdit dès décembre 1809, il a joué un rôle essentiel dans la diffusion des idées nationales et libérales en Allemagne.}}, dont les membres se font appeler des \teiHi[rend={italic}]{Wissende}, c’est-à-dire, indique Dumas en note, « Qui savent, qui sont du secret\teiNote[xmlid={ftn65-2},n={14},place={foot}]{\teiP[xmlid={p-099}]{Alexandre Dumas, \teiHi[rend={italic}]{Le Capitaine Richard} \lbrack{}1854\rbrack{}, Paris, Dufour, Mulat et Boulanger, 1858, p. 27.}} ». La réunion a lieu dans les ruines du château d’Abensberg, non loin de Ratisbonne, et son but est de désigner le futur assassin de Napoléon et par la même occasion le libérateur de l’Allemagne. Il s’agit donc, là encore, de représenter « un héros qui se fait conspirateur pour contrer un pouvoir tyrannique\teiNote[xmlid={ftn66-2},n={15},place={foot}]{\teiP[xmlid={p-100}]{Jean-Noël Tardy, \teiHi[rend={italic}]{L’Âge des ombres}, \teiHi[rend={italic}]{op. cit.}, p. 330.}} ». Mais Dumas utilise un espace tout à fait particulier : lieu « hors de tous les lieux, bien que pourtant \lbrack{}il\rbrack{} soi\lbrack{}t\rbrack{} effectivement localisabl\lbrack{}e\rbrack{} », la ruine n’est pas, à proprement parler, un lieu mais une trace de lieu, un indice de lieu passé, d’autant plus que cette scène se déroule précisément dans la « salle souterraine\teiNote[xmlid={ftn67-2},n={16},place={foot}]{\teiP[xmlid={p-101}]{Dumas, \teiHi[rend={italic}]{Le Capitaine Richard}, \teiHi[rend={italic}]{op. cit.}, p. 30.}} » d’un château en ruines. Nerval et Dumas choisissent le même cadre, une « salle en ruines » du « château de Wurtzbourg\teiNote[xmlid={ftn68},n={17},place={foot}]{\teiP[xmlid={p-102}]{Gérard de Nerval, \teiHi[rend={italic}]{Léo Burckart (1838-1839)}, Sylvie Lécuyer (éd.), Paris, Honoré Champion, « Textes de littérature moderne et contemporaine », 2018, p. 457.}} », pour accueillir la réunion de la société des Voyants, qui conspirent pour renverser le pouvoir établi et instaurer l’unité de l’Allemagne. Il en va différemment de Charles Didier, qui oppose deux sociétés secrètes : celle des carbonari, la Charbonnerie, et celle des sanfédistes. Didier représente, comme Dumas dans \teiHi[rend={italic}]{Le Capitaine Richard}, deux sociétés secrètes ayant réellement existé : la Charbonnerie, célèbre société d’obédience libérale et républicaine\teiNote[xmlid={ftn69},n={18},place={foot}]{\teiP[xmlid={p-103}]{Voir le chapitre que lui consacre Jean-Noël Tardy, dans \teiHi[rend={italic}]{L’Âge des ombres}, \teiHi[rend={italic}]{op. cit.}, p. 47-120.}}, et la \teiHi[rend={italic}]{Santa Fede}, société ultra catholique et anti-républicaine dont les adversaires sont les carbonari et l’empire autrichien. Contrairement aux autres auteurs, Didier ne fait pas le choix du lieu souterrain, mais il est lui aussi marqué par l’héritage gothique et mélodramatique quand il choisit la vieille tour d’Asture pour réunir ses carbonari ; Didier propose ainsi une variation dans l’exploitation de la verticalité symbolique du lieu de réunion. En revanche, le consistoire des sanfédistes tient séance « au couvent de Saint-François-d’Assise, dans le parloir du cardinal de Pétralie\teiNote[xmlid={ftn70},n={19},place={foot}]{\teiP[xmlid={p-104}]{Charles Didier, \teiHi[rend={italic}]{Rome souterraine} \lbrack{}1833\rbrack{}, Paris, Gosselin, 1841, p. 81.}} ». Ces deux lieux ne sont pas \teiHi[rend={italic}]{a priori} des hétérotopies mais plutôt des espaces d’opportunité ; néanmoins, les conditions d’accès à ces lieux participent de la construction des hétérotopies.}
              \teiP[xmlid={p6-3}]{Parmi les critères qui caractérisent les hétérotopies, Foucault distingue « un système d’ouverture et de fermeture qui, à la fois, les isole et les rend pénétrables. En général, on n’accède pas à un emplacement hétérotopique comme dans un moulin. \lbrack{}…\rbrack{} On ne peut y entrer qu’avec une certaine permission et une fois qu’on a accompli un certain nombre de gestes\teiNote[xmlid={ftn71},n={20},place={foot}]{\teiP[xmlid={p-105}]{Foucault, « Des espaces autres », art. cité, p. 1579.}}. » La description de ces gestes, qui peuvent être des paroles, est un passage obligé dans la représentation des sociétés secrètes, car elle inscrit textuellement le passage de l’extérieur vers l’intérieur, elle manifeste l’entrée dans le monde du secret. Chez Hugo, la reconnaissance se fait par un serrage de main et un échange de paroles codées :}
              \begin{teiCit}[xmlid={cit03-3}]

                \begin{teiQuote}
PREMIER CONJURÉ, \teiHi[rend={italic}]{portant seul une torche allumée}.\teiLb{}\teiHi[rend={italic}]{Ad augusta}.\teiLb{}DEUXIÈME CONJURÉ\teiLb{}\teiHi[rend={italic}]{Per angusta}.\teiLb{}PREMIER CONJURÉ\teiLb{}Les saints\teiLb{}\teiLb{}Nous protègent.\teiLb{}TROISIÈME CONJURÉ\teiLb{}Les morts nous servent.\teiLb{}PREMIER CONJURÉ\teiLb{}Dieu nous garde\teiNote[xmlid={ftn72},n={21},place={foot}]{\teiP[xmlid={p-106}]{Hugo, \teiHi[rend={italic}]{Hernani}, IV, 3, p. 630.}}.
\end{teiQuote}
              
\end{teiCit}
              \teiP[xmlid={p7-3}]{Les conjurés suivants, eux, pénètrent dans le tombeau en prononçant uniquement « \teiHi[rend={italic}]{Ad augusta} » et en « échange\lbrack{}ant\rbrack{} des signes de mains avec les autres\teiNote[xmlid={ftn73},n={22},place={foot}]{\teiP[xmlid={p-107}]{\teiHi[rend={italic}]{Ibid}., p. 631.}} ». Hugo met en scène une représentation topique qu’il reprend pour répondre aux attentes du spectateur ; il en va de même des costumes des personnages, « cachés sous leurs manteaux et leurs chapeaux\teiNote[xmlid={ftn74},n={23},place={foot}]{\teiP[xmlid={p-108}]{\teiHi[rend={italic}]{Ibid}., p. 630.}} », ainsi que de l’éclairage de la scène aux flambeaux. Hugo reprend un imaginaire gothique et mélodramatique qu’il adapte à la conjuration contre Don Carlos, qui est pour lui le cœur de la scène : le complot politique est plus important que la recréation fidèle d’une réunion de membres d’une société secrète. Tout au contraire, Nerval et Dumas apportent un soin de détail à la description et à la mise en scène de la réunion ; dans \teiHi[rend={italic}]{Le Capitaine Richard}, Dumas reprend presque mot pour mot une partie de la scène 4 de l’acte IV de \teiHi[rend={italic}]{Léo Burckart}, et l’affiliation de Waldeck à la société des Voyants devient chez Dumas l’affiliation d’un néophyte au \teiHi[rend={italic}]{Tugendbund} :}
              \begin{teiCit}[xmlid={cit04-2}]

                \begin{teiQuote}
— Frère, quelle heure est-il ? — L’heure où le maître veille et où l’esclave dort, répondit le récipiendaire. — Comptez-la. — Je ne l’entends plus depuis qu’elle sonne pour le maître. — Quand l’entendrez-vous ? — Quand elle aura réveillé l’esclave. — Où est le maître ? — À table. — Où est l’esclave ? — À terre. — Que boit le maître ? — Du sang. — Que boit l’esclave ? — Ses larmes. — Que voulez-vous faire de tous les deux ? — Je veux asseoir l’esclave à table, et coucher le maître à terre. — Êtes-vous maître, ou êtes-vous esclave ? — Ni l’un ni l’autre. — Qu’êtes-vous donc ? — Je ne suis rien encore ; mais j’aspire à devenir quelque chose. — Quoi ? — Voyant. — En savez-vous les fonctions ? — Je les apprends. — Qui vous les enseigne ? — Dieu. — Avez-vous des armes ? — J’ai cette corde et ce poignard. — Qu’est-ce que cette corde ? — Le symbole de notre force et de notre union. — Qu’êtes-vous selon ce symbole ? — Je suis un des fils de ce chanvre, que l’union a rapprochés, et que la force a tordus. — Pourquoi avez-vous pris cette corde ? — Pour lier et pour étreindre. — Pourquoi ce poignard ? — Pour couper et pour désunir. — Êtes-vous prêt à jurer que vous ferez usage de cette corde et de ce poignard contre tout condamné dont le nom sera inscrit au livre de sang ? — Oui. — Jurez-le. — Je le jure. — Vous dévouez-vous à la corde et au poignard vous-même, s’il vous arrivait de trahir le serment que vous venez de faire sur le glaive et sur la croix ? — Je m’y dévoue\teiNote[xmlid={ftn75},n={24},place={foot}]{\teiP[xmlid={p-109}]{Dumas, \teiHi[rend={italic}]{Le Capitaine Richard}, \teiHi[rend={italic}]{op. cit.}, p. 32. Voir Nerval, \teiHi[rend={italic}]{Léo Burckart}, \teiHi[rend={italic}]{op. cit.}, p. 465-468.}} !
\end{teiQuote}
              
\end{teiCit}
              \teiP[xmlid={p8-3}]{Ce passage rend compte, comme Dumas l’indique en note, de « la formule exacte de l’affiliation\teiNote[xmlid={ftn76},n={25},place={foot}]{\teiP[xmlid={p-110}]{Dumas, \teiHi[rend={italic}]{ibid}.}} » telle qu’elle a été donnée par Nerval après son étude des sociétés secrètes. Les deux auteurs, contrairement à Hugo, prétendent donner une image fidèlement historique de l’ouverture de la société secrète à ceux qui ont « accompli un certain nombre de gestes ». L’entrée spatiale dans le lieu ne peut ainsi se faire que par l’entrée symbolique dans la société secrète. Le rituel d’initiation et de reconnaissance est un seuil obligatoire sans le franchissement duquel il est impossible de participer à la conjuration. Il en est de même chez les carbonari et les sanfédistes de Didier : des étrangers sont autorisés à entrer dans la tour des premiers au moyen d’un « signe mystérieux », un « attouchement carbonique\teiNote[xmlid={ftn77},n={26},place={foot}]{\teiP[xmlid={p-111}]{Charles Didier, \teiHi[rend={italic}]{Rome souterraine}, \teiHi[rend={italic}]{op. cit.}, p. 39-40.}} » qui accompagne le serrage de main, quand les sanfédistes se reconnaissent par la possession d’une médaille :}
              \begin{teiCit}[xmlid={cit05-2}]

                \begin{teiQuote}
Comme le carbonarisme, le consistoire avait ses statuts, ses grades, ses emblèmes. Chaque initié recevait en pénétrant dans le sanctuaire une médaille d’airain, dont le symbole mystique était destiné à lui rappeler sans cesse le but de l’association : une madone, soutenue par un groupe d’anges, tend d’une main un faisceau de palmes et brandit de l’autre le glaive dont elle vient de frapper un esprit de ténèbres, mort à ses pieds. L’allégorie est diaphane : la madone est l’Italie ; les anges, les sanfédistes ; l’esprit des ténèbres, l’Autriche\teiNote[xmlid={ftn78},n={27},place={foot}]{\teiP[xmlid={p-112}]{\teiHi[rend={italic}]{Ibid}., p. 86.}}.
\end{teiQuote}
              
\end{teiCit}
              \teiP[xmlid={p9-3}]{La nature hétérotopique des lieux de la conjuration politique réside ainsi en partie dans le « système d’ouverture et de fermeture » qui les caractérise : pour participer à la conjuration, il faut d’abord pénétrer dans le lieu de la société secrète qui ourdit le complot. Cette entrée se fait au moyen de serments, de formules et de gestes cryptés qui sont autant de clés pour franchir le seuil du \teiHi[rend={italic}]{locus secretus}.}
              \teiP[xmlid={p10-3}]{Les lieux de la conjuration répondent également à un autre principe de l’hétérotopie, qui est sa nature hétérochronique. Selon Foucault, « les hétérotopies sont liées, le plus souvent, à des découpages du temps, c’est-à-dire qu’elles ouvrent sur ce qu’on pourrait appeler, par pure symétrie, des hétérochronies ; l’hétérotopie se met à fonctionner à plein lorsque les hommes se trouvent dans une sorte de rupture absolue avec leur temps traditionnel\teiNote[xmlid={ftn79},n={28},place={foot}]{\teiP[xmlid={p-113}]{Foucault, « Des espaces autres », art. cité, p. 1578.}} ». Le tombeau de Charlemagne dans \teiHi[rend={italic}]{Hernani} unit le passé, le présent et le futur, il fait le lien entre le temps humain et la divine éternité : comme le pape, l’empereur abrite en son sein un « suprême mystère\teiNote[xmlid={ftn80},n={29},place={foot}]{\teiP[xmlid={p-114}]{Hugo, \teiHi[rend={italic}]{Hernani}, IV, 2, p. 626.}} » qui le relie à l’histoire humaine (l’époque de Pierre et César), donc à la chronologie historique, et à la sphère divine, anhistorique. Il en va de même des ruines de \teiHi[rend={italic}]{Léo Burckart} et du \teiHi[rend={italic}]{Capitaine Richard}. La salle qui accueille les Voyants et l’Union de Vertu est « l’ancienne salle du tribunal secret\teiNote[xmlid={ftn81},n={30},place={foot}]{\teiP[xmlid={p-115}]{Dumas, \teiHi[rend={italic}]{Le Capitaine Richard}, \teiHi[rend={italic}]{op. cit.}, p. 34. On lit, dans \teiHi[rend={italic}]{Léo Burckart} : « Mais n’est-ce pas ici même que se tenaient les séances du tribunal secret ? » (Nerval, \teiHi[rend={italic}]{op. cit.}, p. 458).}} », c’est-à-dire de la Sainte-Vehme, à la fois cour de justice et société secrète apparue en Westphalie au début du \teiHi[rend={small-caps}]{xiii}\teiHi[rend={sup}]{e} siècle\teiNote[xmlid={ftn82},n={31},place={foot}]{\teiP[xmlid={p-116}]{Sur la Sainte-Vehme, voir notre article « Léo Burckart et l’imaginaire du tribunal secret à l’âge romantique », à paraître dans les \teiHi[rend={italic}]{Cahiers Alexandre Dumas}.}}. Inscrit dans l’histoire, ce lieu marque néanmoins une rupture avec le temps humain car il se situe dans un château en ruines ; comme le rappelle Claude Millet à propos de la « poétique des morts » qu’est le romantisme français, « les vestiges, ces signes de la contingence morbide de l’histoire des hommes \lbrack{}…\rbrack{}, les vestiges laissent entrevoir l’éternité, cette vie infinie qui se confond avec la mort\teiNote[xmlid={ftn83},n={32},place={foot}]{\teiP[xmlid={p-117}]{Claude Millet, \teiHi[rend={italic}]{Le Romantisme. Du bouleversement des lettres dans la France postrévolutionnaire}, Paris, LGF, « Références », 2007, p. 138.}} ». Nerval et Dumas proposent un lieu où le temps semble suspendu, où les strates temporelles se superposent jusqu’à se confondre avec l’éternité. On observe, dans \teiHi[rend={italic}]{Rome souterraine}, deux autres configurations, qui s’expliquent par l’opposition idéologique qui structure le roman de Didier, républicain et révolutionnaire. La tour d’Asture est d’emblée replacée dans la géographie historique et mythologique de Rome : alors qu’Anselme et Marius se trouvent non loin du vallon où eut lieu le viol de Lucrèce qui donna naissance à la république romaine, le premier s’exclame : « Dieu veuille \lbrack{}…\rbrack{} que la tour d’Asture soit pour nos descendants ce qu’est pour nous ce vallon sacré\teiNote[xmlid={ftn84},n={33},place={foot}]{\teiP[xmlid={p-118}]{Charles Didier, \teiHi[rend={italic}]{Rome souterraine}, \teiHi[rend={italic}]{op. cit.}, p. 9.}} ! » Le lieu de la conjuration se caractérise non pas par une « rupture absolue », mais par la répétition virtuelle d’un moment qui mêle histoire et mythologie. En revanche, le parloir du cardinal de Pétralie n’accueille les sanfédistes que parce que le pape, véritable chef de la \teiHi[rend={italic}]{Santa Fede}, ne peut organiser les séances au Vatican : « \teiHi[rend={italic}]{Alter ego} du souverain pontife, qui décemment ne pouvait paraître, le cardinal de Pétralie l’y représentait, et partant y occupait le premier rang\teiNote[xmlid={ftn85},n={34},place={foot}]{\teiP[xmlid={p-119}]{\teiHi[rend={italic}]{Ibid}., p. 86.}}. » La construction des lieux de la conjuration dans \teiHi[rend={italic}]{Rome souterraine} participe ainsi de l’imaginaire politique qui innerve le roman.}
              \teiP[xmlid={p11-3}]{Le \teiHi[rend={italic}]{locus secretus} est donc, dans nos œuvres, un \teiHi[rend={italic}]{locus politicus} dont les caractéristiques en font une hétérotopie. Et comme toute hétérotopie, il a, « par rapport à l’espace restant, une fonction\teiNote[xmlid={ftn86},n={35},place={foot}]{\teiP[xmlid={p-120}]{Foucault, « Des espaces autres », art. cité, p. 1580.}} », essentiellement politique.}
            
\end{teiDiv}
            \begin{teiDiv}[type={section1},xmlid={div02-2}]

              \teiHead[xmlid={fonction-politique-du-locus-secretus-une-heterotopie-de-compensation-001}]{Fonction politique du \teiHi[rend={italic}]{locus secretus }: une hétérotopie de compensation}
              \teiP[xmlid={p12-3}]{Dans son essai sur les complots, Jean-Noël Tardy rappelle en introduction que « la conspiration est fondamentalement une construction à laquelle participent de nombreux acteurs aux objectifs différents » mais qui tous « adhèrent à une idée : le pouvoir en place n’est pas légitime\teiNote[xmlid={ftn87},n={36},place={foot}]{\teiP[xmlid={p-121}]{Jean-Noël Tardy, \teiHi[rend={italic}]{L’Âge des ombres}, \teiHi[rend={italic}]{op. cit.}, p. 18-19.}} ». Les lieux hétérotopiques permettent en effet de réunir des hommes qui, dans le \teiHi[rend={italic}]{locus publicus}, n’auraient pas pu se rencontrer ou se réunir. Ainsi, lorsque Don Ricardo donne l’identité des conjurés à Don Carlos, tous ont une raison différente de comploter contre lui : le duc de Gotha « veut un Allemand d’Allemagne à l’empire », Hohenbourg préférerait « l’enfer avec François \lbrack{}I\teiHi[rend={sup}]{er}\rbrack{} que le ciel avec \lbrack{}lui\rbrack{}\teiNote[xmlid={ftn88},n={37},place={foot}]{\teiP[xmlid={p-122}]{Hugo, \teiHi[rend={italic}]{Hernani}, IV, 2, p. 620.}} », Giron veut se venger de l’amant de sa femme, etc. Comme la majorité des conjurations politiques, celle qui vise Don Carlos réunit des hommes aux motivations diverses (volonté politique, détestation personnelle, vengeance domestique, etc.), mais qui considèrent nécessaire d’empêcher Don Carlos de devenir empereur : le projet politique commun permet d’unir, dans cette hétérotopie qu’est le tombeau de Charlemagne, une société secrète indissociable de son lieu de réunion. Les ruines d’Abensberg jouent le même rôle que le tombeau dans \teiHi[rend={italic}]{Hernani }: les Voyants s’opposent sur de nombreux points, qu’il s’agisse de leurs origines sociales (populaires, bourgeoises, nobles), de la classe dirigeante qu’ils souhaitent mettre au pouvoir (« C’est une querelle à vider plus tard entre vainqueurs\teiNote[xmlid={ftn89},n={38},place={foot}]{\teiP[xmlid={p-123}]{Nerval, \teiHi[rend={italic}]{Léo Burckart}, \teiHi[rend={italic}]{op. cit.}, p. 474.}} », déclare un voyant), ou bien encore de ceux qu’ils représentent. L’homme masqué, au début de la séance, rappelle que « de même qu’au congrès, chaque ministre représente un roi, de même ici chacun de nous représente un peuple\teiNote[xmlid={ftn90},n={39},place={foot}]{\teiP[xmlid={p-124}]{\teiHi[rend={italic}]{Ibid}., p. 463.}} ». Dans \teiHi[rend={italic}]{Léo Burckart} comme dans \teiHi[rend={italic}]{Le Capitaine Richard}, les conjurés sont désignés par le nom d’un royaume ou d’une région :}
              \begin{teiCit}[xmlid={cit06-2}]

                \begin{teiQuote}
— Frères, continua le même orateur, n’oublions pas que, de même qu’au congrès chaque ministre représente un roi, de même, ici, chacun de nous représente un peuple. Veilleur, appelez les noms.\teiLb{}Le veilleur prononça l’un après l’autre les noms suivants :\teiLb{}— Baden, Nassau, Hesse, Wurtemberg, Westphalie, Autriche, Italie, Hongrie, Bohême, Espagne, Tyrol, Saxe, Luxembourg, Hanovre, Holstein, Mecklembourg, Bavière\teiNote[xmlid={ftn91},n={40},place={foot}]{\teiP[xmlid={p-125}]{Dumas, \teiHi[rend={italic}]{Le Capitaine Richard}, \teiHi[rend={italic}]{op. cit.}, p. 31.}}.
\end{teiQuote}
              
\end{teiCit}
              \teiP[xmlid={p13-3}]{L’ampleur géographique de cette liste symbolise l’hétérogénéité des membres présents, comme chez Nerval et Hugo. Les carbonari de Didier sont eux aussi d’abord présentés dans leur diversité régionale :}
              \begin{teiCit}[xmlid={cit07}]

                \begin{teiQuote}
— Je récuse Septime, répliqua Grimaldi ; sa neutralité m’est suspecte : il ne m’est contraire que parce qu’il est Piémontais et moi Génois.\teiLb{}— Il faut pourtant bien, dit Côme, le Toscan, un arbitre entre un Génois et un Vénitien ; les plaies du Bosphore et de Chiozza saignent encore.\teiLb{}— En qualité de Lombard, s’écria Cavalcabo, je suis neutre. Fils de Saint-Marc et de Gênes, je vous cite à ma barre, venez plaider votre cause\teiNote[xmlid={ftn92},n={41},place={foot}]{\teiP[xmlid={p-126}]{Charles Didier, \teiHi[rend={italic}]{Rome souterraine}, \teiHi[rend={italic}]{op. cit.}, p. 14. Et citation suivante.}}.
\end{teiQuote}
              
\end{teiCit}
              \teiP[xmlid={p14-2}]{Didier met ainsi en évidence les inimitiés entre les différents peuples italiens, dont les affrontements passés motivent les rivalités et les méfiances contemporaines. Le jeune Azzo intervient alors pour mettre fin à l’altercation entre ses compagnons : « Rien, même en jouant, ne doit réveiller ici les rivalités de l’Italie ; tout, au contraire, doit nous rappeler que nous sommes les enfants d’une mère commune, et que nous ne formons qu’une seule famille. » L’emploi du déictique \teiHi[rend={italic}]{ici} manifeste textuellement le rôle-clé que joue le lieu hétérotopique dans la conjuration : alors que le \teiHi[rend={italic}]{locus publicus} est un lieu de division, le \teiHi[rend={italic}]{locus secretus} est un vecteur d’union. De la même manière, dans \teiHi[rend={italic}]{Hernani}, à la question « Qui frappera ? », tous répondent : « Nous tous\teiNote[xmlid={ftn93},n={42},place={foot}]{\teiP[xmlid={p-127}]{Hugo, \teiHi[rend={italic}]{Hernani}, IV, 3, p. 633.}} ! » Il est frappant de constater que le lieu de réunion des sanfédistes, à propos duquel nous avons vu qu’il n’est pas tout à fait une hétérotopie, ne permet pas de véritable union d’éléments par trop hétérogènes :}
              \begin{teiCit}[xmlid={cit08}]

                \begin{teiQuote}
Ces ennemis \lbrack{}des carbonari\rbrack{}, quels sont-ils ? L’autel qui les excommunie, le trône qui les décime, Rome et César ont fait alliance contre eux. Pour mieux trancher les mille têtes de l’hydre populaire qui les enlace dans l’ombre de ses immenses replis, ils ont pactisé, ils ont sacrifié sur l’autel de la peur, Rome ses prétentions guelfes, César ses inimitiés gibelines. Ce sacrifice forcé n’est qu’apparent ; les vieux levains couvent, fermentent au fond des âmes. On veut bien des deux côtés la destruction de l’ennemi commun, mais c’est à condition d’avoir pour soi, et pour soi tout seul, sa dépouille et les fruits de sa victoire. L’orgueil d’aucun des deux ne veut plier ; leur ambition s’irrite à l’idée seule d’un partage. Mais une guerre d’existence les rallie ; on dissimule, on combat ensemble, ensemble on les tue… après la victoire, on verra. De là de nouvelles complications, de nouveaux mystères ; car si l’on s’embrasse au grand jour, on se déchire dans les ténèbres\teiNote[xmlid={ftn94},n={43},place={foot}]{\teiP[xmlid={p-128}]{Charles Didier, \teiHi[rend={italic}]{Rome souterraine}, \teiHi[rend={italic}]{op. cit.}, p. 84.}}.
\end{teiQuote}
              
\end{teiCit}
              \teiP[xmlid={p15}]{Contrairement aux carbonari, les sanfédistes sont incapables de faire taire leurs volontés et ambitions personnelles, qui s’expriment notamment dans la répartition géographique des futurs royaumes d’Italie : « C’étaient quatre Italies en une\teiNote[xmlid={ftn95},n={44},place={foot}]{\teiP[xmlid={p-129}]{\teiHi[rend={italic}]{Ibid}., p. 94.}}. » Le parloir du cardinal de Pétralie ne doit donc pas être considéré comme une hétérotopie, mais comme le lieu d’une conjuration qui repose sur un morcellement fondamental de ses membres et de ses buts. Ce \teiHi[rend={italic}]{locus politicus} de la \teiHi[rend={italic}]{Realpolitik} se distingue profondément des autres \teiHi[rend={italic}]{loci secreti} de nos œuvres, qui sont, eux, des hétérotopies de compensation.}
              \teiP[xmlid={p16}]{Foucault distingue deux « pôles extrêmes » entre lesquels se déploie la fonction de l’hétérotopie. Il existe, selon lui, l’hétérotopie d’illusion, « qui dénonce comme plus illusoire encore tout l’espace réel », et l’hétérotopie de compensation, qui crée « un autre espace, un autre espace réel, aussi parfait, aussi méticuleux, aussi bien arrangé que le nôtre est désordonné, mal agencé et brouillon\teiNote[xmlid={ftn96},n={45},place={foot}]{\teiP[xmlid={p-130}]{Foucault, « Des espaces autres », art. cité, p. 1580.}} ». Dans les œuvres qui nous intéressent ici, les lieux de la conjuration se présentent comme une version améliorée du \teiHi[rend={italic}]{locus publicus}. Alors que Léo Burckart intervient pour défendre le trône, l’accusateur lui rappelle qu’il n’est plus dans le monde tyrannique : « Frère, tu oublies que nous sommes ici au-dessus des fictions politiques et légales. Les princes de la terre ne sont pas nos princes, à nous\teiNote[xmlid={ftn97},n={46},place={foot}]{\teiP[xmlid={p-131}]{Nerval, \teiHi[rend={italic}]{Léo Burckart}, \teiHi[rend={italic}]{op. cit.}, p. 474.}} ! » Il s’agit ainsi, si l’on repense à l’analyse de Jean-Noël Tardy, de substituer à la fiction du \teiHi[rend={italic}]{locus publicus} la fiction compensatrice du \teiHi[rend={italic}]{locus secretus}, qui fonctionne comme « l’envers de l’endroit\teiNote[xmlid={ftn98},n={47},place={foot}]{\teiP[xmlid={p-132}]{Alexandra Delcamp, « Les sociétés secrètes dans La Comédie humaine, “l’endroit” de l’envers ? », \teiHi[rend={italic}]{L’Imaginaire des sociétés secrètes dans la littérature du }\teiHi[rend={ small-caps italic}]{xix}\teiHi[rend={italic sup}]{e}\teiHi[rend={italic}]{ siècle}, \teiHi[rend={italic}]{op. cit.}, \teiRef[target={https://serd.hypotheses.org/files/2021/05/6.-Alexandra-Delcamp.pdf}]{\teiHi[rend={underline}]{https://serd.hypotheses.org/files/2021/05/6.-Alexandra-Delcamp.pdf}} (consulté le 28 août 2023).}} ». Parmi les fictions politiques présentes dans les œuvres, il en est une qui occupe une place essentielle et que l’hétérotopie réalise : l’unité nationale. Chez Didier, Nerval et Dumas, l’hétérotopie a pour fonction de réaliser dans sa dimension minimale l’unité rêvée pour la nation. Après qu’Azzo a mis fin à la querelle de ses camarades, les carbonari s’exclament : « Romains et Lombards, Vénitiens ou Génois, nous sommes tous frères. Embrassons-nous donc, et vive la grande République Ausonienne ! » Et le narrateur de commenter : « Il y eut à ces mots une effusion générale, et l’on sentit que le sang italien coulait dans toutes ces veines, versait dans tous ces cœurs l’amour sacré du pays natal\teiNote[xmlid={ftn99},n={48},place={foot}]{\teiP[xmlid={p-133}]{Charles Didier, \teiHi[rend={italic}]{Rome souterraine}, \teiHi[rend={italic}]{op. cit.}, p. 14.}}. » Les membres de l’Union de Vertu se battent également pour l’unité de leur nation, que Napoléon a juré de détruire :}
              \begin{teiCit}[xmlid={cit09}]

                \begin{teiQuote}
J’accuse l’empereur Napoléon de tenter le plus grand crime qui existe aux yeux d’un Allemand, c’est-à-dire de vouloir détruire la nationalité de l’Allemagne. C’est pour détruire la nationalité de l’Allemagne qu’il a nommé son beau-frère Murat grand-duc de Berg ; c’est pour détruire la nationalité de l’Allemagne qu’il a nommé son frère Jérôme roi de Westphalie ; c’est pour détruire la nationalité de l’Allemagne qu’il veut détrôner l’empereur François II et mettre à sa place son frère Joseph, dont ne veulent pas les Espagnols ; enfin, c’est pour détruire la nationalité de l’Allemagne qu’il fait battre aujourd’hui la Bavière contre l’Autriche, la confédération du Rhin contre l’Empire, amis contre amis, Allemands contre Allemands, frères contre frères\teiNote[xmlid={ftn100},n={49},place={foot}]{\teiP[xmlid={p-134}]{Dumas, \teiHi[rend={italic}]{Le Capitaine Richard}, \teiHi[rend={italic}]{op. cit.}, p. 37.}} !
\end{teiQuote}
              
\end{teiCit}
              \teiP[xmlid={p17}]{Ce passage rend compte de l’importance des sociétés secrètes dans l’élaboration du sentiment national allemand\teiNote[xmlid={ftn101},n={50},place={foot}]{\teiP[xmlid={p-135}]{Voir, en particulier, Abel Douay et Gérard Hertault, \teiHi[rend={italic}]{Franc-maçonnerie et sociétés secrètes contre Napoléon. Naissance de la nation allemande}, Paris, Nouveau Monde et Fondation Napoléon, 2005.}}, que l’accusateur de la société des Voyants rappelait déjà dans \teiHi[rend={italic}]{Léo Burckart }: « En 1806, les princes d’Allemagne vinrent à nous, et nous dirent : Peuples et noblesse, nous avons un maître qui nous pèse, venez en aide à notre puissance, et nous irons en aide à votre liberté\teiNote[xmlid={ftn102},n={51},place={foot}]{\teiP[xmlid={p-136}]{Nerval, \teiHi[rend={italic}]{Léo Burckart}, \teiHi[rend={italic}]{op. cit.}, p. 473.}}. » Si cette promesse de libération nationale n’a pas été tenue par les princes, elle a pris forme au sein même de cette hétérotopie qu’est la salle du château de Wurtzbourg où se réunissent symboliquement tous les peuples allemands. Comme l’écrit Laurent Colantonio à propos de la nation romantique, « le sentiment national repose moins sur l’adhésion rationnelle à des principes que sur l’appartenance à une communauté d’origine – qui prédispose l’individu dès sa naissance –, enracinée depuis le fond des âges sur un territoire où elle est destinée à se perpétuer\teiNote[xmlid={ftn103},n={52},place={foot}]{\teiP[xmlid={p-137}]{Laurent Colantonio, « Nation », Alain Vaillant (dir.), \teiHi[rend={italic}]{Dictionnaire du Romantisme}, Paris, CNRS Éditions, 2012, p. 499.}} ». Aux fictions politiques et légales du \teiHi[rend={italic}]{locus publicus}, l’hétérotopie du lieu de la conjuration substitue une fiction politique : celle d’une nation enfin réunie, libre et juste.}
            
\end{teiDiv}
          
\end{teiElement}
        
\end{teiElement}
        \begin{teiElement}[name={group},type={chapter},data-page-title={La politique dans le caveau},data-page-subtitle={Le chronotope gothique d’Une ténébreuse affaire},data-page-authors={Kathia Huynh@@Université d’Orléans, Polen, Orléans, France},data-include-href={../XML/chap\_4\_La\_politique\_dans\_le\_caveau\_Locus\_politicus.xml},xmlid={chapter-003}]

          \begin{teiElement}[name={front}]

            \begin{teiDiv}[type={titlePage},xmlid={titlepage-004}]

              \teiP[rend={title-main},xmlid={p-138}]{La politique dans le caveau}
              \teiP[rend={title-sub},xmlid={p-139}]{Le chronotope gothique d’\teiHi[rend={italic}]{Une ténébreuse affaire}}
              \teiP[rend={author-aut},xmlid={p-140}]{Kathia \teiHi[rend={small-caps}]{Huynh}}
              \teiP[rend={authority\_affiliation},xmlid={p-141}]{Université d’Orléans, Polen, Orléans, France}
            
\end{teiDiv}
          
\end{teiElement}
          \begin{teiElement}[name={body},xmlid={body-004}]

            \teiP[xmlid={p1-4}]{Condensé « d’obscurité et de complication » livrant le « tableau mouvant de politique et d’histoire\teiNote[xmlid={ftn1-4},n={1},place={foot}]{\teiP[xmlid={p-142}]{Alain, \teiHi[rend={italic}]{Balzac}, Paris, Gallimard, 1999, p. 38 et 37.}} » d’une période trouble, \teiHi[rend={italic}]{Une ténébreuse affaire} (1841) doit sa redoutable complexité à sa disparité narrative, à l’origine d’une difficulté de lecture analogue à la confusion éprouvée face à l’« opacité absolue de la politique\teiNote[xmlid={ftn35-3},n={2},place={foot}]{\teiP[xmlid={p-143}]{Jacques-David Ebguy, « \teiHi[rend={italic}]{Une ténébreuse affaire}, un roman politique ? Souveraineté, société et dissensus », Romuald Fonkoua, Pierre Hartmann et Éléonore Reverzy (dir.),\teiHi[rend={italic}]{ Les Fables du politique des Lumières à nos jours}, Strasbourg, Presses universitaires de Strasbourg, 2012, p. 86.}} » au \teiHi[rend={small-caps}]{xix}\teiHi[rend={sup}]{e }siècle. Ce roman historique retrace le déclin de l’Ancien Régime à travers cinq aristocrates séditieux (Laurence de Cinq-Cygne, ses cousins les jumeaux Simeuse et leurs alliés, les frères d’Hauteserre), de leur implication dans la conspiration de Cadoudal et Pichegru en 1803 jusqu’à leur disparition quasi-totale à l’épilogue, situé en 1833, dans le salon de Diane de Cadignan ; l’ultime apparition de Laurence, veuve et isolée, fait office de chant du cygne de cette vieille noblesse. Ce proto-roman policier\teiNote[xmlid={ftn36-3},n={3},place={foot}]{\teiP[xmlid={p-144}]{André Vanoncini, « Balzac et la ténébreuse naissance du roman policier », \teiHi[rend={italic}]{Romanische Studien}, n\teiHi[rend={sup}]{o} 3, 2016, republié dans \teiHi[rend={italic}]{Balzac, roman, histoire, philosophie}, Paris, Honoré Champion, 2019, p. 45-55.}} entend ensuite « peindre la police politique aux prises avec la vie privée\teiNote[xmlid={ftn37-3},n={4},place={foot}]{\teiP[xmlid={p-145}]{Honoré de Balzac, préface d’\teiHi[rend={italic}]{Une ténébreuse affaire }(1841), Pierre-Georges Castex (éd.), \teiHi[rend={italic}]{La Comédie humaine}, Paris, Gallimard, 1978, t. VIII, p. 492. Toutes les références à \teiHi[rend={italic}]{La Comédie humaine} renverront à cette édition. }} » \teiHi[rend={italic}]{via} l’espion Corentin, bras droit de Fouché, lesquels tirent dans l’ombre les ficelles d’une action dont les vrais ressorts restent longtemps cachés. Classé parmi les \teiHi[rend={italic}]{Scènes de la vie politique}, le récit relate enfin le procès intenté contre les nobles condamnés à tort pour l’enlèvement du sénateur Malin de Gondreville\teiNote[xmlid={ftn38-3},n={5},place={foot}]{\teiP[xmlid={p-146}]{L’épisode est inspiré de l’enlèvement réel du sénateur Clément de Ris. Voir Pierre Laubriet, « Autour d’\teiHi[rend={italic}]{Une ténébreuse affaire} », \teiHi[rend={italic}]{L’Année balzacienne}, 1968, p. 267-282.}}, éclairé \teiHi[rend={italic}]{a posteriori }lors de l’analepse finale assumée par Henri de Marsay, nommé Premier ministre. Là seulement est dévoilée l’implication du ministre de la Police et de son séide dans cette affaire qui aura servi à tirer profit de la destruction secrète de documents compromettants pour venger Corentin de l’humiliant coup de cravache donné par Laurence lors de l’enquête au début du roman.}
            \teiP[xmlid={p2-4}]{Cette opacité vient aussi de ce que chaque séquence – de la conspiration au complot, de l’enquête policière au procès politique – repose sur la démultiplication des secrets, liés dans le roman à la sphère du politique et du pouvoir\teiNote[xmlid={ftn39-3},n={6},place={foot}]{\teiP[xmlid={p-147}]{Pour ne citer qu’un exemple : « Cet homme au pâle visage \lbrack{}Fouché\rbrack{}, élevé dans les dissimulations monastiques, qui possédait les secrets des Montagnards auxquels il appartint, et ceux des royalistes auxquels il finit par appartenir, avait lentement et silencieusement étudié les hommes, les choses, les intérêts de la scène politique ; il pénétra les secrets de Bonaparte \lbrack{}…\rbrack{} » (VIII, p. 552-553).}}. Cette tendance est d’abord partagée par l’ensemble d’un personnel romanesque extrêmement hétéroclite, fracturé entre aristocrates épris d’un héroïsme aussi noble qu’anachronique, bourgeois opportunistes prêts à changer de principes comme de régime, et espions capables de toutes les infamies. Elle a en outre la particularité de s’incarner en un seul et même \teiHi[rend={italic}]{locus secretus }: un ancien caveau dissimulé sous les profondeurs de l’obscure forêt de Nodesme, utilisé d’abord par le clan aristocratique dans la première partie pour soustraire les conjurés à la surveillance de la police, par les sbires de Fouché ensuite, dans la troisième partie, pour séquestrer à l’insu de tous, en particulier des premiers propriétaires des lieux, le sénateur Malin de Gondreville.}
            \teiP[xmlid={p3-4}]{Au centre d’\teiHi[rend={italic}]{Une ténébreuse affaire}, point de fuite de deux complots symétriques, trône un \teiHi[rend={italic}]{locus gothicus}, qui vaut moins pour sa dimension onirique ou psychanalytique\teiNote[xmlid={ftn40-3},n={7},place={foot}]{\teiP[xmlid={p-148}]{Voir Annie Le Brun, \teiHi[rend={italic}]{Les Châteaux de la subversion}, Paris, Garnier frères, 1982.}} que pour sa valeur politique et historique\teiNote[xmlid={ftn41-3},n={8},place={foot}]{\teiP[xmlid={p-149}]{Cette préséance de l’historique et du politique est aussi présente dans le souterrain, autre chronotope gothique. Voir Thomas Conrad, « Les souterrains au \teiHi[rend={small-caps}]{xix}\teiHi[rend={sup}]{e} siècle : des images du temps », \teiHi[rend={italic}]{Communications}, n\teiHi[rend={sup}]{o} 105, 2019/2, p. 55-69 et « Souterrains politiques verniens. Variations sur un motif, de \teiHi[rend={italic}]{L’Île mystérieuse }à \teiHi[rend={italic}]{Face au drapeau} », \teiHi[rend={italic}]{Romanesques}, hors-série, 2022, p. 99-117.}} ; moins miroir de la psyché, donc, que chronotope faisant l’objet d’un traitement original. Lieu d’une « corrélation essentielle des rapports spatio-temporels\teiNote[xmlid={ftn42-3},n={9},place={foot}]{\teiP[xmlid={p-150}]{Mikhaïl Bakhtine, \teiHi[rend={italic}]{Esthétique et théorie du roman} \lbrack{}1978\rbrack{}, Dara Olivier (trad.), Paris, Gallimard, 1987, p. 237.}} », le chronotope figure le temps historique\teiNote[xmlid={ftn43-3},n={10},place={foot}]{\teiP[xmlid={p-151}]{\teiHi[rend={italic}]{Ibid. }: « Ici, le temps se condense, devient compact, visible pour l’art, tandis que l’espace s’intensifie, s’engouffre dans le mouvement du temps, du sujet, de l’Histoire. Les indices du temps se découvrent dans l’espace, celui-ci est perçu et mesuré d’après le temps. »}} tel qu’il est conçu et représenté à une époque donnée. Le \teiHi[rend={italic}]{devenir} historique, quant à lui, transparaît au fil de la démonstration de Bakhtine : c’est dans la reprise transséculaire d’un même chronotope (l’univers pastoral de l’idylle, détruit aux \teiHi[rend={small-caps}]{xviii}\teiHi[rend={sup}]{e} et \teiHi[rend={small-caps}]{xix}\teiHi[rend={sup}]{e} siècles, par exemple) que se dégagent, d’une comparaison d’œuvre en œuvre, les nouvelles relations au temps et à l’histoire, ainsi que le processus de mutation historique. La mesure du devenir historique dépend donc d’un dialogue intertextuel maintenu dans la durée d’une histoire littéraire pluriséculaire.}
            \teiP[xmlid={p4-4}]{Or \teiHi[rend={italic}]{Une ténébreuse affaire} resserre ce phénomène dans le cours du récit : un même chronotope, le caveau gothique, est récupéré dans le cadre de deux complots politiques distincts, dont les ressorts et les modalités témoignent de la concurrence de deux systèmes historiques coprésents puis consécutifs. Le remplacement de l’un par l’autre figure le mouvement d’un \teiHi[rend={small-caps}]{xix}\teiHi[rend={sup}]{e} siècle en devenir. D’un complot l’autre, le chronotope du roman gothique configure un passage de relais historique, représentatif d’une bascule entre ancien et nouvel ordre du monde, caractérisée par un glissement de la place et de la visée du secret dans le cadre de l’action politique.}
            \teiP[xmlid={p5-4}]{Après avoir exposé les traits de cet avatar du chronotope gothique, il s’agira de faire état de l’usage différencié du caveau dans le cadre des complots. Théâtre du politique, ou plutôt de conceptions divergentes du politique, le caveau devient aussi l’indicateur d’une dynamique d’ordre historique. On verra comment les manières d’occuper ou de posséder ce chronotope sont révélatrices d’une certaine gestion du passé, qui conditionne le caractère inopérant ou au contraire efficace de l’exercice du pouvoir au présent.}
            \begin{teiDiv}[type={section1},xmlid={div01-3}]

              \teiHead[xmlid={un-caveau-en-foret-espace-et-recit-du-secret-001}]{Un caveau en forêt : espace et récit du secret}
              \teiP[xmlid={p6-4}]{\teiHi[rend={italic}]{Topos} consacré par Ann Radcliffe (\teiHi[rend={italic}]{L’Italien ou le Confessionnal des pénitents noirs}), Lewis (\teiHi[rend={italic}]{Le Moine}) ou Maturin (\teiHi[rend={italic}]{Melmoth ou l’Homme errant}) dont Balzac a été lecteur et critique\teiNote[xmlid={ftn44-3},n={11},place={foot}]{\teiP[xmlid={p-152}]{Nombreux sont les récits de Balzac qui pastichent les romans noirs, des œuvres de jeunesse (\teiHi[rend={italic}]{Le Centenaire}) à \teiHi[rend={italic}]{La Comédie humaine} (\teiHi[rend={italic}]{Melmoth réconcilié}, \teiHi[rend={italic}]{La Muse du département}). Voir Maurice Bardèche, \teiHi[rend={italic}]{Balzac romancier} \lbrack{}1940\rbrack{}, Paris, Plon, 1943 ; Jacques-David Ebguy, « \teiHi[rend={italic}]{Melmoth réconcilié}, ou le gothique “repris” », Catriona Seth\teiHi[rend={small-caps}]{ (}dir.),\teiHi[rend={italic}]{ Imaginaires gothiques. Aux sources du roman noir français}, Paris, Desjonquères, 2010, p. 193-227 ; Christelle Girard, « Balzac et les romans noirs. L’innutrition gothique comme partage poétique », \teiHi[rend={italic}]{L’Année balzacienne}, 2021, p. 169-191.}}, le caveau fait trois apparitions dans les ex-chapitres VIII « Un coin de forêt », XI « Revanche de la police » et XVIII « Marthe compromise\teiNote[xmlid={ftn45-3},n={12},place={foot}]{\teiP[xmlid={p-153}]{Les chapitres, présents dans la publication en feuilletons dans \teiHi[rend={italic}]{Le Commerce}, ont été supprimés pour l’édition Furne. L’édition Pléiade ne conserve que trois grandes subdivisions : « Les chagrins de la police » ; « Revanche de Corentin » ; « Un procès politique sous l’Empire ».}} », où il présente les attributs du chronotope gothique\teiNote[xmlid={ftn46-3},n={13},place={foot}]{\teiP[xmlid={p-154}]{Pour Bakhtine, le roman gothique a pour chronotope le château, \teiHi[rend={italic}]{Esthétique et théorie du roman}, \teiHi[rend={italic}]{op. cit.}, p. 387. D’autres travaux ont montré que le genre se décline en réalité en un éventail de lieux – forêts, labyrinthes, cloîtres, souterrains, couvents, cachots, ruines, cimetières. Voir Maurice Lévy, \teiHi[rend={italic}]{Le Roman « gothique » anglais, 1764-1824}, Paris, Albin Michel, 1995 ; Simone Bernard\teiHi[rend={small-caps}]{-}Griffiths et Jean Sgard (dir.), \teiHi[rend={italic}]{Mélodrames et romans noirs (1750-1890)}, Toulouse, Presses universitaires du Mirail, 2000 ; Catriona Seth (dir.), \teiHi[rend={italic}]{Imaginaires gothiques. Aux sources du roman noir}, \teiHi[rend={italic}]{op. cit}.}} : la labyrinthique forêt de Nodesme, où se perd Laurence, aidée par le régisseur Michu pour assurer à ses cousins un abri contre la police, cache en son « centre » un « couvent » souterrain (VIII, 564), avec un « caveau voûté », « du genre de ceux qu’on nommait l’\teiHi[rend={italic}]{in-pace}, le cachot des couvents » (VIII, 565), représenté sous « la lueur trompeuse de la lune » (VIII, 562), à « minuit » (VIII, 569). Découvert par Corentin et Peyrade au chapitre \teiHi[rend={small-caps}]{xi}, le caveau fait l’objet d’un second investissement plus noir encore par les ennemis anonymes des Simeuse et des d’Hauteserre, qui y séquestrent Malin.}
              \teiP[xmlid={p7-4}]{Plus noir, car le caveau renoue dans « Marthe compromise » avec l’atmosphère propre à ce chronotope synonyme de claustration et de persécution, en passant du \teiHi[rend={italic}]{locus amoenus} au \teiHi[rend={italic}]{locus horribilis}. L’espace apparaît d’abord sous la « belle lueur d’argent » de la lune dans un « joli coin de forêt » organisé en de « fuyantes perspectives », et bercé du doux « cri des grenouilles » plongées dans une « eau tranquille et ignorée » (VIII, 566). Cette quiétude apollinienne, civilisée, en témoignent les métaphores du « berceau de cave » et du « salon pour causer » (VIII, 566 et 567) employées par Laurence et Michu, s’inverse du tout au tout quand les mystérieux ravisseurs prennent possession du lieu. Lorsque Marthe Michu, mandatée par un faux contrefaisant l’écriture de son mari pour aller alimenter Malin, se rend à son tour au caveau dans une nuit noire d’effroi et de danger, elle le découvre fermé « par un cadenas » (VIII, 651) pour empêcher le prisonnier de s’évader. Les anciennes « idées de liberté » (VIII, 566) et l’espoir de survie évoqués par le lieu sont étouffés sous les barreaux et les angoisses de mort ; l’aménité du refuge disparaît sous la barbarie primitive d’un cachot rendu à sa fonction originelle. Aussi le caveau gothique retrouve-t-il son inquiétante étrangeté et son inhérente violence à l’issue d’une substitution actantielle, on y reviendra.}
              \teiP[xmlid={p8-4}]{Quoique capital dans le « \lbrack{}roman\rbrack{} à cachettes\teiNote[xmlid={ftn47-3},n={14},place={foot}]{\teiP[xmlid={p-155}]{Alain, \teiHi[rend={italic}]{Balzac}, \teiHi[rend={italic}]{op. cit.}, p. 37.}} » qu’est \teiHi[rend={italic}]{Une ténébreuse affaire}, le chronotope fait l’objet de peu de représentations, à rebours de l’esthétique du roman gothique qui se complaît dans le pittoresque horrifique des caveaux et cachots en tous genres. L’espace caché dispose donc d’une présence narrative elle-même dérobée, corrélée à la nature et à la fonction du lieu. Le caveau tient en effet sa valeur de son existence gardée secrète\teiNote[xmlid={ftn48-3},n={15},place={foot}]{\teiP[xmlid={p-156}]{Secret omniprésent au chapitre VIII : le caveau est recouvert « d’épais buissons impénétrables que, depuis 1794, Michu s’était plu à épaissir \lbrack{}…\rbrack{} » (VIII, 564) « pour s’en s’approprier le secret et le rendre impénétrable » (VIII, 565).}}, qui garantit en retour le maintien et l’efficacité des secrets dont il est le réceptacle – c’est pourquoi Michu, une fois levé l’impératif du secret avec la radiation de ses maîtres de la liste des émigrés, ne s’alarme guère de savoir le caveau trouvé par ses ennemis : « la découverte de cette cachette était maintenant sans danger, puisque les gentilshommes redevenaient Français, et avaient recouvré leur liberté » (VIII, 600). Sans secret à abriter, le caveau devient caduc.}
              \teiP[xmlid={p9-4}]{La narration semi-elliptique autour du caveau d’\teiHi[rend={italic}]{Une ténébreuse affaire} se conforme ainsi à la dynamique du secret encodée dans le chronotope, dans une tension entre rétention de l’information à l’égard d’un tiers exclu, et dévoilement de cette rétention, sans l’intuition de laquelle il y aurait vide plutôt que secret. Après avoir initié le lecteur au secret du lieu, le récit jette aussitôt dans l’ombre de la non-représentation cette autre scène souterraine, pour se concentrer sur la face visible des événements, là où la connaissance du caveau est dérobée. Le jeu entre mention et omission narrative figure ainsi le décalage entre savoir à cacher et opacité par là secrétée, le récit collant à l’action telle qu’elle est perçue de l’autre côté du secret, depuis le point de vue partiel de ceux qui ignorent l’existence ou l’implication du caveau dans l’affaire en cours : Corentin et Peyrade d’abord, qui pressentent que le mot et le lieu de l’énigme leur échappent\teiNote[xmlid={ftn49-3},n={16},place={foot}]{\teiP[xmlid={p-157}]{« Vous avez, \teiHi[rend={italic}]{je ne sais comment}, endormi Violette, et vous, votre femme, votre petit gars, vous avez passé la nuit dehors pour avertir Mlle de Cinq-Cygne de notre arrivée et faire sauver ses cousins que vous avez amenés ici, \teiHi[rend={italic}]{je ne sais pas encore où} » (VIII, 594, nous soulignons.)}}, le clan aristocratique ensuite, frappé par le coup « inexplicable » (VIII, 645) de leurs « ennemis cachés dans l’ombre » (VIII, 671). Le rôle nodal du \teiHi[rend={italic}]{locus secretus} dans l’intrigue d’\teiHi[rend={italic}]{Une ténébreuse affaire} commande ainsi une \teiHi[rend={italic}]{fabula secreta} à son image.}
              \teiP[xmlid={p10-4}]{Puisqu’\teiHi[rend={italic}]{Une ténébreuse affaire} postule l’indissociabilité du secret et du pouvoir – ce qu’indique la courbe narrative des « Chagrins de la police » et d’« Un procès politique sous l’Empire » où la répartition des vainqueurs et des perdants recoupe un partage entre initiés et ignorants\teiNote[xmlid={ftn50-3},n={17},place={foot}]{\teiP[xmlid={p-158}]{« Enfin vous nous avez battus » (VIII, 594) concède Corentin, après avoir énuméré les points restés obscurs. Avant le verdict qui prononce la culpabilité de Michu et de ses maîtres, Grandville, l’avocat de la Défense, interpelle les jurés, qui commettent « une fatale erreur que rien ne dissipera », ayant été « le jouet d’une puissance inconnue et machiavélique » (VIII, 671).}} –, le chronotope prend une place centrale dans les complots politiques où celui-ci est plus ou moins ouvertement brigué. D’un complot l’autre, le caveau fait toutefois l’objet d’un usage stratégique très différent, lié à une divergence d’exercice du pouvoir. Dans le passage de relais des nobles aux espions, le caveau enregistre dès lors un glissement d’ordre politico-historique : au \teiHi[rend={small-caps}]{xix}\teiHi[rend={sup}]{e} siècle, l’action et la pensée politiques ne sauraient se concevoir sur le modèle de ce qui, à présent, relève d’un passé révolu.}
            
\end{teiDiv}
            \begin{teiDiv}[type={section1},xmlid={div02-3}]

              \teiHead[xmlid={le-theatre-des-operations-de-la-politique-du-coup-declat-a-la-politique-du-secret-001}]{Le théâtre des opérations : de la politique du coup d’éclat à la politique du secret}
              \teiP[xmlid={p11-4}]{Entre la conspiration de 1803 et l’enlèvement de 1806, le caveau est occupé par des personnages qui en font un usage révélateur de leur conception de l’action politique. Le caveau se situe d’abord dans les marges du complot de Cadoudal et Pichegru, dont la répression par le gouvernement consulaire contraint les conspirateurs revenus sur le territoire au repli tactique pour échapper à la police. Sas d’attente passive où se réfugient ceux qui n’ont d’autre choix que de renoncer à l’assaut armé, le caveau se fait antichambre de la vie politique, qui suit son cours pendant les sept mois passés sous terre, marqués par les purges royalistes puis l’avènement de l’Empire\teiNote[xmlid={ftn51-3},n={18},place={foot}]{\teiP[xmlid={p-159}]{« Ces précautions durèrent autant que le procès Rivière, Polignac et Moreau. Quand le sénatus-consulte qui appelait à l’Empire la famille Bonaparte et nommait Napoléon Empereur fut soumis à l’acceptation du peuple français, M. d’Hauteserre signa sur le registre que vint lui présenter Goulard. Enfin on apprit que le pape viendrait sacrer Napoléon » (VIII, 596).}}. L’inaction à laquelle les Simeuse et les d’Hauteserre sont réduits au caveau indique que c’est sur la scène visible que se joue pour eux l’action politique, envisagée comme coup d’éclat, dont le secret est la condition, non la nature profonde. Laurence et ses alliés veulent « détruire Bonaparte et ramener les Bourbons » (VIII, 541) suivant un modèle actif structuré autour de l’attentat, de l’agression : « \lbrack{}Laurence\rbrack{} le visait, du fond de sa vallée et de ses forêts, avec une fixité terrible, elle voulait parfois aller le tuer aux environs de Saint-Cloud ou de la Malmaison » ; « plusieurs méditaient, comme Laurence, la mort de cet homme, sans s’effrayer du mot assassinat » (VIII, 538). Dans cette perspective, il est significatif que l’épisode du caveau soit omis par une ellipse ou survolé lors d’un sommaire de menues tâches prosaïques (« nourrir les quatre gentilshommes et \lbrack{}…\rbrack{} veiller à ce qu’ils ne manquassent de rien », VIII, 596) : le creux narratif, vide d’actions, traduit à ce moment le recul du politique, le vide de l’action.}
              \teiP[xmlid={p12-4}]{Si pour les aristocrates, le caveau est un espace de néant d’où l’action politique ne saurait être menée, il en est tout autrement pour leurs mystérieux adversaires. Dans l’affaire de l’enlèvement, le caveau où le sénateur est séquestré est la scène principale et cachée d’une opération qui prend une ampleur politique\teiNote[xmlid={ftn52-3},n={19},place={foot}]{\teiP[xmlid={p-160}]{L’extension politique de l’affaire est explicite : « Un pareil attentat contre un membre de son Sénat excita la colère de l’Empereur \lbrack{}…\rbrack{} », qui voit là « un attentat contre ses institutions, un fatal exemple de résistance aux effets de la Révolution, une atteinte à la grande question des biens nationaux, et un obstacle à cette fusion des partis qui fut la constante occupation de sa politique intérieure » (VIII, 639).}} : annexe contingente dans les menées royalistes, le lieu est au contraire indispensable au second complot dont il devient indissociable. De solution de repli dans la conspiration ratée, il devient la condition même du complot politique réussi. Ce rôle éclate au moment du procès, lors de la comparution de Malin, mystérieusement libéré après que le plaidoyer de Grandville a failli sauver ses clients : la révélation de sa détention dans un lieu dont « l’existence \lbrack{}…\rbrack{} n’était connue que de Marthe, de son fils, de Michu, des quatre gentilshommes et de Laurence » (VIII, 650) fait du caveau le chaînon manquant entre le délit et les inculpés, alors à tort jugés coupables d’avoir attenté à la vie d’un représentant officiel. Loin du refuge passif qu’il était, le caveau devient un foyer actif où se trament les fils de la machination, notamment la complicité postiche de Marthe, reconnue par Malin, qui entraîne la culpabilité tout aussi fausse de ses alliés. L’efficacité inédite du caveau tel qu’il est inséré dans le coup monté révèle l’émergence d’une nouvelle tactique politique où l’exécution de l’action passe de la visibilité de la scène à l’obscurité des coulisses, ce qui témoigne d’une relation nouvelle entre secret et politique : au secret comme simple voile et condition d’une action éclatante en attente de se produire au grand jour succède le secret comme levier et essence même du politique – d’où la promotion des figures de l’ombre, Fouché et Corentin, mais aussi Malin et Louis \teiHi[rend={small-caps}]{XVIII}\teiNote[xmlid={ftn53-4},n={20},place={foot}]{\teiP[xmlid={p-161}]{Allié du Premier Consul et de Fouché, Malin est aussi « dans les secrets de la maison de Bourbon » (VIII, 524). Le Roi, « muet sur les causes \lbrack{}du\rbrack{} désastre » causé par « le secret de l’enlèvement du sénateur » (VIII, 684), fait selon de Marsay bon usage d’un secret bien gardé : « l’histoire de cette nuit a été terrible ; je vous la dis, parce que moi seul la sais, parce que Louis \teiHi[rend={small-caps}]{XVIII} ne l’a pas dite à la pauvre M\teiHi[rend={sup}]{me} de Cinq-Cygne, et qu’il est indifférent au gouvernement actuel qu’elle le sache » (VIII, 689).}}.}
              \teiP[xmlid={p13-4}]{Le réinvestissement du caveau, des chevaliers d’antan à la police moderne de Fouché, illustre ainsi « la fin de la politique ancienne\teiNote[xmlid={ftn54-4},n={21},place={foot}]{\teiP[xmlid={p-162}]{Jacques-David Ebguy, « \teiHi[rend={italic}]{Une ténébreuse affaire}, un roman politique ?… », art. cité, p. 86.}} », déplorée par Laurence qui comprend « qu’il ne s’agissait plus ici de renverser un homme ou le pouvoir à l’aide de gens dévoués, de sympathies fanatiques enveloppées dans les ombres du mystère » (VIII, 647-648). Triomphe alors une « politique absolument opaque\teiNote[xmlid={ftn55-4},n={22},place={foot}]{\teiP[xmlid={p-163}]{Alain, \teiHi[rend={italic}]{Balzac}, \teiHi[rend={italic}]{op. cit.}, p. 37.}} », écrite dans la topographie d’\teiHi[rend={italic}]{Une ténébreuse affaire }: « les forteresses prises d’assaut des temps anciens sont remplacées par les espaces labyrinthiques de la campagne champenoise, les combats à visage découvert se déroulent à présent la nuit et dans des lieux souterrains\teiNote[xmlid={ftn56-4},n={23},place={foot}]{\teiP[xmlid={p-164}]{Jacques-David Ebguy, « \teiHi[rend={italic}]{Une ténébreuse affaire}, un roman politique ?... », art. cité, p. 86.}}. » Encodée dans le transfert du caveau, cette substitution d’une politique du secret à une politique du coup se voit aussi dans le fait que les personnages sonores et solaires d’\teiHi[rend={italic}]{Une ténébreuse affaire} (les Simeuse et Laurence, résumés par leur cri de guerre ; Napoléon, allégorisé par le « soleil » liminaire, VIII, 501) incarnent des systèmes politiques en voie de disparition (la féodalité, l’Ancien Régime, mais aussi l’Empire, dont l’aurore annonce déjà le crépuscule\teiNote[xmlid={ftn57-4},n={24},place={foot}]{\teiP[xmlid={p-165}]{« L’automne de l’année 1803 fut un des plus beaux de la première période de ce siècle que nous nommons l’Empire. \lbrack{}…\rbrack{} Aussi le peuple commença-t-il à établir entre le ciel et Bonaparte, alors déclaré consul à vie, une entente à laquelle cet homme a dû l’un de ses prestiges ; et, chose étrange ! le jour où, en 1812, le soleil lui manqua, ses prospérités cessèrent » (VIII, 501).}}). À l’inverse, les êtres qui font corps avec la nuit et le secret, silencieux et « impénétrables au moment où ils jouent » (VIII, 523), Malin, Corentin et Fouché deviennent les piliers de tous les régimes : Malin, homme du « jeu double », voire « triple » (VIII, 524), se range derrière Bonaparte, Fouché, les Bourbons et les d’Orléans ; Corentin est promu dans \teiHi[rend={italic}]{Splendeurs et misères des courtisanes} chef de la police à la Restauration ; ministre de la Police sous l’Empire, Fouché reprendra son titre, quoique de façon éphémère, en 1815, sous Louis \teiHi[rend={small-caps}]{XVIII}.}
              \teiP[xmlid={p14-3}]{La question de la violence du chronotope gothique peut maintenant être reprise. \teiHi[rend={italic}]{Locus amoenus} en marge de l’action, le caveau tel qu’investi par les émigrés ne saurait récupérer la violence inhérente au roman noir parce que leur conception de la violence déroge aux modalités du chronotope, le secret et le souterrain, lié à l’horizontalité, au monde d’en bas. La représentation décalée du \teiHi[rend={italic}]{topos} révèlerait une décorrélation entre violence et secret dans le cadre de l’action politique envisagée comme coup d’éclat et justifiée par une valeur venue d’en haut : pour le clan aristocratique, l’affrontement politique prend la forme d’un duel épique\teiNote[xmlid={ftn58-4},n={25},place={foot}]{\teiP[xmlid={p-166}]{Voir Jacques-David Ebguy, « \teiHi[rend={italic}]{Une ténébreuse affaire}, un roman politique ?... », art. cité, p. 82-87.}} qui tire sa légitimité d’un principe transcendant, qu’il s’agisse de l’éthique chevaleresque, du roi ou de Dieu. Dans ce cadre, tout attentat, forme de la violence politique, doit être revendiqué \teiHi[rend={italic}]{par}, dirigé \teiHi[rend={italic}]{contre} et mené \teiHi[rend={italic}]{au nom de}, comme le révèlent en creux le soulagement des parents d’Hauteserre, qui « regardèrent comme un bonheur que Bonaparte eût échappé à ce danger, \teiHi[rend={italic}]{en croyant que les Républicains étaient les auteurs de cet attentat} » (VIII, 548, nous soulignons.), ou la résolution de Michu, dévoué à la noblesse, qui refuse de désavouer son intention de « tuer Malin » (VIII, 612). Ce \teiHi[rend={italic}]{locus secretus} qui est aussi \teiHi[rend={italic}]{locus amoenus}, c’est-à-dire qu’il en expulse la violence en dépit des lois du genre, est à l’image d’une pensée politique où le secret n’est pas le terreau d’origine et d’exercice d’une violence au contraire revendiquée, valant pour \teiHi[rend={italic}]{ceux à qui} et \teiHi[rend={italic}]{ce à quoi} elle se rapporte.}
              \teiP[xmlid={p15-2}]{À l’inverse, le retour du \teiHi[rend={italic}]{topos} de la séquestration signe l’alliance du secret et du monde d’en bas à une violence ainsi anonymisée, sans origine, désidéologisée même. Dans la nuit du caveau, la violence, à la faveur du secret, n’a plus ni visage (il faut attendre les dernières pages pour attribuer le complot à Fouché et à Corentin), ni couleur (pour défendre ses intérêts, Fouché séquestre sans scrupule Malin, pourtant représentant de l’Empire comme lui, et ancien allié dans la conspiration avortée contre le Premier Consul à la veille de Marengo), ni classe (Corentin, qui se venge du coup de cravache de Laurence, rétorque à Peyrade, inquiet au sujet de sa « position » sociale : « est-ce que je distingue ? Tout est poisson dans la mer ! », VIII, 590) – sans oublier que le crime sera, à l’issue du procès, mal attribué et mal justifié, mis sur le compte de la morgue anti-bourgeoise des aristocrates en réalité victimes d’usurpation d’identité. Le secret assure ainsi à la violence d’être sans auteur ni principe, à une époque où celle-ci s’est démocratisée (la vraie menace vient des « gens de bas étage », VIII, 612, à savoir Malin, Corentin et Fouché, issus du peuple, plus que du grand homme\teiNote[xmlid={ftn59-4},n={26},place={foot}]{\teiP[xmlid={p-167}]{Enterré dans la forêt, le caveau relève, en partie du moins dans \teiHi[rend={italic}]{Une ténébreuse affaire}, de l’« imaginaire des profondeurs » concomitant à l’avènement dans le champ social de « l’homme d’en bas », identifié chez Hugo et Sue par Pierre Macherey, dans \teiHi[rend={italic}]{Pour une théorie de la production littéraire}, Paris, Maspero, 1966, p. 75-96. Le passage de relais qui se joue au niveau du caveau, de la noblesse à l’humanité inférieure, relève bien du changement de paradigme socio-historique que Macherey voit à l’œuvre au milieu du \teiHi[rend={small-caps}]{xix}\teiHi[rend={sup}]{e} siècle, à la différence près que l’idéologie d’un Balzac, qui voit d’un mauvais œil la démocratisation du champ social, n’a rien à voir avec le socialisme d’un Sue ou d’un Hugo.}}) et où elle a cessé d’être l’apanage du politique qui la légitime au nom des plus hauts intérêts de l’État\teiNote[xmlid={ftn60-4},n={27},place={foot}]{\teiP[xmlid={p-168}]{On renvoie à la légitimation du massacre de la Saint-Barthélemy par Catherine de Médicis et de la Terreur par Robespierre dans \teiHi[rend={italic}]{Sur Catherine de Médicis}.}}, pour servir les intérêts privés les plus vils et l’édification du pouvoir personnel\teiNote[xmlid={ftn61-4},n={28},place={foot}]{\teiP[xmlid={p-169}]{Le narrateur commente ainsi la venue secrète de Malin lorsque la conspiration est éventée : celle-ci « lui donnait un air de zèle aux yeux de Bonaparte, tandis qu’au lieu de s’agir de l’État, il ne s’agissait que de lui-même » (VIII, 523).}}. Le retour du caveau, choisi comme théâtre de la cruauté rendu à son inquiétante noirceur, marque l’avènement d’une violence sans transcendance et perpétrée en secret, occupée à cacher ses auteurs et ses buts réels, au fur et à mesure que ceux-ci, ayant perdu leur caractère public et avouable pour devenir intéressés et égoïstes, se privatisent. Autrement dit, le caveau se fait l’emblème d’une violence horizontalisée mise au service d’acteurs invisibles. Par le recours au souterrain gothique, ils jettent ainsi le voile du secret sur les motivations strictement personnelles de leurs actions.}
              \teiP[xmlid={p16-2}]{À travers les rapports des protagonistes au caché et au secret, la situation du caveau éclaire la nature de l’action politique, telle que chaque clan, qui recoupe un système socio-historique, la pense. Dans la mesure où le chronotope gothique est lié au « passé historique\teiNote[xmlid={ftn62-4},n={29},place={foot}]{\teiP[xmlid={p-170}]{Mikhaïl Bakhtine, \teiHi[rend={italic}]{Esthétique et théorie du roman}, \teiHi[rend={italic}]{op. cit.}, p. 387.}} », les façons de l’habiter, de se le (ré)approprier, révèlent une certaine \teiHi[rend={italic}]{gestion} du passé, qui dans \teiHi[rend={italic}]{Une ténébreuse affaire} conditionne l’impuissance ou l’efficacité de l’action politique menée au présent. C’est cette relation du présent au passé cristallisée dans l’usage nostalgique ou stratégique du chronotope qu’il s’agira d’examiner.}
            
\end{teiDiv}
            \begin{teiDiv}[type={section1},xmlid={div03-2}]

              \teiHead[xmlid={du-refuge-au-piege-le-caveau-miroir-du-devenir-historique-001}]{Du refuge au piège : le caveau, miroir du devenir historique}
              \teiP[xmlid={p17-2}]{Dans \teiHi[rend={italic}]{Une ténébreuse affaire}, le caveau renoue avec l’expression du temps archaïque. Vestige d’un « monastère pris, saccagé, démoli » et en « ruines » datant des \teiHi[rend={small-caps}]{xii}\teiHi[rend={sup}]{e} et \teiHi[rend={small-caps}]{xiii}\teiHi[rend={sup}]{e} siècles (VIII, 564), souvenir d’une époque où le scepticisme de la raison n’avait pas dissipé le sentiment religieux, il « appartenait à la plus vieille coupe de la forêt » (VIII, 565), que Michu découvre en 1794 grâce aux indications du feu marquis de Simeuse, père des jumeaux (VIII, 564), en quelque sorte posé comme le premier patron des lieux.}
              \teiP[xmlid={p18}]{Aussi, lorsque les Simeuse et les d’Hauteserre descendent au caveau, c’est dans un passé prérévolutionnaire avec lequel ils s’identifient (« Nous sommes de vrais chevaliers du Moyen Âge », VIII, 620), et dans les entrailles de leurs ancêtres qu’ils se réfugient, comme l’explicite Michu : « leur pauvre père avait peut-être une vision en me mettant sur la piste de cette cachette, il a pressenti que ses fils s’y sauveraient » (VIII, 567). Double de l’hôtel de Cinq-Cygne, « démoli avec une inexplicable rapidité » (VIII, 523) à la Révolution, le caveau sert de refuge archaïque et serein de la passivité ; il miroite « l’isolement politique et l’anachronisme de ces figures\teiNote[xmlid={ftn63-4},n={30},place={foot}]{\teiP[xmlid={p-171}]{Jacques-David Ebguy, « \teiHi[rend={italic}]{Une ténébreuse affaire}, un roman politique ?... », art. cité, p. 84.}} », ainsi que leur déni de l’histoire présente au profit d’un repli aveugle dans le passé.}
              \teiP[xmlid={p19}]{Il faut revenir sur les sept mois passés sous terre, au cours desquels les Simeuse et les d’Hauteserre ne voient pas les prémices de l’Empire. Cette réclusion volontaire engage un traitement spécifique du récit, par l’ellipse et le sommaire, qui occultent, sinon nivellent tout relief événementiel. Ces deux formes de repli, topographique et narratif, traduisent une manière de se rapporter au temps et à l’événementialité historique : l’occupation du caveau est emblématique d’un désenracinement par rapport à l’histoire présente, vécue sur un tout autre plan, et d’un aveuglement au devenir, dont fera preuve le clan aristocratique d’un bout à l’autre d’\teiHi[rend={italic}]{Une ténébreuse affaire}, en évoluant dans une bulle anachronique coupée du dehors. Comme au caveau, l’histoire advient sans eux : tandis que « le monde entier se préoccupait du dénouement » des campagnes impériales (VIII, 608), « les habitants de Cinq-Cygne ne firent aucune attention au couronnement de l’empereur Napoléon » (VIII, 607-608), Laurence « n’apprit l’étonnant triomphe d’Austerlitz » que « deux mois après » (VIII, 609) ; le passé, comme le caveau, est un ailleurs qui refoule le présent sans l’empêcher d’advenir : admis sur le territoire, les ex-émigrés voient dans Gondreville, qu’ils « regard\lbrack{}ent\rbrack{} toujours \lbrack{}…\rbrack{} comme à \lbrack{}eux\rbrack{} » (VIII, 611), une terre idyllique où s’adonner à l’amour et à la chasse comme il y a « deux siècles » (VIII, 616), niant par leur \teiHi[rend={italic}]{ethos} aristocratique les acquis révolutionnaires et les lois impériales en arpentant, « malgré les arrêtés de préfecture » (VIII, 619), une terre qui ne leur appartient plus. La manière d’occuper le caveau a donc une valeur allégorique : la retraite dans ces ruines souterraines traduit l’appartenance de ces personnages, en cela fossiles\teiNote[xmlid={ftn64-3},n={31},place={foot}]{\teiP[xmlid={p-172}]{Jacques-David Ebguy emploie l’étiquette de « revenant » : il est celui qui veut « retrouver sa place dans la société par un mouvement de négation ou de suspens de l’Histoire », \teiHi[rend={italic}]{Le Héros balzacien. Balzac et la question de l’héroïsme}, Saint-Cyr-sur-Loire, Christian Pirot, 2010, p. 209.}}, à une strate temporelle déjà recouverte par le mouvement de l’histoire.}
              \teiP[xmlid={p20}]{Ironie du sort, ce caveau censé être un sanctuaire, mi-arche mi-berceau, devant selon Michu sauver ses maîtres pour les réintégrer au champ social à restaurer, sera en réalité leur tombeau, à l’issue d’un détournement stratégique du passé incarné par le chronotope. Le piège de Fouché et Corentin repose en effet sur un retournement pervers du caveau depuis la question de la propriété : les véritables coupables s’emparent du lieu pour faire jouer son appartenance première, confirmée lors des débats où « il fut acquis \lbrack{}…\rbrack{} que cette cachette, découverte par Michu, n’était connue que de lui, de Laurence et des quatre gentilshommes » (VIII, 667). Par une rétorsion ironique, les héros crispés sur l’idée qu’ils se font de la propriété\teiNote[xmlid={ftn65-3},n={32},place={foot}]{\teiP[xmlid={p-173}]{On rappelle que les habitants de Cinq-Cygne refusent de céder Gondreville. Pour racheter leur terre, ils se lancent dans une expédition destinée à récupérer leur or enterré dans la forêt. Celle-ci sera néanmoins récupérée par l’Accusation, qui verra dans cette échappée nocturne la preuve de leur responsabilité dans l’enlèvement perpétré le même soir. À l’issue du procès, le désir de propriété se retourne contre lui-même : loin de récupérer Gondreville, les Simeuse et les d’Hauteserre, condamnés à s’engager dans l’armée impériale pour y trouver l’exil et la mort, en sont définitivement bannis.}}, liée sous l’Ancien Régime à la transmission d’un nom et d’une lignée, périssent par où ils ont péché. La réappropriation du caveau dirait alors quelque chose de la gestion du passé au présent à travers le prisme de la propriété, bouleversée à la Révolution par la vente des biens nationaux, qui hante \teiHi[rend={italic}]{Une ténébreuse affaire}, avant sa régulation par le Code civil. S’en faisant la chambre d’échos, le caveau servirait deux fins : sanctionner un passé périmé ; décrire le présent de Juillet.}
              \teiP[xmlid={p21}]{L’ancienne conception de la propriété, fondée sur l’évidence d’une fusion de l’être et de la terre, à travers notamment la perpétuation en elle du nom et de la lignée, est d’abord périmée par la facilité avec laquelle les sbires de Fouché attribuent le caveau à ses anciens propriétaires, qui, à ce moment de la diégèse, n’y résident plus, pas plus qu’ils n’y songent. Il suffit d’un faux imitant l’écriture de Michu pour que ses maîtres soient impliqués et que le crime commis en ce lieu fasse d’eux des criminels. La façon dont ces signes contrefaits parviennent à faire passer pour réelle une propriété fictive accuse au fond l’artifice du lien de propriété, sans consubstantialité ni nécessité aucune, mais de l’ordre plutôt de la convention, voire de l’imaginaire. Le succès du coup monté qui retourne contre elle-même cette conception datée et monolithique de la propriété sanctionne la défaite d’une crispation anachronique au passé, inadaptée à la nouvelle cartographie foncière du présent né de la Révolution ; s’il est possible, en réactionnaire, de déplorer cet état de fait, il serait peu réaliste de l’ignorer. Ce désaveu pervers d’une conception qu’on dirait essentialiste de la propriété rejoint d’une certaine manière l’invitation pragmatique du marquis de Chargebœuf, parent des héros, qui les enjoint à admettre que les « temps \lbrack{}sont\rbrack{} bien changés » et à se « soumettre aux événements » en allant « acheter \teiHi[rend={italic}]{quelque} belle terre dans un \teiHi[rend={italic}]{autre} coin de la France » (VIII, 614, nous soulignons) : dans les deux cas, c’est l’arbitraire et la contingence d’une propriété fabriquée par des signes (par l’argent, par le contrat, par les faux aussi), qui se révèlent dans la valse de réappropriations et désappropriations dont les lieux font l’objet.}
              \teiP[xmlid={p22}]{On notera enfin que l’ultime locataire du caveau est Malin. C’est lui qui, en dernière instance, sort gagnant de ce séjour sous terre jouant de façon grinçante avec les motifs de l’anabase et de la résurrection : aux héros sortis du caveau temporairement restaurés dans leurs droits civils pour mieux mourir succède un antagoniste enterré sénateur d’Empire aux velléités nobiliaires (« Malin voudrait être comte de Gondreville », VIII, 615) d’abord revenu « nommé \lbrack{}…\rbrack{} comte de Gondreville » (VIII, 668), et finalement sacré propriétaire des biens nationaux à l’issue du procès. Ce caveau qui expulse des nobles vivants-morts pour accoucher d’un bourgeois triomphant illustrerait l’édification, sur les ruines de l’Ancien Régime, d’un nouveau corps politique qui fait du passé un usage stratégique, c’est-à-dire qui sait s’en défaire au moment précis où il devient \teiHi[rend={italic}]{passé}, après avoir « pressur\lbrack{}é\rbrack{} les événements pour son compte » (VIII, 523) dans l’intervalle où ceux-ci portaient encore le nom de \teiHi[rend={italic}]{présent}. Ainsi s’éclaire le \teiHi[rend={italic}]{girouettisme} de Malin, introduit dès les premières pages en plein exercice de « politique expectante » (VIII, 526), tâchant de déterminer si l’avenir, et par contrecoup le passé, sera impérial ou royal, afin de décider qui, de Bonaparte ou des Bourbons, il faut soutenir. Sous la forme d’une frise chronologique dépliant son devenir sur le long \teiHi[rend={small-caps}]{xix}\teiHi[rend={sup}]{e} siècle, son portrait à l’épilogue le pose ainsi comme}
              \begin{teiCit}[xmlid={cit01-4}]

                \begin{teiQuote}
l’ancien clerc venu d’Arcis, l’ancien Représentant du Peuple, l’ancien Thermidorien, l’ancien tribun, l’ancien conseiller d’État, l’ancien comte de l’Empire et sénateur, l’ancien pair de Louis XVIII, le nouveau pair de Juillet \lbrack{}…\rbrack{} (VIII, 687).
\end{teiQuote}
              
\end{teiCit}
              \teiP[xmlid={p23}]{L’anaphore d’« ancien » fait de Malin l’homme de tous les passés, dont l’ancrage successif mais éphémère dans toutes les strates temporelles – et non la fossilisation dans une seule, à l’instar de ses rivaux –, lui assure d’être toujours « nouveau », c’est-à-dire de ressusciter à \teiHi[rend={italic}]{son} présent. Dans cette perspective, Malin émergeant victorieux du caveau figurerait un usage opportuniste du passé qui, pour être efficace dans la quête du pouvoir politique, n’en est pas bon pour autant. Sur ce point, le roman de Balzac, entre ses nobles défaits et son exécrable bourgeois, reste aporétique ; sans doute parce que l’arène politique est le lieu d’une dissolution de l’éthique qui trouve dans le chronotope gothique, espace de « nuit morale\teiNote[xmlid={ftn66-3},n={33},place={foot}]{\teiP[xmlid={p-174}]{Annie Le Brun, \teiHi[rend={italic}]{Les Châteaux de la subversion}, \teiHi[rend={italic}]{op. cit.}, p. 14.}} », son reflet.}
              \teiP[xmlid={p24}]{À travers son chronotope gothique, \teiHi[rend={italic}]{Une ténébreuse affaire} figure le point de bascule entre deux époques, l’Ancien Régime et le monde qu’il faudrait qualifier de moderne, qui s’accompagne d’une modification du paradigme politique, observable dans la conception et le rapport au secret, inscrit dans l’usage du lieu.}
              \teiP[xmlid={p25}]{Avec ses deux complots et ses deux humanités, cet espace de coexistence et de rivalité politico-historique qu’est le caveau renvoie au double visage de Balzac romancier lui-même : la cohabitation en un \teiHi[rend={italic}]{même} lieu dans le cours de l’intrigue résume le travail de Balzac-archéologue du présent, conscient que celui-ci est formé de multiples « strates\teiNote[xmlid={ftn67-3},n={34},place={foot}]{\teiP[xmlid={p-175}]{Aude Déruelle, « Balzac penseur de l’histoire ? », Francesco Spandri (dir.), \teiHi[rend={italic}]{Balzac penseur}, Paris, Classiques Garnier, 2019, p. 106.}} » temporelles en coprésence ; son occupation \teiHi[rend={italic}]{successive }et \teiHi[rend={italic}]{différenciée} dans le temps du récit annonce en revanche le souci de Balzac-historien du \teiHi[rend={small-caps}]{xix}\teiHi[rend={sup}]{e} siècle de capter, selon le mot de l’« Avant-propos », le vertigineux « mouvement » (I, 12) d’une société révolutionnée désormais livrée au devenir.}
            
\end{teiDiv}
          
\end{teiElement}
        
\end{teiElement}
        \begin{teiElement}[name={group},type={chapter},data-page-title={Un cabinet de toilette bien politique},data-page-subtitle={Les dessous du Second Empire dans trois romans d’Émile Zola},data-page-authors={Émilie Bauduin@@Université du Québec à Montréal, Figura, Montréal, Canada@@Université Sorbonne Nouvelle, CRP 19, EA 3423, Paris, France},data-include-href={../XML/chap\_5\_Un\_cabinet\_de\_toilette\_Locus\_politicus.xml},xmlid={chapter-004}]

          \begin{teiElement}[name={front}]

            \begin{teiDiv}[type={titlePage},xmlid={titlepage-005}]

              \teiP[rend={title-main},xmlid={p-176}]{Un cabinet de toilette bien politique}
              \teiP[rend={title-sub},xmlid={p-177}]{Les dessous du Second Empire dans trois romans d’Émile Zola}
              \teiP[rend={author-aut},xmlid={p-178}]{Émilie \teiHi[rend={small-caps}]{Bauduin}}
              \teiP[rend={authority\_affiliation},xmlid={p-179}]{Université du Québec à Montréal, Figura, Montréal, Canada}
              \teiP[rend={authority\_affiliation},xmlid={p-180}]{Université Sorbonne Nouvelle, CRP 19, EA 3423, Paris, France}
            
\end{teiDiv}
          
\end{teiElement}
          \begin{teiElement}[name={body},xmlid={body-005}]

            \teiP[xmlid={p1-5}]{Dans la France post-révolutionnaire, où désormais « \lbrack{}l\rbrack{}a vie publique postule la transparence\teiNote[xmlid={ftn1-5},n={1},place={foot}]{\teiP[xmlid={p-181}]{Michelle Perrot, « Avant et ailleurs », Philippe Ariès et Georges Duby (dir.), \teiHi[rend={italic}]{Histoire de la vie privée. Tome IV. De la Révolution à la Grande Guerre}, Paris, Éditions du Seuil, 1999, p. 15.}} » et où l’on « soupçonne les “intérêts privés ou particuliers” de contrarier le changement et d’abriter l’intrigue et le complot\teiNote[xmlid={ftn35-4},n={2},place={foot}]{\teiP[xmlid={p-182}]{Michelle Perrot, « Public, privé et rapports de sexes », Jacques Chevalier (dir.), \teiHi[rend={italic}]{Public/privé}, Paris, PUF, 1995, p. 67.}} », intime et social, privé et public, individuel et collectif s’entremêlent fréquemment. Cela est particulièrement sensible au sein des habitations domestiques, où « se matérialisent les visées du pouvoir, les rapports interpersonnels et la quête de soi\teiNote[xmlid={ftn36-4},n={3},place={foot}]{\teiP[xmlid={p-183}]{Michelle Perrot, « Manières d’habiter », Philippe Ariès et Georges Duby (dir.),\teiHi[rend={italic}]{ Histoire de la vie privée}, \teiHi[rend={italic}]{op. cit.}, p. 296.}} » selon l’historienne Michelle Perrot. Aussi les romanciers du \teiHi[rend={small-caps}]{xix}\teiHi[rend={sup}]{e} siècle donnent-ils à lire les enjeux qui animent leur époque en dépeignant les lieux de l’intimité : pour ces auteurs, représenter l’espace privé équivaut à « montrer les notions politiques qui le sous-tendent, les conditions économiques qui le produisent, les relations sociales qu’il engendre, les rapports humains qu’il permet et les croyances, savoirs et discours qu’il met en scène\teiNote[xmlid={ftn37-4},n={4},place={foot}]{\teiP[xmlid={p-184}]{Jean-François Richer\teiHi[rend={italic}]{, Les Boudoirs dans l’œuvre d’Honoré de Balzac : surveiller, mentir, désirer, mourir}, Montréal, Nota Bene, 2012, p. 11.}} ».}
            \teiP[xmlid={p2-5}]{Balzac et Flaubert pavent la voie à cette tendance dans la première moitié du siècle ; puis, dans la seconde, c’est au tour des écrivains naturalistes, et notamment d’Émile Zola, de s’emparer de ce nouvel objet. Le maître de Médan accorde en effet une place privilégiée à la représentation de l’espace domestique dans ses textes. Elle lui offre l’occasion de se livrer à la « vaste entreprise de dévoilement\teiNote[xmlid={ftn38-4},n={5},place={foot}]{\teiP[xmlid={p-185}]{Colette Becker, \teiHi[rend={italic}]{Zola. Le Saut dans les étoiles}, Paris, Presses de la Sorbonne Nouvelle, 2002, p. 52.}} » du cycle des \teiHi[rend={italic}]{Rougon-Macquart} (1871-1893), qui vise à révéler tant les secrets des alcôves que les rouages souterrains de la vie publique. Parmi les pièces qui composent l’habitation, l’un des espaces privilégiés par l’écrivain pour procéder à ce dévoilement conjoint de l’intime et du social est le cabinet de toilette.}
            \teiP[xmlid={p3-5}]{« \lbrack{}P\rbrack{}etite pièce où l’on s’habille\teiNote[xmlid={ftn39-4},n={6},place={foot}]{\teiP[xmlid={p-186}]{Émile Littré, \teiHi[rend={italic}]{Dictionnaire de la langue française. Tome IV. Q-Z}, Paris, Hachette, 1874, p. 2241.}} » que l’on ne trouve, jusqu’à la toute fin du siècle, que dans les habitations bourgeoises les plus luxueuses\teiNote[xmlid={ftn40-4},n={7},place={foot}]{\teiP[xmlid={p-187}]{Comme le remarque Georges Vigarello, « \lbrack{}c\rbrack{}e sont surtout les espaces intimes de l’élite qui sont transformés par les nouveaux outils d’embellissement » qui apparaissent au \teiHi[rend={small-caps}]{xix}\teiHi[rend={sup}]{e} siècle (\teiHi[rend={italic}]{Histoire de la beauté. Le corps et l’art d’embellir de la Renaissance à nos jours}, Paris, Éditions du Seuil, 2004, p. 179).}}, le cabinet de toilette est doté sous le Second Empire de fonctions paradoxales qui le placent à la frontière des sphères publique et privée. D’une part, il s’agit toujours à cette époque d’un lieu de sociabilité où la dame de la maison – puisque c’est surtout à elle qu’est attribuée cette pièce selon le dimorphisme sexué de l’espace qui a cours durant cette période – reçoit ses invités les plus chers, comme c’était le cas au \teiHi[rend={small-caps}]{xvIII}\teiHi[rend={sup}]{e} siècle, lors du rituel de la « seconde toilette », « sorte d’imitation hybride entre la ruelle des précieuses et une version intimiste du lever royal d’autrefois\teiNote[xmlid={ftn41-4},n={8},place={foot}]{\teiP[xmlid={p-188}]{Philippe Perrot, \teiHi[rend={italic}]{Le Travail des apparences, ou Les Transformations du corps féminin, }\teiHi[rend={ small-caps italic}]{xv III}\teiHi[rend={italic sup}]{e}\teiHi[rend={italic}]{-}\teiHi[rend={ small-caps italic}]{xix}\teiHi[rend={italic sup}]{e}\teiHi[rend={italic}]{ siècles}, Paris, Éditions du Seuil, 1984, p. 42. Sous l’Ancien Régime, la toilette se fait donc en deux temps. « Après \lbrack{}de\rbrack{} prosaïques préparations de coulisse, faites d’ablutions, d’aspersions, d’immersions, de lavements et de frottements \lbrack{}…\rbrack{}, place au spectacle », à la « toilette de simulacre » : « \lbrack{}l\rbrack{}a maîtresse de céans \lbrack{}…\rbrack{} reçoit \lbrack{}dans son cabinet de toilette\rbrack{} son public visiteur, entourée d’un décor où l’art s’est emparé du moindre objet » (\teiHi[rend={italic}]{ibid.}, p. 42, 45 et 40).}} ». « Sanctuaire de la femme », où celle-ci « imprime sa marque particulière\teiNote[xmlid={ftn42-4},n={9},place={foot}]{\teiP[xmlid={p-189}]{Blanche Augustine Angèle Soyer Staffe, \teiHi[rend={italic}]{Le Cabinet de toilette}, Paris, Victor-Havard, 1891, p. 1.}} », le cabinet de toilette est ainsi « un salon ou tout au moins un boudoir où la maîtresse de maison peut se tenir le matin et même recevoir parfois ses amies intimes\teiNote[xmlid={ftn43-4},n={10},place={foot}]{\teiP[xmlid={p-190}]{\teiHi[rend={italic}]{La Construction moderne}, 22 septembre 1900, cité dans Monique Eleb et Anne Debarre, \teiHi[rend={italic}]{L’Invention de l’habitation moderne. Paris 1880-1914}, Paris-Bruxelles, F. Hazan-Archives d’architecture moderne, 1995, p. 225. Pour ce qui est de la toilette masculine, Monique Eleb et Anne Debarre indiquent que « selon les rares traités de savoir-vivre qui abordent cette question, \lbrack{}…\rbrack{} \lbrack{}la plupart des hommes\rbrack{} utilisent le cabinet de toilette de leur épouse, le matin, avant elle » (\teiHi[rend={italic}]{ibid.})\teiHi[rend={italic}]{.}}} ». Mais, d’autre part, dans la foulée des avancées du discours hygiéniste et des transformations techniques et urbaines qui caractérisent les premières décennies de la seconde moitié du \teiHi[rend={small-caps}]{xix}\teiHi[rend={sup}]{e} siècle, il devient aussi progressivement un espace de toilette corporelle\teiNote[xmlid={ftn44-4},n={11},place={foot}]{\teiP[xmlid={p-191}]{« Avec des stases, des inerties, des rétroactions, des régressions temporaires, des disparités géographiques et sociales considérables, l’élan \lbrack{}…\rbrack{} est donné, la pompe amorcée : du tub à la baignoire, du “cabinet de toilette” à la “salle de bains”, lentement mais sûrement, on passe de la friction humide ou de la brève aspersion à la longue et complète immersion, d’une eau parcimonieuse répandue à la verticale à une eau pléthorique goûtée à l’horizontale » (Philippe Perrot, \teiHi[rend={italic}]{Le Travail des apparences}, \teiHi[rend={italic}]{op. cit.}, p. 137).}}. Jusqu’aux années 1880, y trônent parfois une baignoire et souvent des cuvettes, des brocs, des peignes, des pommades et des savons – ce qui implique un certain, si ce n’est un complet dénudement des corps\teiNote[xmlid={ftn45-4},n={12},place={foot}]{\teiP[xmlid={p-192}]{Dans les années\teiHi[rend={italic}]{ }1880, une période marquée par la révolution pastorienne et la découverte des « microbes », « le cabinet de toilette se double, dans les appartements luxueux, d’une salle de bains », laquelle est spécifiquement dédiée aux soins hygiéniques intimes. Monique Eleb et Anne Debarre, \teiHi[rend={italic}]{L’Invention de l’habitation moderne}, \teiHi[rend={italic}]{op. cit.}, p. 219. Au sujet des développements de l’hygiène au \teiHi[rend={small-caps}]{xix}\teiHi[rend={sup}]{e} siècle et de ses effets sur les pratiques de propreté, voir Georges Vigarello, \teiHi[rend={italic}]{Le Propre et le sale}, Paris, Éditions du Seuil, 1985, p. 191-229 et p. 230-234.}}. Ces nouvelles fonctions hygiéniques ne privatisent pas immédiatement la pièce ; à l’époque que décrit Zola, celle-ci rassemble plutôt les deux formes de toilette, privée et publique\teiNote[xmlid={ftn46-4},n={13},place={foot}]{\teiP[xmlid={p-193}]{Philippe Perrot précise notamment que si une baignoire est installée dans le cabinet de toilette, « il la faut dissimulée “derrière une riche tenture orientale”, cet ameublement demeurant trop démonstratif d’une hygiène intime, évoquant trop crûment la nudité du corps qu’il implique, la liberté des gestes qu’il autorise » (Philippe Perrot, \teiHi[rend={italic}]{Le Travail des apparences}, \teiHi[rend={italic}]{op. cit.}, p. 133). Selon Georges Vigarello, « \lbrack{}c\rbrack{}’est à la fin du \teiHi[rend={small-caps}]{xix}\teiHi[rend={sup}]{e} siècle que se systématise une injonction : fermer rigoureusement les accès des cabinets de toilette et des salles de bains » (\teiHi[rend={italic}]{Le Propre et le sale}, \teiHi[rend={italic}]{op. cit.}, p. 246).}}. Dans ses romans, l’écrivain fait alors sienne la rencontre de l’individuel et du collectif rendue possible par le cabinet de toilette, afin de « mettre à nu » les personnages qu’il représente, et ce, en particulier lorsqu’il est question de critiquer le règne impérial et ses représentants. Ce lieu de l’intimité est notamment celui où des personnages liés à la sphère politique de Napoléon \teiHi[rend={small-caps}]{III }trahissent leur camp (\teiHi[rend={italic}]{La Fortune des Rougon}, 1871), mènent leurs campagnes (\teiHi[rend={italic}]{Son Excellence Eugène Rougon}, 1876) et signent leur défaite (\teiHi[rend={italic}]{La Débâcle}, 1892), l’écrivain tissant des liens étroits entre les attributs de cette pièce et les dérives du pouvoir sous le Second Empire. Dans cette réflexion, il s’agira de montrer comment, en dépeignant le cabinet de toilette et les usages qui en sont faits par certains personnages liés à la sphère politique, Zola expose les dessous (et les tares) du règne du prince-président.}
            \begin{teiDiv}[type={section1},xmlid={div01-4}]

              \teiHead[xmlid={ablutions-et-trahison-dans-la-fortune-des-rougon-001}]{Ablutions et trahison dans\teiHi[rend={italic}]{ La Fortune des Rougon}}
              \teiP[xmlid={p4-5}]{\teiHi[rend={italic}]{La Fortune des Rougon} met en scène la rébellion infructueuse des républicains, à la suite du coup d’État du 2 décembre 1851, dans la ville fictive de Plassans. Dans ce roman, les caractéristiques hygiéniques des personnages permettent de les situer sur l’échelle sociale de la sous-préfecture. Comme le rappelle l’historien Philippe Perrot, dans la société postrévolutionnaire :}
              \begin{teiCit}[xmlid={cit01-5}]

                \begin{teiQuote}
\lbrack{}P\rbrack{}ar ce qu’elle implique en effet de matière première (l’eau), d’infrastructure (espace, équipement, savon, linge, etc.), d’habitudes (c’est-à-dire d’éducation) et de temps « perdu », l’opération de se laver demeure un privilège d’élite dont le résultat visible constitue un élément crucial de sa panoplie signalétique\teiNote[xmlid={ftn47-4},n={14},place={foot}]{\teiP[xmlid={p-194}]{Philippe Perrot, \teiHi[rend={italic}]{Le Travail des apparences}, \teiHi[rend={italic}]{op. cit.}, p. 107.}}.
\end{teiQuote}
              
\end{teiCit}
              \teiP[xmlid={p5-5}]{En régime zolien, les attributs hygiéniques des personnages donnent également l’occasion au romancier de les situer sur une autre échelle, axiologique cette fois. Ces attributs, de même que ceux associés au cabinet de toilette, sont en effet utiles à la charge critique portée par Zola contre le règne de Napoléon \teiHi[rend={small-caps}]{III}, comme le révèlent la description du personnage d’Antoine Macquart et celle de sa visite impromptue au cabinet de toilette.}
              \teiP[xmlid={p6-5}]{Antoine Macquart est l’un des enfants illégitimes d’Adélaïde Fouque dont Pierre Rougon et sa femme, Félicité, ont usurpé l’héritage. Ancien soldat subalterne devenu un vannier que « la moindre besogne épouvant\lbrack{}e\rbrack{}\teiNote[xmlid={ftn48-4},n={15},place={foot}]{\teiP[xmlid={p-195}]{Émile Zola, \teiHi[rend={italic}]{La Fortune des Rougon} \lbrack{}1871\rbrack{}, dans \teiHi[rend={italic}]{Œuvres complètes, t. IV. La guerre et la Commune. 1870-1871}, Paris, Nouveau Monde, 2003, p. 115. Désormais, les références à cette œuvre seront indiquées entre parenthèses dans le texte par le sigle \teiHi[rend={italic}]{F.}}} », il peine à se tailler une position dans la sous-préfecture. La description de son apparence en témoigne : selon la narration, Macquart ne « jou\lbrack{}e\rbrack{} au monsieur \lbrack{}que\rbrack{} tant qu’il \lbrack{}a\rbrack{} de l’argent en poche » (\teiHi[rend={italic}]{F.}, p. 119). Lorsqu’il vit aux crochets de ses proches, il est ainsi décrit comme « \lbrack{}s\rbrack{}oigneusement rasé » (\teiHi[rend={italic}]{ibid.}). Mais, dès que le personnage est privé des ressources pécuniaires extorquées à sa famille, il ne peut plus procéder à de tels soins, ce qui le rend moins respectable : « Il ne fut plus “monsieur” Macquart, cet ouvrier rasé et endimanché tous les jours, qui jouait au bourgeois ; il redevint le grand diable malpropre qui avait spéculé jadis sur ses haillons » (\teiHi[rend={italic}]{F.}, p. 138) Dans ce portrait, c’est le caractère trompeur de l’apparence qui domine. Le terme de « jeu » est employé à deux reprises pour qualifier les aspirations bourgeoises d’Antoine, ce qui en indique la théâtralité et la fausseté. L’oscillement du personnage entre des débuts malpropres et une réussite sociale trompeuse établie sous le sceau de la propreté préfigure le contexte de sa traîtrise politique dans la suite du texte : tout comme il affecte d’être un bourgeois, Antoine se fait passer pour un républicain lorsqu’il croit cette situation préférable, jusqu’à ce que de singulières ablutions exposent sa vraie nature.}
              \teiP[xmlid={p7-5}]{La description morale du personnage indique sa nature duplice. C’est son égoïsme qui le pousse à prendre parti pour la République :}
              \begin{teiCit}[xmlid={cit02-4}]

                \begin{teiQuote}
Antoine Macquart, rongé d’envie et de haine, rêvant des vengeances contre la société entière, accueillit la République comme une ère bienheureuse où il lui serait permis d’emplir ses poches dans la caisse du voisin, et même d’étrangler le voisin, s’il témoignait le moindre mécontentement. \lbrack{}…\rbrack{} Ce qui fit surtout de lui un républicain féroce, ce fut l’espérance de se venger enfin des Rougon, qui se rangeaient franchement du côté de la réaction (\teiHi[rend={italic}]{F.}, p. 120-121).
\end{teiQuote}
              
\end{teiCit}
              \teiP[xmlid={p8-5}]{Il est opposé en cela à son neveu, Silvère Mouret, un jeune ouvrier dont les aspirations républicaines sont qualifiées de « purement morales » (\teiHi[rend={italic}]{F.}, p. 130). Macquart initie Silvère aux valeurs de la République, lesquelles sont, tant pour le jeune homme que pour Zola, associées aux valeurs viriles prônées par le siècle\teiNote[xmlid={ftn49-4},n={16},place={foot}]{\teiP[xmlid={p-196}]{Au sujet de ces valeurs, parmi lesquelles on trouve notamment la force, le courage et le patriotisme, voir Alain Corbin, Jean-Jacques Courtine et Georges Vigarello (dir.), \teiHi[rend={italic}]{Histoire de la virilité, t. II. Le Triomphe de la virilité. Le }\teiHi[rend={ small-caps italic}]{xix}\teiHi[rend={italic sup}]{e}\teiHi[rend={italic}]{ siècle}, Paris, Éditions du Seuil, 2011. Au sujet de la critique que fait Zola de la « perte de “virilité” des mâles » sous le Second Empire, voir Nadia Beaudoin, « Émile Zola et la décadence. Les motifs décadents chez “Le père du Naturalisme” », \teiHi[rend={italic}]{Québec français}, n\teiHi[rend={sup}]{o} 13, printemps 1999, p. 76.}}. Silvère émet ainsi un discours très élogieux à l’égard d’Antoine lorsqu’il le croit dévoué à la cause républicaine. Comme son oncle semble valoriser avec lui la démocratie sociale, il fait notamment partie des « vrais » hommes selon lui. C’est du moins ce que suggère le rapport qu’il établit, lorsqu’il observe les contingents s’approchant de Plassans, entre son oncle et les insurgés, dont il reconnaît l’énergie et la vigueur viriles : « Le contingent de Chavanoz ! dit-il. Il n’y a que huit hommes, mais ils sont solides ; l’oncle Antoine les connaît… » (\teiHi[rend={italic}]{F.}, p. 43). Dans la pensée de Silvère, cette caractérisation distingue le vannier des bourgeois conservateurs, dont il considère l’apparence soignée non comme un signe de respectabilité, mais bien comme la manifestation de leur supériorité méprisante face au peuple. Chez lui, la constatation de la bonne tenue de l’apparence est une insulte : il espère pouvoir « di\lbrack{}r\rbrack{}e leur fait \lbrack{}…\rbrack{} à ces beaux messieurs » (\teiHi[rend={italic}]{F.}, p. 35), une fois la victoire des républicains advenue. Mais, loin de se présenter comme un sincère défenseur de la République, Macquart est au fond un personnage hypocrite et lâche qui cherche tout au long du récit à revêtir l’apparence des bourgeois qu’il prétend exécrer. C’est ainsi lors de son emprisonnement dans le cabinet de toilette de M. Garçonnet, un lieu étroitement associé au féminin (et dont le notable propriétaire porte, de surcroît, un nom peu viril), qu’il trahit les intérêts républicains.}
              \teiP[xmlid={p9-5}]{Macquart est enfermé dans le cabinet de toilette par son demi-frère lorsque ce dernier tente de réprimer les rébellions républicaines dans la ville. La description de son enfermement établit une dialectique entre ses convictions républicaines (qui désignent la sphère publique et les valeurs viriles dans la pensée zolienne) et son désir de luxe (qui, sous la plume de Zola, renvoie plutôt au privé, à l’individualisme et à la dégénérescence bonapartiste). Après deux jours de détention, Antoine « éprouv\lbrack{}e\rbrack{} des envies de briser la porte, à la pensée que son frère se carr\lbrack{}e\rbrack{} dans la pièce voisine » et « se prome\lbrack{}t\rbrack{} de l’étrangler de ses propres mains lorsque les insurgés viendr\lbrack{}o\rbrack{}nt le délivrer » (\teiHi[rend={italic}]{F.}, p. 235). Ces pensées, qui correspondent à l’idéal viril par leur référence à la force brutale et au désir de domination du personnage, témoignent de son penchant, semble-t-il toujours intact, pour la cause républicaine. Mais, peu à peu, Antoine est rasséréné par le cabinet de toilette en raison de ses attributs :}
              \begin{teiCit}[xmlid={cit03-4}]

                \begin{teiQuote}
Il y respirait \teiHi[rend={italic}]{une odeur douce}, un sentiment de bien-être qui détendait ses nerfs. \lbrack{}...\rbrack{} \lbrack{}L\rbrack{}e divan était \teiHi[rend={italic}]{moelleux et tiède} ; des parfums, des pommades, des savons garnissaient le lavabo de marbre, et le jour pâlissant tombait du plafond avec \teiHi[rend={italic}]{des voluptés molles}, pareil aux lueurs d’une lampe pendue dans une \teiHi[rend={italic}]{alcôve} \lbrack{}…\rbrack{} (\teiHi[rend={italic}]{F.}, p. 235, nous soulignons).
\end{teiQuote}
              
\end{teiCit}
              \teiP[xmlid={p10-5}]{L’« odeur douce », la mollesse des meubles et l’atmosphère voluptueuse qui caractérisent le cabinet renvoient aux attributs propres à l’hédonisme féminin qui lui est associé dans l’imaginaire du \teiHi[rend={small-caps}]{xix}\teiHi[rend={sup}]{e} siècle. L’adéquation progressive d’Antoine à cet espace hautement féminin annonce sa dévirilisation prochaine, laquelle est étroitement liée au vacillement de ses idées politiques. Le luxe du cabinet de toilette, les promesses d’aisance et de bien-être que font miroiter son mobilier et ses accessoires contaminent en effet les convictions politiques du personnage et l’incitent à changer de camp. L’« air musqué, fade et assoupi, qui traîne dans les cabinets de toilette » porte d’abord Antoine à se dire « que ces diables de riches “étaient bien heureux tout de même” » (\teiHi[rend={italic}]{F.}, p. 235) ; puis, alors qu’il examine sa prison de fortune, cette observation le mène à des réflexions peu démocratiques :}
              \begin{teiCit}[xmlid={cit04-3}]

                \begin{teiQuote}
Il se disait que jamais il n’aurait un pareil coin pour se débarbouiller. Le lavabo surtout l’intéressait ; ce n’était pas malin, pensait-il, de se tenir propre, avec tant de petits pots et tant de fioles. Cela le fit penser amèrement à sa vie manquée. L’idée lui vint qu’il avait peut-être fait fausse route ; on ne gagne rien à fréquenter les gueux ; il aurait dû ne pas faire le méchant et s’entendre avec les Rougon (\teiHi[rend={italic}]{F.}, p. 236).
\end{teiQuote}
              
\end{teiCit}
              \teiP[xmlid={p11-5}]{Ce renversement se développe tout au long de la scène jouée au cabinet de toilette : bien que Macquart rejette d’abord l’idée de trahir la bande insurrectionnelle, « les tiédeurs, les souplesses du divan continu\lbrack{}ent\rbrack{} à l’adoucir, à lui donner un regret vague » (\teiHi[rend={italic}]{F.}, p. 236). Il se fait ensuite la réflexion qu’« \lbrack{}a\rbrack{}près tout, les insurgés l’abandonn\lbrack{}ent\rbrack{}, \lbrack{}qu’\rbrack{}ils se f\lbrack{}on\rbrack{}t battre comme des imbéciles », avant de conclure « que la République \lbrack{}est\rbrack{} une duperie \lbrack{}…\rbrack{} \lbrack{}et que\rbrack{} \lbrack{}d\rbrack{}écidément, il aurait dû se vendre à la réaction » (\teiHi[rend={italic}]{F.}, p. 236). C’est l’espace dans lequel il est confiné qui le pousse à de telles conclusions antirépublicaines, puisque, « \lbrack{}e\rbrack{}n pensant cela, il lorgn\lbrack{}e\rbrack{} le lavabo, pris d’une grande envie d’aller se laver les mains avec une certaine poudre de savon contenue dans une boîte de cristal » (\teiHi[rend={italic}]{F.}, p. 236). Sous l’influence du luxueux décor, Antoine cède à son envie :}
              \begin{teiCit}[xmlid={cit05-3}]

                \begin{teiQuote}
La tentation devint trop forte ; Macquart s’installa devant le lavabo. Il se lava les mains, la figure ; il se coiffa, se parfuma, fit une toilette complète. Il usa de tous les flacons, de tous les savons, de toutes les poudres. Mais sa plus grande jouissance fut de s’essuyer avec les serviettes du maire ; elles étaient souples, épaisses. Il y plongea sa figure humide, y respira béatement toutes les senteurs de la richesse (\teiHi[rend={italic}]{F.}, p. 236).
\end{teiQuote}
              
\end{teiCit}
              \teiP[xmlid={p12-5}]{La description, qui procède par une énumération et une accumulation des pratiques de propreté toutes deux inscrites sous le signe de la totalisation (« toilette complète », usage de « tous » les produits) met en lumière le caractère excessif de ces ablutions. Cette démesure évoque les excès de richesse de l’Empire à naître, lesquels, semble-t-il, sont les arguments décisifs pour la conversion d’Antoine à la cause bonapartiste. Assiégé par les odeurs de la prospérité escomptée, le républicain opportuniste bascule dans la traîtrise politique : « Puis, quand il fut pommadé, quand il sentit bon de la tête aux pieds, il revint s’étendre sur le divan, rajeuni, porté aux idées conciliantes. Il éprouvait un mépris encore plus grand pour la République, depuis qu’il avait mis le nez dans les fioles de M. Garçonnet » (\teiHi[rend={italic}]{F.}, p. 236). Lors de ce nettoyage, Macquart se « lave » donc de ses convictions républicaines, tandis que son aspect prétendument viril est évincé de la description : d’ouvrier solide et patriote, du moins aux dires de Silvère, Antoine succombe dans ce lieu à sa nature voluptueuse et lâche, ce qui le place aux antipodes de l’idéal viril masculin dans la pensée de Zola.}
              \teiP[xmlid={p13-5}]{Dans l’économie du roman, le cabinet de toilette devient ainsi le lieu où se révèle la véritable nature du faux républicain, puisque c’est là que celui-ci peut goûter au luxe bourgeois qui détermine chez lui la nécessité de la trahison. En situant dans l’espace domestique un changement de camp aussi décisif, l’auteur souligne combien, sous Napoléon \teiHi[rend={small-caps}]{III}, les intérêts privés régissent la vie publique. Comme le cabinet de toilette est caractérisé par les mêmes attributs que ceux que Zola déplore chez les représentants du Second Empire (excès de luxe, de jouissance\teiNote[xmlid={ftn50-4},n={17},place={foot}]{\teiP[xmlid={p-197}]{À propos de ces motifs et des liens qui les unissent sous la plume de Zola, voir Ye Young chun, « La “dépense” chez Zola », \teiHi[rend={italic}]{Littérature}, vol. 1, n\teiHi[rend={sup}]{o} 153, 2009, p. 79-84.}} et féminité dévirilisante\teiNote[xmlid={ftn51-4},n={18},place={foot}]{\teiP[xmlid={p-198}]{À ce sujet, voir Nadia Beaudoin, « Émile Zola et la décadence… », art. cité, p. 76.}}, notamment), il est également l’endroit idéal où consacrer, dans l’économie romanesque, l’avènement du régime : c’est en effet la trahison de Macquart qui détermine, à Plassans, la réussite du guet-apens antirépublicain orchestré par les Rougon. C’est aussi lors d’une scène jouée dans le cabinet de toilette que sera annoncée la chute imminente de l’Empire, dans le pénultième volume du cycle ; mais, avant que d’en étudier sa fin, il convient d’observer l’usage que font de cette pièce les artisans les plus actifs du régime bonapartiste à son apogée.}
            
\end{teiDiv}
            \begin{teiDiv}[type={section1},xmlid={div02-4}]

              \teiHead[xmlid={le-cabinet-de-toilette-dans-son-excellence-eugene-rougon-incursion-dans-le-deshabille-de-l-a-politique-001}]{Le cabinet de toilette dans \teiHi[rend={italic}]{Son Excellence Eugène Rougon} : incursion dans « le déshabillé de l\lbrack{}a\rbrack{} politique »}
              \teiP[xmlid={p14-4}]{\teiHi[rend={italic}]{Son Excellence Eugène Rougon} met en scène les ascensions et les chutes du ministre Eugène Rougon, notamment au gré des luttes auxquelles se livrent ce personnage et une jeune intrigante en puissance, Clorinde Balbi, sous le règne du prince-président\teiNote[xmlid={ftn52-4},n={19},place={foot}]{\teiP[xmlid={p-199}]{Zola, \teiHi[rend={italic}]{Son Excellence Eugène Rougon} \lbrack{}1876\rbrack{}, \teiHi[rend={italic}]{Œuvres complètes, t. VII. La République en marche. 1875-1876}, Paris, Nouveau Monde, 2003, p. 236. Désormais, les références à cette œuvre seront indiquées entre parenthèses dans le texte par le sigle \teiHi[rend={italic}]{SE}.}}. Dans ce roman, Zola prolonge l’association mise en place dans \teiHi[rend={italic}]{La Fortune} entre la toilette des représentants de l’Empire et les tares de ce régime, cette fois dans le but d’en critiquer les députés. C’est en outre au cabinet de toilette que Clorinde mène ses batailles les plus décisives, ce qui permet au texte de dépeindre, en en montrant les limites, la place de la femme en politique sous le Second Empire.}
              \teiP[xmlid={p15-3}]{Lorsque les personnages du roman débattent de graves enjeux diplomatiques, ils sont fréquemment occupés à de minutieuses pratiques hygiéniques et cosmétiques. Ainsi d’Eugène Rougon qui, assailli dans son bureau par l’un de ses protégés, va se rafraîchir à la cuvette dissimulée dans la pièce, tandis qu’un autre de ses collègues, M. d’Escorailles, « ayant fini de classer sa correspondance, \lbrack{}tire\rbrack{} de sa poche une petite lime à manche d’écaille et se travaill\lbrack{}e\rbrack{} les ongles, délicatement » (\teiHi[rend={italic}]{SE}, p. 364). Rougon, dont l’influence est sollicitée par nombre de partis, affirme même qu’« \lbrack{}o\rbrack{}n ne le laiss\lbrack{}e\rbrack{} plus respirer ; \lbrack{}et que\rbrack{} la veille encore, on l’avait relancé jusque dans son cabinet de toilette, pendant qu’il se faisait la barbe » (\teiHi[rend={italic}]{SE}, p. 373), une protestation qui montre bien la proximité des activités politiques du ministre avec celles, plus intimes, liées au cabinet de toilette. Il semble pourtant que ce soit Rougon lui-même qui provoque cette insertion des affaires publiques dans sa vie privée. C’est en effet lors de sa toilette qu’à Niort, il se renseigne sur la préfecture : « Tout en se débarbouillant, le ministre demandait au préfet des détails sur le pays, les histoires des fonctionnaires, les besoins des uns, les vanités des autres » (\teiHi[rend={italic}]{SE}, p. 384). Les gestes de toilette sont ainsi toujours associés aux enjeux politiques – lesquels sont étroitement liés aux intérêts privés –, ce qui permet à Zola de remettre en question le sérieux des artisans de la politique impériale. Ceci se révèle clairement lorsque la narration expose leur intérêt exacerbé pour les questions hygiéniques à la fin du récit. Lors des ablutions minutieuses de M. de Combelot dans le couloir aux lavabos du Corps législatif, le désintérêt de celui-ci pour la bibliothèque adjacente est opposé à son souci pour sa toilette selon une juxtaposition narrative moqueuse :}
              \begin{teiCit}[xmlid={cit06-3}]

                \begin{teiQuote}
M. de Combelot, les mains plongées au fond d’une grande cuvette, les y frottait doucement, en souriant à leur blancheur. \lbrack{}…\rbrack{} Et il prit le temps de s’éponger longuement les mains, à l’aide d’une serviette chaude qu’il remit ensuite dans l’étuve, aux portes de cuivre. Même il alla, à l’extrémité du couloir, devant une haute glace, peigner sa belle barbe noire, avec un petit peigne de poche.\teiLb{}La bibliothèque était vide. Les livres dormaient dans leurs casiers de chêne ; toutes nues, les deux grandes tables étalaient la sévérité de leurs tapis verts ; aux bras des fauteuils, rangés en bon ordre, les pupitres mécaniques se repliaient, gris d’une légère poussière. Et, au milieu de ce recueillement, dans l’abandon de la galerie où traînait une odeur de papiers, M. La Rouquette dit tout haut, en faisant claquer la porte :\teiLb{}« Il n’y a jamais personne là-dedans ! » (\teiHi[rend={italic}]{SE}, p. 449)
\end{teiQuote}
              
\end{teiCit}
              \teiP[xmlid={p16-3}]{En soulignant l’importance que les hommes politiques accordent à la propreté corporelle au détriment du savoir juridique, cette scène expose leur vanité et les discrédite. Le lexique intime convoqué pour décrire la bibliothèque dévoile de surcroît l’ambiguïté qui règne dans l’espace ministériel entre le public et le privé. Dans ce lieu, les livres « dorm\lbrack{}ent\rbrack{} » ; les tables sont « toutes nues » et elles « étal\lbrack{}ent\rbrack{} » – geste langoureux – leurs tapis verts ; les pupitres sont « aux bras » des fauteuils, le tout dans une atmosphère de « recueillement » et « d’abandon » correspondant mieux à des lieux comme le boudoir ou la chambre qu’à un espace de réflexion politique ou intellectuelle. Malgré le caractère accueillant de ce lieu d’intelligence et de savoir, qui s’« offre » littéralement à tous les membres du Corps législatif, ce dernier reste vide : c’est dire le désintérêt des artisans de la politique impériale à l’égard de ces deux sphères, sur lesquelles M. La Rouquette va jusqu’à « claquer la porte ». Par cette description moqueuse des espaces gouvernementaux et de leur (in)occupation par les députés, le roman critique les dérives du pouvoir sous Napoléon \teiHi[rend={small-caps}]{III}, dont les représentants, aveugles aux affaires importantes, ne consultent que leur intérêt personnel pour décider du sort du pays.}
              \teiP[xmlid={p17-3}]{La rencontre entre vie publique et vie privée mise en lumière dans ce passage est essentielle à l’implication de Clorinde en politique. À l’inverse de la bibliothèque déserte du Corps législatif, le cabinet de toilette de la jeune femme, qu’elle instrumentalise à des fins diplomatiques, est fort couru.}
              \teiP[xmlid={p18-2}]{Comme le note Michelle Perrot, « la Révolution \lbrack{}…\rbrack{} différencie les rôles sexuels en opposant hommes politiques et femmes domestiques\teiNote[xmlid={ftn53-5},n={20},place={foot}]{\teiP[xmlid={p-200}]{Michelle Perrot, « Avant et ailleurs », Philippe Ariès et Georges Duby (dir.), \teiHi[rend={italic}]{Histoire de la vie privée}, \teiHi[rend={italic}]{op. cit.}, p. 15.}} ». Les femmes sont ainsi non seulement évincées de la sphère publique, mais aussi écartées de la vie politique. Dans le roman de 1876, où c’est une femme qui se fait l’instigatrice des luttes de pouvoir, il pouvait donc paraître naturel que Clorinde fasse du privé son champ de bataille, de son intimité un atout, de son corps une arme et qu’elle se « garde comme un argument irrésistible » (\teiHi[rend={italic}]{SE}, p. 342), afin de mener sa campagne. Dans ce contexte, le cabinet de toilette acquiert une importance nouvelle : il devient la tribune d’où le personnage féminin peut faire valoir ses « arguments », à défaut de se faire entendre à la Chambre.}
              \teiP[xmlid={p19-2}]{La première description de l’ameublement du cabinet de toilette de Clorinde révèle la nature incongrue de cette pièce. Là où l’on s’attendrait à trouver une table de toilette, des divans, ou même une baignoire, ce sont en effet plutôt « un grand bureau d’acajou déverni », « un fauteuil de cuir » et un « cartonnier » (\teiHi[rend={italic}]{SE}, p. 298) qui inaugurent la liste des meubles et accessoires de la pièce. Le texte mentionne que « \lbrack{}d\rbrack{}es paperasses \lbrack{}y\rbrack{} traîn\lbrack{}ent\rbrack{} sous une épaisse couche de poussière » (\teiHi[rend={italic}]{SE}, p. 298), l’accumulation de dossiers et de saleté montrant bien la longue utilisation de cet espace à des fins bureaucratiques plutôt que cosmétiques ou hygiéniques\teiNote[xmlid={ftn54-5},n={21},place={foot}]{\teiP[xmlid={p-201}]{La malpropreté du cabinet de toilette de Clorinde fait écho à sa malpropreté physique, laquelle participe à son érotisation dans le texte. Comme le remarque Véronique Cnockaert, « \lbrack{}p\rbrack{}our éviter la fixité, Zola favorise dans chacun de ses détails portrographiques l’expression d’un certain désordre, l’inscription d’un tremblement dans le convenu. L’érotisation du fragment qui rejaillit sur l’ensemble de la figure passe par le débraillé : les déesses zoliennes ont souvent des cheveux en broussailles et des allures de batteuses de pavé » (« Du marbre et du chiffonné. Propos sur la beauté dans l’œuvre d’Émile Zola », \teiHi[rend={italic}]{Études françaises}, vol. 39, n\teiHi[rend={sup}]{o} 2, 2003, p. 38).}}. Le caractère suspect des activités menées en ce lieu est lui aussi souligné, car le cabinet est comparé aux bureaux d’« un huissier louche » (\teiHi[rend={italic}]{SE}, p. 298). Il s’agit là du premier indice quant à la nature de la participation de Clorinde à la vie politique, une participation souterraine rapidement confirmée dans la suite de la scène : lorsque la jeune femme reçoit Rougon, elle est affairée à la rédaction de lettres destinées à de grands magistrats italiens (\teiHi[rend={italic}]{SE}, p. 300). Le cabinet de toilette est donc présenté comme un espace hautement politisé. Pour autant, la politique n’y supplante pas l’intime. Au contraire, comme le remarque Rougon, missives, dossiers et encriers se mêlent confusément aux jupes, aux savons et aux corsets :}
              \begin{teiCit}[xmlid={cit07-2}]

                \begin{teiQuote}
Sur le cartonnier, il lut, comme chez les hommes d’affaires : \teiHi[rend={italic}]{Quittances }; \teiHi[rend={italic}]{Lettres à classer} ; \teiHi[rend={italic}]{Dossiers A}. Il sourit en apercevant, au milieu des paperasses du bureau, un corset qui traînait, usé, craqué à la taille. Il y avait encore un savon dans la coquille de l’encrier, et des bouts de satin bleu à terre, les rognures de quelque raccommodage de jupe, qu’on avait oublié de balayer (\teiHi[rend={italic}]{SE}, p. 301, l’auteur souligne).
\end{teiQuote}
              
\end{teiCit}
              \teiP[xmlid={p20-2}]{Les accessoires nécessaires à l’entretien du corps apparaissent ainsi aux côtés des documents. Leur cohabitation met en lumière le seul pouvoir effectif dont est dotée la femme sous le régime impérial, du moins dans la pensée de Zola : son corps.}
              \teiP[xmlid={p21-2}]{Clorinde est bien consciente du pouvoir de séduction qu’elle peut tirer de son physique. Elle utilise stratégiquement son cabinet de toilette – lieu de sociabilité où les dévoilements sont permis – pour mener sa campagne. Elle y reçoit les hommes influents afin d’obtenir les plus récentes informations diplomatiques autant qu’elle en profite pour exposer, par vanité, son implication étendue dans la politique étrangère (\teiHi[rend={italic}]{SE}, p. 412). En échange des secrets diplomatiques qu’elle soutire à ses proches, la jeune femme doit cependant montrer son corps. Lorsque Clorinde, à peine sortie de la baignoire, accueille par exemple dans son intimité M. de Reuthlinguer, l’un de ses informateurs, la narration le montre dans une posture singulière qui suggère le voyeurisme (toléré) auquel se livre le visiteur lors de son passage au cabinet : « Elle restait à demi nue, la chemise glissée des épaules. Le baron, de ses lèvres pâles, trouva un sourire d’indulgence ; et il se tint debout près d’elle, les yeux froids et clairs, penché dans un salut d’extrême politesse » (\teiHi[rend={italic}]{SE}, p. 415). Clorinde ne se formalise de sa nudité qu’à sa sortie du cabinet :}
              \begin{teiCit}[xmlid={cit08-2}]

                \begin{teiQuote}
Quand il la quitta \lbrack{}…\rbrack{} \lbrack{}e\rbrack{}lle l’accompagna jusqu’à la porte. Mais, tout d’un coup, elle croisa ses bras sur sa poitrine, très rouge, en s’écriant :\teiLb{}« Ah ! bien moi qui m’en vais comme ça avec vous ! » (\teiHi[rend={italic}]{SE}, p. 415)
\end{teiQuote}
              
\end{teiCit}
              \teiP[xmlid={p22-2}]{Le cabinet de toilette, s’il est d’abord présenté comme un bureau politique, retrouve dans \teiHi[rend={italic}]{Son Excellence Eugène Rougon} la fonction de théâtre du corps qu’il possédait sous l’Ancien Régime : Clorinde s’y livre à une « toilette de prétexte \lbrack{}…\rbrack{} qui favorise le développement de mille attraits cachés ou non encore aperçus\teiNote[xmlid={ftn55-5},n={22},place={foot}]{\teiP[xmlid={p-202}]{Philippe Perrot, \teiHi[rend={italic}]{Le Travail des apparences}, \teiHi[rend={italic}]{op. cit.}, p. 45.}} », ce qui lui permet d’entretenir, avec les hommes de pouvoir de son entourage, des relations d’intimité (corporelle et sexuelle) qui seraient en tout autre lieu considérées inconvenantes – et qui sont d’ailleurs remises en cause dès que les personnages en sortent. La représentation du cabinet de toilette et de l’usage qu’en fait le personnage féminin met ainsi en lumière son ingéniosité tout en révélant les limites spatiales et sociales que lui impose sa condition.}
              \teiP[xmlid={p23-2}]{Puisque Clorinde ne peut agir que depuis la sphère intime, le texte rappelle qu’elle ne joue pas sur le même terrain que ses rivaux masculins. L’intrigante parvient en effet à influencer le jeu politique, mais seulement d’une manière occulte. Aussi savamment échafaudés que soient les projets tramés dans son cabinet de toilette, la jeune femme ne peut aspirer à une véritable position dans la sphère publique mise en scène par Zola. Henri Mitterand rappelle à cet égard que}
              \begin{teiCit}[xmlid={cit09-2}]

                \begin{teiQuote}
\lbrack{}l\rbrack{}a vérité politique voulait \lbrack{}…\rbrack{} que Clorinde Balbi ne fit pas campagne pour elle-même, mais pour un \lbrack{}…\rbrack{} homme. Car une femme d’intrigue, en 1860, peut tout au plus accéder au lit de l’Empereur et devenir, quelque temps, une favorite ; elle peut, par les protections acquises et par son intuition des objectifs, des enjeux et des chances, orienter dans un sens favorable la carrière de l’homme qu’elle soutient ou qu’elle pousse, ou la cause à laquelle elle s’est vouée. \lbrack{}…\rbrack{} Mais il n’est pas question qu’elle s’élève au-dessus de ce rôle d’agent, d’entremetteuse des intérêts politiques\teiNote[xmlid={ftn56-5},n={23},place={foot}]{\teiP[xmlid={p-203}]{Henri Mitterand, \teiHi[rend={italic}]{Le Regard et le signe. Poétique du roman réaliste et naturaliste}, Paris, PUF, 1987, p. 197.}}.
\end{teiQuote}
              
\end{teiCit}
              \teiP[xmlid={p24-2}]{Clorinde n’en est pas moins puissante. Il semble même que le rapport qu’elle entretient avec l’intime étend son influence. Lorsqu’elle s’attaque à la conquête de Rougon, les résultats de cette séduction sont par exemple souvent décrits au moyen de qualificatifs liés à l’usage qu’elle fait du cabinet de toilette. Puisque c’est là qu’elle prend soin de son corps, la pensée de la jeune femme est notamment rapprochée des effets délassants du bain dans l’imaginaire du ministre :}
              \begin{teiCit}[xmlid={cit10}]

                \begin{teiQuote}
Il pensait à elle, comme s’il l’avait déjà quittée depuis longtemps ; et, les yeux sur le parquet, il descendait dans des pensées à demi formulées, fort douces, dont il goûtait le chatouillement intérieur. \teiHi[rend={italic}]{Il lui semblait sortir d’un bain tiède, avec une langueur de membres délicieuse}. \teiHi[rend={italic}]{Une odeur particulière, d’une rudesse presque sucrée, le pénétrait. Cela lui aurait paru bon, de se coucher sur un des canapés et de s’y endormir, dans cette odeur} (\teiHi[rend={italic}]{SE}, p. 270, nous soulignons).
\end{teiQuote}
              
\end{teiCit}
              \teiP[xmlid={p25-2}]{Cette réflexion, où perce le désir sexuel de Rougon, a sur lui des effets pathologiques, ce qui souligne la puissance reconnue au personnage féminin (et en particulier à son corps) dans le texte : « Rougon se trouva sur l’avenue des Champs-Élysées, il demeura un moment étourdi à respirer l’air frais qui soufflait des hauteurs de l’Arc de Triomphe. \lbrack{}…\rbrack{} Il venait d’avoir comme un coup de sang, il se passait les mains sur la face » (\teiHi[rend={italic}]{SE}, p. 274). Un phénomène similaire se produit lorsque le ministre séjourne à Compiègne, avec Clorinde et son mari. Lors d’un banquet donné par l’empereur, la salle à manger est comparée à une immersion dans un bain exhalant les mêmes effluves que ceux privilégiés par Clorinde dans l’intimité\teiNote[xmlid={ftn57-5},n={24},place={foot}]{\teiP[xmlid={p-204}]{L’indiscrétion d’Auguste Jobelin, dans la deuxième partie du roman, dévoile les parfums privilégiés par la jeune femme lors de ses ablutions, qui sont les mêmes que ceux employés pour la préparation du gibier à Compiègne : « À ce moment, un bruit étrange vint du cabinet de toilette, une sorte de barbotement doux et continu. Le colonel se hâta d’aller voir, et il trouva Auguste très intéressé par la baignoire qu’Antonia avait oublié de vider. \teiHi[rend={italic}]{Des ronds de citron, dont Clorinde s’était servie pour ses ongles, flottaient} » (\teiHi[rend={italic}]{SE}, p. 419, nous soulignons).}} : « C’était une approche presque tendre, une arrivée gourmande dans un milieu de luxe, de lumière et de tiédeur, \teiHi[rend={italic}]{comme un bain sensuel où les odeurs musquées des toilettes se mêlaient à un léger fumet de gibier}, \teiHi[rend={italic}]{relevé d’un filet de citron} » (\teiHi[rend={italic}]{SE}, p. 323, nous soulignons). Ces émanations affectent singulièrement Rougon, qui est d’autant plus piqué de jalousie que Clorinde se prête à des jeux de séduction irritants avec l’un de ses rivaux lors de ce repas (\teiHi[rend={italic}]{SE}, p. 325). Le texte, grâce aux rapprochements qu’il établit entre les effets de la présence de Clorinde et ceux d’un passage au cabinet de toilette, fait ainsi de l’intrigante une adversaire redoutable, capable d’obtenir les plus précieux secrets diplomatiques et d’amollir les hommes d’influence sans même avoir accès à leur monde, par le seul emploi de son corps et de son espace intime comme d’un instrument politique. Il en va tout autrement de la représentation que fait Zola des subterfuges cosmétiques de l’empereur et de sa visite au cabinet de toilette dans \teiHi[rend={italic}]{La Débâcle}, qui visent tous deux à signaler l’étiolement de son pouvoir à la veille de la chute de l’Empire.}
            
\end{teiDiv}
            \begin{teiDiv}[type={section1},xmlid={div03-3}]

              \teiHead[xmlid={portrait-de-lempereur-au-cabinet-de-toilette-dans-la-debacle-001}]{Portrait de l’empereur au cabinet de toilette dans \teiHi[rend={italic}]{La Débâcle}}
              \teiP[xmlid={p26}]{Lorsque Silvère est mis à mort, au terme de \teiHi[rend={italic}]{La Fortune}, la narration indique que cette scène force « \lbrack{}u\rbrack{}n petit bourgeois propret \lbrack{}à\rbrack{} se retir\lbrack{}er\rbrack{}, en déclarant que s’il restait davantage, ça l’empêcherait de dîner » (\teiHi[rend={italic}]{F.}, p. 266). Sous la plume de Zola, la « propreté » de ce figurant, à l’instar de la corruption au milieu des savons d’Antoine Macquart et des frais de toilette récurrents des ministres du Second Empire, est un repère signalétique de l’ère d’hypocrisie vaniteuse dans laquelle est entrée la France après la victoire bonapartiste. En ce siècle où chacun aspire à se rendre « accessible à la morale de la propreté\teiNote[xmlid={ftn58-5},n={25},place={foot}]{\teiP[xmlid={p-205}]{Alain Corbin, « Coulisses », Philippe Ariès et Georges Duby (dir.), \teiHi[rend={italic}]{Histoire de la vie privée}, \teiHi[rend={italic}]{op. cit.}, p. 414.}} », l’écrivain révèle en effet dans tous les romans de son cycle la « boue humaine\teiNote[xmlid={ftn59-5},n={26},place={foot}]{\teiP[xmlid={p-206}]{Zola, \teiHi[rend={italic}]{L’Argent} \lbrack{}1891\rbrack{}, Paris, Gallimard, 1980, p. 211.}} » qui se cache sous les belles apparences des représentants de l’Empire. Quelque vingt ans après la parution du premier volume de la série, l’écrivain emploie toujours la même stratégie : lorsqu’il représente la chute du règne de Napoléon \teiHi[rend={small-caps}]{III}, il le fait au moyen de motifs hygiéniques. Dans \teiHi[rend={italic}]{La Débâcle}, qui relate la défaite française à Sédan, ce sont tant les vaines pratiques cosmétiques de l’empereur que sa souffrance intime, dévoilée au cabinet de toilette, qui annoncent la déchéance du régime.}
              \teiP[xmlid={p27}]{Le portrait donné de Napoléon \teiHi[rend={small-caps}]{III} au champ de bataille insiste sur le maquillage qu’il arbore dans l’espoir de cacher son piètre état de santé\teiNote[xmlid={ftn60-5},n={27},place={foot}]{\teiP[xmlid={p-207}]{La mention de ce fard a d’ailleurs fait controverse lors de la publication du roman en raison de discussions au sujet de sa véracité historique. À ce propos, voir Bernadette C. Lintz, « L’Empereur fardé : Napoléon \teiHi[rend={small-caps}]{III} des “Châtiments” à “La Débâcle” », \teiHi[rend={italic}]{Nineteenth-Century French Studies}, vol. 35, n\teiHi[rend={sup}]{o} 3-4, printemps-été 2007, p. 610. En réponse à ses détracteurs, Zola répond : « Moi, je le trouve superbe, ce fard, digne d’un des grands héros de Shakespeare, haussant la figure de Napoléon \teiHi[rend={small-caps}]{III} à une mélancolie tragique d’une infinie grandeur » (\teiHi[rend={italic}]{Les Rougon-Macquart}, Henri Mitterand éd., Paris, Gallimard, 1961, t. V, p. 1455).}} :}
              \begin{teiCit}[xmlid={cit11}]

                \begin{teiQuote}
C’était bien Napoléon \teiHi[rend={small-caps}]{III}, qui lui apparaissait plus grand, à cheval, et les moustaches si fortement cirées, les joues si colorées, qu’il le jugea tout de suite rajeuni, fardé comme un acteur. Sûrement, il s’était fait peindre, pour ne pas promener, parmi son armée, l’effroi de son masque blême, décomposé par la souffrance, au nez aminci, aux yeux troubles. Et, averti dès cinq heures qu’on se battait à Bazeilles, il était venu, de son air silencieux et morne de fantôme, aux chairs ravivées de vermillon\teiNote[xmlid={ftn61-5},n={28},place={foot}]{\teiP[xmlid={p-208}]{Zola, \teiHi[rend={italic}]{La Débâcle} \lbrack{}1892\rbrack{}, Paris, Gallimard, 1984, p. 214-215. Désormais, les références à cette œuvre seront indiquées entre parenthèses dans le texte par le sigle \teiHi[rend={italic}]{D.}}}.
\end{teiQuote}
              
\end{teiCit}
              \teiP[xmlid={p28}]{Dans ce passage, la répétition de l’adverbe « si » souligne l’intensité de l’usage que fait l’empereur des produits de toilette censés lui redonner une apparence de santé. Cependant, le lexique employé révèle le caractère vain de cette mise en scène cosmétique alors que sonne le glas du régime : l’empereur est comparé à « un acteur », ce qui souligne la théâtralité de son maquillage et la nature factice de la vie qu’il affecte, une vie si fausse qu’il est en outre décrit comme « un fantôme », animé par le seul effet du fard. La disparition progressive de ce lourd maquillage dévoile son déclin : alors qu’on le présente d’abord « peign\lbrack{}é\rbrack{} » et « bichonn\lbrack{}é\rbrack{}, avec toutes sortes d’histoires sur la figure » (\teiHi[rend={italic}]{D.}, p. 247), les traces de sa toilette minutieuse s’effacent rapidement de son visage à la suite des défaites subies sur le champ de bataille : « Sous la sueur d’angoisse de cette marche au travers de la défaite, le fard s’en était allé des joues, les moustaches cirées s’étaient amollies, pendantes, la face terreuse avait pris l’hébétement douloureux d’une agonie » (\teiHi[rend={italic}]{D.}, p. 260). Cette déchéance cosmétique évoque le déclin du régime impérial. Dans la pensée zolienne, le règne de Napoléon \teiHi[rend={small-caps}]{III} est en effet lui aussi caractérisé par sa démesure et son caractère trompeur : montrer que ces deux attributs faillent à son souverain, c’est donc insinuer qu’ils failliront bientôt aussi à l’Empire. Le texte souligne par la même occasion, en signant l’échec des « histoires » peintes sur le visage de l’empereur, le caractère dérisoire des apparences menteuses prônées sous son règne, au moment où celui-ci connaît une fin violente. Cet échec des apparences est plus spécifiquement mis en lumière, dans l’économie romanesque, au sein de la pièce vouée au paraître : le cabinet de toilette.}
              \teiP[xmlid={p29}]{Napoléon \teiHi[rend={small-caps}]{III} souffre dans le roman, outre d’une affection des voies urinaires\teiNote[xmlid={ftn62-5},n={29},place={foot}]{\teiP[xmlid={p-209}]{« Décidément, il a une sale pierre dans son sac », affirme le major Bouroche (\teiHi[rend={italic}]{D.}, p. 84).}}, d’une grave crise de dysenterie\teiNote[xmlid={ftn63-5},n={30},place={foot}]{\teiP[xmlid={p-210}]{« La présence de sa mère et de sa femme l’empêchait de désigner plus clairement la dysenterie dont l’empereur souffrait depuis le Chesne et qui le forçait à s’arrêter ainsi dans les fermes, le long de la route » (\teiHi[rend={italic}]{D.}, p. 183).}}. C’est au cabinet de toilette qu’il se réfugie pour laisser libre cours à sa souffrance, comme le relate Rose, la fille de la concierge de la sous-préfecture où séjourne l’empereur :}
              \begin{teiCit}[xmlid={cit12}]

                \begin{teiQuote}
Non, vous ne pouvez pas savoir ce qu’il souffre ! ... Imaginez-vous qu’hier soir j’étais montée pour aider à donner du linge. Alors voilà qu’en passant dans la pièce qui touche au cabinet de toilette, j’ai entendu des gémissements, oh ! des gémissements, comme si quelqu’un était en train de mourir. Et je suis restée tremblante, le cœur glacé, en comprenant que c’était l’empereur… Il paraît qu’il a une maladie affreuse qui le force à crier ainsi. Quand il y a du monde, il se retient ; mais, dès qu’il est seul, c’est plus fort que sa volonté, il crie, il se plaint, à vous faire dresser les cheveux sur la tête (\teiHi[rend={italic}]{D.}, p. 246-247).
\end{teiQuote}
              
\end{teiCit}
              \teiP[xmlid={p30}]{Présenté dans ce passage comme un haut lieu de la vie privée, où l’empereur peut exprimer sa faiblesse sans craindre d’être entendu – à l’exception des oreilles indiscrètes des quelques femmes qui le servent à la sous-préfecture –, le cabinet de toilette, ainsi que le révèle le discours de Rose, n’offre cependant qu’une intimité relative au prince-président. Si la jeune fille prend connaissance de ses douleurs grâce à une indiscrétion de sa part (elle écoute à sa porte) et qu’elle ne tarde pas à en faire le compte rendu à Henriette « sans \lbrack{}même\rbrack{} attendre d’être questionnée » (\teiHi[rend={italic}]{D.}, p. 246), les souffrances de l’empereur sont en effet \teiHi[rend={italic}]{déjà }connues de tous, comme le montre la référence à la rumeur populaire contenue dans le syntagme « il paraît que ». C’est dire l’inefficacité de la mise en scène impériale lors du conflit de 1870, tant sur le plan individuel que collectif : le maquillage et les précautions de l’empereur à l’égard de ses souffrances intimes ne dupent personne et ses apparitions sur le champ de bataille se résument, pour reprendre la formule de Bernadette C. Lintz, à un « \lbrack{}à un\rbrack{} artifice qui ne fait que pointer le simulacre, le leurre et le néant\teiNote[xmlid={ftn64-4},n={31},place={foot}]{\teiP[xmlid={p-211}]{Bernadette C. Lintz, « L’Empereur fardé : Napoléon \teiHi[rend={small-caps}]{III} des “Châtiments” à “La Débâcle” », art. cité, p. 616.}} » de cette campagne militaire et du régime qui l’a lancée. Lors de ce passage au cabinet de toilette, corps public et corps privé se confondent donc, cette fois afin de révéler l’ampleur de leur corruption : « lieu où le corps de l’empereur se décompose, se vide, se liquéfie, rappelle Bernadette C. Lintz, ce “petit coin où il se cache pour crier” \lbrack{}…\rbrack{} est le lieu où s’opère symboliquement la liquidation d’un régime en train de se dissoudre sous l’effet de sa propre corruption, “dans la sanie et l’excrément\teiNote[xmlid={ftn65-4},n={32},place={foot}]{\teiP[xmlid={p-212}]{\teiHi[rend={italic}]{Ibid.}, p. 617-618. La citation est de Maurice Descotes, \teiHi[rend={italic}]{Le Personnage de Napoléon} \teiHi[rend={ small-caps italic}]{III}\teiHi[rend={italic}]{ dans }Les Rougon-Macquart, Paris, Minard, 1970, p. 57.}}”. » Dans l’espace romanesque même où triomphait auparavant une politique secrète du pouvoir fondée sur les apparences stratégiques et sur les exploitations théâtrales du corps, les masques, définitivement, tombent.}
              \teiP[xmlid={p31}]{Pièce au luxe excessif et aux attributs féminins dévirilisants ; espace de toilette qui tantôt discrédite les hommes de pouvoir qui abusent de son usage, tantôt souligne l’ingéniosité des femmes qui l’instrumentalisent ; lieu privilégié de la révélation des tares individuelles et sociales et de la mise en lumière des rouages secrets du régime fondé sur le mensonge, les intérêts privés et les apparences de Napoléon \teiHi[rend={small-caps}]{III}, le cabinet de toilette, dans la pensée de Zola, cristallise nombre d’enjeux propres à sa critique du Second Empire. Le régime en vigueur, au moment où Zola rédige son cycle, n’est pas pour autant présenté comme la voie du salut. L’avènement de la Troisième République, dans la \teiHi[rend={italic}]{Débâcle}, s’amorce en effet avec les mêmes pratiques hygiéniques dissimulatrices que celles ayant donné naissance à l’Empire : au lendemain de la Commune, « Maurice s’aperç\lbrack{}oit\rbrack{} que les rangs \lbrack{}des Communards\rbrack{} s’éclairciss\lbrack{}ent\rbrack{} » et que « des camarades fil\lbrack{}ent\rbrack{} sans bruit, \lbrack{}vont\rbrack{} \teiHi[rend={italic}]{se laver les mains}, mettre une blouse, dans la terreur des représailles » (\teiHi[rend={italic}]{D.}, p. 547, nous soulignons). Bien qu’il s’agisse d’une référence à un fait historique – au lendemain de la semaine sanglante, les Versaillais ont procédé à une inspection des mains des Communards afin d’y déceler des traces de poudre et de les arrêter pour leur implication dans l’insurrection, entraînant par là-même la nécessité de les tenir propres\teiNote[xmlid={ftn66-4},n={33},place={foot}]{\teiP[xmlid={p-213}]{Voir, notamment, la gravure « Souvenirs de la Commune. Inspection des mains des insurgés à Belleville », parue dans \teiHi[rend={italic}]{L’Univers illustré }du 16 septembre 1871 et présentée lors de l’Exposition \teiHi[rend={italic}]{La Commune de Paris}, Paris, Hôtel de ville, 18 mars-28 mai 2011.}} – et bien que cet épisode puisse sembler tenir du document, celui-ci s’inscrit dans la droite ligne de la description que fait Zola de l’avènement de l’Empire depuis les débuts de son cycle romanesque. En effet, dans \teiHi[rend={italic}]{La Curée} (1872), la dissimulation permise par l’hygiène et par les produits cosmétiques exemplifiait déjà les commencements de l’ère d’hypocrisie et de mensonge que représente le règne de Napoléon \teiHi[rend={small-caps}]{III} aux yeux du romancier :}
              \begin{teiCit}[xmlid={cit13}]

                \begin{teiQuote}
C’était l’heure où les aventuriers du 2 Décembre, après avoir payé leurs dettes, jetaient dans les égouts leurs bottes éculées, leurs redingotes blanchies aux coutures, rasaient leurs barbes de huit jours, et devenaient des hommes comme il faut. Saccard entrait enfin dans la bande, il se nettoyait les ongles et ne se lavait plus qu’avec des poudres et des parfums inestimables\teiNote[xmlid={ftn67-4},n={34},place={foot}]{\teiP[xmlid={p-214}]{Zola, \teiHi[rend={italic}]{La Curée} \lbrack{}1872\rbrack{}, Paris, Le Livre de Poche, 1996, p. 96.}}.
\end{teiQuote}
              
\end{teiCit}
              \teiP[xmlid={p32}]{Par l’itération des mêmes gestes de propreté trompeurs, le texte souligne les risques que l’histoire se répète après le règne de Napoléon \teiHi[rend={small-caps}]{III} : à la critique du Second Empire inscrite dans la représentation du cabinet de toilette et de son usage s’ajoute donc une seconde critique, moins visible, mais néanmoins signifiante, celle de la Troisième République.}
            
\end{teiDiv}
          
\end{teiElement}
        
\end{teiElement}
      
\end{teiElement}
      \begin{teiElement}[name={group},type={section1},xmlid={section1-002}]

        \teiHead[xmlid={discours-exhiber-le-secret-politique-001}]{Discours : exhiber le secret politique}
        \begin{teiElement}[name={group},type={chapter},data-page-title={Énonciation référentielle et révélation des secrets du pouvoir chez Pontis, Sirot et Lenet},data-page-authors={Cécile Leduc@@Sorbonne Université, STIH, EA 4509, Paris, France},data-include-href={../XML/chap\_6\_Enonciation\_referentielle\_Locus\_politicus.xml},xmlid={chapter-005}]

          \begin{teiElement}[name={front}]

            \begin{teiDiv}[type={titlePage},xmlid={titlepage-006}]

              \teiP[rend={title-main},xmlid={p-215}]{Énonciation référentielle et révélation des secrets du pouvoir chez Pontis, Sirot et Lenet}
              \teiP[rend={author-aut},xmlid={p-216}]{Cécile \teiHi[rend={small-caps}]{Leduc}}
              \teiP[rend={authority\_affiliation},xmlid={p-217}]{\teiHi[rend={small-caps}]{S}orbonne Université\teiHi[rend={small-caps}]{, }STIH, EA 4509, Paris, France}
            
\end{teiDiv}
          
\end{teiElement}
          \begin{teiElement}[name={body},xmlid={body-006}]

            \teiP[xmlid={p1-6}]{L’époque moderne est le théâtre d’un « basculement inédit », qui substitue à « une culture médiévale du secret – entendu comme mystère – \lbrack{}…\rbrack{} des pratiques politiques du secret devenues techniques de cryptage et d’occultation\teiNote[n={1},place={foot},xmlid={ftn1-6}]{\teiP[xmlid={p-218}]{Bernard Darbord et Agnès Delage, « Présentation », dans \teiHi[rend={italic}]{idem} (dir.), \teiHi[rend={italic}]{Le Partage du secret. Cultures du dévoilement et de l’occultation en Europe, du Moyen-Âge à l’époque moderne}, Paris, Armand Colin, 2013, p. 7-12, ici p. 8.}} » : la lecture des Mémoires du \teiHi[rend={small-caps}]{xvii}\teiHi[rend={sup}]{e} siècle nous permet d’observer ce mouvement de près. Le secret y est en effet représenté comme « technique de contrôle et de dissimulation\teiNote[n={2},place={foot},xmlid={ftn23}]{\teiP[xmlid={p-219}]{Introduction à la « Partie III : les spectacles du secret », \teiHi[rend={italic}]{ibid.}, p. 329-332, ici p. 329.}} » par un pouvoir évoluant vers l’absolutisme\teiNote[n={3},place={foot},xmlid={ftn24}]{\teiP[xmlid={p-220}]{Voir la « Présentation » de Christian Renoux à la partie « Fabriques du secret », dans Laure Depretto, Christian Renoux, Christophe Speroni \teiHi[rend={italic}]{et al.} (dir.), \teiHi[rend={italic}]{Cultures du secret à l’époque moderne. Raisons, espaces, paradoxes, fabriques}, Paris, Classiques Garnier, 2023.}}. Les mémorialistes, qui accèdent aux secrets du pouvoir et qui connaissent le pouvoir de ces secrets, les dévoilent, s’inscrivant ainsi en faux de l’Histoire officielle\teiNote[n={4},place={foot},xmlid={ftn25}]{\teiP[xmlid={p-221}]{Frédéric Charbonneau fait du refus de l’Histoire officielle au profit d’une Histoire secrète l’origine commune de ce corpus hétérogène (\teiHi[rend={italic}]{Les Silences de l’Histoire. Les Mémoires français du }\teiHi[rend={small-caps italic}]{xvii}\teiHi[rend={italic sup}]{e}\teiHi[rend={italic}]{ siècle}, Paris, Hermann, 2016).}}. Cette entreprise de dévoilement sert une critique des ministres qui sont représentés comme les véritables détenteurs du savoir et du pouvoir.}
            \teiP[xmlid={p2-6}]{Il s’agira ici d’analyser les textes de trois mémorialistes, victimes ou complices des secrets et donc – selon l’étymologie du terme – inclus ou exclus du pouvoir : Claude de Létouf, baron de Sirot, disgracié par Mazarin ; Louis de Pontis, exclu du pouvoir par Richelieu mais proche de Louis \teiHi[rend={small-caps}]{XIII }; Pierre Lenet, au cœur des négociations avec Mazarin durant la Fronde. Le secret ne prenant vie et forme que dans les discours, nous analyserons uniquement des scènes de discours rapportés entre le mémorialiste et le ministre ou le roi. Ce qui est par essence un non-savoir constitue paradoxalement dans ces scènes le centre du savoir du mémorialiste et la pierre de touche de son plaidoyer.}
            \teiP[xmlid={p3-6}]{C’est l’énonciation spécifique au régime factuel, appelée référentielle, qui retiendra notre attention. Contrairement aux textes fictionnels, l’énonciateur des Mémoires est un individu réel ayant vécu ce qu’il raconte. Le terme \teiHi[rend={italic}]{narrateur }ne sera pas employé : forgé pour la fiction, il gomme justement l’identité entre la personne ayant historiquement existé et l’instance énonciative. Ce récit d’un passé vécu implique la scission de l’instance énonciative entre un \teiHi[rend={italic}]{je} narrant et un \teiHi[rend={italic}]{je }narré\teiNote[n={5},place={foot},xmlid={ftn26}]{\teiP[xmlid={p-222}]{Pour une étude précise de ce dispositif énonciatif dans les Mémoires, voir Delphine de Garidel, \teiHi[rend={italic}]{Poétique de Saint-Simon. Cours et détours du récit historique dans} Les Mémoires, Paris, Honoré Champion, 2005, et Emmanuèle Lesne, \teiHi[rend={italic}]{La Poétique des Mémoires (1650-1685)}, Paris, Honoré Champion, 1996.}}, sans qu’il soit toujours possible de mesurer l’influence du premier sur le second. Or l’étude du secret dans les Mémoires s’avère particulièrement adaptée à l’analyse de cette scission : le secret peut impliquer un décalage de connaissances entre les deux instances. Le \teiHi[rend={italic}]{je }narrant éclairerait de son savoir rétrospectif les secrets de l’expérience du \teiHi[rend={italic}]{je }narré. Même si les Mémoires de Pontis n’ont pas été rédigés par Pontis lui-même, ils reprennent la même configuration énonciative\teiNote[n={6},place={foot},xmlid={ftn27}]{\teiP[xmlid={p-223}]{On est sur ce point en désaccord avec René Démoris qui postule un « narrateur doublement fictif » dans ces Mémoires\teiHi[rend={italic}]{ }: un « Pontis-narrateur » et un « Pontis mondain », mis en scène par le premier, tous deux différant du « Pontis de la retraite », le « Pontis réel » (\teiHi[rend={italic}]{Le Roman à la première personne du classicisme aux Lumières}, Genève, Droz, 2002, p. 107). La référentialité énonciative annule à nos yeux l’instance « Pontis-narrateur » ; nous considérons que l’origine énonciative des Mémoires est bien le « Pontis réel » racontant son existence à Du Fossé, et qui se scinde, comme toute énonciation référentielle, entre le \teiHi[rend={italic}]{je }narrant (correspondant au « Pontis de la retraite ») et le \teiHi[rend={italic}]{je }narré (« le Pontis mondain »).}}, puisque Thomas Du Fossé n’a fait que retranscrire les propos de Pontis qui « parlait comme avec lui-même devant cet ami dont il ne sait pendant longtemps qu’il couche par écrit sa relation\teiNote[n={7},place={foot},xmlid={ftn28}]{\teiP[xmlid={p-224}]{Andrée Villard, « Présentation », dans Louis de Pontis, \teiHi[rend={italic}]{Mémoires (1676). Avec la totalité des modifications de 1678}, recueillis par Pierre Thomas Du Fossé et Andrée Villard (éd.), Paris, Honoré Champion, 2000, p. 7-78, ici p. 29. Toutes nos références renverront à cette édition.}} ». Il ne s’agira pas de présumer la dimension intentionnelle de la narration faite par le \teiHi[rend={italic}]{je }narrant, encore moins sa dimension esthétique, mais bien plutôt d’étudier de près la manière dont les deux instances énonciatives s’articulent dans la narration et les effets rhétoriques qui en découlent. En quoi l’énonciation référentielle permet-elle d’associer représentation et révélation des secrets du pouvoir ? Quelles en sont les conséquences ? On verra tout d’abord comment l’expérience réelle des arcanes du pouvoir conditionne la représentation du secret dans les Mémoires, avant d’observer les implications du dispositif énonciatif des Mémoires dans la représentation et la révélation des secrets.}
            \begin{teiDiv}[type={section1},xmlid={div01-5}]

              \teiHead[xmlid={le-je-et-le-pouvoir-une-representation-des-coulisses-001}]{Le \teiHi[rend={italic}]{je }et le pouvoir : une représentation des coulisses}
              \begin{teiDiv}[type={section2},xmlid={div02-5}]

                \teiHead[xmlid={montage-narratif-un-je-entre-inclusion-et-exclusion-001}]{Montage narratif : un \teiHi[rend={italic}]{je} entre inclusion et exclusion}
                \teiP[xmlid={p4-6}]{Dans les Mémoires, les faits relatés sont conditionnés par l’existence du témoin, c’est de son accès réel aux coulisses du pouvoir que dépend sa connaissance des secrets. Il s’agira ici d’observer les manifestations narratives de la position du mémorialiste sur l’échiquier politique.}
                \teiP[xmlid={p5-6}]{Le deuxième volume des \teiHi[rend={italic}]{Mémoires}\teiNote[n={8},place={foot},xmlid={ftn29}]{\teiP[xmlid={p-225}]{\teiHi[rend={italic}]{Mémoires et la vie de messire Claude de Letouf, chevalier baron de Sirot, lieutenant général des camps et armée du roy, \&c. Sous les règnes des rois Henri }\teiHi[rend={small-caps italic}]{iv}\teiHi[rend={italic}]{, Loüis }\teiHi[rend={small-caps italic}]{XIII}\teiHi[rend={italic}]{, et Loüis }\teiHi[rend={small-caps italic}]{XIV}, Paris, Claude Barbin, 1683, vol. 2. Toutes nos références renverront à cette édition. L’orthographe sera modernisée dans les citations.}} de Sirot est scandé par ses tentatives répétées auprès de Mazarin pour obtenir des faveurs royales. La narration met clairement en valeur ses demandes qui suivent toutes le même schéma : requête de Sirot, différentes promesses de Mazarin débouchant soit sur un refus total, soit sur un accord d’une faveur bien en deçà de ce qui lui avait été promis. Ces épisodes répétés sont toujours relatés avec le même vocabulaire : les « protestations » (p. 117 et 138) de Mazarin les ouvrent, puis suivent les « espérances » (p. 134 et 259) de Sirot, et ils sont clos par les refus de Mazarin, introduits par la conjonction de coordination \teiHi[rend={italic}]{mais}\teiNote[n={9},place={foot},xmlid={ftn30}]{\teiP[xmlid={p-226}]{« \lbrack{}…\rbrack{} mais il refusa tous ces expédients » (p. 121), « \lbrack{}…\rbrack{} mais tout cela n’aboutit à rien » (p. 132), « \lbrack{}…\rbrack{} mais l’aversion naturelle qu’il avait pour moi, s’opposait trop fortement à tout ce qui pouvait contribuer à ma satisfaction » (p. 154), « \lbrack{}…\rbrack{} mais \lbrack{}…\rbrack{} je trouvai le cardinal entièrement changé pour mes intérêts » (p. 259).}} qui souligne bien le revirement narratif. N’opposant que son bon service de soldat à la duplicité du ministre, Sirot est contraint de rester en-dehors du jeu, ce que souligne le montage narratif par les multiples répétitions d’un même canevas.}
                \teiP[xmlid={p6-6}]{Soldat lui aussi, Pontis a néanmoins accès à une partie des \teiHi[rend={italic}]{arcana imperii }en tant qu’intime de Louis \teiHi[rend={small-caps}]{XIII}. La narration confronte à trois reprises une scène secrète et intime, entre Pontis et le roi, à une scène officielle et hostile, entre Pontis et Richelieu. Est tout d’abord opposée une « petite farce » (p. 314) jouée par Pontis dans la chambre de Louis \teiHi[rend={small-caps}]{XIII }avec pour seul autre spectateur Saint-Simon, à « une audience extraordinaire » devant la porte du roi « en présence des cardinaux, des princes et des seigneurs de la cour » (p. 315). Une disgrâce officielle ordonnée par Richelieu contraste ensuite avec un rendez-vous secret donné par le roi dans une galerie de Saint-Germain. Enfin, deux interrogatoires, le premier mené par Richelieu et le second par De Noyers, dans lesquels Pontis est « tourn\lbrack{}é\rbrack{} et \lbrack{}…\rbrack{} retourn\lbrack{}é\rbrack{} en toutes manières » (p. 540), alternent avec deux scènes intimes avec Louis \teiHi[rend={small-caps}]{XIII}, la première consistant en une véritable confession royale, dans la garde-robe même du souverain\teiNote[n={10},place={foot},xmlid={ftn31}]{\teiP[xmlid={p-227}]{Suivant Furetière, la garde-robe est le lieu secret par excellence au \teiHi[rend={small-caps}]{xvii}\teiHi[rend={sup}]{e} siècle : « S\teiHi[rend={small-caps}]{ecret}, est aussi une epithete particuliere du lieu. Un lieu secret, c’est le privé, la garderobbe. \teiHi[rend={italic}]{Un escalier secret} ou desrobé, est un petit escalier par où on se coule sans bruit, ou sans être veu » (\teiHi[rend={italic}]{Dictionnaire universel contenant généralement tous les mots françois, tant vieux que modernes, et les termes de toutes les sciences et des arts}, 1690, article « \teiHi[rend={small-caps}]{S}ecret »). On notera que lors de la première entrevue, Pontis monte dans la chambre du roi grâce à l’escalier secret de la garde-robe : avec ces notations spatiales, on est en plein dans l’imaginaire du secret.}}. La confrontation directe de ces scènes de confiance et de défiance met nettement en valeur la complicité qui unit Louis \teiHi[rend={small-caps}]{XIII} et Pontis, dépositaire des secrets du roi.}
                \teiP[xmlid={p7-6}]{Négociateur pour le prince de Condé, Lenet est un expert des secrets du pouvoir. Ses \teiHi[rend={italic}]{Mémoires}\teiNote[n={11},place={foot},xmlid={ftn32}]{\teiP[xmlid={p-228}]{\teiHi[rend={italic}]{Mémoires de Monsieur L***, conseiller d’État, contenant l’Histoire des guerres civiles des années 1649 et suivantes ; principalement celles de Guienne et autres provinces}, \lbrack{}Paris, Marin et Guérin\rbrack{}, 1729, 2 volumes. Toutes nos références renverront à cette édition. L’indication du volume sera abrégée par un chiffre romain avant le numéro de page. L’orthographe sera modernisée dans les citations. En l’absence d’une édition scientifique moderne, nous avons en effet choisi l’édition \teiHi[rend={italic}]{princeps.}}} consignent les très longs entretiens qu’il a eus avec le cardinal Mazarin. Les deux adversaires s’y jaugent, agissent en fonction d’intérêts cachés et essaient de découvrir les secrets de l’autre pour le devancer. Lors de leur première rencontre dans les Mémoires, Lenet et Mazarin tentent de rester alliés, mais Mazarin se rapproche en même temps secrètement de la maison Vendôme, rivale des Condé, tandis que Lenet réclame pour Condé des terres indispensables s’il devait engager une guerre civile. Les deux autres entretiens confidentiels qui suivent relatent les négociations de Mazarin et Lenet pour s’allier afin de permettre la libération des princes, alors même qu’ils négocient chacun secrètement avec les frondeurs. Lenet a une place de choix et un vrai pouvoir durant les tractations diplomatiques de la Fronde. La narration en rend parfaitement compte par ces longs dialogues rapportant les diverses tentatives de manipulation opérées par les deux camps.}
                \teiP[xmlid={p8-6}]{On constate ainsi un double mouvement : les récits sont informés par chaque expérience singulière du pouvoir, mais en retour les narrations s’ajustent pour en rendre compte le plus fidèlement possible. Les diverses formes prises par le secret dans les dialogues dépendent de même des liens existant entre le mémorialiste et le pouvoir.}
              
\end{teiDiv}
              \begin{teiDiv}[type={section2},xmlid={div03-4}]

                \teiHead[xmlid={visages-discursifs-du-secret-un-je-acteur-ou-victime-001}]{Visages discursifs du secret : un \teiHi[rend={italic}]{je} acteur ou victime}
                \teiP[xmlid={p9-6}]{Naissant d’une chose tue ou dite autrement, le secret modèle les discours. On analysera ici ses différentes manifestations et la manière dont il s’y transforme en une arme maniée ou subie par les acteurs de l’interaction.}
                \teiP[xmlid={p10-6}]{Les raisons de l’action de Mazarin demeurent partiellement opaques à Sirot qui, trop crédule pour la duplicité régnante, croit autant aux signes extérieurs de Mazarin qu’à ses promesses. Il rapporte ainsi les « promesses que son Éminence \lbrack{}lui\rbrack{} avait faites » (p. 120), insiste sur la charge que « le cardinal \lbrack{}lui\rbrack{} avait promis\lbrack{}e\rbrack{} » (p. 261). La duplicité du pouvoir trouve son expression linguistique dans cette manipulation du performatif. Dans cette dernière citation, l’aspect accompli du verbe indique que l’action est envisagée après son terme, comme étant achevée. Mais l’aspect accompli du performatif, la promesse faite, ne s’accompagne pas de sa réalisation, l’attribution de la charge. Le décalage introduit entre la performativité de ces verbes et leur irréalisation confère au secret le visage de la tromperie, et la répétition de cette parole vide insiste sur ce dérèglement dont Sirot est victime.}
                \teiP[xmlid={p11-6}]{Le secret chez Lenet concerne non seulement les intentions d’autrui, mais aussi les propres intentions du mémorialiste, dont les propos masquent ses desseins véritables. Comme l’échange suivant le montre, toutes les paroles de Lenet et de Mazarin sont aiguisées d’implicite, destiné à faire agir l’adversaire contre son gré :}
                \begin{teiCit}[xmlid={cit01-6}]

                  \begin{teiQuote}
Si Votre Éminence savait ce qui se brasse de toutes parts contre elle, elle ferait de sérieuses réflexions sur tout ce que je lui dis. Aussi fais-je, me repartit-il ; je sais bien que j’ai beaucoup d’ennemis, mais j’espère d’en venir à bout, comme j’ai fait jusqu’à présent ; j’ai de la résolution, des amis, et la protection de la Reine : demain nous nous reverrons (II, p. 486-487).
\end{teiQuote}
                
\end{teiCit}
                \teiP[xmlid={p12-6}]{Lenet commence par discréditer son adversaire tout en le menaçant. La capacité de Mazarin à comprendre sa propre situation est en effet remise en cause par la subordonnée hypothétique qui envisage cette compréhension comme contraire à l’état actuel des choses. Le verbe semi-factif\teiNote[n={12},place={foot},xmlid={ftn33-2}]{\teiP[xmlid={p-229}]{Voir Axelle Vatrican, « Savoir que et la notion de présupposition », \teiHi[rend={italic}]{Langages}, n\teiHi[rend={sup}]{o} 186, 2012/2, p. 69-84.}} \teiHi[rend={italic}]{savoir}, qui présuppose la vérité de la complétive qu’il régit, introduit une menace pour Mazarin qui est contraint de croire ces propos. L’acte illocutoire n’a pas l’effet perlocutoire escompté : Mazarin – peut-être en réalité déstabilisé – parvient en apparence à prendre l’ascendant dans la discussion. Il réemploie le verbe \teiHi[rend={italic}]{savoir }pour désamorcer la menace de son adversaire et montrer qu’il maîtrise la situation ; puis il menace à son tour Lenet, en employant des substantifs généraux (« des ennemis », « des amis ») qui peuvent laisser penser à Lenet qu’il est lui aussi trahi. Il reprend enfin définitivement le pouvoir en décidant de la fin de l’échange. Le secret ici est manipulation, et ces longs échanges rapportés avec Mazarin nous montrent la grande habileté de Lenet négociateur.}
                \teiP[xmlid={p13-6}]{Dans les Mémoires de Pontis, le secret consiste clairement en un « jeu du tiers inclus-exclu\teiNote[n={13},place={foot},xmlid={ftn34-2}]{\teiP[xmlid={p-230}]{Expression de Louis Marin qui souligne ce paradoxe du secret : il exige qu’« un des joueurs \teiHi[rend={italic}]{en tant que joueur }ne fa\lbrack{}sse\rbrack{} pas partie précisément \teiHi[rend={italic}]{pour que }le jeu ait lieu » (« Logiques du secret », \teiHi[rend={italic}]{Lectures traversières}, Paris, Albin Michel, 1992, p. 247-257, ici p. 251).}} » : il oppose Pontis et Louis \teiHi[rend={small-caps}]{XIII} à Richelieu. Après avoir reçu Pontis en secret et lui avoir pardonné son acte d’insubordination, Louis \teiHi[rend={small-caps}]{XIII} met en scène la rencontre officielle et y participe, jouant la carte de la distance pour dissimuler la mise en scène. La comparaison du début des deux demandes de grâce de Pontis met en évidence la dimension factice d’une parole déterminée par le secret :}
                \begin{teiCit}[xmlid={cit02-5}]

                  \begin{teiQuote}
Sire, je ne puis assez \teiHi[rend={italic}]{remercier} Votre Majesté de la grâce et de l’honneur qu’elle me fait de vouloir bien que je lui rende compte de mes actions \lbrack{}…\rbrack{} (p. 308).\teiLb{}Sire, je viens me jeter aux pieds de Votre Majesté pour implorer sa miséricorde. \lbrack{}…\rbrack{} Mais je \teiHi[rend={italic}]{supplie} très humblement votre bonté, Sire, de vouloir auparavant m’accorder la grâce de m’entendre \lbrack{}…\rbrack{} (p. 314, nous soulignons).
\end{teiQuote}
                
\end{teiCit}
                \teiP[xmlid={p14-5}]{Lors de la première entrevue, secrète, le verbe performatif \teiHi[rend={italic}]{remercier }est opérant. Or, lors de l’entrevue officielle, la performativité des verbes \teiHi[rend={italic}]{implorer }et \teiHi[rend={italic}]{supplier }est annulée par le secret partagé entre le souverain et Pontis. Ce dernier demande en effet ce qu’il a déjà obtenu lors de la première entrevue. C’est l’ordre des verbes performatifs qui montre la mise en scène et rend impossible la performativité des deux verbes : on supplie puis on remercie et non l’inverse. Comme la scène elle-même qui est pure mise en scène, la parole est scindée entre être et paraître. Ces deux extraits soulignent bien la connivence des deux interactants, et ce de manière directe dans la chambre du souverain, et de manière indirecte devant Richelieu.}
                \teiP[xmlid={p15-4}]{L’expérience des coulisses du pouvoir informe la narration et les discours des Mémoires, qui les retranscrivent et les façonnent en retour afin d’en souligner les spécificités. Si cette narration est un véritable plaidoyer \teiHi[rend={italic}]{pro domo }chez Sirot, le plaidoyer est plus ténu chez Lenet qui rend compte de son rôle dans l’histoire de Condé. Quant à Pontis, il défend bien sa cause comme il l’aurait fait au moment des faits mais, lorsqu’il raconte son histoire à Du Fossé, il est détaché des vanités mondaines et ne désire plus convaincre que de la nécessité de se tourner vers Dieu\teiNote[n={14},place={foot},xmlid={ftn35-5}]{\teiP[xmlid={p-231}]{Ajoutons que dans le cas spécifique des Mémoires de Pontis, cette entreprise du \teiHi[rend={italic}]{je }narrant double celle du transcripteur qui entend par ce récit faire « remarquer les voies différentes par lesquelles Dieu se plaît de conduire ceux qu’il veut enfin attirer à son service » (« Avertissement », \teiHi[rend={italic}]{op. cit.}, p. 83).}}. Nous voyons mieux ainsi les liens qui unissent \teiHi[rend={italic}]{je }narrant et \teiHi[rend={italic}]{je }narré pour ce qui est des événements rapportés ; regardons à présent la manière dont ces deux instances articulent représentation et révélation des secrets.}
              
\end{teiDiv}
            
\end{teiDiv}
            \begin{teiDiv}[type={section1},xmlid={div04}]

              \teiHead[xmlid={je-narre-et-je-narrant-le-probleme-de-la-connaissance-des-secrets-001}]{\teiHi[rend={italic}]{Je }narré et \teiHi[rend={italic}]{je }narrant : le problème de la connaissance des secrets}
              \begin{teiDiv}[type={section2},xmlid={div05}]

                \teiHead[xmlid={je-narre-et-je-narrant-representation-et-revelation-des-secrets-001}]{\teiHi[rend={italic}]{Je} narré et \teiHi[rend={italic}]{je} narrant : représentation et révélation des secrets}
                \teiP[xmlid={p16-4}]{En littérature factuelle, le secret implique une connaissance. Il a pu être découvert immédiatement, entre le temps de l’événement et celui de l’écriture, ou bien être encore inconnu. Il est alors intéressant d’observer si la narration respecte la chronologie de connaissance du secret.}
                \teiP[xmlid={p17-4}]{Dans les Mémoires de Pontis, \teiHi[rend={italic}]{je }narré et \teiHi[rend={italic}]{je }narrant rendent compte de ce processus de découverte. On peut définir un premier cas de figure dans lequel le \teiHi[rend={italic}]{je }narré connaissait le secret et le révèle ; le \teiHi[rend={italic}]{je }narrant ne fait alors que le reprendre. L’origine secrète de la haine de Richelieu pour Pontis est ainsi dévoilée une première fois par le biais du \teiHi[rend={italic}]{je }narré, puis reprise par le \teiHi[rend={italic}]{je }narrant dans une intervention proleptique : « mais je ressentis bientôt les effets de la mauvaise volonté de Monsieur l’évêque d’Amiens qui réussit si habilement à irriter le cardinal contre moi, que je me vis en un instant dépouillé de tout » (p. 498). Second cas de figure, le \teiHi[rend={italic}]{je }narrant a appris une information ultérieurement ; la narration le signale alors clairement, par exemple au moyen d’une proposition subordonnée qui permet le décrochage énonciatif : « le roi connut aussitôt le déguisement et l’artifice de cette demande, ainsi qu’il me le témoigna lui-même depuis » (p. 534). En l’occurrence, les circonstances dans lesquelles Louis \teiHi[rend={small-caps}]{XIII} a rapporté au \teiHi[rend={italic}]{je }narré la demande manipulatrice de Richelieu sont exposées quelques pages plus loin, ce qui renforce la cohérence narrative. La reprise par le \teiHi[rend={italic}]{je }narrant des secrets déjà révélés par le \teiHi[rend={italic}]{je }narré et l’intégration dans le \teiHi[rend={italic}]{continuum }narratif de données détenues par le \teiHi[rend={italic}]{je }narrant seul, font de la défiance de Richelieu un fil rouge du récit : celle-ci apparaît comme la clé des malheurs de Pontis. Ces interventions pourraient être liées à la genèse particulière du texte : elles seraient une trace de la posture de conteur adoptée par Pontis relatant sa vie à Du Fossé, ou un indice de l’intervention du scripteur dans la mise en récit.}
                \teiP[xmlid={p18-3}]{Chez Lenet, le \teiHi[rend={italic}]{je }narrant vient confirmer les dires du \teiHi[rend={italic}]{je }narré, qui sait immédiatement énoncer les ressorts cachés de l’événement. Le \teiHi[rend={italic}]{je }narré peut ainsi expliquer son hypothèse sur le comportement de Mazarin, élaborée au moment des faits, et le \teiHi[rend={italic}]{je }narrant la confirmer, au présent d’énonciation : « je m’imaginai, et je crois que je n’eus pas tort, que le Cardinal qui allait à ses fins comme nous allions aux nôtres voulait que je ne parlasse à Mademoiselle que de l’affaire de Montrond » (II, p. 480). Le \teiHi[rend={italic}]{je }narrant peut aussi introduire une confirmation factuelle des interprétations du \teiHi[rend={italic}]{je }narré\teiHi[rend={italic}]{ }: « j’ai jugé \lbrack{}…\rbrack{} qu’il pouvait bien avoir insinué au duc de Rohan de proposer cela au prince pour voir comment il le recevrait, et pour s’en servir après contre lui, comme il ne manqua pas de faire » (I, p. 41). Le secret n’est pas élucidé entre les faits et leur narration, mais cet intervalle de temps est important : les événements qui ont suivi ont conforté les opinions du \teiHi[rend={italic}]{je }narré. Celui-ci révèle les intentions cachées des actions de Mazarin et le recul du \teiHi[rend={italic}]{je }narrant crédibilise ces révélations parce que le temps est venu les confirmer. La clairvoyance immédiate du négociateur n’en est que davantage soulignée.}
                \teiP[xmlid={p19-3}]{Le \teiHi[rend={italic}]{je }narré dans les Mémoires de Sirot est exclu des secrets du pouvoir, toute révélation devrait donc être imputable au \teiHi[rend={italic}]{je }narrant. Or il est très difficile de reconstituer la chronologie de la découverte du secret et la narration présente des contradictions. Alors que Sirot sollicite le paiement de sa rançon, est d’abord révélée par le \teiHi[rend={italic}]{je }narré l’origine supposée des refus de Mazarin : des espions lui ont appris ses entrevues avec Condé, et le ministre craint « un attachement secret entre \lbrack{}eux\rbrack{} qui pouvait lui être préjudiciable » (p. 120). Mais cette explication confirmée par un discours rapporté de Mazarin et par le \teiHi[rend={italic}]{je }narré (p. 123 et 127-128) est contredite plus loin dans le récit. Un monologue auto-rapporté\teiNote[n={15},place={foot},xmlid={ftn36-5}]{\teiP[xmlid={p-232}]{Selon la terminologie proposée par Dorrit Cohn dans \teiHi[rend={italic}]{La Transparence intérieure. Modes de représentation de la vie psychique dans le roman}, Alain Bony (trad.), Paris, Éditions du Seuil, 1981, p. 185-190.}} expose les doutes du \teiHi[rend={italic}]{je }narré au sujet du comportement du ministre : « mais ne trouverait-il point \lbrack{}…\rbrack{} à redire, aux déférences et à l’obéissance que je rends à Messieurs le Prince de Condé, et le duc d’Enghien son fils. Cela ne peut être » (p. 136). Puis le \teiHi[rend={italic}]{je }narré évoque « une certaine aversion secrète, dont \lbrack{}il\rbrack{} ne savai\lbrack{}t\rbrack{} pas l’origine » (p. 153). Que penser alors de la première explication alléguée ? Quand Sirot a-t-il appris que des espions avaient informé le ministre de ses entrevues avec Condé ? L’a-t-il su immédiatement, ce qui est contredit par la suite de la narration ? L’a-t-il su durant l’intervalle entre les faits et leur mise en récit et réintroduit \teiHi[rend={italic}]{a posteriori} sur le fil chronologique de la narration ? Dans ce cas, les deux confirmations par le récit qui suivent sont-elles à traiter sur le même plan – le \teiHi[rend={italic}]{je }narrant réinterprétant les événements pour leur donner une cohérence – ou sont-elles des explications déjà détenues par le \teiHi[rend={italic}]{je }narré ? On voit qu’un secret connu du seul \teiHi[rend={italic}]{je }narrant peut être intégré à la narration sans que le processus de connaissance du secret soit signalé. La manipulation des informations par le \teiHi[rend={italic}]{je }narrant – non nécessairement consciente et intentionnelle – viendrait ici renforcer l’entreprise de dénonciation d’un ministre défiant et ambitieux.}
                \teiP[xmlid={p20-3}]{Chaque texte articule donc différemment \teiHi[rend={italic}]{je }narré et \teiHi[rend={italic}]{je }narrant pour révéler les secrets : le dévoilement des intentions secrètes de Richelieu par les deux \teiHi[rend={italic}]{je }chez Pontis permet de montrer que le \teiHi[rend={italic}]{je }narré n’était pas dupe des manœuvres du ministre en renforçant la critique émanant du \teiHi[rend={italic}]{je }narrant ; les interprétations immédiates du \teiHi[rend={italic}]{je }narré chez Lenet valorisent la sagacité du négociateur ; et le savoir rétrospectif du \teiHi[rend={italic}]{je }narrant serait potentiellement utilisé à des fins de manipulations narratives chez Sirot, afin d’accentuer ce qui est présenté comme acharnement de la part du ministre sur le \teiHi[rend={italic}]{je }narré. Toutes ces révélations sont tout de même étonnantes dans un récit factuel puisqu’elles concernent les intentions cachées d’autrui. Examinons donc les moyens déployés par la narration pour les représenter avec vraisemblance.}
              
\end{teiDiv}
              \begin{teiDiv}[type={section2},xmlid={div06}]

                \teiHead[xmlid={limmixtion-in-vraisemblable-du-je-dans-les-secrets-dautrui-001}]{L’immixtion (in)vraisemblable du \teiHi[rend={italic}]{je} dans les secrets d’autrui}
                \teiP[xmlid={p21-3}]{Pour nos mémorialistes, le véritable \teiHi[rend={italic}]{locus politicus} est ce \teiHi[rend={italic}]{locus secretus }qu’est la psyché opaque des ministres. Afin de mettre au jour ce qu’ils pensent être leurs intentions secrètes, ils emploient fréquemment des compléments circonstanciels de but en \teiHi[rend={italic}]{pour }(désormais CCB)\teiHi[rend={italic}]{. }Ces CCB présentent deux caractéristiques intéressantes pour notre propos. Premièrement, ils sont toujours de l’ordre de l’interprétation, pourtant ils présentent cette interprétation de manière objective. Dans un CCB, c’est en effet « l’activité de synthèse mentale du sujet\teiNote[n={16},place={foot},xmlid={ftn37-5}]{\teiP[xmlid={p-233}]{Pierre Cadiot, \teiHi[rend={italic}]{De la grammaire à la cognition : la préposition }pour, Paris, CNRS Éditions, 1991, p. 261-262.}} » qui est responsable de la mise en relation d’une action avec une intention\teiHi[rend={italic}]{. }Or le CCB met en place un « dispositif de coénonciation\teiNote[n={17},place={foot},xmlid={ftn38-5}]{\teiP[xmlid={p-234}]{\teiHi[rend={italic}]{Ibid.}, p. 260.}} » qui présente l’intention du locuteur du « même point de vue\teiNote[n={18},place={foot},xmlid={ftn39-5}]{\teiP[xmlid={p-235}]{\teiHi[rend={italic}]{Ibid}., p. 259.}} » que l’action. Secondement, ces interprétations peuvent être aussi bien le fait du \teiHi[rend={italic}]{je }narré, qui analyse la situation au moment des faits, que du \teiHi[rend={italic}]{je }narrant qui réanalyse la situation des années après. Si l’intervention du \teiHi[rend={italic}]{je }narrant n’est pas précisée, elle est invisible et seul le décalage de connaissances entre les deux instances peut nous l’indiquer. On étudiera donc les conséquences de l’utilisation de ces CCB dans la révélation des secrets.}
                \teiP[xmlid={p22-3}]{Commençons par Lenet qui rapporte toutes les interprétations qu’il a faites du comportement de Mazarin au moment même des faits. Tout est clair : le \teiHi[rend={italic}]{je }narré connaissait une partie des secrets de son adversaire, et la narration distingue nettement la transcription des actions de Mazarin de leurs interprétations par Lenet. Un verbe d’opinion signale souvent au lecteur l’activité interprétative du \teiHi[rend={italic}]{je }narré, qui élabore une hypothèse sur l’intention cachée de l’interlocuteur : « j’ai jugé que le cardinal lui préparait dès lors ce piège pour brouiller \lbrack{}Condé\rbrack{} avec le duc d’Orléans » (I, p. 41). Le verbe d’opinion expose explicitement ce que suppose implicitement le CCB. Comme précédemment, le recul du \teiHi[rend={italic}]{je }narrant autorise ces révélations, mais on peut supposer la connaissance immédiate de ces intentions cachées ; la dimension interprétative et subjective du \teiHi[rend={italic}]{je} narré est de plus clairement indiquée.}
                \teiP[xmlid={p23-3}]{Au contraire, les Mémoires de Sirot ne rendent jamais compte de cette activité interprétative\teiHi[rend={italic}]{ }; on peut les soupçonner de se servir du recul du \teiHi[rend={italic}]{je }narrant pour expliquer les tribulations du \teiHi[rend={italic}]{je }narré. Dans l’exemple suivant : « pour me persuader ce qu’il me disait, \lbrack{}Mazarin\rbrack{} me mena chez la Reine et m’y présenta » (p. 117), l’intention de Mazarin n’apparaît pas comme une interprétation du \teiHi[rend={italic}]{je }mais bien comme un but naturel de l’action. Cette intention peut même être présupposée, comme par exemple dans une subordonnée, ce qui empêche sa réfutation : « cette Éminence qui ne demandait qu’à gagner du temps pour éluder les propositions qu’on lui faisait, lui dit, qu’il tâcherait de persuader à la reine d’écouter cette nouvelle proposition » (p. 131-132). La narration contraint ainsi le lecteur à croire au caractère manipulateur des actions de Mazarin. On peut vraisemblablement penser que la disgrâce finale du \teiHi[rend={italic}]{je} narré oriente ces interprétations, le CCB permet alors au \teiHi[rend={italic}]{je }narrant de s’introduire dans le déroulement des faits sans être vu. Cette invisibilisation est un élément de plus à mettre au compte de la critique de Mazarin.}
                \teiP[xmlid={p24-3}]{Chez Pontis, les savoirs respectifs du \teiHi[rend={italic}]{je }narrant et du \teiHi[rend={italic}]{je }narré sont bien distingués, mais c’est dans ce texte que l’effacement de l’énonciation au profit de la révélation des pensées secrètes est la plus saisissante. Comme chez Lenet, le \teiHi[rend={italic}]{je }narré peut connaître le secret et le révéler. Mais il est aussi fait mention de l’origine du savoir du \teiHi[rend={italic}]{je }narré : c’est parce que ce dernier a appris de Louis \teiHi[rend={small-caps}]{XIII} que Richelieu avait essayé de manipuler le roi, en se servant de lui, qu’il est fondé à dire que le cardinal « s’avisa \lbrack{}…\rbrack{} de se servir de \lbrack{}s\rbrack{}on nom pour emporter sur son ennemi ce qu’il pouvait souffrir qu’il lui enlevât » (p. 533). La narration objective rend ainsi plausible la révélation des intentions cachées de Richelieu. Dans le passage suivant, qui détaille la réaction de Richelieu venant d’apprendre que le roi s’est opposé à ses desseins en confiant un gouvernement à Cinq-Mars, un pas de plus est franchi dans ce dévoilement :}
                \begin{teiCit}[xmlid={cit03-5}]

                  \begin{teiQuote}
Le cardinal, \teiHi[rend={italic}]{qui savait bien que cet autre était Monsieur le grand écuyer,} regarda comme le dernier affront qu’il pût recevoir, de ce que celui qui était sa créature, et qui depuis était devenu son ennemi, avait pu emporter sur lui ce gouvernement. \teiHi[rend={italic}]{Comme il s’était persuadé qu’il était maître de tout, et qu’il croyait même s’être rabaissé d’avoir demandé une chose qui dépendait de son pouvoir}, il fut piqué très sensiblement du refus du roi ; et jugeant que ce ne pouvait être qu’un effet de la mauvaise volonté de ses ennemis qui l’avaient mis mal dans l’esprit de Sa Majesté, il commença à entrer dans quelque appréhension de voir bientôt renverser toute sa fortune. \teiHi[rend={italic}]{Car il savait}, \teiHi[rend={bold}]{comme j’ai dit auparavant}, \teiHi[rend={italic}]{qu’il se formait de puissantes cabales contre lui }\lbrack{}…\rbrack{} (p. 532-533).
\end{teiQuote}
                
\end{teiCit}
                \teiP[xmlid={p25-3}]{Nous avons laissé en caractères romains la séquence qui relève d’une observation purement factuelle de la scène, correspondant à la mention de la colère de Richelieu. C’est cette réaction qui est expliquée par les pensées du ministre. En italiques, on trouve les savoirs ou les croyances vraisemblables que Pontis attribue à Richelieu et, soulignés, deux brefs basculements en focalisation interne, introduits par trois verbes qui dénotent un processus mental. La mention des croyances vraisemblables de Richelieu permet le glissement insensible vers la focalisation interne. En gras, enfin, nous signalons l’intervention du \teiHi[rend={italic}]{je }narrant, qui vient justifier l’exposition des pensées d’autrui. Cette justification, qui consiste notamment en un renvoi à la narration antérieure, renforce la cohérence narrative et objective les composantes interprétatives de la scène. Tous ces éléments confèrent de la vraisemblance au coup de force qu’est, en régime factuel, l’exposition des pensées d’autrui\teiNote[n={19},place={foot},xmlid={ftn40-5}]{\teiP[xmlid={p-236}]{Dorrit Cohn fait de la focalisation interne un critère de distinction entre la fiction et la factualité (\teiHi[rend={italic}]{Le Propre de la fiction}, Claude Hary-Schaeffer (trad.), Paris, Éditions du Seuil, 2001), suivant en cela les analyses de Käte Hamburger (\teiHi[rend={italic}]{Logique des genres littéraires}, Pierre Cadiot trad., Paris, Éditions du Seuil, 1986). Marc Hersant, identifiant un passage de focalisation interne dans les Mémoires de Saint-Simon, a discuté ce critère et a conclu qu’il n’était pas « décisif pour “trancherˮ » du caractère factuel ou fictionnel du texte (\teiHi[rend={italic}]{Le Discours de vérité dans les “Mémoires” du duc de Saint-Simon}, Paris, Honoré Champion, 2009, p. 524).}}. Ce passage précède une tentative de manipulation de Louis \teiHi[rend={small-caps}]{XIII }par Richelieu : il met efficacement au jour les intentions secrètes de ce dernier pour s’accaparer le pouvoir. La psyché du ministre n’est plus \teiHi[rend={italic}]{locus secretus}, mais elle est exposée au grand jour et dénoncée par un \teiHi[rend={italic}]{je }non mentionné dans le passage.}
                \teiP[xmlid={p26-2}]{La révélation des secrets du pouvoir prend ainsi diverses formes qui laissent plus ou moins apparaître leur part subjective et interprétative. L’origine référentielle est plus marquée chez Lenet par la présence du \teiHi[rend={italic}]{je }narrant et la mention de l’activité interprétative du \teiHi[rend={italic}]{je }narré ; elle l’est moins chez Sirot qui emploie aussi le \teiHi[rend={italic}]{je }narrant mais intègre ses interprétations à la narration sans les signaler ; elle l’est à la fois plus et moins chez Pontis : l’emploi régulier du \teiHi[rend={italic}]{je }narrant cohabite avec l’effacement temporaire de la première personne dans un passage en focalisation interne. Cette invisibilisation de l’origine référentielle contribue activement à l’efficacité rhétorique du texte, puisque qu’elle permet de présenter comme une vérité objective ce qui tient en fait de l’interprétation subjective.}
                \teiP[xmlid={p27-2}]{On observe finalement que la fusion des deux instances énonciatives est aussi importante que leur scission dans la poétique du secret des Mémoires. Les montages narratifs et les manifestations discursives du secret sont en effet déterminés par l’expérience du \teiHi[rend={italic}]{je} narré, mais leur mise en forme finale est bien imputable au \teiHi[rend={italic}]{je} narrant, qui rend visibles les spécificités de son expérience des coulisses royales. Si chez Pontis et Lenet, le \teiHi[rend={italic}]{je }narrant vient confirmer ce que savait déjà le \teiHi[rend={italic}]{je }narré, on a souligné chez Sirot les potentielles manipulations narratives d’un \teiHi[rend={italic}]{je }narrant qui incorpore à sa narration des secrets inconnus du \teiHi[rend={italic}]{je }narré. Mais c’est au sein des CCB que cette fusion des deux instances énonciatives est la plus forte car la plus invisible. Ils constituent en cela un outil précieux de cohérence narrative parce qu’ils permettent au \teiHi[rend={italic}]{je} narrant de donner sens à l’expérience du \teiHi[rend={italic}]{je} narré par la révélation d’intentions cachées. Ces révélations servent en même temps la critique efficace d’un pouvoir duplice. Les Mémoires de Pontis constituent sur ce point un hapax : la cohérence narrative et l’accusation d’un ministre machiavélique semblent en effet autoriser l’usage de la focalisation interne, problématique en régime factuel. Toutes nos observations sur cet auteur – alternance de scènes de confiance et de défiance, récits théâtraux de demande de grâce, prolepses narratives et focalisation interne – semblent converger vers une plus grande intervention du \teiHi[rend={italic}]{je} narrant dans la mise en récit de l’existence du \teiHi[rend={italic}]{je} narré. On peut légitimement se demander dans quelle mesure l’oralité primitive de ce récit et sa retranscription par un tiers auraient partie liée avec ces conclusions.}
                \teiP[xmlid={p28-2}]{Il convient d’interpréter cette poétique du secret à l’aune de l’acte rhétorique qui fonde l’écriture des Mémoires. Quel que soit le moment où le secret vient à être connu, il est remarquable qu’il soit toujours révélé au début des épisodes narrés sans jamais constituer un ressort de dramatisation de l’action. En désactivant la possibilité d’un récit à intrigue, ces révélations entravent l’immersion du lecteur. En cela, nos textes factuels se distinguent nettement des récits fictionnels qui usent du secret comme d’« une clé d’accès \lbrack{}…\rbrack{} pour favoriser l’immersion du lecteur dans l’espace textuel\teiNote[n={20},place={foot},xmlid={ftn41-5}]{\teiP[xmlid={p-237}]{Guiomar Hautcœur, \teiHi[rend={italic}]{Roman et secret. Essai sur la lecture à l’époque moderne (}\teiHi[rend={small-caps italic}]{xvi}\teiHi[rend={italic sup}]{e}\teiHi[rend={italic}]{-}\teiHi[rend={small-caps italic}]{xviii}\teiHi[rend={italic sup}]{e}\teiHi[rend={italic}]{ siècles)}, Paris, Classiques Garnier, 2019, p. 25.}} ». Au contraire, suivant Raphaël Baroni, la fonction « des comptes rendus historiques \lbrack{}…\rbrack{} est davantage d’expliquer le passé que d’intriguer le lecteur\teiNote[n={21},place={foot},xmlid={ftn42-5}]{\teiP[xmlid={p-238}]{Raphaël Baroni, « Le rôle des personnages dans les rouages de l’intrigue », \teiHi[rend={italic}]{Letras de Hoje}, vol. 52, 2017, p. 156-166, ici p. 157.}} », ce qui amène dès lors à considérer « l’absence de dispositifs intrigants au sein de certains récits \lbrack{}non\rbrack{} comme un défaut, mais au contraire comme une qualité assurant le succès de l’acte de langage\teiNote[n={22},place={foot},xmlid={ftn43-5}]{\teiP[xmlid={p-239}]{\teiHi[rend={italic}]{Ibid.}, p. 159.}} ». L’attention du lecteur est de ce fait attirée sur les enjeux et les conséquences du secret dans la destinée du mémorialiste. On constate ainsi l’utilisation de la spécificité énonciative des Mémoires à des fins différentes : elle sert d’une part la force narrative des textes (raconter sa propre expérience), et d’autre part leur force rhétorique (convaincre de sa vérité sur l’histoire), deux forces qui entrent en tension et se complètent.}
              
\end{teiDiv}
            
\end{teiDiv}
          
\end{teiElement}
        
\end{teiElement}
        \begin{teiElement}[name={group},type={chapter},data-page-title={Un secret de polichinelle},data-page-subtitle={La révélation de la puissance de l’argent dans le Code des gens honnêtes (1825) de Balzac},data-page-authors={Marie-Agathe Tilliette@@Université du Littoral Côte d’Opale, HLLI, UR 4030, Boulogne-sur-Mer, France},data-include-href={../XML/chap\_7\_Un\_secret\_de\_polichinelle\_Locus\_politicus.xml},xmlid={chapter-006}]

          \begin{teiElement}[name={front}]

            \begin{teiDiv}[type={titlePage},xmlid={titlepage-007}]

              \teiP[rend={title-main},xmlid={p-240}]{Un secret de polichinelle}
              \teiP[rend={title-sub},xmlid={p-241}]{La révélation de la puissance de l’argent dans le \teiHi[rend={italic}]{Code des gens honnêtes} (1825) de Balzac}
              \teiP[rend={author-aut},xmlid={p-242}]{Marie-Agathe \teiHi[rend={small-caps}]{Tilliette}}
              \teiP[rend={authority\_affiliation},xmlid={p-243}]{Université du Littoral Côte d’Opale, HLLI, UR 4030, Boulogne-sur-Mer, France}
            
\end{teiDiv}
          
\end{teiElement}
          \begin{teiElement}[name={body},xmlid={body-007}]

            \teiP[xmlid={p1-7}]{Dans \teiHi[rend={italic}]{La Comédie humaine}, l’ordre social et politique trouve ses fondements dans des espaces cachés, qui se manifestent notamment à travers l’action des sociétés secrètes, depuis les Treize jusqu’aux Frères de la Consolation, en passant par divers avatars d’organisations clandestines comme les usuriers réunis au café Thémis dans \teiHi[rend={italic}]{Gobseck} ou encore le Cénacle des \teiHi[rend={italic}]{Illusions perdues}. Cette prégnance des réseaux souterrains dans la représentation du politique, qui correspond à la fameuse distinction opérée par Vautrin entre « l’Histoire officielle, menteuse qu’on enseigne » et « l’Histoire secrète » que seuls les initiés connaissent\teiNote[xmlid={ftn1-7},n={1},place={foot}]{\teiP[xmlid={p-244}]{Cette distinction éminemment balzacienne est formulée par Vautrin (\teiHi[rend={italic}]{alias} l’abbé Carlos Herrera) à l’intention de Lucien de Rubempré tenté par le suicide. Honoré de Balzac, \teiHi[rend={italic}]{Illusions perdues}, Paris, Gallimard, « Bibliothèque de la Pléiade », 1977, p. 695.}}, a fait l’objet de nombreuses études, liant approche herméneutique et enjeux poétiques dans \teiHi[rend={italic}]{La Comédie humaine}\teiNote[xmlid={ftn38-6},n={2},place={foot}]{\teiP[xmlid={p-245}]{Voir notamment Lyon-Caen Boris, \teiHi[rend={italic}]{Balzac et la comédie des signes : essai sur une expérience de pensée}, Saint-Denis, Presses universitaires de Vincennes, 2006 ; Chantal Massol, \teiHi[rend={italic}]{Une poétique de l’énigme : le récit herméneutique balzacien}, Genève, Droz, « Histoire des idées et critique littéraire », 2006 ; Anne-Marie Baron, \teiHi[rend={italic}]{Balzac occulte. Alchimie, magnétisme, sociétés secrètes}, Lausanne, L’Âge d’homme, 2012, p. 143-212 ; Rebecca Sugden, « Vers une poétique de la conspiration chez Balzac », Nicolas Aude et Marie-Agathe Tilliette (dir.), \teiHi[rend={italic}]{L’Imaginaire des sociétés secrètes dans la littérature du }\teiHi[rend={ small-caps italic}]{xix}\teiHi[rend={italic sup}]{e}\teiHi[rend={italic}]{ siècle}, actes de l’atelier de la Société des études romantiques et dix-neuviémistes, 28 février 2020, Bibliothèque de l’Arsenal (BNF), \teiRef[target={https://serd.hypotheses.org/8154}]{\teiHi[rend={underline}]{https://serd.hypotheses.org/8154}} (consulté le 4 août 2023).}}. C’est à partir d’un texte antérieur de Balzac, œuvre du publiciste plutôt que de l’écrivain, que l’on voudrait réfléchir aux lieux du secret politique : le \teiHi[rend={italic}]{Code des gens honnêtes} (1825). Ce code parodique, dont l’attribution à Balzac a parfois fait débat, doit être replacé dans l’entreprise collective d’un groupe de jeunes journalistes, menés par Horace Raisson, pour composer une série de détournements de la nouvelle pièce maîtresse de l’ordre social français, le Code civil ou Code Napoléon (1804)\teiNote[xmlid={ftn39-6},n={3},place={foot}]{\teiP[xmlid={p-246}]{Le \teiHi[rend={italic}]{Code des gens honnêtes }paraît en mars 1825, sans nom d’auteur, chez le libraire-éditeur parisien Jean-Nicolas Barba. Il fait l’objet d’une seconde édition au mois de juillet, augmentée seulement d’un frontispice, puis d’une troisième édition, légèrement modifiée et cette fois signée par Horace Raisson, en 1829 chez Nicolas Roret, sous le titre \teiHi[rend={italic}]{Code pénal, manuel complet des honnêtes gens}. Cette réédition a jeté le doute sur l’identité du ou des auteurs, mais l’on s’accorde aujourd’hui pour l’attribuer à Honoré de Balzac. Voir Balzac, \teiHi[rend={italic}]{Œuvres diverses}, Roland Chollet, Christiane Guise et René Guise (éd.), Paris, Gallimard, « Bibliothèque de la Pléiade », 1966, t. II, p. 1335-1338.}}. Sous-titré \teiHi[rend={italic}]{L’art de ne pas être dupe des fripons}, le code balzacien se présente comme un manuel enseignant aux « gens honnêtes » à échapper aux vols dont ils risquent quotidiennement d’être victimes, étant autant menacés par les voleurs et autres émanations des bas-fonds que par des figures prétendument honorables comme les notaires ou les agents de change. Car ce n’est pas la puissance souterraine des sociétés secrètes qui est ici révélée, mais une puissance plus fondamentale encore, nécessaire à l’existence des premières : celle de l’argent, dont la peinture parcourt, comme on le sait, toute l’œuvre balzacienne\teiNote[xmlid={ftn40-6},n={4},place={foot}]{\teiP[xmlid={p-247}]{Selon Alexandre Péraud, dans \teiHi[rend={italic}]{La Comédie humaine}, Balzac « montre comment l’argent est désormais devenu – pour le pire et pour le meilleur – la condition et la modalité du rapport que l’individu moderne entretient avec le réel, la cause et la conséquence de l’agir individuel, le moteur et le poison de la société. Bref, un principe, et non simplement un outil, puissamment intériorisé par l’\teiHi[rend={italic}]{homo œconomicus} en train d’advenir » (Alexandre Péraud dir., \teiHi[rend={italic}]{La Comédie (in)humaine de l’argent}, Lormont, Le Bord de l’eau, 2013, p. 13).}}, mais qui est ici traitée sous un jour ludique, comme secret de polichinelle. L’argent dont parle le \teiHi[rend={italic}]{Code des gens honnêtes}, ce n’est pas seulement la monnaie sonnante et trébuchante, mais plus largement toute possession matérielle susceptible d’être transformée en capital ; ce n’est pas seulement non plus l’argent des individus, mais plus largement toute richesse, jusqu’au niveau de l’État même. L’argent, nous dit le \teiHi[rend={italic}]{Code}, définit le jeu politique, à la fois au sens large \teiHi[rend={italic}]{du} politique (de l’organisation de la cité) et au sens plus précis de \teiHi[rend={italic}]{la} politique (des institutions et des gouvernants), et peut-être vaut-il mieux en rire qu’en pleurer\teiNote[xmlid={ftn41-6},n={5},place={foot}]{\teiP[xmlid={p-248}]{Sur le traitement comique du motif de l’argent au \teiHi[rend={small-caps}]{xix}\teiHi[rend={sup}]{e} siècle, voir le volume dirigé par Florence Fix et Marie-Ange Fougère, \teiHi[rend={italic}]{L’Argent et le rire de Balzac à Mirbeau}, Rennes, PUR, « Interférences », 2012. Le premier article de ce volume fait l’étude d’un autre texte parodique de Balzac, publié dix ans plus tard, qui décrit le règne de l’argent, « Melmoth réconcilié » (\teiHi[rend={italic}]{Le Livre des conteurs}, 1835) : Émilie Pézard\teiHi[rend={small-caps}]{, }« Melmoth à la banque : la réécriture ironique du roman noir dans \teiHi[rend={italic}]{Melmoth réconcilié} de Balzac », \teiHi[rend={italic}]{L’Argent et le rire de Balzac à Mirbeau}, Florence Fix et Marie-Ange Fougère (dir.), \teiHi[rend={italic}]{op. cit.}, p. 21-34.}}.}
            \begin{teiDiv}[type={section1},xmlid={div01-6}]

              \teiHead[xmlid={rire-de-la-puissance-de-largent-001}]{Rire de la puissance de l’argent}
              \teiP[xmlid={p2-7}]{Le \teiHi[rend={italic}]{Code des gens honnêtes} s’ouvre par un avant-propos qui clame haut et fort la toute-puissance de l’argent et les menaces pesant sur un objet aussi désirable :}
              \begin{teiCit}[xmlid={cit01-7}]

                \begin{teiQuote}
L’argent, par le temps qui court, donne le plaisir, la considération, les amis, les succès, les talents, l’esprit même ; ce doux métal doit donc être l’objet constant de l’amour et de la sollicitude des mortels de tout âge, de toute condition, depuis les rois jusqu’aux grisettes, depuis les propriétaires jusqu’aux émigrés.\teiLb{}Mais cet argent, source de tous les plaisirs, origine de toutes les gloires, est aussi le but de toutes les tentatives.\teiLb{}La vie peut être considérée comme un combat perpétuel entre les riches et les pauvres. Les uns sont retranchés dans une place forte à murs d’airain, pleine de munitions ; les autres tournent, virent, sautent, attaquent, rongent les murailles ; et malgré les ouvrages à cornes que l’on bâtit, en dépit des portes, des fossés, des batteries, il est rare que les assiégeants, ces cosaques de l’état social, n’emportent pas quelque avantage\teiNote[xmlid={ftn42-6},n={6},place={foot}]{\teiP[xmlid={p-249}]{Balzac, \teiHi[rend={italic}]{Code des gens honnêtes}, dans \teiHi[rend={italic}]{Œuvres diverses}, t. II, \teiHi[rend={italic}]{op. cit.}, p. 147.}}.
\end{teiQuote}
              
\end{teiCit}
              \teiP[xmlid={p3-7}]{Le registre héroï-comique qui ouvre le \teiHi[rend={italic}]{Code} en révèle toute l’ambivalence : d’un côté, il se place très explicitement sous le signe de l’humour par l’exagération des énumérations et l’expression de la totalité ; de l’autre, il aborde de front les enjeux de la marchandisation de la vie moderne et des inégalités sociales, dont la peinture, plus souvent tragique que comique, sera au cœur de \teiHi[rend={italic}]{La Comédie humaine}. L’évocation de la vie comme « combat perpétuel entre les riches et les pauvres » n’a d’ailleurs pas manqué d’être relevée par la critique balzacienne et parfois traduite en termes de lutte des classes\teiNote[xmlid={ftn43-6},n={7},place={foot}]{\teiP[xmlid={p-250}]{Jean-Hervé Donnard, \teiHi[rend={italic}]{Les Réalités économiques et sociales dans }La Comédie humaine, Paris, Armand Colin, 1961, p. 32-33.}}. Comme l’écrit Pierre Barbéris, le \teiHi[rend={italic}]{Code des gens honnêtes} est sans doute le premier texte de Balzac qui « pose directement le problème moderne » qu’est la domination de l’argent, devenu « fin en soi » plutôt que moyen\teiNote[xmlid={ftn44-5},n={8},place={foot}]{\teiP[xmlid={p-251}]{Pierre Barbéris, \teiHi[rend={italic}]{Balzac et le mal du siècle. Contribution à une physiologie du monde moderne}, tome I :\teiHi[rend={italic}]{ 1799-1829, Une expérience de l’absurde : aliénations et prises de conscience}, Paris, Gallimard, « NRF », 1970, p. 702 et 703.}}.}
              \teiP[xmlid={p4-7}]{Toutefois, l’humour reste la règle principale de la composition éclectique du \teiHi[rend={italic}]{Code des gens honnêtes} : les observations sociales et la description minutieuse des escroqueries se font par le biais d’anecdotes humoristiques parfois tendancieuses, de conseils ironiques et d’aphorismes héroï-comiques. L’éclectisme de la composition du \teiHi[rend={italic}]{Code} tient certes à la rapidité d’une rédaction désinvolte, également révélée par l’inexactitude presque systématique des renvois internes, mais il permet de suggérer et même d’incarner dans toute sa mystérieuse diversité la puissance de l’argent. Ses articles varient de la simple phrase à de véritables petits canevas romanesques ou dramatiques. Pour n’en donner qu’un exemple, l’avertissement suivant, « Songez que souvent vous pourrez être le centre de toute une intrigue, et que deux, trois ou quatre acteurs différents joueront leurs rôles pour faire sortir de votre magasin ou de votre poche cette précieuse panacée, \teiHi[rend={italic}]{l’argent}\teiNote[xmlid={ftn45-5},n={9},place={foot}]{\teiP[xmlid={p-252}]{Balzac, \teiHi[rend={italic}]{Code des gens honnêtes}, \teiHi[rend={italic}]{op. cit.}, p. 175.}} ! », est suivi d’une illustration détaillée, avec descriptions et dialogues : un Anglais met en dépôt chez une lingère d’excellents crayons à papier, dont la vente devrait rapporter un fort pourcentage ; un jeune homme fort bien mis les achète mais n’a pas l’argent sur lui et la lingère avance les frais ; ni l’Anglais ni le jeune homme ne reparaissent jamais et les crayons entreposés s’avèrent défectueux ; la lingère est volée\teiNote[xmlid={ftn46-5},n={10},place={foot}]{\teiP[xmlid={p-253}]{\teiHi[rend={italic}]{Ibid}., p. 175-177.}}. Ce canevas est encadré de conseils mi-sérieux mi-humoristiques, dont on peut se demander s’ils font l’éloge des voleurs ou la critique de la naïveté des dupes, les uns comme les autres étant menés par une forme de cupidité.}
              \teiP[xmlid={p5-7}]{Les trois livres du \teiHi[rend={italic}]{Code des gens honnêtes} se construisent en \teiHi[rend={italic}]{crescendo} : on part des vols les plus concrets, vols à la tire et cambriolages, pour arriver aux manipulations psychologiques plus subtiles du parasite, de l’« ami de collège dans le malheur\teiNote[xmlid={ftn47-5},n={11},place={foot}]{\teiP[xmlid={p-254}]{\teiHi[rend={italic}]{Ibid.}, p. 220.}} » ou encore du sacristain, avant de conclure par la virtuosité mise en œuvre « dans ce labyrinthe que l’on nomme palais\teiNote[xmlid={ftn48-5},n={12},place={foot}]{\teiP[xmlid={p-255}]{\teiHi[rend={italic}]{Ibid}., p. 264.}} » par les avoués, notaires et avocats, grands maîtres ès escroqueries. Les premiers vols peuvent faire l’objet de conseils pratiques (« Rien n’est si utile que de garder toute la nuit une lumière vive dans son appartement\teiNote[xmlid={ftn49-5},n={13},place={foot}]{\teiP[xmlid={p-256}]{\teiHi[rend={italic}]{Ibid}., p. 198.}} »), qui incitent à prendre toutes les précautions possibles et à choisir la prudence, fût-ce au prix du ridicule (« D’honorables personnes mettent leurs mouchoirs dans leurs chapeaux\teiNote[xmlid={ftn50-5},n={14},place={foot}]{\teiP[xmlid={p-257}]{\teiHi[rend={italic}]{Ibid}., p. 165.}} »). Les suivants relèvent davantage d’une logique de révélation, que ce soit du mensonge de qui habille son intérêt égoïste des oripeaux d’une noble cause, ou des incompréhensibles procédures des avoués, dont l’unique objectif est de pressurer leurs clients. Le rédacteur du \teiHi[rend={italic}]{Code} se targue alors, par le détour humoristique d’un vers de \teiHi[rend={italic}]{Bajazet}, de posséder de précieuses connaissances d’initié : « Nourri dans le sérail, j’en connais les détours\teiNote[xmlid={ftn51-5},n={15},place={foot}]{\teiP[xmlid={p-258}]{\teiHi[rend={italic}]{Ibid}., p. 243. Jean Racine, \teiHi[rend={italic}]{Bajazet}, acte IV, scène 8, v. 1424.}} », affirme-t-il, à l’instar du rusé politique qu’est le grand vizir Acomat.}
              \teiP[xmlid={p6-7}]{Le sème spatial propre au secret (présent dans son origine latine \teiHi[rend={italic}]{secretum}, « lieu écarté ») s’incarne dans une série de lieux de l’ombre, suggérés plus que décrits. Le dédale du Palais de justice, comparé au « sérail », en est le plus explicite, mais l’ensemble du \teiHi[rend={italic}]{Code} repose sur une obscure profondeur spatiale en amont et en aval des escroqueries. Il y a d’abord les lieux où les vols s’élaborent, puis ceux où les richesses et les voleurs disparaissent dans un laps de temps remarquablement bref. Un des enjeux du \teiHi[rend={italic}]{Code} est de faire pénétrer plus ou moins profondément les lectrices et lecteurs dans ce monde mystérieux.}
              \teiP[xmlid={p7-7}]{Ainsi, on trouve dans le \teiHi[rend={italic}]{Code des gens honnêtes} le procédé de révélation qui est au fondement même du réalisme balzacien et qui fait des sphères du pouvoir économique et politique une série de lieux cachés reliés en réseaux mystérieux, mis en lumière par la narration. Toutefois, comme l’a montré Chantal Massol, la révélation ne se fait jamais de manière univoque dans \teiHi[rend={italic}]{La Comédie humaine} et il en va de même dans le \teiHi[rend={italic}]{Code des gens honnêtes} où le traitement ludique de l’adoration du veau d’or et de ses avatars monétaires induit un jeu avec le secret : de même que les narrateurs de \teiHi[rend={italic}]{La Comédie humaine} construisent et déconstruisent à différents niveaux la position de domination des figures détentrices du secret (depuis les chefs des sociétés secrètes jusqu’au romancier lui-même\teiNote[xmlid={ftn52-5},n={16},place={foot}]{\teiP[xmlid={p-259}]{Chantal Massol, « “Rois inconnus”. De l’usage du secret dans le roman balzacien », Sylvie Triaire et Alain Vaillant (dir.), \teiHi[rend={italic}]{Écritures du pouvoir et pouvoirs de la littérature}, Montpellier, Presses universitaires de la Méditerranée, « Collection des littératures », 2001, p. 233-255, DOI : 10.4000/books.pulm.193. Chantal Massol résume ainsi la poétique balzacienne (§ 2) : « le roman balzacien pose inlassablement la société comme énigme, pour mieux la déchiffrer ; il y aperçoit sans relâche des secrets, pour mieux les révéler. »}}), le rédacteur du \teiHi[rend={italic}]{Code} se divertit à faire de l’argent un fondement social et politique à la fois caché et évident, agissant dans l’ombre et éclatant au grand jour. Que l’argent mène le monde, ce n’est guère une sensationnelle révélation, c’est un secret connu de tous, un secret de polichinelle.}
              \teiP[xmlid={p8-7}]{Le \teiHi[rend={italic}]{Code} semble donc osciller entre deux objectifs qui paraissent contradictoires dans leur rapport au secret : constater ce que chacun peut voir, et révéler ce qui est caché. Les premiers paragraphes de l’avant-propos, que l’on a déjà cités, surjouent la révélation, mais quelques pages plus loin, l’écrivain avertit ses lecteurs et lectrices de ne pas tomber dans les poncifs qu’il vient lui-même d’employer : « On dit toujours : “Actuellement l’argent est tout, celui qui a de l’argent est maître de tout.” Ah ! gardez-vous de répéter ces phrases banales, vous auriez l’air d’un niais\teiNote[xmlid={ftn53-6},n={17},place={foot}]{\teiP[xmlid={p-260}]{Balzac, \teiHi[rend={italic}]{Code des gens honnêtes}, \teiHi[rend={italic}]{op. cit.}, p. 153.}}. » Certes, ce qu’il nie dans la suite du texte, en en profitant pour se donner une place dans la lignée valorisante des grands satiristes latins, ce n’est pas tant le règne de l’argent que le caractère inédit de son étendue\teiNote[xmlid={ftn54-6},n={18},place={foot}]{\teiP[xmlid={p-261}]{« Celui qui a estropié Juvénal, Horace, et les auteurs de toutes les nations, doit savoir que, de tout temps, l’argent a été chéri et recherché avec une ardeur égale », \teiHi[rend={italic}]{ibid}.}}, mais les premières phrases du \teiHi[rend={italic}]{Code}, affirmant l’omnipotence de l’argent « par le temps qui court », lui donnent malgré tout, \teiHi[rend={italic}]{a posteriori}, « l’air d’un niais » – ou d’un polichinelle. Même en considérant qu’il y a bien révélation dans le détail de certains procédés employés, celle-ci n’en reste pas moins parodique par le double recours constant à l’exagération et à l’ironie.}
              \teiP[xmlid={p9-7}]{De plus, les protagonistes de ce secret de comédie ne sont pas les mystérieux affiliés d’une loge maçonnique ou de ténébreux conspirateurs, mais tout un chacun, « depuis les rois jusqu’aux grisettes, depuis les propriétaires jusqu’aux émigrés\teiNote[xmlid={ftn55-6},n={19},place={foot}]{\teiP[xmlid={p-262}]{\teiHi[rend={italic}]{Ibid}., p. 147.}} ». Toute la galerie de types que met en scène le \teiHi[rend={italic}]{Code}, porteurs d’un nom topique ou d’une simple initiale, constitue un ample panorama de la société française de la Restauration\teiNote[xmlid={ftn56-6},n={20},place={foot}]{\teiP[xmlid={p-263}]{Comme le souligne Boris Lyon-Caen en étudiant l’équivocité générique du \teiHi[rend={italic}]{Code des gens honnêtes}, celui-ci ressort à certains égards à la littérature panoramique, que l’on peut considérer comme « proto-sociologique » (« Du sérieux à l’essai : lecture du \teiHi[rend={italic}]{Code des gens honnêtes} (Balzac) », \teiHi[rend={italic}]{Romantisme}, vol. 2, n\teiHi[rend={sup}]{o} 164, 2014, p. 41-49).}}, et on s’amuse de leur confrontation, car les plus manipulateurs ne sont pas ceux qu’on croit (« Entre la parole d’honneur d’un avoué et celle d’une actrice, n’hésitez pas : croyez l’actrice\teiNote[xmlid={ftn57-6},n={21},place={foot}]{\teiP[xmlid={p-264}]{Balzac, \teiHi[rend={italic}]{Code des gens honnêtes}, \teiHi[rend={italic}]{op. cit.}, p. 220.}} »). Le ton de la révélation adopté par le rédacteur du \teiHi[rend={italic}]{Code} fait ainsi partie de la parodie : non pas détournement juridique ici, mais paradoxalement parodie du réalisme balzacien lui-même, avant même qu’il ne prenne toute son ampleur dans \teiHi[rend={italic}]{La Comédie humaine}.}
            
\end{teiDiv}
            \begin{teiDiv}[type={section1},xmlid={div02-6}]

              \teiHead[xmlid={collusion-du-politique-et-de-leconomique-001}]{Collusion du politique et de l’économique}
              \teiP[xmlid={p10-7}]{La sphère politique à proprement parler ne fait pas exception au règne de l’argent décrit par le \teiHi[rend={italic}]{Code des gens honnêtes} et si aucun « titre\teiNote[xmlid={ftn58-6},n={22},place={foot}]{\teiP[xmlid={p-265}]{Un « titre » est l’une des subdivisions du texte de loi, que reprend la parodie du \teiHi[rend={italic}]{Code}.}} » ne lui est réservé, c’est que le Code civil parodié n’a pas pour vocation de statuer sur les institutions politiques, mais d’organiser les relations sociales. Le \teiHi[rend={italic}]{locus politicus} qui transparaît dans ce tableau panoramique est lui aussi gangréné par le pouvoir inédit de l’argent, comme l’illustre cette considération du titre II du Livre I, intitulé « Escroqueries » :}
              \begin{teiCit}[xmlid={cit02-6}]

                \begin{teiQuote}
Un jour, en France, un Écossais, nommé Law, se trouva en train d’escroquer tout le royaume. Ce Law passe généralement pour le plus grand homme produit par la classe des coquins ; mais aujourd’hui les politiques avouent qu’il est le fondateur du système des banques et du crédit. Vous voyez qu’il y a des temps où une escroquerie est mieux reçue qu’en d’autres. Par le temps qui court, l’Écossais serait peut-être ministre inamovible\teiNote[xmlid={ftn59-6},n={23},place={foot}]{\teiP[xmlid={p-266}]{Balzac, \teiHi[rend={italic}]{Code des gens honnêtes}, \teiHi[rend={italic}]{op. cit.}, p. 191-192.}}.
\end{teiQuote}
              
\end{teiCit}
              \teiP[xmlid={p11-7}]{De fait, John Law, dont le système monétaire mis en place entre 1715 et 1720 fragilisa les finances de l’État français, est bien l’un des précurseurs de la finance moderne. Le rédacteur du \teiHi[rend={italic}]{Code des gens honnêtes} y voit l’une des preuves de la toute-puissance de l’argent et estime avec ironie que celui qu’il peint comme le plus grand des escrocs pourrait prétendre aux plus hautes dignités du monde contemporain : la sphère politique ainsi suggérée est dominée par les intérêts économiques de quelques-uns.}
              \teiP[xmlid={p12-7}]{Une telle soumission de la sphère politique à l’intérêt économique est également représentée dans le chapitre liminaire du \teiHi[rend={italic}]{Code des gens honnêtes}, intitulé « Considérations morales, politiques, littéraires, philosophiques, législatives, religieuses et budgétaires, sur la compagnie des voleurs ». Là encore, le registre est celui de la révélation parodique, dans la mesure où il s’agit de mettre au jour le pouvoir des voleurs, dépeints comme une corporation ou même une micro-société tentaculaire. Si une telle révélation ne déparerait pas les descriptions des bas-fonds de \teiHi[rend={italic}]{Splendeurs et misères des courtisanes}, elle est ici placée sous le signe de l’exagération. Le narrateur, se posant en révélateur d’un secret connu de lui seul, se fabrique un \teiHi[rend={italic}]{ethos} de bonimenteur plutôt que d’historien des mœurs contemporaines :}
              \begin{teiCit}[xmlid={cit03-6}]

                \begin{teiQuote}
Nous avons \lbrack{}…\rbrack{} cru nécessaire, avant de tenter de dévoiler les ruses des voleurs privilégiés comme non privilégiés de toutes les classes, de nous livrer à des considérations impartiales sur les voleurs ; nous seuls, peut-être, pouvions les examiner sous toutes leurs faces avec sang-froid ; et certes, on ne nous accusera pas de vouloir les défendre, nous qui leur coupons les vivres, et signalons toutes leurs opérations, en élevant dans ce livre un phare qui les domine\teiNote[xmlid={ftn60-6},n={24},place={foot}]{\teiP[xmlid={p-267}]{\teiHi[rend={italic}]{Ibid}., p. 150. Le registre publicitaire revient à plusieurs reprises dans le \teiHi[rend={italic}]{Code}, comme par exemple dans cette rassurante affirmation sur la rentabilité de la dépense faite : « c’est un conseil qui vaut plus que le petit écu avec lequel vous aurez acheté ce livre » (\teiHi[rend={italic}]{ibid}., p. 252).}}.
\end{teiQuote}
              
\end{teiCit}
              \teiP[xmlid={p13-7}]{Le registre héroï-comique se poursuit dans la suite des « Considérations » qui font un éloge contradictoire et hyperbolique des voleurs, « classe industrieuse et commerçante\teiNote[xmlid={ftn61-6},n={25},place={foot}]{\teiP[xmlid={p-268}]{\teiHi[rend={italic}]{Ibid}., p. 155.}} » indispensable à la survie de toute société, avant de les décrire comme un État dans l’État, une « république qui a ses lois et ses mœurs\teiNote[xmlid={ftn62-6},n={26},place={foot}]{\teiP[xmlid={p-269}]{\teiHi[rend={italic}]{Ibid.}}} ». Les fantaisistes minutes du parlement des voleurs londoniens, mieux organisés – nous informe le rédacteur du \teiHi[rend={italic}]{Code} – que leurs confrères parisiens, poussent à son comble la parodie du pouvoir politique. Cette représentation carnavalesque n’a rien de bien inquiétant dans son ensemble, au contraire, elle renvoie à un imaginaire ancien de la gueuserie comme monde à l’envers, plutôt qu’à des inquiétudes contemporaines\teiNote[xmlid={ftn63-6},n={27},place={foot}]{\teiP[xmlid={p-270}]{À propos de l’imaginaire de la gueuserie, voir Bronisław Geremek, \teiHi[rend={italic}]{Les Fils de Caïn. L’image des pauvres et des vagabonds dans la littérature européenne du }\teiHi[rend={ small-caps italic}]{xv}\teiHi[rend={italic sup}]{e}\teiHi[rend={italic}]{ au }\teiHi[rend={ small-caps italic}]{xvii}\teiHi[rend={italic sup}]{e}\teiHi[rend={italic}]{ siècle} \lbrack{}1980\rbrack{}, Paris, Flammarion, 1991. Sur les contre-sociétés dans l’imaginaire social des bas-fonds au \teiHi[rend={small-caps}]{xix}\teiHi[rend={sup}]{e} siècle, voir Dominique Kalifa, \teiHi[rend={italic}]{Les Bas-fonds. Histoire d’un imaginaire}, Paris, Éditions du Seuil, « L’univers historique », 2013, en particulier p. 61-67. Louis Chevalier mentionne lui aussi le \teiHi[rend={italic}]{Code des gens honnêtes} en soulignant que la criminalité qui y est peinte relève d’une forme ancienne des « classes dangereuses », qui ne se confondent pas encore avec les « classes laborieuses ». Louis Chevalier, \teiHi[rend={italic}]{Classes laborieuses et classes dangereuses à Paris pendant la première moitié du }\teiHi[rend={ small-caps italic}]{xix}\teiHi[rend={italic sup}]{e}\teiHi[rend={italic}]{ siècle }\lbrack{}1958\rbrack{}, Paris, Plon, 2002, p. 59-64.}}. Toutefois, quelques remarques ponctuelles donnent une profondeur supplémentaire à cette description ludique, à l’instar d’une allusion aux tensions politiques et sociales que connaît la voisine Angleterre :}
              \begin{teiCit}[xmlid={cit04-4}]

                \begin{teiQuote}
Lorsque les barrières dont les lois entourent le bien d’autrui sont franchies, il faut reconnaître un invincible besoin, une fatalité ; car enfin la société ne donne même pas du pain à tous ceux qui ont faim ; et, quand ils n’ont aucun moyen d’en gagner, que voulez-vous qu’ils fassent ? Mais, bien plus, le jour où la masse des malheureux sera plus forte que la masse des riches, l’état social sera tout autrement établi ; et en ce moment l’Angleterre est menacée d’une révolution de ce genre\teiNote[xmlid={ftn64-5},n={28},place={foot}]{\teiP[xmlid={p-271}]{Balzac, \teiHi[rend={italic}]{Code des gens honnêtes}, \teiHi[rend={italic}]{op. cit.}, p. 152.}}.
\end{teiQuote}
              
\end{teiCit}
              \teiP[xmlid={p14-6}]{On voit que l’exagération comique n’est finalement pas si éloignée de réelles inquiétudes sociales, bien justifiées en ces décennies marquées par les bouleversements de la première industrialisation.}
              \teiP[xmlid={p15-5}]{En outre, au-delà de la parodie, le fait que la « corporation » des voleurs puisse être décrite comme une entité politique met en relief la collusion, pour le rédacteur du \teiHi[rend={italic}]{Code}, des domaines économique et politique, non pas tant par ce qui pourrait être la dénonciation d’une forme d’oligarchie ou de ploutocratie que par l’identité parfaite des enjeux des deux sphères. L’\teiHi[rend={italic}]{homo politicus }est un \teiHi[rend={italic}]{homo œconomicus} (c’est-à-dire un individu dont l’action rationnelle est fondée sur la « maximisation de la satisfaction\teiNote[xmlid={ftn65-5},n={29},place={foot}]{\teiP[xmlid={p-272}]{Une telle conception commence à émerger chez David Hume et Adam Smith et se trouve au fondement de l’économie politique au \teiHi[rend={small-caps}]{xix}\teiHi[rend={sup}]{e} siècle. À ce propos, voir Pierre Demeulenaere, \teiHi[rend={italic}]{Homo œconomicus. Enquête sur la constitution d’un paradigme} \lbrack{}1996\rbrack{}, Paris, PUF, « Quadrige », 2003.}} ») et tel est peut-être le grand enjeu sérieux du \teiHi[rend={italic}]{Code des gens honnêtes}. Il n’y a pas seulement soumission de la sphère politique aux enjeux économiques, mais confusion des deux domaines. Les lieux du pouvoir politique réel se révèlent être ceux du pouvoir économique, là où s’échange, s’accumule, se multiplie l’argent.}
            
\end{teiDiv}
            \begin{teiDiv}[type={section1},xmlid={div03-5}]

              \teiHead[xmlid={un-rire-prophylactique-001}]{Un rire prophylactique}
              \teiP[xmlid={p16-5}]{En riant de l’emprise de l’argent sur les esprits, en s’amusant de la cupidité généralisée qui est le véritable fondement de l’ordre social et politique moderne, le rédacteur du \teiHi[rend={italic}]{Code des gens honnêtes} se pose en satiriste et son texte n’est pas dépourvu d’une fonction morale. \teiHi[rend={italic}]{Castigat ridendo mores }: il décrit certes les efforts acharnés et incessants des voleurs de tous types pour s’emparer des possessions d’autrui, mais il n’épargne pas non plus les « gens honnêtes » auxquels il s’adresse, en peignant les petits calculs mesquins des parents cherchant à marier leur fille ou l’amour-propre aveugle du bourgeois qui se croit auteur. Dans le canevas résumé plus haut, l’avidité de la lingère n’est finalement peut-être pas plus grande que celle des escroqueurs qui, en un sens, méritent leur gain par leur inventivité. Le rapport à l’argent est le même et si certains restent à l’intérieur du cadre légal, ce n’est pas tant par honnêteté que par manque d’imagination ou du fait de leur situation privilégiée. À cet égard, le \teiHi[rend={italic}]{Code des gens honnêtes} est une sorte de physiologie universelle du rapport à la possession, dans laquelle l’humour du rédacteur fait émerger un univers moral en demi-teinte. Avant la longue thérapie curative de \teiHi[rend={italic}]{La Comédie humaine}, qui analyse les causes et les conséquences du règne de l’argent dans la société contemporaine, le \teiHi[rend={italic}]{Code des gens honnêtes} apparaît comme un traitement prophylactique, dont l’un des plus efficaces ferments est le détournement ludique du paradigme du secret.}
              \teiP[xmlid={p17-5}]{De plus, s’il s’adresse avant tout aux individus, il étend peu à peu la portée de son propos jusqu’aux gouvernements dans le troisième et dernier livre. Alors que le premier chapitre de ce livre énumère les accommodements louches et les malversations des notaires et des avoués, le second élargit le sujet aux institutions et pratiques publiques que sont les monts-de-piété, la loterie royale, les maisons de jeux et les dettes publiques. Toutes quatre sont dénoncées, avec souvent plus de véhémence que d’humour, comme des escroqueries légales. La critique faite du mont-de-piété, nommé « trafic infâme » et même « brigandage », dépasse celle que l’on trouvera dans les pages des \teiHi[rend={italic}]{Mystères de Paris}\teiNote[xmlid={ftn66-5},n={30},place={foot}]{\teiP[xmlid={p-273}]{Eugène Sue, \teiHi[rend={italic}]{Les Mystères de Paris} \lbrack{}1842-1843\rbrack{}, Paris, Robert Laffont, « Bouquins », 1989, voir en particulier la note 1 de l’auteur p. 1091, qui cite les analyses d’Alphonse Esquiros.}}. En outre, respectant la logique graduée du reste du \teiHi[rend={italic}]{Code}, son rédacteur est plus virulent encore contre la loterie et les maisons de jeux, et il présente ainsi ces dernières : « À Paris, dans cette capitale du monde civilisé, dans ce centre de sociabilité, de commerce, d’industrie, \lbrack{}…\rbrack{} il existe des maisons où l’usure et le vol sont autorisés ; où la ruine, le désespoir, le suicide sont affermés, et rapportent au gouvernement des sommes immenses que l’on peut appeler le prix du sang\teiNote[xmlid={ftn67-5},n={31},place={foot}]{\teiP[xmlid={p-274}]{Balzac, \teiHi[rend={italic}]{Code des gens honnêtes}, \teiHi[rend={italic}]{op. cit.}, p. 270. Le terme « affermer » appartient au lexique du droit fiscal et désigne le droit de percevoir des taxes. Quelques lignes plus loin, le rédacteur fustige « l’immoralité de ce revenu flétrissant du fisc » (\teiHi[rend={italic}]{ibid}., p. 271).}}. » En autorisant et en tirant des revenus de telles institutions, le gouvernement fait passer son intérêt financier devant toute considération de morale publique ; il est désigné par cette ironique périphrase : « l’autorité qui respecte la morale publique quand son intérêt ne lui conseille pas de la violer\teiNote[xmlid={ftn68-2},n={32},place={foot}]{\teiP[xmlid={p-275}]{\teiHi[rend={italic}]{Ibid.}, p. 270.}}. » En d’autres termes, le pouvoir politique est autant rongé par la cupidité que les individus dont le \teiHi[rend={italic}]{Code} a longuement peint les agissements.}
              \teiP[xmlid={p18-4}]{Par conséquent, tout comme les individus risquent de devenir les dupes de leur propre convoitise, le gouvernement est susceptible d’être escroqué et c’est à ce danger que se consacre le dernier paragraphe du \teiHi[rend={italic}]{Code}, intitulé « Emprunts. Dettes publiques ». Le ton, d’abord très ironique pour vanter la disparition de toutes les inégalités grâce au système des emprunts, devient plus sérieux en rappelant les menaces qui pèsent sur l’État :}
              \begin{teiCit}[xmlid={cit05-4}]

                \begin{teiQuote}
Combien d’exemples cependant devraient mettre en garde contre cette avidité de chances. Les gouvernements font banqueroute comme les particuliers ; et les gouvernements ne craignent pas les galères. La France a jadis ruiné ses sujets en les réduisant aux deux tiers ; aujourd’hui même elle est à la veille de faire une banqueroute d’une espèce nouvelle, en convertissant les rentes cinq pour cent en trois pour cent. On trouve cependant toujours des gens disposés à acheter, à brocanter, sur ces valeurs idéales qui rapportent moins qu’une maison, qu’une terre, et que l’on ne peut cependant faire assurer contre la grêle et l’incendie\teiNote[xmlid={ftn69-2},n={33},place={foot}]{\teiP[xmlid={p-276}]{\teiHi[rend={italic}]{Ibid}., p. 272.}}.
\end{teiQuote}
              
\end{teiCit}
              \teiP[xmlid={p19-4}]{Le souvenir de l’échec du système de Law, évoqué à nouveau dans cet ultime paragraphe, doit servir d’avertissement à la fois général et particulier, et souligne les dangers de l’argent devenu virtuel, « idéal » au sens littéral. Quoique l’annonce récurrente d’un quatrième livre et l’absence de conclusion donnent l’impression d’une rédaction désinvolte plus liée aux contraintes de librairie qu’à un dessein réfléchi, peut-être doit-on voir dans cette fin abrupte un réel point d’aboutissement : l’argent s’est autonomisé dans un monde virtuel jusqu’à présenter une menace pour la seule autorité qui aurait pu maîtriser la puissance de ce golem moderne. La révélation reprend alors tout son sens : le secret de polichinelle qu’est la puissance de l’argent est tellement omniprésent, tellement incontournable, tellement évident, qu’on en vient à l’oublier. Le secret politique de la société de la Restauration, qui est fondamentalement un secret économique, est d’autant mieux caché dans ses détails et ses procédés qu’il est exhibé dans ses principes.}
              \teiP[xmlid={p20-4}]{À l’issue de cette réflexion, le \teiHi[rend={italic}]{Code des gens honnêtes} apparaît bien comme un texte charnière dans la production du jeune Balzac, non seulement en ce qu’il contient de nombreux motifs, thèmes et personnages types qui seront repris et développés dans \teiHi[rend={italic}]{La Comédie humaine}\teiNote[xmlid={ftn70-2},n={34},place={foot}]{\teiP[xmlid={p-277}]{Pierre Barbéris, \teiHi[rend={italic}]{op. cit.}, p. 702 ; Boris Lyon\teiHi[rend={small-caps}]{-}Caen, art. cité, p. 43-44.}}, mais aussi, plus profondément, en ce qu’il expérimente paradoxalement une forme de détournement parodique de la logique réaliste, fondée sur la révélation d’un secret, avant même la pleine expression de celle-ci. Le secret – la puissance de l’argent – est en fait connu de tous et il n’est pas certain que cette non-révélation vaille « le petit écu » qui a servi à acheter le livre\teiNote[xmlid={ftn71-2},n={35},place={foot}]{\teiP[xmlid={p-278}]{Balzac, \teiHi[rend={italic}]{Code des gens honnêtes}, \teiHi[rend={italic}]{op. cit.}, p. 252.}}. C’est dans un second temps que l’attitude herméneutique prend sens au-delà de la parodie, dans la mesure où les enjeux économiques sont si considérables qu’ils en viennent à se confondre avec l’ordre social et politique. Le secret de polichinelle devient alors la « lettre volée » de la nouvelle d’Edgar Allan Poe : si voyant qu’il est invisible. Le rappel d’une telle évidence, d’abord adressé aux individus, intéresse l’État lui-même qui se comporte en \teiHi[rend={italic}]{homo œconomicus}.}
              \teiP[xmlid={p21-4}]{Les lieux du secret politique se confondent avec les lieux, de plus en plus virtuels, où l’argent prend son autonomie, où il établit ses propres lois. Si la réunion du parlement des voleurs londoniens n’est déjà que la partie émergée de leurs activités (et encore est-elle bien sûr réservée aux seuls membres de la corporation), \teiHi[rend={italic}]{a fortiori} les lieux du politique ne sont que la couverture de ce qui se joue dans d’autres lieux, ces lieux mystérieux où s’élaborent les vols, les escroqueries légales et les institutions publiques comme les monts-de-piété ou la loterie, tous placés sur le même plan par le rédacteur du \teiHi[rend={italic}]{Code}. On a affaire, dans ce texte bigarré et dont le statut auctorial demeure ambigu, à un récit détourné du secret politique, composé à une époque où « l’omnipotence, l’omniscience, l’omniconvenance\teiNote[xmlid={ftn72-2},n={36},place={foot}]{\teiP[xmlid={p-279}]{Balzac, \teiHi[rend={italic}]{La Maison Nucingen}, Paris, Gallimard, « Bibliothèque de la Pléiade », 1977, t. VI, p. 331.}} » de l’argent renouvelle les lieux de l’exercice du pouvoir. Les différents niveaux de parodie du Code civil correspondent bien à la rouerie des « fripons » et de leur « art », et ils introduisent un ultime jeu de miroir, en ce que le rédacteur lui-même, publiciste et bonimenteur, n’est pas si éloigné des escrocs beaux parleurs qu’il met en scène. Les lecteurs et lectrices qui, rappelons-le, ne sont pas épargnés par la satire du \teiHi[rend={italic}]{Code}, semblent placés dans un palais des glaces où l’on en vient à ne plus savoir ce qui est l’image reflétée et ce qui est la réalité\teiNote[xmlid={ftn73-2},n={37},place={foot}]{\teiP[xmlid={p-280}]{L’attraction foraine qu’est le palais des glaces ou labyrinthe de miroirs trouverait son origine dans la galerie des Glaces de Versailles et on peut considérer que l’infini jeu des reflets, tellement visibles qu’ils troublent la perception de la réalité, est une représentation fidèle des fastes du pouvoir politique et de sa démocratisation au \teiHi[rend={small-caps}]{xix}\teiHi[rend={sup}]{e} siècle.}} : le secret, ce n’est pas ce qui est caché, mais ce qu’on ne sait plus reconnaître, et ses lieux ne sont plus seulement impossibles à situer, ils sont impossibles à comprendre.}
            
\end{teiDiv}
          
\end{teiElement}
        
\end{teiElement}
        \begin{teiElement}[name={group},type={chapter},data-page-title={L’échange épistolaire, lieu du secret politique dans La Fortune des Rougon},data-page-authors={Étienne Poirier@@Université McGill, Montréal, Canada},data-include-href={../XML/chap\_8\_L\_echange\_epistolaire\_Locus\_politicus.xml},xmlid={chapter-007}]

          \begin{teiElement}[name={front}]

            \begin{teiDiv}[type={titlePage},xmlid={titlepage-008}]

              \teiP[rend={title-main},xmlid={p-281}]{L’échange épistolaire, lieu du secret politique dans \teiHi[rend={italic}]{La Fortune des Rougon}}
              \teiP[rend={author-aut},xmlid={p-282}]{Étienne \teiHi[rend={small-caps}]{Poirier}}
              \teiP[rend={authority\_affiliation},xmlid={p-283}]{Université McGill, Montréal, Canada}
            
\end{teiDiv}
          
\end{teiElement}
          \begin{teiElement}[name={body},xmlid={body-008}]

            \teiP[xmlid={p1-8}]{C’est dans un climat d’incertitude politique que s’ouvre le roman inaugural du cycle des \teiHi[rend={italic}]{Rougon-Macquart} d’Émile Zola, \teiHi[rend={italic}]{La Fortune des Rougon}. Depuis la sous-préfecture fictive de Plassans, dans le Midi, les personnages reçoivent les échos dispersés des événements qui agitent la capitale à la suite du coup d’État du 2 décembre 1851 et cherchent à se ranger dans le clan qui sortira victorieux. Tandis que le jeune couple de Silvère et Miette choisit de rejoindre les républicains qui préparent la résistance à l’Empire, celui formé par Pierre et Félicité Rougon se rapproche plutôt des élites municipales de Plassans afin de tirer profit de l’instabilité pour s’emparer du pouvoir de la ville : « la révolution de 1848 trouva donc tous les Rougon sur le qui-vive, exaspérés par leur mauvaise chance et disposés à violer la fortune, s’ils la rencontraient jamais au détour d’un sentier. C’était une famille de bandits à l’affût, prêts à détrousser les événements\teiNote[xmlid={ftn1-8},n={1},place={foot}]{\teiP[xmlid={p-284}]{Zola, \teiHi[rend={italic}]{La Fortune des Rougon}, dans \teiHi[rend={italic}]{Les Rougon-Macquart}, vol. I, Armand Lanoux (dir.), Paris, Gallimard, « Bibliothèque de la Pléiade », 1960, p. 72. Désormais, les renvois à cette œuvre seront suivis du numéro de page, dans le corps du texte.}}. »}
            \teiP[xmlid={p2-8}]{Il est cependant bien difficile de « détrousser les événements » qui se déroulent à Paris depuis Plassans : à la distance géographique qui isole la sous-préfecture de la métropole se superpose une distance temporelle, les nouvelles des événements n’arrivant que plusieurs heures, voire plusieurs jours après leur issue. Cette double marginalisation provoque une dépendance vis-à-vis de la capitale, dont il s’agit de reproduire les développements en province, comme le remarque Pierluigi Pellini :}
            \begin{teiCit}[xmlid={cit01-8}]

              \begin{teiQuote}
Le travail des personnages (leur héroïsme éventuel) ne consiste donc pas à intervenir activement dans une mêlée qui “ferait l’histoire”, mais plutôt à deviner l’issue des événements parisiens, pour se rallier en temps utile au temps des vainqueurs. L’Histoire est déjà “faite” ailleurs : aux personnages du roman de Zola n’est réservé que le rôle passif de la répéter\teiNote[xmlid={ftn18},n={2},place={foot}]{\teiP[xmlid={p-285}]{Pierluigi Pellini, « “Si je triche un peu” : Zola et le roman historique », \teiHi[rend={italic}]{Les Cahiers naturalistes}, n\teiHi[rend={sup}]{o} 75, 2001, p. 20-21.}}.
\end{teiQuote}
            
\end{teiCit}
            \teiP[xmlid={p3-8}]{Exclus de la grande histoire qui se joue à Paris, les habitants de la province peuvent au mieux tenter de reproduire les mêmes actions, ainsi que l’a fréquemment mentionné la critique, qui décrit la prise de pouvoir de Plassans comme un succédané du coup d’État et Pierre Rougon comme un double grotesque de Napoléon III\teiNote[xmlid={ftn19},n={3},place={foot}]{\teiP[xmlid={p-286}]{Voir notamment Éléonore Reverzy, \teiHi[rend={italic}]{La Chair de l’idée. Politique de l’allégorie dans les }Rougon-Macquart, Genève, Droz, 2007, chap. II, et Corinne Saminadayar\teiHi[rend={small-caps}]{-}Perrin, « \teiHi[rend={italic}]{Storytelling }: Fictions de l’histoire dans \teiHi[rend={italic}]{La Fortune des Rougon} », Émilie Piton\teiHi[rend={small-caps}]{-}Foucault et Henri Mitterand (dir.), \teiHi[rend={italic}]{Lectures de Zola}, La Fortune des Rougon, Rennes, PUR, 2015, p. 21.}}.}
            \teiP[xmlid={p4-8}]{Un tel phénomène de répétition est toutefois impossible à réaliser sans connaître le déroulement des troubles politiques dans la capitale. À ce titre, Pierre et Félicité Rougon bénéficient d’une source d’information privilégiée en la personne de leur fils Eugène, qui écrit régulièrement à son père dans le but de l’aider à prendre la tête de l’élite régionale. La correspondance confidentielle entre le père et le fils devient dès lors un « pivot » du roman\teiNote[xmlid={ftn20},n={4},place={foot}]{\teiP[xmlid={p-287}]{Béatrice Laville, « Le secret dans \teiHi[rend={italic}]{La Fortune des Rougon} », Béatrice Laville et Florence Pellegrini (dir.), La Fortune des Rougon \teiHi[rend={italic}]{d’Émile Zola : lectures croisées}, Bordeaux, Presses universitaires de Bordeaux, 2015, p. 189.}}, dans la mesure où l’échange épistolaire structure les rapports entre la métropole et la province, comme l’explique Olivier Lumbroso :}
            \begin{teiCit}[xmlid={cit02-7}]

              \begin{teiQuote}
Dans l’économie du roman, la lettre ne fait pas seulement parler, agir et réagir les personnages. Elle est aussi l’élément du réseau postal qui couvre le territoire national et participe du maillage spatio-temporel du roman, sinon du cycle. Au sein de cette chronotopie épistolaire, on pourra évoquer en premier lieu la structuration entre Paris et la province au moyen des épîtres échangées entre des personnages “stationnaires” et des personnages “mobiles\teiNote[xmlid={ftn21},n={5},place={foot}]{\teiP[xmlid={p-288}]{Olivier Lumbroso, « Entre réalisme et intimisme : la lettre dans \teiHi[rend={italic}]{Les Rougon-Macquart} de Zola », Colette Becker, Jean-Louis Cabanès et Jean-Marc Hovasse (dir.), \teiHi[rend={italic}]{Écrire l’intime au temps du réalisme et du naturalisme. Mélanges offerts à Pierre-Jean Dufief}, Paris, Honoré Champion, « Romantisme et modernités », 2020, p. 52.}}”.
\end{teiQuote}
            
\end{teiCit}
            \teiP[xmlid={p5-8}]{Si l’on convient assez intuitivement que la correspondance permet de rompre l’écart qui sépare les deux villes, il est aussi possible d’y voir se jouer une réorganisation des lieux associés au pouvoir politique. En s’assimilant à un espace secret, l’échange épistolaire exerce un double rôle dans le roman : il permet aux Rougon d’anticiper le déroulement des événements à Plassans et rétablit du même coup la sphère privée en tant que lieu de pouvoir, au détriment de l’espace public, qui perd toute influence politique.}
            \teiP[xmlid={p6-8}]{La distinction opérée dans cet article entre les différents espaces de pouvoir reprend celle qu’effectue Jürgen Habermas dans son étude sur l’essor de la sphère publique au cours du \teiHi[rend={small-caps}]{xviii}\teiHi[rend={sup}]{e} siècle. Le philosophe distingue l’espace public, lié au pouvoir de l’État, du domaine privé, constitué de la sphère publique, qui émerge avec la création de lieux où peut se réunir la société civile pour échanger, et de la sphère privée, associée au modèle de la maison bourgeoise, limitée à l’intimité familiale\teiNote[xmlid={ftn22},n={6},place={foot}]{\teiP[xmlid={p-289}]{Jürgen Habermas, \teiHi[rend={italic}]{L’Espace public. Archéologie de la publicité comme dimension constitutive de la société bourgeoise}, Marc B. de Launay (trad.), Paris, Payot, 1978, p. 41.}}. Le modèle convoqué par Habermas permet d’illustrer dans quelle mesure le roman présente les événements de décembre 1851 comme le déclencheur d’une redistribution du pouvoir entre les espaces liés à l’État, à la société civile et à la sphère familiale sous le Second Empire. Il s’agira de montrer dans un premier temps comment la marginalisation politique et temporelle de Plassans contribue à entretenir une hiérarchisation des lieux du pouvoir politique dépassée sous un régime démocratique, puis de caractériser l’espace privé que symbolise l’échange épistolaire qui se crée en réaction à cette organisation. C’est de cette manière qu’il sera possible d’évaluer dans quelle mesure la correspondance secrète accomplit une double mission : elle rompt la distance qui oppose Paris et Plassans et, en conséquence, déplace le centre du pouvoir politique vers l’espace privé de la lettre.}
            \begin{teiDiv}[type={section1},xmlid={div01-7}]

              \teiHead[xmlid={plassans-en-retard-sur-le-monde-001}]{Plassans, en retard sur le monde}
              \teiP[xmlid={p7-8}]{C’est en mettant l’accent sur son caractère reculé que s’effectue la description de Plassans en ouverture du deuxième chapitre du roman. La ville, « adossée au nord contre les collines de Garrigues, une des dernières ramifications des Alpes » est plus qu’éloignée, elle semble être à l’extrémité du pays : « comme située au fond d’un cul-de-sac » (p. 36). Cette distance géographique explique pourquoi la sous-préfecture n’a perdu aucune de ses anciennes habitudes datant d’avant la Révolution : « il y a une vingtaine d’années, grâce sans doute au manque de communications, aucune ville n’avait mieux conservé le caractère dévot et aristocratique des anciennes cités provençales » (p. 38). Même lorsqu’elle est comparée aux villes environnantes, Plassans se démarque par son refus de modifier ses pratiques démodées : « tout l’esprit de la ville, fait de poltronnerie, d’égoïsme, de routine, de la haine du dehors et du désir religieux d’une vie cloîtrée, se trouvait dans ces tours de clef donnés aux portes chaque soir \lbrack{}…\rbrack{}. Il n’y a pas de cité, je crois, qui se soit entêtée si tard à s’enfermer comme une nonne » (p. 38). L’ensemble de la description est orienté pour faire valoir le caractère arriéré de la ville, ainsi que le résume Pierre Glaudes : « l’“esprit de la ville” \lbrack{}p. 75\rbrack{}, comme sa situation géographique \lbrack{}…\rbrack{}, la prédisposent à rester à l’écart du mouvement général de son temps\teiNote[xmlid={ftn23-2},n={7},place={foot}]{\teiP[xmlid={p-290}]{Pierre Glaudes, « Le naturel et le social dans \teiHi[rend={italic}]{La Fortune des Rougon} », Pierre Glaudes et Alain Pagès (dir.), \teiHi[rend={italic}]{Relire }La Fortune des Rougon : \teiHi[rend={italic}]{Hommage à David Baguley}, Paris, Classiques Garnier, « Rencontres », 2015, p. 242.}}. »}
              \teiP[xmlid={p8-8}]{La remarque est confirmée dès qu’il est fait mention de la politique. Tandis qu’à Paris s’affrontent républicains et bonapartistes, la domination nobiliaire et ecclésiastique de la ville ramène les débats à ceux du siècle précédent. De même, la vie politique demeure dominée par les secrets et les prises de décisions d’une sphère publique restreinte, limitée aux classes dirigeantes, à la différence de la situation dans la métropole :}
              \begin{teiCit}[xmlid={cit03-7}]

                \begin{teiQuote}
Tout se passe entre le clergé, la noblesse et la bourgeoisie. Les prêtres, très nombreux, donnent le ton à la politique de l’endroit ; ce sont des mines souterraines, des coups dans l’ombre, une tactique savante et peureuse qui permet à peine de faire un pas en avant ou en arrière tous les dix ans. Ces luttes secrètes d’hommes qui veulent avant tout éviter le bruit, demandent une finesse particulière, une aptitude aux petites choses, une patience de gens privés de passions. Et c’est ainsi que les lenteurs provinciales, dont on se moque volontiers à Paris, sont pleines de traîtrises, d’égorgillements sournois, de défaites et de victoires cachées. Ces bonshommes, surtout quand leurs intérêts sont en jeu, tuent à domicile, à coup de chiquenaudes, comme nous tuons à coups de canon, en place publique (p. 73).
\end{teiQuote}
              
\end{teiCit}
              \teiP[xmlid={p9-8}]{Imperméables au mouvement démocratique poussé par la Deuxième République, les élites de la sous-préfecture se réunissent dans le salon jaune des Rougon pour débattre des affaires publiques. La vie politique s’organise ainsi de façon anachronique, les changements causés par l’élan républicain n’ayant pas affecté les mœurs héritées du siècle précédent. La description, par Habermas, des salons qui « faisaient se rencontrer, pour ainsi dire sur le même pied, la noblesse et la grande bourgeoisie des banquiers et des fonctionnaires qui là s’étaient intégrés à l’“intelligentsia\teiNote[xmlid={ftn24-2},n={8},place={foot}]{\teiP[xmlid={p-291}]{Jürgen Habermas, \teiHi[rend={italic}]{L’Espace public}, \teiHi[rend={italic}]{op. cit.}, p. 44.}} ”» correspond tout à fait au type de rassemblement qui se produit chez les Rougon. Sans s’attarder trop longuement sur ce lieu de pouvoir municipal dont plusieurs critiques ont déjà analysé la portée symbolique\teiNote[xmlid={ftn25-2},n={9},place={foot}]{\teiP[xmlid={p-292}]{Daniele Frescaroli, « La description, lieu de mémoire fictionnel. L’espace architectural et l’histoire-mémoire dans \teiHi[rend={italic}]{Les Rougon-Macquart} de Zola », Patrizia Oppici et Susi Pietri (dir.), \teiHi[rend={italic}]{L’Architecture du texte, l’architecture dans le texte}, Macerata, EUM, 2018, p. 75-91.}} et mythique\teiNote[xmlid={ftn26-2},n={10},place={foot}]{\teiP[xmlid={p-293}]{Jacques Perrin, « Le Salon jaune des Rougon ou les métamorphoses d’un mythe », \teiHi[rend={italic}]{Zola, }La Fortune des Rougon \teiHi[rend={italic}]{: figures du pouvoir}, Paris, Ellipses, « Analyses et réflexions », 1994, p. 44-47.}}, il faut souligner que l’endroit est présenté comme étant tout aussi dépassé que la ville. À l’intérieur de la pièce, l’ameublement est démodé et Félicité doit faire preuve de créativité pour « donner un nouveau lustre à toutes ces ruines » (p. 69), ce qui donne au salon une allure démodée : « aux murs étaient pendues six lithographies représentant les grandes batailles de Napoléon. Cet ameublement datait des premières années de l’Empire » (p. 70). La distance géographique et temporelle qui touche la ville est ainsi reproduite dans la description du salon jaune.}
              \teiP[xmlid={p10-8}]{Cette « lenteur provinciale » n’est pas seulement un effet du décalage avec la métropole ; elle témoigne aussi de la manière dont se déroulent les joutes politiques au sein de la ville. À l’image du lieu où elles se mettent en œuvre, les stratégies politiques relèvent de vieilles façons de faire, où les jeux de coulisses et les tractations privées prolifèrent, à l’opposé de la situation parisienne : « que des rois se volent un trône ou que des républiques se fondent, la ville s’agite à peine. On dort à Plassans, quand on se bat à Paris » (p. 73). On ne s’étonne donc pas que le coup d’État désarçonne les élites méridionales, majoritairement royalistes : « les plus fins politiques de Plassans, ceux qui dirigeaient le mouvement réactionnaire, ne flairèrent l’Empire que fort tard. \lbrack{}…\rbrack{} Mais on ne leur laissa pas le temps de prendre un parti ; le coup d’État éclata sur leurs têtes, et ils durent applaudir » (p. 75). La marginalisation géographique et temporelle ne fait pas que maintenir les habitants de Plassans à l’écart de la grande histoire, elle affecte toute l’organisation politique de la ville, qui a conservé ses pratiques d’Ancien Régime.}
            
\end{teiDiv}
            \begin{teiDiv}[type={section1},xmlid={div02-7}]

              \teiHead[xmlid={la-restitution-dune-sphere-privee-001}]{La restitution d’une sphère privée}
              \teiP[xmlid={p11-8}]{Le recours à une correspondance secrète se présente ainsi comme une conséquence de la division anachronique des lieux du pouvoir politique à Plassans. En effet, tandis que les élites politiques continuent d’occuper le salon des Rougon, la fonction de la maison bourgeoise, qu’Habermas associe à la construction d’une sphère privée, se voit transformée. Si le philosophe défend qu’au cours du \teiHi[rend={small-caps}]{xviii}\teiHi[rend={sup}]{e} siècle « le salon s’\lbrack{}est\rbrack{} détaché \lbrack{}de la sphère publique\rbrack{} pour devenir le lieu où se rencontrent les pères de famille bourgeois et leurs épouses\teiNote[xmlid={ftn27-2},n={11},place={foot}]{\teiP[xmlid={p-294}]{Jürgen Habermas, \teiHi[rend={italic}]{L’Espace public}, \teiHi[rend={italic}]{op. cit}., p. 56.}} », force est de constater qu’en raison du décalage qu’accuse Plassans et des nombreux rassemblements qu’elle accueille, la maison des Rougon demeure plutôt associée à la sphère publique. Il s’agit en outre d’un espace où l’influence politique est solidaire des renseignements que l’on détient, mais où la plupart des propos rapportés ne sont que des rumeurs ou des commérages\teiNote[xmlid={ftn28-2},n={12},place={foot}]{\teiP[xmlid={p-295}]{Sur ce point, voir François-Marie Mourad, « Médisances et commérages : les usages de la rumeur dans \teiHi[rend={italic}]{La Fortune des Rougon} », Émilie Piton-Foucault et Henri Mitterand (dir.), \teiHi[rend={italic}]{Lectures de Zola}, \teiHi[rend={italic}]{op. cit.}, p. 87-98.}}. C’est pourquoi le marquis de Carnavant, noble désargenté, parvient à maintenir une forme de supériorité au sein du groupe qui s’y rassemble, en semblant posséder une information privilégiée :}
              \begin{teiCit}[xmlid={cit04-5}]

                \begin{teiQuote}
Il commandait au nom de personnages inconnus, dont il ne livrait jamais les noms. “Ils veulent ceci, ils ne veulent pas de cela” disait-il. Ces dieux cachés, veillant aux destinées de Plassans, au fond de leur nuage, sans paraître directement se mêler aux affaires publiques, devaient être certains prêtres, les grands politiques du pays (p. 79).
\end{teiQuote}
              
\end{teiCit}
              \teiP[xmlid={p12-8}]{Bien que ce type de renseignement puisse causer un grand émoi dans la sphère publique, sa valeur politique est moindre que l’information qui circule dans une sphère plus restreinte, comme l’explique Daniel Lequette :}
              \begin{teiCit}[xmlid={cit05-5}]

                \begin{teiQuote}
L’échelon \lbrack{}du pouvoir\rbrack{} le plus bas est celui où l’information arrive \teiHi[rend={italic}]{sans source ni destinataire ni finalité identifiables }sous la forme massive et insinuante \teiHi[rend={italic}]{de la rumeur, de l’écho, du bruit }: le contenu proprement informatif est presque nul, inversement proportionnel à sa puissance émotionnelle. L’échelon le plus haut se distingue par les caractères inverses : les renseignements sont hiérarchisés, classés, authentifiés et engagent nommément les partenaires des échanges ; c’est l’échelon de la détention de \teiHi[rend={italic}]{secrets} et de la distillation autorisée de \teiHi[rend={italic}]{mensonges}. Eugène est l’exemple du négociant d’information qui par l’effet des circonstances mais aussi de son habileté se trouve à l’intersection des trois axes et peut ainsi faire circuler des renseignements et des rétributions en charges officielles et en honneurs. Il permet à ses parents de dépasser le cercle d’action qui leur était jusqu’alors imparti du fait de leur capital financier\teiNote[xmlid={ftn29-2},n={13},place={foot}]{\teiP[xmlid={p-296}]{Daniel Lequette, « Pouvoir et information dans \teiHi[rend={italic}]{La Fortune des Rougon} », \teiHi[rend={italic}]{Zola, }La Fortune des Rougon \teiHi[rend={italic}]{: Figures du pouvoir}, \teiHi[rend={italic}]{op. cit}., p. 103.}}.
\end{teiQuote}
              
\end{teiCit}
              \teiP[xmlid={p13-8}]{S’il est associé au pouvoir, le salon n’est donc pas associé à l’échelon le plus élevé en termes d’influence politique. Certes, le fait de tenir salon rapproche les Rougon du pouvoir, mais il les prive aussi d’un espace où l’information peut s’échanger librement. En témoigne la d’Eugène en 1849, d’Eugène, qui s’abstient de s’engager dans les discussions du groupe : « sa façon recueillie d’écouter, sa complaisance inaltérable avaient concilié toutes les sympathies. On le jugeait nul, mais bon enfant » (p. 81).}
              \teiP[xmlid={p14-7}]{Dès lors, l’échange épistolaire crée un espace abrité où peuvent s’échanger des informations ayant une plus grande valeur politique. Il n’est d’ailleurs pas anodin que la discussion qui décide de la stratégie épistolaire adoptée se tienne dans un lieu intime, à l’écart du salon jaune, dans la chambre du couple Rougon : « Félicité, restée dans la chambre, essaya vainement d’écouter. Les deux hommes parlaient bas, comme s’ils eussent redouté qu’une seule de leurs paroles pût être entendue de dehors » (p. 84). La réception de la correspondance se déroule au sein de l’espace privé de Pierre, de sa chambre à coucher, où il s’enferme pour lire les lettres d’Eugène, jusqu’au vieux secrétaire où il les garde précieusement. Tous les personnages qui en partagent le secret appartiennent aussi à la sphère intime, puisque, si la correspondance ne touche d’abord que le père et son fils, c’est la famille entière qui doit bénéficier des informations transmises\teiNote[xmlid={ftn30-2},n={14},place={foot}]{\teiP[xmlid={p-297}]{Béatrice Laville, « Le secret dans \teiHi[rend={italic}]{La Fortune des Rougon} », Béatrice Laville et Florence Pellegrini (dir.), La Fortune des Rougon \teiHi[rend={italic}]{d’Émile Zola : lectures croisées},\teiHi[rend={italic}]{ op. cit.}, p. 198.}}. Avant même que la première lettre n’ait été écrite et envoyée, c’est en tant que lieu privé que se conçoit l’échange épistolaire, comme l’annonce Eugène à sa mère, qu’il tient temporairement à l’écart : « d’ailleurs mon père vous instruira, quand l’heure sera venue » (p. 85). De même, au fil du roman, ceux qui apprendront l’existence de cet espace virtuel privé ou qui y auront accès par la lecture des lettres appartiennent à la cellule familiale : Félicité, d’abord, qui vole la clé du secrétaire à son mari endormi pour lire les missives de son fils, puis Aristide, qui apprend le secret de la correspondance par inadvertance, lui-même caché dans un endroit privé de la maison de ses parents « dans un trou noir que formait un petit escalier menant aux combles » (p. 105), lieu topique de découverte de confidences, et enfin le marquis de Carnavant, que l’on soupçonne être le père illégitime de Félicité. Seul Vuillet fait entorse à la règle en violant la confidentialité que son statut de maître des postes lui impose, mais il reconnaît lui-même que c’est en raison du privilège lié à son rôle que cette intrusion est possible : « les mille secrets de Plassans étaient là ; il touchait à l’honneur des femmes, à la fortune des hommes, et il n’avait qu’à briser les cachets, pour en savoir aussi long que le grand vicaire de la cathédrale, le confident des personnes comme il faut de la ville » (p. 261).}
              \teiP[xmlid={p15-6}]{Cette configuration privée de l’échange épistolaire se confirme par la façon dont réagissent les personnages qui accèdent au secret et qui comprennent rapidement qu’il est dans l’intérêt familial de le préserver. Ainsi, bien qu’elle soit en colère contre son mari qui refuse de lui partager l’information qu’Eugène envoie, Félicité se résout tout de même à garder le secret : « jamais complice ne fit moins de bruit et plus de besogne » (p. 95). Aristide, tout en se sentant dupé par sa famille, convient aussi qu’« il fallait mieux attendre et se taire » (p. 106), alors que le marquis, pourtant opposé au retour des Bonaparte, fait le pari de soutenir la cause de Félicité lorsqu’il sait la sienne perdue et garde le silence auprès des autres membres du salon jaune. Même le lecteur est tenu à l’écart de cet espace : jamais le contenu des lettres n’est cité dans le roman, et leur propos est résumé rapidement. C’est lorsque Félicité lit les lettres qu’on a une idée de leur teneur, sans en pouvoir tirer aucun détail : « chacune de ses lettres constatait les progrès de la cause et faisait prévoir un dénouement prochain. Elles se terminaient généralement par l’exposé de la ligne de conduite que Pierre devait tenir à Plassans » (p. 94). Assimilée à un espace privé, la correspondance permet d’assurer un plus grand pouvoir politique à la famille Rougon, à l’abri du « public » qui fourmille dans le salon jaune.}
            
\end{teiDiv}
            \begin{teiDiv}[type={section1},xmlid={div03-6}]

              \teiHead[xmlid={une-prise-de-pouvoir-anticipee-001}]{Une prise de pouvoir anticipée}
              \teiP[xmlid={p16-6}]{Le premier rôle de ces lettres est donné d’emblée : il s’agit de s’assurer que les Rougon sauront agir conformément au pouvoir impérial, lorsque celui-ci prendra Paris, afin de sécuriser leur conquête de Plassans. Or pour y arriver, le clan familial doit combler le décalage qui touche la petite ville et empêche la majorité de ses habitants de prévoir le dénouement du conflit. Ceux qui participent à l’insurrection sont incapables de voir que leur résistance à l’Empire est vaine : « ces hommes, qui marchaient dans l’aveuglement de la fièvre que les événements de Paris avaient mis au cœur des républicains, s’extasiaient au spectacle de cette longue bande de terre toute secouée de révolte » (p. 163). De même, les réguliers du salon jaune souffrent du manque d’accès aux informations provenant de la capitale, ce qui explique l’aveuglement du groupe face au déroulement des événements :}
              \begin{teiCit}[xmlid={cit06-4}]

                \begin{teiQuote}
Les nouvelles les plus contradictoires arrivaient de Paris ; tantôt les républicains l’emportaient, tantôt le parti conservateur écrasait la République. L’écho des querelles qui déchiraient l’Assemblée législative parvenait au fond de la province, grossi un jour, affaibli le lendemain, changé au point que les plus clairvoyants marchaient en pleine nuit. Le seul sentiment général était qu’un dénouement approchait (p. 97).
\end{teiQuote}
              
\end{teiCit}
              \teiP[xmlid={p17-6}]{Même si les plus perspicaces pressentent qu’un bouleversement politique est imminent, ils peinent à deviner lequel. Félicité, « le nez le plus fin de la famille » (p. 76), et le marquis qui « flaira l’Empire » (p. 92) doivent consulter les lettres d’Eugène pour corroborer ce qui, autrement, ne serait que des conjectures : « M. de Carnavant ne s’était pas trompé et ses propres soupçons le confirmaient. Il y avait là une quarantaine de lettres, dans lesquelles \lbrack{}Félicité\rbrack{} put suivre le grand mouvement bonapartiste qui devait aboutir à l’Empire » (p. 94).}
              \teiP[xmlid={p18-5}]{L’espace privé de la correspondance abolit ainsi l’écart entre Paris et Plassans. Si l’amenuisement de la distance géographique est évident par le trajet qu’effectue la lettre, c’est le raccourcissement de l’écart temporel qui est le plus marqué et le plus important. L’échange épistolaire est teinté de la certitude d’Eugène que le changement de régime annoncé est sur le point de se produire, ce qui pousse les Rougon à se presser afin d’être prêts le moment venu : « d’un jour à l’autre, une crise décisive pouvait avoir lieu » (p. 98). Alors que la tendance au sein du salon jaune est plutôt en faveur d’un retour à la monarchie, témoignant une fois de plus de son décalage politique, les Rougon vont le convaincre de changer d’allégeance politique : « Eugène avait la foi. Il parlait à son père du prince Louis Bonaparte comme de l’homme nécessaire et fatal qui seul pouvait dénouer la situation. Il avait cru en lui avant même son retour en France, lorsque le bonapartisme était traité de chimère ridicule » (p. 94). Au moment du coup d’État, les Rougon sont déjà parvenus à établir leur prééminence politique dans la ville : « selon ses vœux, en novembre 1851, le salon jaune était maître de Plassans » (p. 90). L’efficacité de la correspondance est telle que Pierre est mis au courant du coup d’État dès la veille, le 1\teiHi[rend={sup}]{er} décembre. Sophie Guermès résume la situation succinctement : « les Rougon ont une supériorité sur les autres, puisqu’ils savent quelle direction veut suivre le prince-président\teiNote[xmlid={ftn31-2},n={15},place={foot}]{\teiP[xmlid={p-298}]{Sophie Guermès, « Le Hasard et la Nécessité dans \teiHi[rend={italic}]{La Fortune des Rougon} », Émilie Piton\teiHi[rend={small-caps}]{-}Foucault et Henri Mitterand (dir.), \teiHi[rend={italic}]{Lectures de Zola. La Fortune des Rougon}, \teiHi[rend={italic}]{op. cit.}, p. 201.}}. »}
              \teiP[xmlid={p19-5}]{Tandis que, face à la métropole, la correspondance permet aux Rougon de rattraper leur retard sur les événements, à Plassans, elle leur permet plutôt d’anticiper les suites de l’insurrection. Elle leur garantit ainsi une prise de pouvoir sur le temps, posé tout au long du roman comme condition d’accès au pouvoir. L’exemple le plus éloquent de la capacité que procure l’échange épistolaire à anticiper les conséquences du soulèvement pour Plassans se produit dans les jours qui suivent le coup d’État. L’interruption de la correspondance à ce moment désoriente complètement Pierre, incapable de décider de la marche à suivre sans l’information qui lui vient de son fils : « en effet, pourquoi Eugène n’écrivait-il pas à son père ? Après l’avoir tenu si fidèlement au courant des succès de la cause bonapartiste, il aurait dû s’empresser de lui annoncer le triomphe ou la défaite du prince Louis » (p. 258). À l’opposé, Vuillet, après avoir lu la lettre d’Eugène qu’il retient à la poste où il apprend la victoire du parti de l’Empereur à Paris, n’hésite pas à condamner avec véhémence les insurgés républicains, alors qu’il craignait de les dénoncer auparavant. C’est d’ailleurs en raison de ce changement soudain d’allégeance de Vuillet que Félicité parvient à deviner qu’il est en possession d’une lettre d’Eugène et réussit à la récupérer : seul quelqu’un appartenant au même espace privé et au fait des événements à Paris serait capable de prédire l’issue des soulèvements dans la sous-préfecture. La scène confirme aussi la fonction de l’espace privé en tant que centre du pouvoir politique, puisque c’est uniquement après en avoir pris connaissance que Vuillet s’engage désormais dans la sphère publique, en situant son journal en faveur de l’Empire et contre l’insurrection populaire.}
            
\end{teiDiv}
            \begin{teiDiv}[type={section1},xmlid={div04-2}]

              \teiHead[xmlid={une-reconfiguration-du-pouvoir-politique-001}]{Une reconfiguration du pouvoir politique}
              \teiP[xmlid={p20-5}]{En rompant la distance entre la métropole et la province, la lettre engage simultanément un déplacement du centre de pouvoir à Plassans, qui passe aux mains de Pierre Rougon. Associée à la sphère privée, elle confirme aussi la subordination de l’espace public et les intérêts des membres du salon jaune à ceux de la famille Rougon. En ce sens, la lettre convoque une nouvelle hiérarchisation des lieux de pouvoir. Ce ne sont donc pas les institutions de la Deuxième République, appartenant au domaine public, qui constituent le centre du pouvoir politique. Si Pierre Glaudes reconnaît que l’ensemble de l’œuvre cherche à représenter « une hiérarchie spatiale qui associe \lbrack{}la\rbrack{} réussite \lbrack{}des Rougon\rbrack{} à l’intégration graduelle d’espaces supérieurs à ceux qu’ils occupaient précédemment dans la société\teiNote[xmlid={ftn32-2},n={16},place={foot}]{\teiP[xmlid={p-299}]{Pierre Glaudes, « Le naturel et le social dans \teiHi[rend={italic}]{La Fortune des Rougon} », art. cité, p. 243.}} », il faudrait ajouter que cette gradation s’effectue au prix d’une reconfiguration des espaces détenant le plus de pouvoir. Au fil de la correspondance, le centre du pouvoir se déplace, d’abord vers le salon jaune, puis vers la famille elle-même. De bourgeois auparavant confinés à l’espace public, les Rougon ont réussi à faire de leur espace privé, par les secrets qui y sont échangés, le cœur du pouvoir municipal. Cette réorganisation des forces politiques ne transforme toutefois pas de façon radicale la vie publique de la sous-préfecture, qui a toujours rejeté la démocratie : « jusqu’en 1830, le peuple n’a pas compté. Encore aujourd’hui, on agit comme s’il n’était pas » (p. 73). Au mieux, le pouvoir est déplacé d’un espace privé aristocrate à un espace privé bourgeois, celui de la famille Rougon. Or, à Paris, le changement est plus important et symbolise un retour du contrôle qu’exerce la sphère privée sur les institutions de l’État. On ne discute plus des stratégies politiques sur la place publique, mais à nouveau de façon secrète, dans un cercle limité.}
              \teiP[xmlid={p21-5}]{Ce déplacement des lieux du pouvoir politique à Paris, qui fait écho à celui décrit à Plassans, n’est cependant que suggéré dans le roman, d’abord par la description des lettres qu’envoie Eugène : « Félicité comprit que son fils était depuis 1848 un agent secret très actif. Bien qu’il ne s’expliquât pas nettement sur sa situation à Paris, il était évident qu’il travaillait à l’Empire, sous les ordres des personnages qu’il nommait avec une sorte de familiarité » (p. 94). C’est désormais Paris qui semble reproduire les tractations politiques de Plassans, tenues sous le sceau du secret, hors de l’espace public, avec des personnes qui sont incidemment associées à des membres d’une même famille. Ce retour de l’espace privé comme centre du pouvoir politique est aussi confirmé par le jeu de connivence qui s’installe entre l’auteur et le lecteur, comme le remarque Marie-Ange Fougère\teiNote[xmlid={ftn33-3},n={17},place={foot}]{\teiP[xmlid={p-300}]{Marie-Ange Fougère, « \teiHi[rend={italic}]{La Fortune des Rougon, }roman de la connivence », Émilie Piton\teiHi[rend={small-caps}]{-}Foucault et Henri Mitterand (dir.), \teiHi[rend={italic}]{Lectures de Zola. La Fortune des Rougon}, \teiHi[rend={italic}]{op. cit}., p. 84-85.}}. En effet, nul besoin, pour le lecteur au fait de l’histoire du Second Empire, de connaître le contenu des lettres pour deviner la résolution de l’insurrection républicaine dans le Midi, puisque celle-ci relève de l’espace public. Par contre, le détail des discussions politiques qui mènent au coup d’État demeure quant à lui inconnu, puisqu’il relève de la sphère privée dont l’accès est restreint. Si, en régime démocratique, il est attendu que l’information relative aux prises de décisions politiques soit partagée avec le public, sous un régime dictatorial, la diffusion est beaucoup plus contrôlée par les personnes en charge. L’Empire confirme le passage du pouvoir du domaine public à la sphère privée à l’échelle du pays, ne serait-ce que par le retour d’une famille, les Bonaparte, aux rênes de l’État. Au jeu de reproduction, par les Rougon, des événements qui se déroulent dans la capitale s’ajoute, selon un autre effet miroir, le déplacement du centre du pouvoir politique à Paris. En rapprochant les stratégies politiques d’Ancien Régime de la sous-préfecture de celles qu’utilisent les partisans de Napoléon III, le secret accomplit un double objectif au sein du roman. Il confirme non seulement le rapprochement temporel entre Paris et Plassans, mais illustre surtout la reconfiguration du pouvoir politique que cause le coup d’État. Dans cette mesure, le parallèle effectué entre les deux villes démontre comment le secret, contenu dans l’espace privé, est rétabli en tant qu’élément central du pouvoir politique sous le Second Empire.}
              \teiP[xmlid={p22-4}]{La dernière lettre qu’envoie Eugène à sa mère semble confirmer ce retour, pour la métropole, à un système politique similaire à celui de Plassans. La lettre fait suite à une missive de Félicité, la seule du roman à parcourir le chemin inverse, de Plassans vers Paris. La mère y annonce à son fils la victoire de Pierre sur les insurgés et joint les articles des journaux locaux qui racontent ses exploits. Eugène écrit en retour que son père obtiendra la sous-préfecture tel qu’entendu, mais surtout que sa loyauté envers le nouvel empereur lui vaudra la Légion d’honneur. Or plutôt que de choisir d’attendre que l’annonce en soit faite à Paris, les membres du salon jaune insistent pour tenir une cérémonie séance tenante, le capitaine Sicardot décorant Pierre en sa qualité de « vieux soldat de Napoléon » (p. 315). Avec la diffusion de la lettre aux membres du salon jaune, il semble ainsi que tous aient maintenant rattrapé leur retard sur Paris, devançant même la cérémonie officielle qui s’y tiendra. Ce couronnement de leur nouveau chef par les membres du salon jaune dans un espace appartenant au domaine privé est aussi révélateur du nouveau centre de pouvoir politique à Paris : si l’effet de décalage entre la ville et la province semble s’estomper à la fin du roman, c’est parce que la capitale a retrouvé d’anciennes façons de pratiquer la politique avec le retour de l’Empire, où les décisions politiques se prennent à nouveau en secret, dans la sphère privée.}
            
\end{teiDiv}
          
\end{teiElement}
        
\end{teiElement}
      
\end{teiElement}
      \begin{teiElement}[name={group},type={section1},xmlid={section1-003}]

        \teiHead[xmlid={personnages-incarner-le-secret-politique-001}]{Personnages : incarner le secret politique}
        \begin{teiElement}[name={group},type={chapter},data-page-title={Spectacle et dissimulation},data-page-subtitle={Anamorphose du tyran dans les récits de voyages de Bernier (1671) et Tavernier (1676)},data-page-authors={Mathilde Mougin@@Université de Lyon, LabEx COMOD-IHRIM, UMR 5317, Lyon, France},data-include-href={../XML/chap\_9\_Spectacle\_et\_dissimulation\_Locus\_politicus.xml},xmlid={chapter-008}]

          \begin{teiElement}[name={front}]

            \begin{teiDiv}[type={titlePage},xmlid={titlepage-009}]

              \teiP[rend={title-main},xmlid={p-301}]{Spectacle et dissimulation}
              \teiP[rend={title-sub},xmlid={p-302}]{Anamorphose du tyran dans les récits de voyages de Bernier (1671) et Tavernier (1676)}
              \teiP[rend={author-aut},xmlid={p-303}]{Mathilde \teiHi[rend={small-caps}]{Mougin}}
              \teiP[rend={authority\_affiliation},xmlid={p-304}]{Université de Lyon\teiHi[rend={small-caps}]{, }LabEx\teiHi[rend={small-caps}]{ COMOD-IHRIM, UMR 5317, }Lyon, France}
            
\end{teiDiv}
          
\end{teiElement}
          \begin{teiElement}[name={body},xmlid={body-009}]

            \teiP[xmlid={p1-9}]{Dans la lignée d’une tradition machiavélienne qui considère que le prince doit être à la fois « lion » et « renard\teiNote[xmlid={ftn1-9},n={1},place={foot}]{\teiP[xmlid={p-305}]{« Le lion ne peut se défendre des filets, le renard des loups ; il faut donc être renard pour connaître les filets, et lion pour faire peur aux loups », conseille l’auteur au souverain. Ce sont « ces deux natures \lbrack{}…\rbrack{} qu’il faut qu’un prince sache bien imiter » (Machiavel, \teiHi[rend={italic}]{Le Prince et autres textes : suivi d’extraits des Œuvres politiques et d’un choix des Lettres familières}, Paul Veyne éd., Paris, Gallimard, « Folio classique », 1980, p. 107‑108).}} », l’exercice du pouvoir sous l’Ancien Régime semble généralement refléter deux tendances apparemment contradictoires : celle, ostentatoire, du spectacle de la force, qui donne notamment lieu à des cérémonies royales fastueuses ainsi qu’à la mise en scène de supplices violents, et celle de la dissimulation et de la ruse. L’histoire secrète du \teiHi[rend={small-caps}]{xvii}\teiHi[rend={sup}]{e} siècle, où le pouvoir politique conçu comme un produit des passions humaines se décide dans le secret des alcôves, et les mémoires apocryphes tirent toutes les potentialités narratives, esthétiques et morales d’une telle politique du secret\teiNote[xmlid={ftn103-2},n={2},place={foot}]{\teiP[xmlid={p-306}]{Voir par exemple \teiHi[rend={italic}]{Dom Carlos }(1672) de Saint-Réal qui met en scène les rapports difficiles entre Philippe II et son successeur, ou, pour le genre romanesque, \teiHi[rend={italic}]{Les Mémoires de d’Artagnan} (1700) de Courtilz de Sandras, témoignage apocryphe du mousquetaire dans lequel se dessine le personnage-type du cardinal dissimulateur dont se souviendra Dumas.}}. Si les dictionnaires contemporains identifient le pouvoir au premier modèle, celui d’un « pouvoir absolu et despotique\teiNote[xmlid={ftn104},n={3},place={foot}]{\teiP[xmlid={p-307}]{Article « \teiHi[rend={small-caps}]{Pouvoir}. s. m. », Antoine Furetière, \teiHi[rend={italic}]{Dictionnaire universel contenant généralement tous les mots françois tant vieux que modernes et les termes de toutes les sciences et des arts}, La Haye et Rotterdam, Arnout et Reinier Leers, 1690.}} » ou d’un « Grand pouvoir \lbrack{}…\rbrack{} absolu » et « indépendant\teiNote[xmlid={ftn105},n={4},place={foot}]{\teiP[xmlid={p-308}]{Article « \teiHi[rend={small-caps}]{Pouvoir}. s. m. », \teiHi[rend={italic}]{Le Dictionnaire de l’Académie françoise dedié au Roy}, Paris, Chez la veuve Jean-Baptiste Coignard, 1694, vol. 2 (M-Z), p. 301.}} » – comme le confirme le faste de la politique louis-quatorzienne, qui multiplie les fêtes royales et autres démonstrations de puissance –, ils l’associent également à une part de mystère. Le \teiHi[rend={italic}]{Dictionnaire de l’Académie} opère en effet une distinction entre la face visible du pouvoir, et une face qui serait cachée qu’exprime l’idée de « marque », d’« effet\teiNote[xmlid={ftn106},n={5},place={foot}]{\teiP[xmlid={p-309}]{\teiHi[rend={italic}]{Ibid.}}} » du pouvoir. En d’autres termes, seuls les effets du pouvoir sont perceptibles, l’élaboration de décisions s’opérant dans le secret d’une antichambre ou au sein d’un conseil restreint. Bien que les manifestations du pouvoir soient sensibles, ses arcanes en demeurent cachés au plus grand nombre.}
            \teiP[xmlid={p2-9}]{Si le pouvoir monarchique français fascine, celui des monarques orientaux – à l’époque qualifiés de « despotes\teiNote[xmlid={ftn107},n={6},place={foot}]{\teiP[xmlid={p-310}]{Article « \teiHi[rend={small-caps}]{Despote}. s. m. », « C’étoit une dignité dans la Cour des Empereurs d’Orient \lbrack{}…\rbrack{} Les Princes d’Orient sont absolus et \teiHi[rend={italic}]{despotiques} » dans Antoine Furetière, \teiHi[rend={italic}]{Dictionnaire universel…}, \teiHi[rend={italic}]{op. cit.} ; Article « \teiHi[rend={small-caps}]{Despote}, s. m\teiHi[rend={italic}]{.} », « Prince souverain qui dépend de l’Empire Otoman », dans Pierre Richelet, \teiHi[rend={italic}]{Dictionnaire françois : contenant les mots et les choses, plusieurs nouvelles remarques sur la langue françoise, ses expressions propres, figurées et burlesques, la prononciation des mots les plus difficiles, le genre des noms, le régime des verbes. Avec les termes les plus connus des Arts \& des Sciences. Le tout tiré de l’usage et des bons auteurs de la langue françoise}, Genève, Chez Jan Herman Widerhold, 1680, vol. A-L, p. 236.}} » – éveille un intérêt notable dont témoigne la littérature viatique contemporaine. En effet, le système politique des États musulmans suscite la vive curiosité des voyageurs, comme le montre la riche collection de gravures accompagnant le récit de Nicolay (1567-1568), lesquelles représentent précisément les grades des soldats de l’armée turque (capitaine, archer, soldat, laquais, etc.), à l’origine d’une longue tradition qui perdure au \teiHi[rend={small-caps}]{xviii}\teiHi[rend={sup}]{e} siècle avec le \teiHi[rend={italic}]{Recueil }de Ferriol\teiNote[xmlid={ftn108},n={7},place={foot}]{\teiP[xmlid={p-311}]{Nicolas de Nicolay, \teiHi[rend={italic}]{Les Quatre premiers livres des navigations et peregrinations orientales, de N. de Nicolay, Dauphinoys, seigneur d’Arfeuille, varlet de chambre, \& geographe ordinaire du Roy. Avec les figures au naturel tant d’hommes que de femmes selon la diversité des nations, \& de leur port, maintien \& habitz.}, Lyon, Par Guillaume Rouille, 1567 ; Charles de Ferriol, Le Hay, \teiHi[rend={italic}]{Recueil de cent estampes représentant les diverses nations du Levant, tirées d’après nature en 1707 et 1708 par les ordres de M. de Ferriol, ambassadeur du Roy à la Porte et gravées en 1712 et 1713 par les soins de Le Hay}, Paris, chez Le Hay et chez Duchange, 1714.}}. Le despote oriental – ottoman, persan ou moghol – semble particulièrement incarner cette tension entre la politique du spectacle et celle du secret, et fournir au monarque français une forme de miroir amplifiant. Le Grand Turc, le Grand Moghol ou le roi de Perse partagent avec leur homologue français une puissance absolue dont les démonstrations de force ne sont qu’un des versants.}
            \teiP[xmlid={p3-9}]{Il s’agit ici de s’intéresser aux stratégies et aux enjeux de la représentation du pouvoir oriental – et plus spécifiquement turc, persan et indien – par les voyageurs français, principalement dans trois récits : celui de François Bernier (1620-1688), médecin – et disciple de Gassendi\teiNote[xmlid={ftn109},n={8},place={foot}]{\teiP[xmlid={p-312}]{Il publie notamment l’\teiHi[rend={italic}]{Abrégé de la philosophie de Gassendi} en huit tomes en 1674.}} – qui a résidé à la cour moghole pendant une dizaine d’années et qui a publié en 1671 ses \teiHi[rend={italic}]{Voyages}, et ceux de Jean-Baptiste Tavernier (1605-1689), protestant et marchand de pierres précieuses qui a consacré sa vie à de nombreux voyages en Orient et dans les cours persane, moghole et ottomane, dont il livre une description précise dans les \teiHi[rend={italic}]{Six voyages} (1676) et la \teiHi[rend={italic}]{Nouvelle relation de l’intérieur du serrail }(1675). Dans ces récits, la politique orientale est caractérisée par sa dimension spectaculaire, de la mise en scène de supplices à l’architecture même du sérail avec la salle de divan. Toutefois, le sérail est également un lieu de dissimulation et de secrets, dont l’entrée est contrôlée, alimentant les fantasmes des voyageurs et à leur suite des lecteurs. Les auteurs se plaisent ainsi à raconter des histoires secrètes de sérail impliquant les personnages de la famille royale, lesquels fournissent par ailleurs le scénario de nombreuses « turqueries » dramatiques\teiNote[xmlid={ftn110},n={9},place={foot}]{\teiP[xmlid={p-313}]{Pour une étude spécifique de ces tragi-comédies qui mettent en scène des épisodes du règne de Soliman le Magnifique, voir Sylvie Requemora, « Les “turqueries” : une vogue théâtrale en mode mineur », \teiHi[rend={italic}]{Littératures classiques}, n\teiHi[rend={sup}]{o} 51 : « Le théâtre au \teiHi[rend={small-caps}]{xvii}\teiHi[rend={sup}]{e }siècle : pratiques du mineur », 2004, p. 131-151 ; Sylvie Requemora, « Scènes de sérail : la construction d’un Orient théâtral et romanesque au \teiHi[rend={small-caps}]{xvii}\teiHi[rend={sup}]{e} siècle », Anne Duprat et Émilie Picherot (dir.), \teiHi[rend={italic}]{Récits d’Orient dans les littératures d’Europe : }\teiHi[rend={ small-caps italic}]{xvi}\teiHi[rend={italic sup}]{e}\teiHi[rend={italic}]{-}\teiHi[rend={ small-caps italic}]{xvii}\teiHi[rend={italic sup}]{e}\teiHi[rend={italic}]{ siècles}, Paris, PUPS, « Recherches actuelles en littérature comparée », 2008, p. 249-262.}} qui marquent les prémices de l’orientalisme\teiNote[xmlid={ftn111},n={10},place={foot}]{\teiP[xmlid={p-314}]{Sur l’orientalisme au \teiHi[rend={small-caps}]{xvii}\teiHi[rend={sup}]{e} siècle, voir Nicholas Dew, \teiHi[rend={italic}]{Orientalism in Louis XIV’s France}, Oxford et New York, Oxford University Press, « Oxford historical monographs », 2009.}}. Toutefois, loin de n’être que la matière de divertissements, cette représentation du pouvoir oriental constitue aussi un miroir en même temps qu’un avertissement adressés au monarque français pour le préserver de la tyrannie\teiNote[xmlid={ftn112},n={11},place={foot}]{\teiP[xmlid={p-315}]{Aristote explique en effet que la tyrannie est une « déviation » du régime monarchique : tandis que la royauté « a en vue l’avantage commun », « la tyrannie est une monarchie qui vise l’avantage du monarque » (\teiHi[rend={italic}]{Les Politiques}, Pierre Pellegrin trad., Paris, Flammarion, « GF », 2015, livre III, chap. 7, 1279 a, p. 229-230).}}.}
            \begin{teiDiv}[type={section1},xmlid={div01-8}]

              \teiHead[xmlid={le-synoptisme-du-pouvoir-001}]{Le « synoptisme » du pouvoir}
              \teiP[xmlid={p4-9}]{Georges Balandier rappelle à quel point le modèle politique de l’Ancien Régime repose sur un système de représentation, évoquant notamment la pratique des entrées royales de la Renaissance et des fêtes royales louis-quatorziennes\teiNote[xmlid={ftn113},n={12},place={foot}]{\teiP[xmlid={p-316}]{« \lbrack{}L\rbrack{}es techniques dramatiques ne sont pas seulement utilisées au théâtre, mais aussi à la direction de la Cité. Le Prince doit se comporter en acteur politique afin de conquérir et conserver le pouvoir » (Georges Balandier, \teiHi[rend={italic}]{Le Pouvoir sur scènes}, Paris, Éditions Balland, 1980, p. 13, voir également p. 45 et p. 15).}} dans lesquelles le souverain devient lui-même acteur. Gérard Leclerc qualifie de « synoptique » ces différentes modalités du spectacle du pouvoir de cette époque\teiNote[xmlid={ftn114},n={13},place={foot}]{\teiP[xmlid={p-317}]{Il distingue notamment les « spectacles à usage interne, où le pouvoir se présente et se représente lui-même devant un public restreint » (celui de la cour notamment), et les « spectacles à usage externe, destinés aux foules, au peuple, où le pouvoir est présent métaphoriquement », comme les supplices, sans oublier, bien sûr, les espaces de représentation que constituent les palais, « théâtre de cette mise en scène du pouvoir » (Gérard Leclerc, \teiHi[rend={italic}]{Le Regard et le pouvoir}, Paris, PUF, « Sociologie d’aujourd’hui », 2006, p. 115‑116).}}, empruntant à Mathiesen une notion forgée pour décrire le système des médias de la société contemporaine\teiNote[xmlid={ftn115},n={14},place={foot}]{\teiP[xmlid={p-318}]{« \lbrack{}O\rbrack{}n constate, à travers le mass media, le développement d’un \teiHi[rend={italic}]{synoptisme} \lbrack{}…\rbrack{} », c’est-à-dire d’un dispositif « où un grand nombre peut voir et suivre du regard une petite minorité » (Thomas Mathiesen, \teiHi[rend={italic}]{On globalization of control. Towards an integrated surveillance system in Europe}, 1999, cité et traduit par Gérard Leclerc, \teiHi[rend={italic}]{Le regard et le Pouvoir},\teiHi[rend={italic}]{ op. cit.}, p. 26).}}. Les auteurs sont par exemple attentifs au défilé militaire, que Balandier compte parmi les « expressions cérémonielles du dogme et de la pédagogie des gouvernants\teiNote[xmlid={ftn116},n={15},place={foot}]{\teiP[xmlid={p-319}]{Georges Balandier, \teiHi[rend={italic}]{Le Pouvoir sur scènes}, \teiHi[rend={italic}]{op. cit.}, p.\teiHi[rend={italic}]{ }20.}} ». Tavernier raconte que Cha-Abbas II, roi de Perse, fit faire en 1654 une revue générale de sa cavalerie qui dura douze jours\teiNote[xmlid={ftn117},n={16},place={foot}]{\teiP[xmlid={p-320}]{Jean-Baptiste Tavernier, \teiHi[rend={italic}]{Les Six Voyages I}, Paris, Chez Gervais Clouzier et Claude Barbin, 1676, p.\teiHi[rend={italic}]{ }584.}}, et dans sa lettre à Colbert, Bernier détaille sur plus d’une dizaine de pages la composition de l’importante milice entretenue par le Grand Moghol\teiNote[xmlid={ftn118},n={17},place={foot}]{\teiP[xmlid={p-321}]{François Bernier, \teiHi[rend={italic}]{Un libertin dans l’Inde moghole}, Paris, Chandeigne, «\teiHi[rend={italic}]{ }Magellane\teiHi[rend={italic}]{ }», 2008, p.\teiHi[rend={italic}]{ }205‑217.}}.}
              \teiP[xmlid={p5-9}]{Cette mise en scène du pouvoir apparaît également à l’endroit des tortures pratiquées, par exemple, en Perse. Tavernier dresse la longue liste des supplices administrés aux voleurs, selon une logique d’horreur croissante :}
              \begin{teiCit}[xmlid={cit01-9}]

                \begin{teiQuote}
Le supplice le plus cruel de tous est de monter le patient sur un cheval avec un bâton par derriere qui luy tient les bras ouverts. Alors avec un coûteau on le larde en divers endroits du corps de chandeles allumées qui brûlent enfin ce misérable\teiNote[xmlid={ftn119},n={18},place={foot}]{\teiP[xmlid={p-322}]{Jean-Baptiste Tavernier, \teiHi[rend={italic}]{Les Six Voyages I}, \teiHi[rend={italic}]{op. cit.}, p\teiHi[rend={italic}]{. }616.}}.
\end{teiQuote}
              
\end{teiCit}
              \teiP[xmlid={p6-9}]{Cette torture constitue bien un spectacle exhibé dans l’espace public, puisque le voyageur rencontre un jour, sur la route d’Ispahan à Schiras, « trois de ces suppliciez qui demanderent \lbrack{}à lui et ses compagnons\rbrack{} en grâce de hâter leur mort\teiNote[xmlid={ftn120},n={19},place={foot}]{\teiP[xmlid={p-323}]{\teiHi[rend={italic}]{Ibid.}, p.\teiHi[rend={italic}]{ }617.}} ». Cette publicité du châtiment administré revêt une fonction pédagogique destinée au peuple, comme l’ont analysé notamment Michel Foucault\teiNote[xmlid={ftn121},n={20},place={foot}]{\teiP[xmlid={p-324}]{Voir la liste d’exemples de châtiments pratiqués sous l’Ancien Régime qui inaugure \teiHi[rend={italic}]{Surveiller et punir }: Michel Foucault, \teiHi[rend={italic}]{Surveiller et punir. Naissance de la prison}, Paris, Gallimard, «\teiHi[rend={italic}]{ }Tel\teiHi[rend={italic}]{ }», 1975, p.\teiHi[rend={italic}]{ }9‑12.}}, Robert Muchembled ou encore Laura Silverman\teiNote[xmlid={ftn122},n={21},place={foot}]{\teiP[xmlid={p-325}]{Robert Muchembled, \teiHi[rend={italic}]{Le Temps des supplices : de l’obéissance sous les rois absolus, }\teiHi[rend={ small-caps italic}]{xv}\teiHi[rend={italic sup}]{e}\teiHi[rend={italic}]{-}\teiHi[rend={ small-caps italic}]{xviii}\teiHi[rend={italic sup}]{e}\teiHi[rend={italic}]{ siècle}, Paris, Le Grand livre du mois, 2000\teiHi[rend={italic}]{ }; Lisa Silverman, \teiHi[rend={italic}]{Tortured Subjects: Pain, Truth, and the Body in Early Modern France}, Chicago, University of Chicago Press, 2001.}} : le supplice constitue une forme de spectacle dissuasif. C’est bien cette perspective qui conduit Aurangzeb – le Grand Moghol – à punir non pas sa sœur Raushanara Begum, mais les eunuques ayant laissé introduire les amants de cette dernière dans son palais : il fit « un grand et exemplaire châtiment sur les eunuques, parce que c’est une chose qui regarde \lbrack{}…\rbrack{} l’honneur de la maison du roi\teiNote[xmlid={ftn123},n={22},place={foot}]{\teiP[xmlid={p-326}]{François Bernier, \teiHi[rend={italic}]{Un libertin dans l’Inde moghole}, \teiHi[rend={italic}]{op. cit.}, p.\teiHi[rend={italic}]{ }145.}} ».}
              \teiP[xmlid={p7-9}]{Mais le personnage principal de ces « micro-synoptismes » du pouvoir, c’est-à-dire cette « multitude de petits dispositifs où les regards du groupe convergent vers le corps et la personne, le visage et la voix, d’un individu \lbrack{}…\rbrack{} ou d’un petit nombre d’individus\teiNote[xmlid={ftn124},n={23},place={foot}]{\teiP[xmlid={p-327}]{Gérard Leclerc, \teiHi[rend={italic}]{Le Regard et le Pouvoir}, \teiHi[rend={italic}]{op. cit.}, p.\teiHi[rend={italic}]{ }27.}} », est le roi. Filant une métaphore théâtrale, Bernier le présente comme le « héros de la pièce\teiNote[xmlid={ftn125},n={24},place={foot}]{\teiP[xmlid={p-328}]{François Bernier,\teiHi[rend={italic}]{ Un libertin dans l’Inde moghole},\teiHi[rend={italic}]{ op. cit.}, p.\teiHi[rend={italic}]{ }53.}} » dans la violente « tragédie\teiNote[xmlid={ftn126},n={25},place={foot}]{\teiP[xmlid={p-329}]{\teiHi[rend={italic}]{Ibid.}, p.\teiHi[rend={italic}]{ }127.}} » fratricide qui l’oppose à ses frères pour succéder à leur père Shah Jahan. Le souverain est un personnage en représentation, comme le montre la magnificence de la tenue du Grand Moghol :}
              \begin{teiCit}[xmlid={cit02-8}]

                \begin{teiQuote}
Le roi paraissait assis sur son trône dans le fond de la grande salle de l’am-kas, magnifiquement vêtu. Sa veste était d’un satin blanc à petites fleurs et relevée d’une fine broderie d’or et de soie, son turban était de toile d’or et il y avait une aigrette dont le pied était couvert de diamants d’une grandeur et d’un prix extraordinaires, avec une grande topaze orientale qu’on peut dire être sans pareille, qui brillait comme un petit soleil ; un collier de grosses perles lui pendait au col jusques sur l’estomac\teiNote[xmlid={ftn127},n={26},place={foot}]{\teiP[xmlid={p-330}]{\teiHi[rend={italic}]{Ibid.}, p.\teiHi[rend={italic}]{ }264.}}.
\end{teiQuote}
              
\end{teiCit}
              \teiP[xmlid={p8-9}]{Ce « petit soleil » du Levant ne peut manquer d’évoquer celui, contemporain et tout aussi fastueux, de Louis \teiHi[rend={small-caps}]{XIV}, analogie qu’exploitent les voyageurs, nous le verrons. Cette magnificence est également sensible à l’endroit du célèbre Trône du Paon sur lequel il repose, « soutenu par six gros pieds d’or massif et tout semé de rubis, d’émeraudes et de diamants\teiNote[xmlid={ftn128},n={27},place={foot}]{\teiP[xmlid={p-331}]{\teiHi[rend={italic}]{Ibid.}}} », dont la description constitue un \teiHi[rend={italic}]{topos} des \teiHi[rend={italic}]{mirabilia} orientales décrites par les voyageurs\teiNote[xmlid={ftn129},n={28},place={foot}]{\teiP[xmlid={p-332}]{Voir la description qu’en livre Jean-Baptiste Tavernier, \teiHi[rend={italic}]{Les Six Voyages II}, Paris, Chez Gervais Clouzier et Claude Barbin, 1676, p.\teiHi[rend={italic}]{ }241‑242.}}.}
              \teiP[xmlid={p9-9}]{À cette scène moghole fait écho celle de la salle d’audience ottomane décrite par Tavernier, théâtre du pouvoir. En cas d’audience donnée aux ambassadeurs, un tapis de soie et d’or filé\teiNote[xmlid={ftn130},n={29},place={foot}]{\teiP[xmlid={p-333}]{Jean-Baptiste Tavernier, \teiHi[rend={italic}]{Nouvelle relation de l’intérieur du serrail du Grand Seigneur, contenant plusieurs singularitez qui jusqu’icy n’ont point esté mises en lumière. Par J. B. Tavernier Escuyer Baron d’Aubonne}, Paris, Chez Olivier de Varennes, 1675, p.\teiHi[rend={italic}]{ }111.}} est disposé sur le chemin menant à cette salle et un trône « assez riche » y est porté, comportant « huit couvertures tres riches faites exprès pour couvrir ce trône, \& qui viennent pendre à terre de trois côtez, par devant, à droite \& à gauche\teiNote[xmlid={ftn131},n={30},place={foot}]{\teiP[xmlid={p-334}]{Jean-Baptiste Tavernier, \teiHi[rend={italic}]{Nouvelle relation de l’intérieur du serrail du Grand Seigneur…}, \teiHi[rend={italic}]{op. cit.}, p.\teiHi[rend={italic}]{ }109.}} », qui sont couvertes de perles, de pierreries (rubis, émeraudes) et de broderies (« brocart d’or »). La richesse du trône est telle que l’auteur l’assimile à une « maniere d’Autel », suggérant par cette comparaison toute la dimension cérémoniale de l’événement, régi par une scénographie précise. En outre, l’ambassadeur est contraint de porter la veste que lui envoie le Grand Seigneur, et un « Maistre des ceremonies\teiNote[xmlid={ftn132},n={31},place={foot}]{\teiP[xmlid={p-335}]{\teiHi[rend={italic}]{Ibid.}, p.\teiHi[rend={italic}]{ }111.}} » l’accompagne devant le souverain. « Toute cette action se passe dans un grand silence\teiNote[xmlid={ftn133},n={32},place={foot}]{\teiP[xmlid={p-336}]{\teiHi[rend={italic}]{Ibid.}, p.\teiHi[rend={italic}]{ }112.}} » et les personnes présentes assistent tels des spectateurs à la scène. De plus, des notations détaillent les attitudes à la manière de didascalies : « le Prince qui doit faire le serment est debout comme j’ay dit, les mains étenduës l’une contre l’autre \& élevées à la hauteur des épaules, pour recevoir le livre de l’Alcoran de celles du \teiHi[rend={italic}]{Kapi-Aga}, qui a esté le prendre de sur le carreau l’ayant baisé \& fait toucher à sa teste\teiNote[xmlid={ftn134},n={33},place={foot}]{\teiP[xmlid={p-337}]{Jean-Baptiste Tavernier, \teiHi[rend={italic}]{Nouvelle relation de l’intérieur du serrail du Grand Seigneur…}, \teiHi[rend={italic}]{op. cit.}, p.\teiHi[rend={italic}]{ }115.}}. »}
              \teiP[xmlid={p10-9}]{Ce sens de la mise en scène apparaît aussi à l’endroit du corps même du souverain. En effet, alors qu’il était atteint d’une « violente fièvre » qui lui « faisait perdre quelquefois le jugement » et engendrait une paralysie lui « ôt\lbrack{}ant\rbrack{} presque la parole », Aurangzeb « se fit porter dans l’assemblée des \teiHi[rend={italic}]{omrahs} pour se faire voir », « dans le fort de sa maladie », « afin de désabuser ceux qui pourraient croire qu’il serait mort ». « Et, ce qui est quasi incroyable », continue l’auteur, « après être revenu d’un évanouissement qui avait fait dire par toute la ville qu’il était mort, il fit entrer deux ou trois des plus grands \teiHi[rend={italic}]{omrahs} \lbrack{}…\rbrack{} pour leur faire voir qu’il était vivant\teiNote[xmlid={ftn135},n={34},place={foot}]{\teiP[xmlid={p-338}]{François Bernier, \teiHi[rend={italic}]{Un libertin dans l’Inde moghole}, \teiHi[rend={italic}]{op. cit.}, p.\teiHi[rend={italic}]{ }137‑139.}} ». Le corps de son frère Dara fait également l’objet d’une mise en scène synoptique, puisqu’après l’avoir fait capturer, il l’expose dans une tenue et un équipage à même de signifier sa déchéance de son statut de prince :}
              \begin{teiCit}[xmlid={cit03-8}]

                \begin{teiQuote}
Ce n’était plus un de ces superbes éléphants de Ceylan ou de Pegu qu’il avait accoutumé de monter, avec des harnais dorés et des couvertures en broderie, et des sièges avec leurs dais tout peints et dorés pour se parer du soleil ; ce n’était qu’un vieil et misérable animal tout sale et tout vilain, avec une vieille couverture toute déchirée et un pauvre siège tout découvert ; on ne lui voyait plus ce collier de grosses perles que les princes ont accoutumé de porter au col et ces riches turbans \lbrack{}…\rbrack{} ou vestes en broderie ; il n’avait pour tout vêtement qu’une veste de grosse toile blanche toute sale et un turban de même, avec un misérable châle ou écharpe de Cachemire sur la tête comme un simple valet\teiNote[xmlid={ftn136},n={35},place={foot}]{\teiP[xmlid={p-339}]{\teiHi[rend={italic}]{Ibid.}, p.\teiHi[rend={italic}]{ }121.}}.
\end{teiQuote}
              
\end{teiCit}
              \teiP[xmlid={p11-9}]{La déchéance, tout autant que le pouvoir qui lui est symétriquement opposé, fait l’objet d’une mise en scène savamment orchestrée. Le Grand Moghol est particulièrement maître dans cette dramaturgie du pouvoir, comme il apparaît encore avec la fastueuse cérémonie de la pesée – lors de laquelle le roi monte sur une balance dont les poids sont en or massif –, où la puissance du souverain est indexée à son poids qui augmente d’année en année\teiNote[xmlid={ftn137},n={36},place={foot}]{\teiP[xmlid={p-340}]{\teiHi[rend={italic}]{Ibid.}, p.\teiHi[rend={italic}]{ }265.}}. Pourtant, il organise le spectacle d’une pénitence expiatrice des crimes qu’il a commis pour s’emparer du pouvoir en s’astreignant à une ascèse sévère. Tavernier – qui lui aussi raconte la guerre de succession moghole – confirme le caractère calculé de cette dernière :}
              \begin{teiCit}[xmlid={cit04-6}]

                \begin{teiQuote}
Dés le moment qu’Aureng-zeb prit possession du trône il ne voulut plus manger de pain de froment, ni de viande, ni de poisson. Il ne se nourrit que de pain d’orge, d’herbages \& de confitures, \& ne boit aucune forte liqueur. C’est une penitence qu’il s’est luy-méme imposée pour tant de crimes qu’il a commis ; mais son ambition \& le desir de régner durent toûjours, \& c’est à quoy apparemment il n’a pas dessein de renoncer de sa vie\teiNote[xmlid={ftn138},n={37},place={foot}]{\teiP[xmlid={p-341}]{Jean-Baptiste Tavernier, \teiHi[rend={italic}]{Les Six Voyages II}, \teiHi[rend={italic}]{op. cit.}, p.\teiHi[rend={italic}]{ }236.}}.
\end{teiQuote}
              
\end{teiCit}
              \teiP[xmlid={p12-9}]{Toutefois, l’auteur explique qu’il mange en cachette\teiNote[xmlid={ftn139},n={38},place={foot}]{\teiP[xmlid={p-342}]{« \lbrack{}…\rbrack{} personne ne voit manger le Roy que ses femmes \& ses Eunuques », \teiHi[rend={italic}]{ibid.}, p.\teiHi[rend={italic}]{ }245.}}. Sa méthode de gouvernement consiste principalement à « jeter \lbrack{}de la\rbrack{} poudre aux yeux du peuple\teiNote[xmlid={ftn140},n={39},place={foot}]{\teiP[xmlid={p-343}]{François Bernier, \teiHi[rend={italic}]{Un libertin dans l’Inde moghole}, \teiHi[rend={italic}]{op. cit.}, p.\teiHi[rend={italic}]{ }95.}} ». Il va même jusqu’à faire « paraître toutes les marques de douleur qu’un fils peut avoir de la perte de son père\teiNote[xmlid={ftn141},n={40},place={foot}]{\teiP[xmlid={p-344}]{\teiHi[rend={italic}]{Ibid.}, p.\teiHi[rend={italic}]{ }195.}} » lors de la mort de Shah Jahan, dont il a pourtant usurpé le pouvoir, et auprès de qui il affiche une déférence jusqu’aux dernières minutes de son règne. En effet, il lui manifeste « mille belles protestations d’affection et de soumission » avant de le capturer, prenant « secrètement » les dispositions nécessaires à son emprisonnement\teiNote[xmlid={ftn142},n={41},place={foot}]{\teiP[xmlid={p-345}]{\teiHi[rend={italic}]{Ibid.}, p.\teiHi[rend={italic}]{ }90‑92.}}.}
              \teiP[xmlid={p13-9}]{Ainsi, les voyageurs s’accordent à dépeindre le pouvoir oriental comme un dispositif spectaculaire, synoptique, dont l’objectif est d’impressionner le plus grand nombre, que ce soit par la mise en scène de tortures publiques ou par le déploiement de cérémonies fastueuses. Les émotions mêmes du souverain – sa douleur, sa contrition, etc. – font l’objet d’un spectacle insincère. Cependant, il ne s’agit là que de la « marque » ou de l’« effet\teiNote[xmlid={ftn143},n={42},place={foot}]{\teiP[xmlid={p-346}]{\teiHi[rend={italic}]{Le Dictionnaire de l’Académie françoise dedié au Roy}, \teiHi[rend={italic}]{op. cit.}, p.\teiHi[rend={italic}]{ }301.}} » du pouvoir, surface visible d’une face dissimulée, dans une politique du secret.}
            
\end{teiDiv}
            \begin{teiDiv}[type={section1},xmlid={div02-8}]

              \teiHead[xmlid={dissimulations-et-dispositifs-panoptiques-de-surveillance-001}]{Dissimulations et dispositifs panoptiques de surveillance}
              \teiP[xmlid={p14-8}]{L’ostentation du pouvoir va paradoxalement de pair avec le motif de la dissimulation, dans une tradition toute machiavélienne du secret\teiNote[xmlid={ftn144},n={43},place={foot}]{\teiP[xmlid={p-347}]{Machiavel déclare qu’il faut « être un grand simulateur et dissimulateur », alléguant l’exemple d’Alexandre\teiHi[rend={ small-caps italic}]{ }\teiHi[rend={small-caps}]{VI} qui « ne fit jamais rien que piper le monde » (\teiHi[rend={italic}]{Le Prince et autres textes}, \teiHi[rend={italic}]{op. cit.}, p.\teiHi[rend={italic}]{ }108).}}, réinvestie dans les miroirs des princes à l’âge classique\teiNote[xmlid={ftn145},n={44},place={foot}]{\teiP[xmlid={p-348}]{Bernard Teyssandier développe notamment l’exemple du texte de Jean Héroard destiné à Louis\teiHi[rend={ small-caps italic}]{ }\teiHi[rend={small-caps}]{XIII} (\teiHi[rend={italic}]{De l’institution du prince}, 1609). Bernard Teyssandier, « Figures du voile et morales du secret dans l’éducation du prince au grand siècle », Françoise Gevrey, Alexis Levrier et Bernard Teyssandier (dir.), \teiHi[rend={italic}]{Éthique, poétique et esthétique du secret de l’Ancien Régime à l’époque contemporaine}, Leuven, Peeters, 2015, p.\teiHi[rend={italic}]{ }31.}}. Or cette culture du secret imprègne également le récit que les voyageurs font de la politique orientale. Les cas de dissimulation abondent dans les textes, comme dans cet épisode à la cour persane que rapporte Tavernier :}
              \begin{teiCit}[xmlid={cit05-6}]

                \begin{teiQuote}
Cha-Abas, pour sçavoir ce qui se passoit dans son Royaume, \& le sçavoir par lui-même sans se trop confier au rapport de ses Ministres, il se déguisoit souvent, \& alloit par la ville comme un simple habitant sous prétexte de vendre \& d’acheter, tâchant de découvrir si le marchand ne faisoit point faux poids \& fausse mesure\teiNote[xmlid={ftn146},n={45},place={foot}]{\teiP[xmlid={p-349}]{Jean-Baptiste Tavernier, \teiHi[rend={italic}]{Les Six Voyages I}, \teiHi[rend={italic}]{op. cit.}, p.\teiHi[rend={italic}]{ }530.}}.
\end{teiQuote}
              
\end{teiCit}
              \teiP[xmlid={p15-7}]{Bernier, dans son récit de la guerre de succession moghole, érige la dissimulation en éthique du pouvoir : « qui ne sait pas dissimuler ne sait pas régner\teiNote[xmlid={ftn147},n={46},place={foot}]{\teiP[xmlid={p-350}]{François Bernier, \teiHi[rend={italic}]{Un libertin dans l’Inde moghole}, \teiHi[rend={italic}]{op. cit.}, p.\teiHi[rend={italic}]{ }55.}} », déclare-t-il, qualifiant Aurangzeb de « rusé et dissimulé au possible\teiNote[xmlid={ftn148},n={47},place={foot}]{\teiP[xmlid={p-351}]{\teiHi[rend={italic}]{Ibid.}, p.\teiHi[rend={italic}]{ }48.}} » et son père Shah Jahan de maître dans l’art de la tromperie. C’est effectivement beaucoup plus souvent par l’exercice de la ruse que de la force brute qu’ils parviennent à leurs fins, comme lorsque Aurangzeb enivre son frère Murad Baksh pour le destituer des insignes de son pouvoir\teiNote[xmlid={ftn149},n={48},place={foot}]{\teiP[xmlid={p-352}]{\teiHi[rend={italic}]{Ibid.}, p.\teiHi[rend={italic}]{ }97.}}.}
              \teiP[xmlid={p16-7}]{En outre, la violence qui caractérise les supplices orientaux peut également faire l’objet d’une dissimulation. En témoigne l’exemple éclatant de l’amant échaudé de Begum Saheb, autre sœur d’Aurangzeb, qui se rend elle aussi coupable d’accueillir un amant dans ses appartements lors du règne de Shah Jahan son père. Elle a à peine le temps de cacher son amant que son père, très possessif et jaloux – et connu pour être attaché à sa fille par un amour presque incestueux\teiNote[xmlid={ftn150},n={49},place={foot}]{\teiP[xmlid={p-353}]{« \lbrack{}L\rbrack{}e bruit courait même qu’il l’aimait jusques à un point qu’on a de la peine à s’imaginer », déclare François Bernier (\teiHi[rend={italic}]{ibid.}, p.\teiHi[rend={italic}]{ }49). Plus explicite, Tavernier précise qu’elle était « tout ensemble sa fille \& sa femme » (\teiHi[rend={italic}]{Les Six Voyages II}, \teiHi[rend={italic}]{op. cit.}, p.\teiHi[rend={italic}]{ }235).}} –, l’entretient un moment et « commanda fort sévèrement qu’on mît le feu à l’heure même sous la chaudière et ne voulut point partir de là que les eunuques ne lui eussent fait comprendre que le misérable était expédié\teiNote[xmlid={ftn151},n={50},place={foot}]{\teiP[xmlid={p-354}]{François Bernier, \teiHi[rend={italic}]{Un libertin dans l’Inde moghole}, \teiHi[rend={italic}]{op. cit.}, p.\teiHi[rend={italic}]{ }51.}} ». La lente agonie de la victime ne donne lieu à aucune exhibition physique ou textuelle, mais n’en est pas moins terrible. La politique du secret semble aller de pair avec une poétique de la litote et de la suggestion à la grande puissance évocatoire, comme dans le récit que Tavernier fait de cet épisode. Il nimbe de mystère la disparition de cet amant : « Depuis ce temps-là on n’en a point oüy parler, \& il n’est pas mal-aisé de croire qu’il se passe d’etranges choses dans l’enclos où ces femmes \& ces filles sont enfermées\teiNote[xmlid={ftn152},n={51},place={foot}]{\teiP[xmlid={p-355}]{Jean-Baptiste Tavernier, \teiHi[rend={italic}]{Les Six Voyages II}, \teiHi[rend={italic}]{op. cit.}, p.\teiHi[rend={italic}]{ }239.}}. » L’approximation de la source de la rumeur (« on ») et l’incertitude du sort du jeune homme laissent toute latitude à l’imagination du lecteur, qui peut se figurer les pires supplices.}
              \teiP[xmlid={p17-7}]{De nombreuses autres anecdotes viatiques correspondant à ce scénario de l’intrusion illicite d’un amant dans l’appartement des femmes – comme celle des deux amants de Raushanara Begum\teiNote[xmlid={ftn153},n={52},place={foot}]{\teiP[xmlid={p-356}]{François Bernier, \teiHi[rend={italic}]{Un libertin dans l’Inde moghole}, \teiHi[rend={italic}]{op. cit.}, p.\teiHi[rend={italic}]{ }145.}} ou encore celle de l’« amourette » entre une femme et un eunuque « coupé tout ras\teiNote[xmlid={ftn154},n={53},place={foot}]{\teiP[xmlid={p-357}]{\teiHi[rend={italic}]{Ibid.}, p.\teiHi[rend={italic}]{ }143‑144.}} » – et participent à constituer le sérail en un lieu de dissimulation et de mystères abondamment exploité par la fiction dramatique et romanesque, avec laquelle le récit de voyage entretient des rapports d’influences réciproques, ainsi que l’a analysé Sylvie Requemora\teiNote[xmlid={ftn155},n={54},place={foot}]{\teiP[xmlid={p-358}]{Voir à ce propos Sylvie Requemora, « Des anecdotes tragi-comiques. L’intertextualité romanesque dans quelques récits de voyage du \teiHi[rend={small-caps}]{xvii}\teiHi[rend={sup}]{e} siècle », dans Sophie Linon\teiHi[rend={small-caps}]{-}Chipon, Véronique Magri\teiHi[rend={small-caps}]{-}Mourgues et Sarga Moussa (dir.), \teiHi[rend={italic}]{Miroirs de textes : récits de voyage et intertextualité}, Nice, Publications de la Faculté des lettres, arts et sciences humaines de Nice, « Publications de la Faculté des lettres, arts et sciences humaines de Nice » n\teiHi[rend={sup}]{o} 49, 1998. Plus généralement, sur la poétique du récit de voyage, voir Sylvie Requemora, \teiHi[rend={italic}]{Voguer vers la modernité. Le voyage à travers les genres au }\teiHi[rend={ small-caps italic}]{xvii}\teiHi[rend={italic sup}]{e}\teiHi[rend={italic}]{ siècle}, Paris, PUPS, « Imago mundi », 2012. À ce motif de l’intrusion dans le sérail fait écho celui de l’intrusion dans le couvent. Voir, par exemple, le conte « Les Lunettes » de La Fontaine (1674), qui rapporte l’intrusion d’un « blondin » dans un couvent, ou encore les \teiHi[rend={italic}]{Lettres portugaises} (1669) de Guilleragues, roman épistolaire qui rapporte le désespoir amoureux d’une religieuse séduite par un officier.}}. Tavernier décrit son architecture comme compliquée, marquée par le secret : « au milieu de la grande galerie on voit une niche pratiquée dans le mur, où le Roy se rend de son Haram par un petit escalier dérobé\teiNote[xmlid={ftn156},n={55},place={foot}]{\teiP[xmlid={p-359}]{Jean-Baptiste Tavernier, \teiHi[rend={italic}]{Les Six Voyages II}, \teiHi[rend={italic}]{op. cit.}, p. 67.}}. » Le harem, partie du sérail où les femmes sont enfermées et soustraites à la vue de tous hormis des eunuques noirs « affreux de visage, \& coupez à net\teiNote[xmlid={ftn157},n={56},place={foot}]{\teiP[xmlid={p-360}]{\teiHi[rend={italic}]{Ibid.}, vol. I, p. 636.}} » chargés de les surveiller, incarne d’ailleurs le paroxysme de cette architecture de la dissimulation. Si le lieu du sérail est un des principaux objets de curiosité que s’emploient à décrire les voyageurs en Orient, il constitue rapidement une surface de projection de leurs désirs, dès le récit de Nicolas de Nicolay\teiNote[xmlid={ftn158},n={57},place={foot}]{\teiP[xmlid={p-361}]{« Avec Nicolay, le microcosme turc revêt une autre fonction : déjà s’y dessine ce lieu fantasmatique où l’Occidental projette et reconnaît d’inavouables désirs » (Frank Lestringant, « Guillaume Postel et l’“obsession turque” », \teiHi[rend={italic}]{Écrire le monde à la Renaissance : quinze études sur Rabelais, Postel, Bodin et la littérature géographique}, Caen, Paradigme, « Collection varia » n\teiHi[rend={sup}]{o} 6, 1994, p. 219).}}, précisément parce qu’il se dérobe à leur vue, comme l’ont montré Frank Lestringant et Michèle Longino\teiNote[xmlid={ftn159},n={58},place={foot}]{\teiP[xmlid={p-362}]{« Plus encore, Tavernier n’a jamais vu de ses propres yeux l’ensemble des appartements et des espaces reculés du sérail, ni observé par lui-même les comportements et les usages dans les lieux qui lui étaient interdits » (Michèle Longino, « Le voyageur, les eunuques et le sérail : l’oculaire par procuration », \teiHi[rend={italic}]{Littératures classiques}, vol. 82, n\teiHi[rend={sup}]{o} 3, 2013, p. 262).}}. Cette dimension fantasmatique prime bientôt sur la réalité du sérail, érigé en « \teiHi[rend={italic}]{topos }obligé de toutes les relations de voyage en Orient », analyse Alain Grosrichard. Les descriptions des sérails ottomans, persans et indiens se ressemblent toutes : « on se répète, on se copie \lbrack{}…\rbrack{} parce que l’image stéréotypée du sérail, composée au début du \teiHi[rend={small-caps}]{xvii}\teiHi[rend={sup}]{e} siècle, répond sans doute exactement à ce qu’on attend\teiNote[xmlid={ftn160},n={59},place={foot}]{\teiP[xmlid={p-363}]{Alain Grosrichard, \teiHi[rend={italic}]{Structure du sérail : la fiction du despotisme asiatique dans l’Occident classique}, Paris, Éditions du Seuil, « Champ freudien », 1979, p. 155.}}. » Le sérail devient alors dans l’imaginaire collectif un lieu qui « enserre » – idée que reflète l’orthographe fautive du mot étudiée par Jean Dubu\teiNote[xmlid={ftn161},n={60},place={foot}]{\teiP[xmlid={p-364}]{Jean Dubu explique qu’à l’époque circulait une fausse étymologie du mot \teiHi[rend={italic}]{sérail}, alors considéré comme dérivant du verbe « serrer » (Jean Dubu, « \teiHi[rend={italic}]{Bajazet }: “serrail’’ et transgression », Pierre Ronzeaud (dir.), \teiHi[rend={italic}]{Racine. La Romaine, la Turque, la Juive}, Aix-en-Provence, PUP, 1986, p. 99-113).}} –, et isole ses occupants du monde extérieur. Il constitue un mystère largement exploité dans le genre des turqueries dramatiques, où il fait figure d’« archi-lieu devenu mythique\teiNote[xmlid={ftn162},n={61},place={foot}]{\teiP[xmlid={p-365}]{Sylvie Requemora, « Scènes de sérail : la construction d’un Orient théâtral et romanesque au \teiHi[rend={small-caps}]{xvii}\teiHi[rend={sup}]{e} siècle », dans Anne Duprat et Émilie Picherot (dir.), \teiHi[rend={italic}]{Récits d’Orient dans les littératures d’Europe…}, \teiHi[rend={italic}]{op. cit.}, p. 261. Elle définit le genre des turqueries – tragi-comédies mettant en scène Soliman le Magnifique, Roxane et ses fils – comme « la théâtralisation d’une histoire secrète, se déroulant dans un palais oriental au sens large et confrontant des héros musulmans à la fois sauvages et linguistiquement “naturalisés”, en proie à la fureur amoureuse » (\teiHi[rend={italic}]{op. cit.}, p. 254).}} », explique Sylvie Requemora, et rencontre aussi un succès dans la fiction romanesque\teiNote[xmlid={ftn163},n={62},place={foot}]{\teiP[xmlid={p-366}]{Voir par exemple le roman de Madeleine de Scudéry, \teiHi[rend={italic}]{Ibrahim ou l’illustre Bassa }(1665), composé avant la pièce de son frère qu’il inspire, et au \teiHi[rend={small-caps}]{xviii}\teiHi[rend={sup}]{e} siècle, \teiHi[rend={italic}]{Les Lettres persanes }(1721) de Montesquieu, qui témoigne bien de la charge fantasmatique attachée à ce lieu d’enfermement dans lequel se révoltent Roxane et les autres femmes d’Usbek.}}.}
              \teiP[xmlid={p18-6}]{Si plusieurs opérations de pouvoir sont soustraites à la vue du plus grand nombre, elles ne sont pas inconnues du sultan qui contrôle tout ce qui se passe dans son royaume, et surtout dans son sérail, lieu de dissimulation mais aussi de surveillance où l’œil est un instrument de maîtrise des sujets. La galerie secrète lui permettant d’assister au conseil du Divan en témoigne. Sans être vu, il admire le spectacle des décisions politiques prises par un cercle restreint d’acteurs du pouvoir :}
              \begin{teiCit}[xmlid={cit06-5}]

                \begin{teiQuote}
Le Grand Seigneur y assiste le plus souvent \lbrack{}au conseil du Divan\rbrack{}, mais san estre veu ; \& c’est ce qui tient, toûjours en crainte, \& le grand Vizir \& les autres Officiers. Il peut se rendre de son appartement par une gallerie couverte à une fenestre qui répond dans la Sale du Divan, \& qui est\teiHi[rend={italic}]{ toûjours cachée }par un rideau de velours, qu’il tire quand il luy plaist, \& quand il void qu’on n’a pas rendu bonne justice\teiNote[xmlid={ftn164},n={63},place={foot}]{\teiP[xmlid={p-367}]{Jean-Baptiste Tavernier, \teiHi[rend={italic}]{Nouvelle relation de l’intérieur du serrail du Grand Seigneur…}, \teiHi[rend={italic}]{op. cit.}, p.\teiHi[rend={italic}]{ }88‑89.}}.
\end{teiQuote}
              
\end{teiCit}
              \teiP[xmlid={p19-6}]{Par sa seule vue, le sultan inspire de la crainte à ceux qui se savent de lui observés : « la vision se révèle être un instrument du pouvoir \lbrack{}…\rbrack{} crucial\teiNote[xmlid={ftn165},n={64},place={foot}]{\teiP[xmlid={p-368}]{Michèle Longino, « Le voyageur, les eunuques et le sérail : l’oculaire par procuration », \teiHi[rend={italic}]{Littératures classiques}, \teiHi[rend={italic}]{op. cit.}, p.\teiHi[rend={italic}]{ }262. Elle explique que le « rideau devient lui-même un œil » ; Roland Barthes, \teiHi[rend={italic}]{Sur Racine}, Paris, Éditions du Seuil, «\teiHi[rend={italic}]{ }Essais\teiHi[rend={italic}]{ }», 1963, p.\teiHi[rend={italic}]{ }17.}} », analyse Michèle Longino à propos de ce passage, à la suite d’Alain Grosrichard qui commente la configuration du sérail « organisé en vue et en fonction de la jouissance du despote \lbrack{}…\rbrack{} qui commence par le privilège scopique\teiNote[xmlid={ftn166},n={65},place={foot}]{\teiP[xmlid={p-369}]{Alain Grosrichard, \teiHi[rend={italic}]{Structure du sérail : la fiction du despotisme asiatique dans l’Occident classique}, \teiHi[rend={italic}]{op. cit.}, p.\teiHi[rend={italic}]{ }157.}} ». Tout, dans le sérail, sert cet « oculaire par procuration » à l’usage du despote, y compris le voile, qui « devient lui-même un œil\teiNote[xmlid={ftn167},n={66},place={foot}]{\teiP[xmlid={p-370}]{Michèle Longino, « Le voyageur, les eunuques et le sérail : l’oculaire par procuration », art. cité, p. 269.}} » ou, selon l’élégante formule de Roland Barthes, « paupière, symbole du Regard masqué\teiNote[xmlid={ftn168},n={67},place={foot}]{\teiP[xmlid={p-371}]{Barthes, \teiHi[rend={italic}]{Sur Racine}, \teiHi[rend={italic}]{op. cit.}, p.\teiHi[rend={italic}]{ }17.}} ». La configuration du sérail ne peut ainsi manquer d’annoncer le projet d’architecture carcéral du \teiHi[rend={italic}]{panopticon} proposé par Jeremy Bentham, caractéristique de la société disciplinaire organisant une « surveillance généralisée\teiNote[xmlid={ftn169},n={68},place={foot}]{\teiP[xmlid={p-372}]{Foucault, \teiHi[rend={italic}]{Surveiller et punir. Naissance de la prison}, \teiHi[rend={italic}]{op. cit.}, p.\teiHi[rend={italic}]{ }244. Voir Jérémie Bentham, \teiHi[rend={italic}]{Panoptique. Mémoire sur un nouveau principe pour construire des maisons d’inspection, et nommément des maisons de force}, Étienne Dumont (trad.), Paris, De l’Imprimerie nationale, 1791.}} » qu’étudie Michel Foucault. Ce régime de surveillance excède d’ailleurs largement les portes du sérail. Tavernier déclare par exemple qu’il est impossible de fuir la Turquie à cause des espions :}
              \begin{teiCit}[xmlid={cit07-3}]

                \begin{teiQuote}
Pour les moyens de prendre la fuite à qui auroit quelque pressentiment de sa perte, il est inutile d’y penser. Tous les Officiers \& les esclaves que les Bachas ont à leur service sont autant d’espions \& de gens qui les éclairent, \& il leur est impossible de s’en cacher. Il seroit dangereux de confier son secret à aucun d’eux, ce sont des ames basses \& incapables d’aucune belle action ; joint que les ports \& les passages sont également fermez pour les uns \& pour les autres\teiNote[xmlid={ftn170},n={69},place={foot}]{\teiP[xmlid={p-373}]{Jean-Baptiste Tavernier, \teiHi[rend={italic}]{Nouvelle relation de l’intérieur du serrail du Grand Seigneur…}, \teiHi[rend={italic}]{op. cit.}, p.\teiHi[rend={italic}]{ }175‑176.}}.
\end{teiQuote}
              
\end{teiCit}
              \teiP[xmlid={p20-6}]{De même, en Perse, des « Radars sont tellement portez dans toute la Perse » qu’« il est en effet comme impossible de se sauver de la Perse tous les passages estant bien gardez\teiNote[xmlid={ftn171},n={70},place={foot}]{\teiP[xmlid={p-374}]{Jean-Baptiste Tavernier, \teiHi[rend={italic}]{Les Six Voyages I}, \teiHi[rend={italic}]{op. cit.}, p.\teiHi[rend={italic}]{ }615. Yasmine Atlas a analysé la difficulté du passage de frontières dont témoigne le voyageur dans son récit, dans « Une frontière qualifiante : Représentation du passage en Perse chez quelques voyageurs français », \teiHi[rend={italic}]{Dix-septième siècle}, n\teiHi[rend={sup}]{o} 1, 2018, p. 49‑62.}} ». Ce régime de contrôle généralisé érige ces pays en États-sérails\teiNote[xmlid={ftn172},n={71},place={foot}]{\teiP[xmlid={p-375}]{Je remercie les directrices du volume pour m’avoir suggéré cette formule qui a émergé de nos échanges préparatoires.}}.}
              \teiP[xmlid={p21-6}]{Les voyageurs décrivent donc bien un exercice oxymoronique du pouvoir qui conjugue l’ostentation au secret, dans une dialectique contradictoire qui semble emblématique des états orientaux, comme l’indique le sous-titre du récit de Bernier : « l’or \& l’argent après avoir circulé dans le monde passent dans l’Hindoustan, d’où ils ne reviennent plus\teiNote[xmlid={ftn173},n={72},place={foot}]{\teiP[xmlid={p-376}]{Le trésor ottoman est dépeint de la même manière : « Mais pour le Tresor particulier \& secret, qui est une voûte sous terre qui ne s’ouvre qu’en la presence du Grand Seigneur, c’est une mer que je puis comparer à la mer Caspienne, où il entre plusieurs rivières que l’on n’en void point sortir » (Jean-Baptiste Tavernier, \teiHi[rend={italic}]{Nouvelle relation de l’intérieur du serrail du Grand Seigneur…}, \teiHi[rend={italic}]{op. cit.}, p.\teiHi[rend={italic}]{ }130).}}. » La machine politique y est dépeinte sous les traits d’une gueule d’enfer engloutissant pour toujours les richesses qui participent au spectacle du pouvoir, selon des opérations mystérieuses dérobées à l’intelligence du voyageur qui entreprend toutefois d’en rendre compte dans son récit.}
            
\end{teiDiv}
            \begin{teiDiv}[type={section1},xmlid={div03-7}]

              \teiHead[xmlid={le-spectacle-oxymoronique-du-pouvoir-oriental-un-anti-miroir-des-princes-001}]{Le spectacle oxymoronique du pouvoir oriental : un anti-miroir des princes}
              \teiP[xmlid={p22-5}]{Ce mystère du pouvoir parfois inquiétant ne laisse pas de fasciner les voyageurs, au même titre que son versant spectaculaire. Toutefois, est-ce par sa charge d’exotisme et d’altérité que ce pouvoir fascine, ou, au contraire, par sa troublante proximité avec le pouvoir français ? Si le pouvoir oriental se distingue du pouvoir français par sa violence et ses excès, il s’en rapproche par de nombreux aspects. Michèle Longino remarque à juste titre que « les mêmes gestes d’ostentation caractérisent et Versailles et la Sublime Porte, qui renvoient l’un à l’autre dans un effet spéculaire\teiNote[xmlid={ftn174},n={73},place={foot}]{\teiP[xmlid={p-377}]{Michèle Longino, « Le voyageur, les eunuques et le sérail : l’oculaire par procuration », art. cité, p.\teiHi[rend={italic}]{ }266.}} ». De même, le régime de surveillance des États ottomans est tout à fait analogue à celui « que pratiquaient les nombreux espions qui rôdaient dans les coulisses de Versailles » « durant cette période où se forment la culture et la politique françaises\teiNote[xmlid={ftn175},n={74},place={foot}]{\teiP[xmlid={p-378}]{\teiHi[rend={italic}]{Ibid.}, p.\teiHi[rend={italic}]{ }268.}} ». Les auteurs semblent aussi glisser certains indices mettant le lecteur sur la voie de cette équivalence : ainsi, le « petit soleil » qui irradie de la tenue du Grand Moghol\teiNote[xmlid={ftn176},n={75},place={foot}]{\teiP[xmlid={p-379}]{François Bernier, \teiHi[rend={italic}]{Un libertin dans l’Inde moghole}, \teiHi[rend={italic}]{op. cit.}, p.\teiHi[rend={italic}]{ }264.}} évoque Louis \teiHi[rend={small-caps}]{XIV}, avec qui l’association entre l’astre solaire et la fonction royale atteint son point culminant\teiNote[xmlid={ftn177},n={76},place={foot}]{\teiP[xmlid={p-380}]{Voir à ce propos, par exemple, l’ouvrage de Peter Burke, \teiHi[rend={italic}]{Louis }\teiHi[rend={ small-caps italic}]{XIV}\teiHi[rend={italic}]{ : les stratégies de la gloire}, Paul Chelma (trad.), Paris, Éditions du Seuil, 1995.}}. En outre, n’est-il pas ironique que Louis \teiHi[rend={small-caps}]{xiv} arbore comme symbole de puissance le diamant bleu indien qu’il achète à Tavernier\teiNote[xmlid={ftn178},n={77},place={foot}]{\teiP[xmlid={p-381}]{Voir la gravure, Jean-Baptiste Tavernier, \teiHi[rend={italic}]{Les Six Voyages II}, \teiHi[rend={italic}]{op. cit.}, p. 336‑337. L’histoire de ce diamant, dont a été découverte en 2007 la parenté avec le diamant bleu Hope, a donné lieu à plusieurs études. Voir Pierre Ménard, \teiHi[rend={italic}]{Le Chasseur de diamants. Les fabuleuses aventures de Jean-Baptiste Tavernier}, Paris, Tallandier, 2023.}} ?}
              \teiP[xmlid={p23-4}]{Cette proximité entre les deux souverains est abondamment exploitée par la littérature libertine et pamphlétaire, comme la \teiHi[rend={italic}]{Cour de France turbanisée}\teiNote[xmlid={ftn179},n={78},place={foot}]{\teiP[xmlid={p-382}]{Monsr L.B.D.E.D.E., \teiHi[rend={italic}]{La Cour de France turbanisée et les trahisons démasquées : en trois parties}, À Cologne, Chez Pierre Marteau, 1686.}} ou encore dans un entretien fictif entre Mahomet et Colbert où ce dernier déclare au prophète qu’« il n’est presentement rien qui ressemble plus à la \teiHi[rend={small-caps}]{Porte}, que le \teiHi[rend={small-caps}]{Louvre}\teiNote[xmlid={ftn180},n={79},place={foot}]{\teiP[xmlid={p-383}]{Anonyme, \teiHi[rend={italic}]{Entretiens dans le royaume des ténèbres entre Mahomet et Colbert}, Cologne, D. Hartman, 1683, p. 5. Pamphlet cité par Pierre Ronzeaud, « Modifications de l’image royale en situation intertextuelle polémique (1670-1685) », Noémi Hepp et Madeleine Bertaud (dir.), \teiHi[rend={italic}]{L’Image du souverain dans les lettres françaises : des guerres de religion à la révocation de l’Édit de Nantes}, Paris, Klincksieck, 1985, p.\teiHi[rend={italic}]{ }234.}} ». Dans ces textes comme dans les récits de voyage se loge un discours à charge contre les dérives autoritaires du règne du monarque français. Stéphane Van Damme, dans un article sur le libertinage, considère que la matière orientale fournit aux voyageurs le moyen d’exprimer un contenu critique, à l’instar du phénomène à l’œuvre dans la pastorale :}
              \begin{teiCit}[xmlid={cit08-3}]

                \begin{teiQuote}
Le recours aux savoirs lointains, en particulier sur l’Orient, et le dépaysement de la critique fondent, à la manière de la pastorale, une double distance critique qui est à la fois un déplacement géographique et un dialogue avec le passé\teiNote[xmlid={ftn181},n={80},place={foot}]{\teiP[xmlid={p-384}]{Stéphane Van Damme, « Curiosité subversive : « orientalisation » du libertinage et géographie morale », \teiHi[rend={italic}]{Études Épistémè. Revue de littérature et de civilisation (}\teiHi[rend={ small-caps italic}]{xvi}\teiHi[rend={italic sup}]{e}\teiHi[rend={italic}]{-}\teiHi[rend={ small-caps italic}]{xviii}\teiHi[rend={italic sup}]{e}\teiHi[rend={italic}]{ siècles)}, n\teiHi[rend={sup}]{o} 26, 2014, DOI : \teiRef[target={https://doi.org/10.4000/episteme.327}]{https://doi.org/10.4000/episteme.327}. Voir aussi Stéphane Van Damme, \teiHi[rend={italic}]{Les Voyageurs du doute : l’invention d’un altermondialisme libertin (1620-1820)}, Paris, Fayard, 2023. Pour une étude spécifique des récits de voyage libertins au \teiHi[rend={small-caps}]{xviii}\teiHi[rend={sup}]{e} siècle, voir la récente thèse de Mathilde Morinet, \teiHi[rend={italic}]{Arpenter le monde en esprit fort : les écrits de voyage libertins au }\teiHi[rend={ small-caps italic}]{xvii}\teiHi[rend={italic sup}]{e}\teiHi[rend={italic}]{ siècle}, sous la direction de Sylvie Requemora et Marie-Christine Pioffet, Aix-Marseille Université et York University, 2022.}}.
\end{teiQuote}
              
\end{teiCit}
              \teiP[xmlid={p24-4}]{L’Orient est en quelque sorte allégorisé, et permet aux voyageurs de dissimuler une critique du pouvoir qu’ils ne pourraient assumer ouvertement, avant les \teiHi[rend={italic}]{Lettres persanes }(1721) qui utilise le potentiel satirique de la fiction du regard étranger. Si Bernier jouit par exemple d’une certaine liberté de ton dans sa très libertine lettre sur la religion hindoue adressée à son ami Chapelain\teiNote[xmlid={ftn182},n={81},place={foot}]{\teiP[xmlid={p-385}]{Pour une étude de la critique libertine des superstitions hindoues, voir les nombreuses études de Frédéric Tinguely, comme \teiHi[rend={italic}]{Le Fakir et le Taj Mahal : l’Inde au prisme des voyageurs français du }\teiHi[rend={ small-caps italic}]{xvii}\teiHi[rend={italic sup}]{e}\teiHi[rend={italic}]{ siècle}, Chêne-Bourg, La Baconnière-arts, « Belles pages de la Bibliothèque de Genève », n\teiHi[rend={sup}]{o} 7, 2011 ; « Un paradis sans miracles : le Cachemire de François Bernier », Frédéric Tinguely et Adrien Paschoud (dir.), \teiHi[rend={italic}]{Voyage et libertinage : }\teiHi[rend={ small-caps italic}]{xvii}\teiHi[rend={italic sup}]{e}\teiHi[rend={italic}]{-}\teiHi[rend={ small-caps italic}]{xviii}\teiHi[rend={italic sup}]{e}\teiHi[rend={italic}]{ siècle}, Lausanne, Faculté des Lettres, « Études de lettres », n\teiHi[rend={sup}]{o} 3, 2006 ; \teiHi[rend={italic}]{La Lecture complice : culture libertine et geste critique}, Genève, Droz, « Les seuils de la modernité », vol. 20, 2016 ; « Différence et dissimulation : la figure du “fakir” dans le voyage en Inde au \teiHi[rend={small-caps}]{xvii}\teiHi[rend={sup}]{e} siècle », \teiHi[rend={italic}]{ASDIWAL, revue genevoise d’anthropologie et d’histoire des religions}, n\teiHi[rend={sup}]{o} 4, 2009, p. 94-110.}}, celle qu’il adresse à Colbert sur la description économique et politique de l’Inde le contraint à la prudence.}
              \teiP[xmlid={p25-4}]{Dominique Carnoy analyse précisément les stratégies mobilisées par les auteurs français – comme Tavernier, Chardin ou Thévenot – pour formuler ces critiques : « bien des pages de nos voyageurs s’appliqueraient à Louis \teiHi[rend={small-caps}]{XIV} plutôt qu’au sultan, pour peu qu’on en supprimât quelques adjectifs ou adverbes d’intensité », tels que « despotique » ou « tyrannique ». « Le pouvoir mahométan cumule en fait les craintes que pouvait nourrir la bourgeoisie sous Louis \teiHi[rend={small-caps}]{XIV} face à la démesure d’un absolutisme d’année en année plus affirmé\teiNote[xmlid={ftn183},n={82},place={foot}]{\teiP[xmlid={p-386}]{Dominique Carnoy, \teiHi[rend={italic}]{Représentations de l’islam dans la France du }\teiHi[rend={ small-caps italic}]{xvii}\teiHi[rend={sup}]{e} \teiHi[rend={italic}]{siècle : la ville des tentations}, Paris, L’Harmattan, 1998, p. 278‑281.}}. » L’Orient devient alors « une scène théâtrale attachée à l’Europe\teiNote[xmlid={ftn184},n={83},place={foot}]{\teiP[xmlid={p-387}]{Edward W. Said, \teiHi[rend={italic}]{L’orientalisme. L’Orient créé par l’Occident}, Catherine Malamoud (trad.), Paris, Éditions du Seuil, 2005, p. 80.}} », selon la formule d’Edward W. Saïd, sur laquelle sont interrogés des aspects de la sociétés française\teiNote[xmlid={ftn185},n={84},place={foot}]{\teiP[xmlid={p-388}]{La portée critique de la représentation de la politique orientale dans le corpus spécifiquement dramatique – celui des turqueries – a fait l’objet d’une étude par Sylvie Requemora, « Les “turqueries” : une vogue théâtrale en mode mineur », art. cité.}}. Le spectacle de la politique orientale offrirait aux lecteurs français la virtualité tyrannique contenue dans la monarchie absolue telle que l’exerce Louis \teiHi[rend={small-caps}]{XIV}.}
              \teiP[xmlid={p26-3}]{Toutefois, n’est-ce qu’une différence de degré qui sépare la monarchie française du despotisme oriental, conçu comme une version intensifiée et dégénérée de la première ? Les descriptions fascinées que les voyageurs font de certains aspects du cérémonial étranger sont l’occasion d’interroger l’essence même du pouvoir, et plus particulièrement son régime de symbolisation. Il en est ainsi de la cérémonie de la pesée évoquée par Bernier, où il apparaît que la puissance du souverain augmente à mesure que s’alourdit sa masse corporelle. Cette identification du pouvoir politique à l’état physique du souverain semble nier le régime de symbolisation qui distingue les deux corps du roi, notamment caractéristique de la monarchie anglaise, selon Ernst Kantorowicz\teiNote[xmlid={ftn186},n={85},place={foot}]{\teiP[xmlid={p-389}]{Voir Ernst Kantorowicz, \teiHi[rend={italic}]{Les Deux Corps du roi. Essai sur la théologie politique au Moyen Âge }\lbrack{}1957\rbrack{}, Jean-Philippe et Nicole Genet (trad.), Paris, Gallimard, « Bibliothèque des Histoires », 2019. Ralph E. Giesey applique un certain nombre de ces analyses à la monarchie française (\teiHi[rend={italic}]{Le Roi ne meurt jamais. Les obsèques royales dans la France de la Renaissance}, Dominique Ebnöther trad., Paris, Flammarion, 1987).}}, et en vertu de laquelle le roi, même malade, continue de renvoyer au corps symbolique de sa fonction. Cette corporéité du pouvoir est également manifeste à l’endroit de la maladie de Shah Jahan, qui fragilise son pouvoir politique et prépare les conditions de sa ruine en allumant une guerre fratricide entre ses quatre fils\teiNote[xmlid={ftn187},n={86},place={foot}]{\teiP[xmlid={p-390}]{« \lbrack{}I\rbrack{}l était tombé dans une grande maladie dont on ne croyait pas qu’il dût jamais relever, ce qui avait mis de la division entre ces quatre frères \lbrack{}Dara, Sultan Shuja, Aurangzeb et Murad Bakhsh\rbrack{} qui prétendaient tous à l’empire » (François Bernier, \teiHi[rend={italic}]{Un libertin dans l’Inde moghole}, \teiHi[rend={italic}]{op. cit.}, p.\teiHi[rend={italic}]{ }45).}}. De même, frappé par la maladie, Aurangzeb tente de conjurer sa faiblesse physique qui menace son trône en défilant devant son peuple\teiNote[xmlid={ftn188},n={87},place={foot}]{\teiP[xmlid={p-391}]{\teiHi[rend={italic}]{Ibid.}, p.\teiHi[rend={italic}]{ }139.}}. Or cette confusion du pouvoir et de la force est au fondement des régimes tyranniques, d’après la conception aristotélicienne du pouvoir politique, et contraint le souverain à établir un dispositif de surveillance généralisé pour « faire en sorte que les habitants soient toujours sous l’œil du tyran et passent leur temps à sa porte\teiNote[xmlid={ftn189},n={88},place={foot}]{\teiP[xmlid={p-392}]{Aristote, \teiHi[rend={italic}]{Les Politiques}, \teiHi[rend={italic}]{op. cit.}, livre V, chap. 11, 1313 a, p.\teiHi[rend={italic}]{ }397.}} ». La structure panoptique du sérail ou encore les grandes armées, qui suscitent la crainte fascinée des Européens et dont la description constitue un invariant des récits de voyage au Levant\teiNote[xmlid={ftn190},n={89},place={foot}]{\teiP[xmlid={p-393}]{Voir par exemple le récit de Nicolas de Nicolay qui détaille les différents grades de l’armée ottomane et les gravures qu’il intègre dans son récit, la lettre de Bernier à Colbert dans laquelle il décrit sur plus de dix pages la composition de la milice moghole (François Bernier, \teiHi[rend={italic}]{Un Libertin dans l’Inde Moghole, op. cit.}, p.\teiHi[rend={italic}]{ }205-216), ou encore le chapitre que Tavernier consacre à la description de l’ordre des soldats persans (\teiHi[rend={italic}]{Six Voyages I}, \teiHi[rend={italic}]{op. cit.}, p. 579-582).}}, sont donc autant d’expressions de la nature tyrannique du pouvoir oriental qu’a notamment analysée Silvia Murr\teiNote[xmlid={ftn191},n={90},place={foot}]{\teiP[xmlid={p-394}]{« Le tyran a pour principe (\teiHi[rend={italic}]{archè}) son propre intérêt, ce qui fait de lui l’ennemi des gouvernés, et implique que le principe de son pouvoir soit la force armée, donc la crainte de la part de sujets, d’où leur haine, donc sa propre insécurité. \lbrack{}…\rbrack{} tout le régime du Mogol repose sur la force, donc d’une part sur la disponibilité d’argent pour financer la force, d’autre part sur l’aptitude actuelle du Prince à manœuvrer cette force et à mettre en œuvre toutes les “finesses’’ de l’art militaire et de l’art politique pour rendre inefficace celle de ses adversaires \lbrack{}…\rbrack{} Bernier montre que le grand Mogol ne règne que tant qu’il est capable de mener en personne son armée » (Sylvia Murr, « Le politique “au Mogol” selon Bernier : appareil conceptuel, rhétorique stratégique, philosophie morale », \teiHi[rend={italic}]{Purusartha}, n\teiHi[rend={sup}]{o} 13, 1990, p. 266).}}. Le pouvoir tyrannique serait par conséquent moins l’effet d’une intensification de certains traits de la monarchie – de ce fait pervertie – que le résultat de la négation du caractère symbolique d’un pouvoir alors purement phénoménal qui rend le monarque « esclave de sa “montre”, c’est-à-dire de l’étalage de sa puissance financière\teiNote[xmlid={ftn192},n={91},place={foot}]{\teiP[xmlid={p-395}]{\teiHi[rend={italic}]{Ibid.}, p. 269.}} » ainsi que de sa force. Les nombreux spectacles du pouvoir oriental reflèteraient, paradoxalement, sa fragilité et sa précarité, et constitueraient un avertissement pour le monarque français.}
              \teiP[xmlid={p27-3}]{Si le souverain oriental est dépourvu de transcendance, il semble également dépourvu d’intériorité, ramené là encore à une logique purement phénoménale, au rebours d’une politique du secret dans le cadre de laquelle chaque décision royale serait le fruit d’une délibération complexe. Le prince oriental, lui, agit par caprice ou « fantaisie » : par exemple Sultan Shuja, autre frère d’Aurangzeb, « augmentait ou retranchait \lbrack{}les\rbrack{} pensions \lbrack{}des \teiHi[rend={italic}]{omrahs}\rbrack{} par fantaisie\teiNote[xmlid={ftn193},n={92},place={foot}]{\teiP[xmlid={p-396}]{François Bernier, \teiHi[rend={italic}]{Un libertin dans l’Inde moghole}, \teiHi[rend={italic}]{op. cit.}, p.\teiHi[rend={italic}]{ }47.}} ». En outre, de nombreuses scènes illustrent l’arbitraire du pouvoir du roi de Perse dans le récit de Tavernier, comme l’histoire de l’architecte chargé de construire au milieu d’Ispahan une tour avec la tête des bêtes exécutées. S’étant acquitté de son ouvrage, l’architecte vient trouver Cha-Sefi, roi de Perse, pour lui dire qu’il ne manque plus qu’une dernière tête pour constituer le sommet de la tour :}
              \begin{teiCit}[xmlid={cit09-3}]

                \begin{teiQuote}
Ce Prince dans la débauche, \& pour faire voir peut-estre aux Ambassadeurs comme il est absolu sur ses sujets, se tournant brusquement vers l’Architecte, Tu as raison, luy dit-il, \& on ne sçauroit trouver pour cela de teste plus propre que la tienne. Il fallut que le mal-heureux Architecte donnât sa teste, \& l’ordre du Roy fut executé en mesme temps\teiNote[xmlid={ftn194},n={93},place={foot}]{\teiP[xmlid={p-397}]{Jean-Baptiste Tavernier, \teiHi[rend={italic}]{Les Six Voyages I}, \teiHi[rend={italic}]{op. cit.}, p.\teiHi[rend={italic}]{ }383‑384.}}.
\end{teiQuote}
              
\end{teiCit}
              \teiP[xmlid={p28-3}]{Le despote oriental agit ici en tyran, c’est-à-dire « selon son bon vouloir\teiNote[xmlid={ftn195},n={94},place={foot}]{\teiP[xmlid={p-398}]{Aristote, \teiHi[rend={italic}]{Les Politiques}, \teiHi[rend={italic}]{op. cit.}, III, 16, 1287 a, p.\teiHi[rend={italic}]{ }266.}} », sans anticiper sa décision. Ce même roi poignarde la sultane reine en rentrant d’un soir de débauche, ivre, sans réfléchir à son geste\teiNote[xmlid={ftn196},n={95},place={foot}]{\teiP[xmlid={p-399}]{« Le Roy s’estant éveillé \& ne voyant point encore la Reine, demanda en colère pourquoi elle ne venoit pas. La Sultane mere du Roy, laquelle comme j’ay dit, n’estoit qu’une esclave Georgienne, qui haïssoit mortellement la jeune Sultane Reine qui estoit fille du Roy de Georgie, parce qu’elle s’en voyoit fort dédaignée, prit occasion de la mettre mal dans l’esprit du Roy, \& se trouvant alors dans la chambre lui fit signe de la main que la jeune Reine estoit cachée dans cette niche. Le Roy se levant alors en furie, fut donner à cette pauvre Princesse cinq ou six coups de poignard dans le ventre, \& sans estre émeu de cette action barbare se r’endormit aussi-tost » (Jean-Baptiste Tavernier, \teiHi[rend={italic}]{Les Six Voyages I}, \teiHi[rend={italic}]{op. cit.}, p.\teiHi[rend={italic}]{ }517).}}. Cette méthode de gouvernement semble incarner une anti-politique du secret, dans laquelle les décisions, instantanées, ne sont pas le fruit d’une préméditation comme elles peuvent l’être dans les histoires secrètes contemporaines qui livrent des représentations du pouvoir européen. Dans \teiHi[rend={italic}]{De l’usage de l’histoire }(1671), Saint-Réal présente en effet sa théorie selon laquelle les événements politiques seraient le produit de causes personnelles, et plus spécifiquement des mouvements des passions :}
              \begin{teiCit}[xmlid={cit10-2}]

                \begin{teiQuote}
\lbrack{}S\rbrack{}avoir, c’est connoître les choses par leurs causes ; \lbrack{}…\rbrack{} étudier l’Histoire, c’est étudier les motifs, les opinions \& les passions des hommes, pour en connoître tous les ressorts, les tours \& les détours, enfin toutes les illusions qu’elles savent faire aux esprits, et les surprises qu’elles font aux cœurs\teiNote[xmlid={ftn197},n={96},place={foot}]{\teiP[xmlid={p-400}]{César Vichard de Saint\teiHi[rend={small-caps}]{-}Réal, \teiHi[rend={italic}]{De l’usage de l’histoire}, Paris, Chez Claude Barbin et Estienne Michallet, 1671, p. 4.}}.
\end{teiQuote}
              
\end{teiCit}
              \teiP[xmlid={p29-2}]{La nouvelle\teiHi[rend={italic}]{ Dom Carlos }reflète cette conception, interprétant les relations entre Dom Carlos et son père Philippe II à l’aune de l’amour qu’ils porteraient à la reine Élisabeth de France, d’abord promise au fils, mais que le père épouse. Le secret des alcôves, les lettres interceptées, les entrevues passionnées constituent les principaux ingrédients de cette nouvelle très romanesque. Or, si la vie du sérail est nimbée de mystère pour les auteurs qui n’accèdent pas à toutes les « etranges choses\teiNote[xmlid={ftn198},n={97},place={foot}]{\teiP[xmlid={p-401}]{Jean-Baptiste Tavernier, \teiHi[rend={italic}]{Les Six Voyages II}, \teiHi[rend={italic}]{op. cit.}, p.\teiHi[rend={italic}]{ }239.}} » qui s’y passent, les prises de décisions politiques en elles-mêmes ne sont pas, à la différence du cas de l’histoire secrète, l’effet d’une lente maturation dont les ressorts seraient cachés au plus grand nombre. Elles n’expriment rien d’autre que l’arbitraire et l’inconséquence du despote. La fréquence des anecdotes de ce genre trahit par ailleurs la fascination qu’elles exercent sur les voyageurs qui ne peuvent manquer de tracer des parallèles avec les dérives arbitraires de leur souverain. C’est bien plutôt la proximité entre les régimes français et persan que leur radicale différence que l’on devine dans la prudente définition aux allures de dénégation que Tavernier livre du pouvoir despotique :}
              \begin{teiCit}[xmlid={cit11-2}]

                \begin{teiQuote}
Le Gouvernement de la Perse est purement Despotique, \& le Roy a droit de vie \& de mort sur ses sujets indépendamment d’aucun conseil, ou d’autres procedures accoûtumées dans nostre Europe. Il peut faire mourir de quelque maniere qu’il luy plaît les premiers du Royaume, sans que le corps de l’Estat s’en formalise, ni qu’on ose luy en demander raison, \& l’on peut dire qu’il n’y a point de Souverain au monde plus absolu que le Roy de Perse\teiNote[xmlid={ftn199},n={98},place={foot}]{\teiP[xmlid={p-402}]{\teiHi[rend={italic}]{Ibid.}, vol. I, p.\teiHi[rend={italic}]{ }379.}}.
\end{teiQuote}
              
\end{teiCit}
              \teiP[xmlid={p30-2}]{Tavernier, que son statut de protestant fragilise dans une période de recrudescence des persécutions religieuses, prend un soin tout particulier à distinguer le pouvoir persan du pouvoir européen, qu’il identifie à un régime encadré par des lois et des « procedures » qui empêcheraient le souverain de gouverner seul et de suivre sa fantaisie. Or c’est bien plutôt l’homologie des deux pouvoirs qui est sensible ici, dans un contexte où les sujets s’insurgent contre la centralisation et l’absolutisation croissantes du pouvoir français également taxé d’arbitraire et conçu comme un « droit de mort et pouvoir sur la vie\teiNote[xmlid={ftn200},n={99},place={foot}]{\teiP[xmlid={p-403}]{Foucault, \teiHi[rend={italic}]{Histoire de la sexualité. 1. La volonté de savoir}, Paris, Gallimard, 1976, p.\teiHi[rend={italic}]{ }177.}} ». Derrière ce diagnostic du pouvoir persan affleure celui de Louis \teiHi[rend={small-caps}]{XIV}, de même que l’histoire politique que fait Bernier de l’empire moghol constitue une anamorphose du pouvoir français que l’auteur, libertin, est forcé de dissimuler.}
              \teiP[xmlid={p31-2}]{L’importante proportion qu’occupe la description du pouvoir étranger dans les récits des voyageurs ne témoigne pas seulement de la fascination qu’il suscite chez eux, mais également de l’usage spécifique qu’ils en font. Loin de n’être qu’un objet de curiosité ethnographique, le despotisme oriental constitue pour les voyageurs un terrain de réflexion leur permettant de procéder à une autopsie du pouvoir français ainsi qu’à une évaluation de ses potentielles dérives. Toutefois, le despotisme oriental n’a-t-il pour fonction que de figurer le risque de dégénérescence tyrannique de la monarchie française, deux régimes qui ne seraient distingués que par une question de degré ? La mise en valeur des deux dynamiques apparemment contradictoires du pouvoir despotique, qui allie synoptisme et dissimulation, semble interroger la structure même du pouvoir, et notamment son aptitude à symboliser. En effet, le faste du pouvoir n’est pas un \teiHi[rend={italic}]{signe }de puissance au sens où il participerait à un dispositif de représentation, conformément à la définition qu’en donne le \teiHi[rend={italic}]{Dictionnaire de l’Académie }qui distingue le pouvoir et sa « marque », son « effet\teiNote[xmlid={ftn201},n={100},place={foot}]{\teiP[xmlid={p-404}]{\teiHi[rend={italic}]{Le Dictionnaire de l’Académie françoise dedié au Roy}, \teiHi[rend={italic}]{op. cit.}, p.\teiHi[rend={italic}]{ }301.}} », mais plutôt la manifestation purement phénoménale d’un pouvoir dépourvu de transcendance, ce qui fait échos aux réflexions contemporaines pascaliennes sur le pouvoir\teiNote[xmlid={ftn202},n={101},place={foot}]{\teiP[xmlid={p-405}]{Outre \teiHi[rend={italic}]{Les Pensées}, voir aussi les \teiHi[rend={italic}]{Trois discours sur la condition des grands} (1671) à partir desquels Louis Marin développe une analyse sur la représentation du pouvoir dans \teiHi[rend={italic}]{Le Portrait du roi }(Paris, Éditions de Minuit, 1981). En reliant la force à la tyrannie, notamment dans le fragment 298 (Brunschvicg) des \teiHi[rend={italic}]{Pensées }qui distingue la force et la justice, Pascal propose en effet une réflexion dans laquelle la représentation du pouvoir comprendrait une dimension performative, devenant elle-même un dispositif de pouvoir.}}. Quant aux mystères qui sourdent de la porte du sérail, ils ne révèlent pas d’autres secrets que celui du simple « droit de vie \& de mort\teiNote[xmlid={ftn203},n={102},place={foot}]{\teiP[xmlid={p-406}]{Jean-Baptiste Tavernier, \teiHi[rend={italic}]{Les Six Voyages I}, \teiHi[rend={italic}]{op. cit.}, p.\teiHi[rend={italic}]{ }379.}} » sur les individus, expression de l’arbitraire du despote. Ainsi, ce qu’apprend le pouvoir français à travers ce pouvoir oriental bicéphale, c’est la fragilité constitutive du régime despotique. Miroirs d’un vertige, ces récits allégorisent l’absence de transcendance d’une monarchie que l’épuisement de la puissance de symbolisation réduit à la tyrannie. Ils font entrevoir le paradoxe même d’un pouvoir fort que le processus d’absolutisation menace d’engloutir, à la manière des trésors qui entrent dans les terres orientales sans jamais en ressortir.}
            
\end{teiDiv}
          
\end{teiElement}
        
\end{teiElement}
        \begin{teiElement}[name={group},type={chapter},data-page-title={Observation, exploration et espionnage dans la Relation du voyage de la mer du Sud d’Amédée Frézier (1716)},data-page-authors={Matthias Soubise@@ENS Lyon, IHRIM, UMR 5317, Lyon, France},data-include-href={../XML/chap\_10\_Observation\_exploration\_Locus\_politicus.xml},xmlid={chapter-009}]

          \begin{teiElement}[name={front}]

            \begin{teiDiv}[type={titlePage},xmlid={titlepage-010}]

              \teiP[rend={title-main},xmlid={p-407}]{Observation, exploration et espionnage dans la \teiHi[rend={italic}]{Relation du voyage de la mer du Sud} d’Amédée Frézier (1716)}
              \teiP[rend={author-aut},xmlid={p-408}]{Matthias \teiHi[rend={small-caps}]{Soubise}}
              \teiP[rend={authority\_affiliation},xmlid={p-409}]{ENS Lyon, IHRIM, UMR 5317, Lyon, France}
            
\end{teiDiv}
          
\end{teiElement}
          \begin{teiElement}[name={body},xmlid={body-010}]

            \teiP[xmlid={p1-10}]{Bien qu’on ne dispose d’aucun document officiel informant sur les instructions et la nature de sa mission, il est certain que le voyage d’Amédée Frézier, publié en 1716 sous le titre \teiHi[rend={italic}]{Relation du voyage de la mer du Sud}, était en grande partie un voyage d’espionnage maritime\teiNote[xmlid={ftn1-10},n={1},place={foot}]{\teiP[xmlid={p-410}]{Voir Gaston Arduz Eguía et Hubert Michéa, « Préface », Amédée Frézier, \teiHi[rend={italic}]{Relation du voyage de la mer du Sud aux côtes du Chili et du Pérou }\lbrack{}1716\rbrack{}, Gaston Arduz Eguía et Hubert Michéa (éd.), Paris, Utz, 1995, p. 11.}}. La lecture du récit ne laisse pas vraiment de doute. En effet, il apparaît à plusieurs reprises comme une succession de descriptions des ports, de leurs configurations et de leurs fortifications, accompagnées de la liste des hommes en armes pour les défendre, comme s’il s’agissait de préparer un abordage militaire. Le contexte historique de la guerre de Succession d’Espagne, joint aux intérêts croissants de la marine française pour le commerce avec l’Amérique du Sud, permet de saisir les raisons d’un tel voyage, commencé secrètement en novembre 1711 à Saint-Malo où Amédée Frézier embarque pour le Pérou, avant de revenir en France en août 1714.}
            \teiP[xmlid={p2-10}]{Si l’on comprend bien souvent la littérature viatique du \teiHi[rend={small-caps}]{xviii}\teiHi[rend={sup}]{e} siècle comme l’expression d’une \teiHi[rend={italic}]{libido sciendi} à échelle mondiale, comme en témoigne l’accumulation des explorations scientifiques lancées pendant cette période, il faut rappeler que le voyage scientifique et l’espionnage militaire forment les deux facettes complémentaires de la question des voyages et de la connaissance des espaces extra-européens. Ces deux types d’expédition forment une entreprise totale qui met en œuvre un rapport complexe à l’observation et à la description, pour produire un savoir sur les formes de la nature et les conditions politico-militaires des territoires parcourus\teiNote[xmlid={ftn31-3},n={2},place={foot}]{\teiP[xmlid={p-411}]{Pour une étude de cas stimulante sur l’imbrication entre voyage botanique et espionnage, voir Thibaud Martinetti, « L’Autre Jésuite : civilisation, exploration et espionnage dans le récit de voyage de Fusée-Aublet en Guyane (1762-1764) », \teiHi[rend={italic}]{Nuevo Mundo Mundos Nuevos}, « Débats », 2022.}}.}
            \teiP[xmlid={p3-10}]{Les voyages ont un rôle stratégique (militaire et commercial) important : la \teiHi[rend={italic}]{libido sciendi} des voyageurs scientifiques n’est jamais éloignée de la \teiHi[rend={italic}]{libido dominandi} des explorateurs, des colons ou des administrateurs. Les secrets de ces territoires étrangers, que le voyage se doit de révéler au grand jour, sont donc autant naturels que politiques. La faune, la flore et la géologie doivent être décrites au même titre que les mœurs des habitants et des ethnies indigènes, la topographie des villes et des ports, afin d’être vues et comprises du public français encore largement ignorant de cette partie du monde, certes colonisée de longue date mais jalousement gardée des regards étrangers par les empires ibériques.}
            \teiP[xmlid={p4-10}]{L’édition de 1771 du dictionnaire de Trévoux donne du terme « explorateur » la définition suivante :}
            \begin{teiCit}[xmlid={cit01-10}]

              \begin{teiQuote}
Espion est le terme ordinaire ; mais il y a des mots inusités qui ont quelque chose de noble et de hardi qui plaît d’abord : il semble que l’usage ait tort de ne les pas recevoir. \teiHi[rend={italic}]{Explorateur} paroît assez de ce caractère. Je crois qu’un peu d’adresse à le produire lui feroit faire aisément fortune, et que l’usage, tout tyran qu’il est, se laisserait fléchir en sa faveur. On peut ajouter que le mot d’\teiHi[rend={italic}]{explorateur} paroît annoncer des fonctions plus nobles et plus distinguées que celui d’espion. \lbrack{}…\rbrack{} Il semble qu’on pourroit encore l’appliquer à celui qu’on envoie à la découverte d’un pays, pour en connoître la situation, l’étendue, \&c\teiNote[xmlid={ftn32-3},n={3},place={foot}]{\teiP[xmlid={p-412}]{« Explorateur », \teiHi[rend={italic}]{Dictionnaire universel françois et latin}, Paris, Compagnie des libraires associés, 1771.}}.
\end{teiQuote}
            
\end{teiCit}
            \teiP[xmlid={p5-10}]{Marie-Noëlle Bourguet a montré que, dans cette définition déjà un peu tardive, s’opérait un glissement sémantique de taille : on passe en effet d’une définition militaire à un programme scientifique, certes encore naissant dans la deuxième moitié du siècle\teiNote[xmlid={ftn33-4},n={4},place={foot}]{\teiP[xmlid={p-413}]{Marie-Noëlle Bourguet, « L’Explorateur », Michel Vovelle (dir.), \teiHi[rend={italic}]{L’Homme des Lumières}, Paris, Éditions du Seuil, 1996, p. 289.}}. L’indécision des termes traduit celle des choses. Pendant toute l’époque moderne, il existe des voyageurs, mais ce mot recoupe des réalités bien différentes, du marin embarqué sur un navire marchand à l’homme de sciences parti en expédition, sans oublier les missionnaires\teiNote[xmlid={ftn34-3},n={5},place={foot}]{\teiP[xmlid={p-414}]{Voir François Moureau, \teiHi[rend={italic}]{Le Théâtre des voyages : une scénographie de l’Âge classique}, Paris, Presses de l’université Paris-Sorbonne, 2005, p. 12.}}. Cette pluralité des fonctions du voyageur a des conséquences sur la réception des récits de voyage, auxquels on prête en général peu de crédit, sauf quand l’auteur peut donner des gages de sa bonne foi et ne semble pas chercher à faire de la littérature (et donc mentir\teiNote[xmlid={ftn35-6},n={6},place={foot}]{\teiP[xmlid={p-415}]{Au sujet du \teiHi[rend={italic}]{topos} du voyageur-menteur, voir Anne-Gaëlle Weber, \teiHi[rend={italic}]{À beau mentir qui vient de loin. Savants, voyageurs et romanciers au }\teiHi[rend={ small-caps italic}]{xix}\teiHi[rend={italic sup}]{e}\teiHi[rend={italic}]{ siècle}, Paris, Honoré Champion, 2004.}}). À l’époque de Frézier, le voyage implique des réalités aussi différentes qu’indissociables. Explorateurs, voyageurs et espions ont ceci en commun qu’il s’agit d’apprendre à bien voir et à transformer le regard en discours. L’édition de 1718 du dictionnaire de l’Académie française définit le verbe « épier » comme « Observer secretement \& adroitement les actions, les discours de quelqu’un\teiNote[xmlid={ftn36-6},n={7},place={foot}]{\teiP[xmlid={p-416}]{« Espier », \teiHi[rend={italic}]{Dictionnaire de l’Académie française}, Paris, Chez Jean-Baptiste Coignard, 1718.}} ». Espionnage et voyage scientifique partagent cet impératif du regard compris comme une opération active de l’esprit et comme une faculté entraînée.}
            \teiP[xmlid={p6-10}]{Dans l’épître dédicatoire qu’il destine au duc d’Orléans, régent du royaume, Frézier décrit l’ambition de sa relation de voyage par trois termes :}
            \begin{teiCit}[xmlid={cit02-9}]

              \begin{teiQuote}
C’est un Recueil des Observations que j’ai faites sur la Navigation, sur les erreurs des Cartes, \& sur la situation des Ports \& des Rades où j’ai été. C’est une Description des Animaux, des Plantes, des Fruits, des Métaux, \& de ce que la terre produit de rare dans les plus riches Colonies du Monde. Ce sont des Recherches exactes sur le Commerce, sur les Forces, le Gouvernement, \& les mœurs des Espagnols-Créoles \& des Naturels du Pays, dont j’ai parlé avec tout le respect que je dois à la Verité\teiNote[xmlid={ftn37-6},n={8},place={foot}]{\teiP[xmlid={p-417}]{Amédée Frézier, \teiHi[rend={italic}]{Relation du voyage de la mer du Sud aux côtes du Chili, du Pérou et du Brésil, fait pendant les années 1712, 1713 \& 1714 }\lbrack{}1\teiHi[rend={sup}]{re} édition : Paris, Nyon/Ganeau/Quillau, 1716\rbrack{}, Amsterdam, Pierre Humbert, 1717, t. I, « Épître », p. vi-vii.}}.
\end{teiQuote}
            
\end{teiCit}
            \teiP[xmlid={p7-10}]{« Observations », « descriptions » et « recherches » : par ce vocabulaire, Frézier présente son ouvrage comme le résultat d’une méthode qu’on pourrait qualifier de scientifique, dans la mesure où elle semble reprendre les grands traits de la science expérimentale baconienne, alors qu’une bonne partie de ce qui est observé ne relève pas de l’histoire naturelle, mais des sociétés, des constructions et des richesses de l’Amérique méridionale. Le récit de voyage de Frézier place en son cœur l’idée d’observation, sans jamais véritablement définir son objet ni sa méthode. Ce flottement est utile puisqu’il permet de jouer sur plusieurs tableaux : le voyageur n’est jamais seulement un espion ou un homme de sciences ou un observateur politique, mais tout cela à la fois selon l’itinéraire suivi, les informations collectées et l’organisation du récit. On s’intéressera ici à la manière dont Frézier se met en scène comme homme d’action avant de comprendre de quelle façon il traduit textuellement sa fonction d’observateur. Ce regard du voyageur cherche à percer les secrets d’une Amérique encore mal connue : c’est dans cette volonté de révéler une vérité dissimulée que réside la valeur politique et poétique du récit.}
            \begin{teiDiv}[type={section1},xmlid={div01-9}]

              \teiHead[xmlid={etre-un-explorateur-arpenter-et-affronter-le-territoire-001}]{Être un explorateur : arpenter et affronter le territoire}
              \teiP[xmlid={p8-10}]{L’explorateur construit au début de son récit une image héroïque de lui-même, à la hauteur de la mission qui lui incombe. Contrairement au père Louis Feuillée, qui avait emprunté quelques années auparavant quasiment le même itinéraire\teiNote[xmlid={ftn38-7},n={9},place={foot}]{\teiP[xmlid={p-418}]{Louis Feuillée, \teiHi[rend={italic}]{Journal des observations physiques, mathematiques et botaniques, faites par l’ordre du Roy sur les côtes orientales de l’Amerique meridionale, \& dans les Indes occidentales, depuis l’année 1707 jusques en 1712}, Paris, Pierre Giffart, 1714.}}, Frézier se présente dans la préface de son ouvrage comme « un jeune homme de fatigue », capable de s’attaquer à un dur travail cartographique qui nécessite « d’aller chercher plusieurs stations dans les lieux écartez, couverts, ou de difficile accès », ce que n’aurait pas pu faire le religieux, « point d’un âge à de rudes exercices\teiNote[xmlid={ftn39-7},n={10},place={foot}]{\teiP[xmlid={p-419}]{Pour les trois dernières citations : Amédée Frézier, \teiHi[rend={italic}]{op. cit.}, t. I, « Avertissement », p. xi.}} ». Le père Feuillée avait en effet quarante-trois ans lors de son départ de Marseille en 1703, alors que Frézier n’en a même pas trente en 1712. Il est vrai que le récit de voyage de Feuillée donne une forte impression d’immobilité : on reste la plupart du temps sur le bateau et le récit prend souvent l’aspect d’une suite de tables d’observations astronomiques. Rien de tel dans la relation de Frézier, dans laquelle il se montre plutôt comme un homme d’action. Si les deux voyageurs partagent les mêmes objectifs – accroître la somme des connaissances scientifiques et cartographiques pour être utiles à leur patrie –, le début du récit de Frézier place le voyage sous l’angle de l’assouvissement d’une curiosité ardente, d’une \teiHi[rend={italic}]{libido sciendi} qui le pousse à se lancer dans le monde.}
              \begin{teiCit}[xmlid={cit03-9}]

                \begin{teiQuote}
La structure de l’Univers, qui est naturellement l’objet de notre admiration, a toujours fait aussi le sujet de ma curiosité ; dès l’enfance je faisois mon plus grand plaisir de tout ce qui pouvoit m’en donner la connoissance, les Globes, les Cartes, les Relations des Voyageurs avoient pour moi des attraits singuliers. À peine étois-je en état de voir les choses par moi-même que j’entrepris le Voyage d’Italie. Le prétexte des Études me servit ensuite à parcourir une partie de la France ; mais enfin fixé par l’Emploi que j’ai eu l’honneur d’obtenir au service du Roi : je croyois qu’il ne me restoit plus d’esperance de satisfaire l’inclination que j’avois de voyager, lorsqu’il plut à Sa Majesté de me permettre de profiter de l’occasion qui se presentoit de voir le Chili \& le Pérou\teiNote[xmlid={ftn40-7},n={11},place={foot}]{\teiP[xmlid={p-420}]{\teiHi[rend={italic}]{Ibid}., t. I, p. 1-2.}}.
\end{teiQuote}
              
\end{teiCit}
              \teiP[xmlid={p9-10}]{En comparaison, la justification de l’intérêt de voyager au début du journal de Feuillée met au premier plan l’intérêt pour l’astronomie, dont le voyage ne semble être qu’une mise en pratique\teiNote[xmlid={ftn41-7},n={12},place={foot}]{\teiP[xmlid={p-421}]{« Ayant eu dès ma plus tendre jeunesse une inclination naturelle pour les Mathématiques, je me sentis porté plus particulièrement à l’Astronomie \lbrack{}…\rbrack{}. / Je méditai de mettre en pratique les connoissances que j’avois aquises, \& je conçus le dessein de travailler à la perfection de l’Astronomie, de la Géographie, \& de l’Hidrographie » (Louis Feuillée, \teiHi[rend={italic}]{Journal des observations physiques… }, \teiHi[rend={italic}]{op. cit.}, p. 1).}}. Frézier, quant à lui, semble davantage mu par une curiosité géographique et un désir d’agir pour le royaume. S’il présente néanmoins des prétentions scientifiques – il est, après tout, ingénieur du roi – celles-ci sont toujours secondaires et ancillaires. L’« avertissement » qui précède la relation annonce sans équivoque son refus de plaire et la primauté de l’utilité du voyage pour les navigateurs et le royaume :}
              \begin{teiCit}[xmlid={cit04-7}]

                \begin{teiQuote}
\lbrack{}…\rbrack{} je ne préviens point ici le Lecteur de ce qu’il trouvera de curieux dans cet Ouvrage ; j’avoue qu’il y auroit beaucoup à retrancher pour ceux qui ne se soucient point de Navigation, si l’on devoit pour l’agréable negliger entierement l’utile ; mais il importe plus à la Republique pour le bien du Commerce, qu’on connoisse les saisons, les vents generaux, les courans, les écueils, les bons mouillages, \& les débarquemens, que des choses simplement curieuses \& divertissantes. \lbrack{}…\rbrack{} On doit apporter plus de soin à la conservation des Vaisseaux \& de leurs agrès, \& d’attention au salut de ceux qui travaillent pour la Patrie, qu’à satisfaire la curiosité de ceux qui dans une vie molle jouïssent des avantages que leur procurent les Navigateurs par des travaux infinis, \& en s’exposant à mille dangers\teiNote[xmlid={ftn42-7},n={13},place={foot}]{\teiP[xmlid={p-422}]{Amédée Frézier, \teiHi[rend={italic}]{Relation du voyage…},\teiHi[rend={italic}]{ op. cit.}, t. I, « Avertissement », p. xii-xiii.}}.
\end{teiQuote}
              
\end{teiCit}
              \teiP[xmlid={p10-10}]{Le refus du « curieux » et de l’« agréable » traduit une stratégie d’écriture qui, dès le \teiHi[rend={small-caps}]{xvii}\teiHi[rend={sup}]{e} siècle, cherche à accréditer le discours du voyageur contre les accusations d’affabulation dont il pouvait faire l’objet\teiNote[xmlid={ftn43-7},n={14},place={foot}]{\teiP[xmlid={p-423}]{Voir à ce sujet l’article de Jacques Chupeau, « Les récits de voyage aux lisières du roman », \teiHi[rend={italic}]{Revue d’histoire littéraire de la France}, vol. 77, nᵒ 3-4, 1977, p. 536-553.}}. Mais derrière cette sorte d’anti-\teiHi[rend={italic}]{captatio} apparaît aussi une mise en scène de soi qui insiste, à demi-mots, sur sa mission au service du royaume et l’action concrète qu’elle implique. La mission n’est jamais nommée explicitement dans le récit, mais le voyageur fait état régulièrement d’une inquiétude propre à son statut. Ne pas être connu, ou que son identité ou sa qualité de Français ne soient pas révélées, sont des éléments récurrents du récit, qui traduisent le secret que Frézier doit garder, et donc l’objectif de son voyage\teiNote[xmlid={ftn44-6},n={15},place={foot}]{\teiP[xmlid={p-424}]{Par exemple, expliquant les raisons de son départ précipité de Rio de Janeiro : « Au-reste il ne m’auroit de rien servi que nous y eussions séjourné plus long-tems ; quelques indiscrets de notre Escadre m’ayant fait connoître aux Officiers Portugais pour Ingénieur, il ne me convenoit plus de m’exposer à quelque affront dans un lieu où le souvenir de l’expédition de Rio de Janeiro, encore recent, rendoit la Nation suspecte. En effet on avoit doublé les gardes partout, \& même établi de nouveaux corps de garde, sur ce qu’il se trouvait déjà en rade cinq Navires François, parmi lesquels il y en avoit deux de force, un de 50 \& l’autre de 70 canons » (Amédée Frézier, \teiHi[rend={italic}]{Relation du voyage…},\teiHi[rend={italic}]{ op. cit.}, t. II, p. 474-476).}}.}
              \teiP[xmlid={p11-10}]{Cet \teiHi[rend={italic}]{ethos} de l’explorateur ainsi défini a une conséquence stylistique majeure tout au long du récit, puisqu’il s’agit, dans l’ordre du discours, de rendre visibles et intelligibles des places fortes : il faut que la description soit pratique, qu’elle aide à reconnaître, à se figurer mentalement les lieux et les caractéristiques dans l’éventualité d’un voyage entrepris après-coup. Faire connaître reviendrait avant tout à faire voir.}
            
\end{teiDiv}
            \begin{teiDiv}[type={section1},xmlid={div02-9}]

              \teiHead[xmlid={lobservation-voir-et-savoir-001}]{L’observation : voir et savoir}
              \teiP[xmlid={p12-10}]{« Voir, faire voir et faire savoir, tel sera dès l’origine le programme du voyageur\teiNote[xmlid={ftn45-6},n={16},place={foot}]{\teiP[xmlid={p-425}]{Roland Le Huenen, \teiHi[rend={italic}]{Le Récit de voyage au prisme de la littérature}, Paris, Presses de l’université Paris-Sorbonne, 2015, p. 27.}}. » Ce programme est commun à toutes les formes du voyage, du découvreur des premiers temps de la conquête de l’Amérique jusqu’aux voyageurs naturalistes, en passant par les missionnaires et les espions. Dans l’édition néerlandaise du récit de Frézier, publiée en 1717, est inséré un frontispice (fig. 1) qui se lit comme un résumé visuel de tout \teiHi[rend={italic}]{ce qui a été vu} dans le voyage : animaux étranges, plantes, outils et instruments, peuples indigènes et leurs mœurs, volcan fumant en arrière-plan. Loin de toute vraisemblance, ce frontispice fait comprendre le récit de voyage comme une succession de choses vues, effaçant complètement la mission militaire. Ce qui est vu est aussi ce qui est rapporté en France, en lieu et place des choses elles-mêmes. Espions et naturalistes collectent le monde par les yeux. Dans le cas de Frézier, il s’agit de bien voir pour pouvoir décrire du mieux possible. Cela traduit une compréhension essentiellement utilitaire de l’observation.}
              \begin{teiFigure}[xmlid={figure01}]

                \teiHead[xmlid={figure-1-frontispice-de-la-relation-du-voyage-de-la-mer-du-sud-amsterdam-pierre-humbert-1717-getty-research-institute-los-angeles-courtesy-of-the-getty-research-institute-001}]{Figure 1. Frontispice de la \teiHi[rend={italic}]{Relation du voyage de la mer du Sud}, Amsterdam, Pierre Humbert, 1717. Getty Research Institute, Los Angeles. Courtesy of the Getty Research Institute.}
              
\end{teiFigure}
              \teiP[xmlid={p13-10}]{Décrivant la forteresse de Valparaíso au Chili, Frézier opère une géométrisation du regard, certes propre à l’objet représenté, mais caractéristique du genre d’observation qu’implique le voyage d’espionnage militaire\teiNote[xmlid={ftn46-6},n={17},place={foot}]{\teiP[xmlid={p-426}]{Amédée Frézier, \teiHi[rend={italic}]{Relation du voyage…},\teiHi[rend={italic}]{ op. cit}., t. I, p. 159-163.}}. La forteresse est ainsi décrite selon sa topographie, sa situation et son orientation, la profondeur du fossé et l’épaisseur des murailles. Le voyageur semble vouloir la montrer de la manière la plus impersonnelle possible, se servant pour cela uniquement de tournures passives et d’un « on » qui permet de fondre l’observation personnelle et l’action des Espagnols qui ont édifié la fortification. Mais derrière cette apparente neutralité de l’observation se révèle une appréciation dépréciative de la construction.}
              \begin{teiCit}[xmlid={cit05-7}]

                \begin{teiQuote}
Le côté de la montagne est composé d’une courtine de vingt-six toises \& de deux demi Bastions de vingt toises de face \& onze de flanc, de sorte que la ligne de défense n’est que de quarante-cinq toises ; toute cette partie est bâtie de briques, élevée de vingt-cinq pieds de haut sur une berme ; la profondeur du fossé est d’environ dix pieds, \& sa largeur est de trois toises vers les angles saillans, d’où il tire sa défense à l’angle de l’épaule ; il est creusé dans du rocher pourri, que l’on a un peu escarpé aux deux bouts pour le rendre inaccessible par les coulées. Les parapets n’ont que deux pieds ½ d’épaisseur, \& le reste du contour de la Place n’est que d’une massonnerie de moilon aussi foible\teiNote[xmlid={ftn47-6},n={18},place={foot}]{\teiP[xmlid={p-427}]{\teiHi[rend={italic}]{Ibid.}, p. 161.}}.
\end{teiQuote}
              
\end{teiCit}
              \teiP[xmlid={p14-9}]{La description d’architecture militaire met en évidence un regard sur l’état de la fortification, qui permet à Frézier de conclure ainsi : « l’on voit combien peu est redoutable la Forteresse de Valparaisso, dès qu’on auroit mis pied à terre, comme on le peut faire de beau temps\teiNote[xmlid={ftn48-6},n={19},place={foot}]{\teiP[xmlid={p-428}]{\teiHi[rend={italic}]{Ibid.}, p. 163.}}. » L’observation méticuleuse et informée permet, comme par transparence, de connaître l’état des forces militaires en présence dans le but – hypothétique – d’une attaque, l’œil habile du voyageur faisant apparaître en même temps que l’édifice les secrets militaires de la colonie espagnole. À cela, Frézier ajoute trois gravures qui illustrent son propos et permettent de multiplier les regards sur ce même objet (fig. 2). La description, qui donne accès à la vue d’ensemble, aux détails et à l’essence de l’objet, est complétée par une image qui apparaît comme une représentation objective de la forteresse, et sans doute plus parlante qu’un texte de plusieurs pages pour un marin envoyé en reconnaissance. Le récit de voyage se veut avant tout didactique et pratique : il restitue une vision passée pour guider les futurs voyageurs. La description de cette forteresse, pour banale qu’elle puisse paraître, permet de penser la pluralité des fonctions attribuées à la vue. Le livre, qu’il s’agisse du texte et des images, doit apparaître comme le miroir fidèle des lieux représentés. Dans un même mouvement, à la description se superpose la connaissance des choses. Le « on voit » qui conclut la description ne se rapporte pas à la faculté sensible, mais à l’intelligence. Ce verbe est d’ailleurs complètement absent de la description, Frézier préférant user de tournures impersonnelles qui effacent sa propre présence, et il n’apparaît que pour exprimer un \teiHi[rend={italic}]{savoir} sur l’objet précédemment décrit.}
              \begin{teiFigure}[xmlid={figure02}]

                \teiHead[xmlid={figure-2-profil-du-fort-de-valparaysso-detail-dans-amedee-frezier-relation-du-voyage-de-la-mer-du-sud-paris-nyon-ganeau-quillau-1716-paris-bibliotheque-nationale-de-france-site-gallica-001}]{Figure 2. \teiHi[rend={italic}]{Profil du fort de Valparaysso }(détail), dans Amédée Frézier, \teiHi[rend={italic}]{Relation du voyage de la mer du Sud}, Paris, Nyon/Ganeau/Quillau, 1716. Paris, Bibliothèque nationale de France, site \teiHi[rend={italic}]{Gallica}.}
              
\end{teiFigure}
              \teiP[xmlid={p15-8}]{Cette tension entre le voir et le savoir a une incidence sur la manière même de construire le récit de voyage. Frézier cherche parfois moins à \teiHi[rend={italic}]{faire voir} ce qu’il a vu qu’à dire ce qu’il a compris des lieux parcourus. Cela le fait glisser d’une rhétorique de l’autopsie, fondée sur le témoignage oculaire, à une organisation presque encyclopédique, qui cherche la collecte des savoirs sans nécessairement les adosser à une pratique personnelle du voyage. L’utilisation presque généralisée des tournures impersonnelles et d’un « on voit » à la place d’un « je vois » témoigne de ce glissement. Ce procédé, comme on a pu le remarquer ailleurs, permet de fondre ensemble l’observation individuelle et les observations réalisées par des tiers ou tirées de lectures, s’il s’agit d’événements passés\teiNote[xmlid={ftn49-6},n={20},place={foot}]{\teiP[xmlid={p-429}]{Frézier (comme d’autres voyageurs) utilise notamment ce procédé pour rapporter des observations de tremblements de terre auxquels il n’a pas assisté, créant de ce fait une impression de témoignage de première main. Voir notre article : « “On s’accoutume à tout, même aux tremblements de terre.” Expérience et écriture du séisme dans quelques récits de voyage français en Amérique du Sud au début du \teiHi[rend={small-caps}]{xviii}\teiHi[rend={sup}]{e} siècle », \teiHi[rend={italic}]{Viatica}, nᵒ 9, 2022, notamment les § 24-25.}}. C’est notamment le cas de sa description des mines de Potosi, où il est très peu probable que Frézier soit allé\teiNote[xmlid={ftn50-6},n={21},place={foot}]{\teiP[xmlid={p-430}]{« Potosi est cette Ville si renommée dans tout le Monde, par les immenses richesses qu’on a tirées autrefois, \& qu’on tire encore de la montagne, au pied de laquelle elle est bâtie » (Amédée Frézier, \teiHi[rend={italic}]{Relation du voyage…},\teiHi[rend={italic}]{ op. cit.}, t. I, p. 252).}}. Remontant les côtes du Chili vers le Pérou, Frézier quitte le port de Coquimbo pour rejoindre la ville de Copiapó, passe devant le port de Cobija avant d’arriver à Arica, alors situé dans le vice-royaume du Pérou. Mais la description de Potosi est placée au moment de la reconnaissance de Cobija. Il mentionne la route qui mène jusqu’à Lipes (San Pablo de Lípez) puis à Potosi, qu’il n’a jamais empruntée. Pourtant, Potosi est présentée d’après ce que le voyageur en a lu ou entendu : le nombre de sa population, le rendement de ses mines ou encore la rudesse de son climat. Frézier fait alterner dans son récit des descriptions issues de son expérience personnelle du voyage et des descriptions de seconde main, lissées par une énonciation qui vise à gommer l’expérience individuelle pour se présenter comme une somme de savoirs sur le territoire parcouru.}
              \teiP[xmlid={p16-8}]{L’ambition de connaissance telle que Frézier la pratique semble cependant contredire la mission d’espionnage, du moins dans sa finalité. L’espionnage de Frézier est, on peut le dire sans trop de doute, un échec, du moins si on le comprend telle que serait, pour nous, une mission réussie d’agent secret. On peut s’étonner que la couronne française ait permis la diffusion de toutes les informations et les cartes contenues dans la relation de voyage (tous les ports de l’Amérique du Sud sont décrits, cartographiés et illustrés), car cette propension à vouloir tout dire et tout montrer, certes propre à l’accréditation du discours du voyageur, abolit l’un des traits essentiels de l’espionnage, à savoir le secret inhérent aux informations découvertes\teiNote[xmlid={ftn51-6},n={22},place={foot}]{\teiP[xmlid={p-431}]{Gaston Arduz Eguía et Hubert Michéa, « Préface », Amédée Frézier, \teiHi[rend={italic}]{Relation du voyage…}, \teiHi[rend={italic}]{op. cit.}, p. 18.}}. L’espionnage est certes une découverte, mais une découverte à public restreint. Une grande partie des missions d’explorations n’était pas imprimée, mais circulait auprès de quelques lecteurs choisis sous forme manuscrite, afin que les informations révélées servent d’abord à l’État. Or Frézier, par sa publication et sa volonté affichée de faire voir et comprendre les lieux parcourus, contredit le projet même d’espionnage, au point que son texte a sans doute davantage servi à des puissances étrangères qu’à la marine française. On pourrait donc en tirer la conclusion que c’est sans doute moins le secret de la mission qui intéresse Frézier que la volonté de faire découvrir les secrets d’un territoire encore mal connu. La \teiHi[rend={italic}]{libido sciendi} du voyageur se conjugue mal avec l’impératif de discrétion de l’espion, mais laisse plus de place à l’émergence d’un regard critique sur un monde qu’il convient non seulement de connaître, mais aussi d’interpréter. La découverte de l’explorateur concerne moins le territoire de l’Amérique – Frézier n’en dit rien de vraiment nouveau – que les secrets et les rouages d’un monde et d’une société encore méconnus.}
            
\end{teiDiv}
            \begin{teiDiv}[type={section1},xmlid={div03-8}]

              \teiHead[xmlid={reveler-les-secrets-de-lamerique-001}]{Révéler les secrets de l’Amérique}
              \teiP[xmlid={p17-8}]{La relation de Frézier prend parfois l’allure du \teiHi[rend={italic}]{Diable boiteux} de Lesage publié un peu plus tôt en 1707\teiNote[xmlid={ftn52-6},n={23},place={foot}]{\teiP[xmlid={p-432}]{Alain-René Lesage, \teiHi[rend={italic}]{Le Diable boiteux}, Paris, Veuve Barbin, 1707.}}. Le diable Asmodée y fait découvrir aux yeux du jeune Cléofas les travers moraux qui se cachent sous les toits des maisons de Madrid. Rien ne permet de dire avec certitude que Frézier comptait Lesage dans sa bibliothèque, mais il ne semble pas invraisemblable de situer ce récit de voyage dans la continuité d’un \teiHi[rend={italic}]{topos} fictionnel récurrent au \teiHi[rend={small-caps}]{xviii}\teiHi[rend={sup}]{e} siècle, à savoir la révélation, le plus souvent par un \teiHi[rend={italic}]{autre} (supposé réel ou fantasmé), des secrets et des travers – le plus souvent moraux – d’une société. S’il est généralement admis que ce thème est au cœur de plusieurs dispositifs romanesques, chez Lesage comme chez Montesquieu ou Diderot, il faut aussi admettre les récits de voyage dans ce mouvement résolument critique qui caractérise, pour le dire rapidement, les Lumières\teiNote[xmlid={ftn53-7},n={24},place={foot}]{\teiP[xmlid={p-433}]{Le baron de Lahontan en est le parfait exemple, qui publie les \teiHi[rend={italic}]{Nouveaux voyages de Mr le baron de Lahontan dans l’Amérique septentrionale }(La Haye, Frères L’Honoré, 1703) puis la \teiHi[rend={italic}]{Suite du voyage de l’Amérique ou Dialogue de M. de Baron de Lahontan et d’un sauvage de l’Amérique} (Amsterdam, Veuve Boetman, 1728).}}.}
              \teiP[xmlid={p18-7}]{On s’en tiendra à un seul élément récurrent du récit de Frézier : la critique de la domination des indigènes américains. La \teiHi[rend={italic}]{Brevísima relación de la destrucción de las Indias} de Bartolomé de Las Casas, traduite dès la fin du \teiHi[rend={small-caps}]{xvi}\teiHi[rend={sup}]{e} siècle en français et fréquemment rééditée, a forgé la « légende noire » des exactions espagnoles pendant la conquête des Amériques\teiNote[xmlid={ftn54-7},n={25},place={foot}]{\teiP[xmlid={p-434}]{Bartolomé de Las Casas, \teiHi[rend={italic}]{Tyrannies et cruautez des Espagnols commises ès Indes occidentales qu’on dit le Nouveau monde…}, Jacques de Miggrode (trad.), Rouen, Jacques Cailloué, 1630. Sur la légende noire, voir Jean-Paul Duviols, \teiHi[rend={italic}]{Le Miroir du Nouveau Monde : images primitives}, Paris, Presses de l’université Paris-Sorbonne, 2006.}}. La déploration du sort des « Indiens » est un lieu commun du discours sur l’Amérique ibérique, souvent associé à des interrogations profondes quant aux mœurs et au tempérament des indigènes. Frézier adopte dans son récit un point de vue plus ou moins empathique envers les indigènes, qui lui permet de soutenir une mise en scène de la révélation des mœurs et de la critique du gouvernement espagnol au Pérou et au Chili. Il ne va jamais tout à fait jusqu’à prendre le point de vue de l’autochtone pour le faire parler, comme l’avait fait plus d’un siècle avant lui Jean de Léry, et comme on le trouvera dans des dispositifs fictionnels plus tardivement, chez Diderot ou Françoise de Graffigny\teiNote[xmlid={ftn55-7},n={26},place={foot}]{\teiP[xmlid={p-435}]{Voir Marc-André Bernier, « Amérique », Bronislaw Baczko, Michel Porret et François Rosset (dir.), \teiHi[rend={italic}]{Dictionnaire critique de l’utopie au temps des Lumières}, Genève, Georg éditeur, 2016, p. 30.}}. Il se focalise cependant assez fréquemment sur leurs malheurs, ces derniers servant de révélateurs à la tyrannie des Espagnols et à la corruption morale de leur gouvernement.}
              \teiP[xmlid={p19-7}]{Un lieu permet à Frézier de développer ce point de vue : il s’agit de la mine. Il en visite quelques-unes lors de son passage au Pérou. Espace du secret par excellence, elle l’est d’abord par sa nature même : enfouie dans les profondeurs de la terre, elle doit, pour révéler ses richesses, être découverte. Ensuite, les Espagnols en Amérique sont en grande partie tributaires de la connaissance qu’ont les indigènes du territoire pour découvrir ces mines, et ces derniers sont rarement enclins à dévoiler leur secret. Enfin, elle est un lieu propice aux fantasmes et aux histoires surnaturelles : récits d’apparitions du Diable, enlèvements, spectres, malédictions, sont monnaie courante. Frézier fait du secret des mines un secret résolument politique, dans la mesure où elle sert de révélateur aux relations de domination qui traversent la société coloniale : si les peuples indigènes conservent le secret de l’emplacement des mines, c’est parce que ce lieu est leur tombeau. Révéler le secret des mines ne consiste pas seulement à faire émerger leur localisation ou à comprendre les mystères scientifiques de la formation du minerai, car le principal secret que renferment les mines est celui de la mort annoncée des indigènes asservis pour leur exploitation.}
              \teiP[xmlid={p20-7}]{Frézier a pour les mines un intérêt d’abord économique. La fascination pour l’or et l’argent des Amériques est encore vive en ce début du \teiHi[rend={small-caps}]{xviii}\teiHi[rend={sup}]{e} siècle. Sa description des richesses économiques ne peut donc pas passer sous silence cette manne financière propre aux Espagnols. Outre la description technique et économique, s’attachant à la qualité des sols et aux manières d’exploiter les mines, Frézier saisit l’occasion d’écrire sur les mines pour évoquer la domination des indigènes par les Espagnols. Un \teiHi[rend={italic}]{leitmotiv} de ce discours est la présence, enfoui dans les profondeurs de la terre, d’un trésor caché, qui n’aurait pas encore été révélé aux Espagnols. La raison principale de ce secret généralisé est à trouver dans la conquête même de l’Amérique et l’asservissement des peuples indigènes :}
              \begin{teiCit}[xmlid={cit06-6}]

                \begin{teiQuote}
La haine implacable des Indiens que cette barbare conduite a attiré aux Espagnols, est cause que les tresors cachez, \& les riches mines dont ils se communiquent entre eux la connoissance, demeurent inconnues \& inutiles aux uns \& aux autres ; car les Indiens ne s’en servent pas même pour eux, contens de vivre de leur travail pauvrement, \& dans la dernière misère\teiNote[xmlid={ftn56-7},n={27},place={foot}]{\teiP[xmlid={p-436}]{Amédée Frézier, \teiHi[rend={italic}]{Relation du voyage…}, \teiHi[rend={italic}]{op. cit.}, t. II, p. 474.}}.
\end{teiQuote}
              
\end{teiCit}
              \teiP[xmlid={p21-7}]{La frugalité heureuse des indigènes est un \teiHi[rend={italic}]{topos} de la littérature de voyage aux Amériques, qui ne cesse de fasciner les voyageurs. Elle s’oppose à l’appât du gain qui caractérise les Espagnols et à leur refus de tout travail manuel. Le secret des indigènes se lirait presque comme une sorte de résistance passive au colonisateur. « Il est certain au reste, que les Indiens connoissent plusieurs mines riches qu’ils ne veulent pas déclarer, de peur qu’on ne les y fasse travailler, \& afin que les Espagnols n’en profitent pas\teiNote[xmlid={ftn57-7},n={28},place={foot}]{\teiP[xmlid={p-437}]{\teiHi[rend={italic}]{Ibid.}, p. 475.}}. » Pour les Espagnols, ce secret bien gardé est nimbé d’une sorte de magie, et Frézier se moque volontiers d’eux, qui « croyent qu’ils \lbrack{}les Indiens\rbrack{} les enchantent, \& font plusieurs contes de morts surprenantes arrivées à ceux qui leur en ont voulu découvrir quelques-unes\teiNote[xmlid={ftn58-7},n={29},place={foot}]{\teiP[xmlid={p-438}]{\teiHi[rend={italic}]{Ibid.}, p. 474-475.}} ». À l’aveuglement superstitieux des Espagnols, le voyageur oppose une lecture démystifiée et résolument politique de ce secret.}
              \teiP[xmlid={p22-6}]{La description apparemment technique du fonctionnement de la mine fait émerger un registre pathétique à l’évocation du sort des indigènes. Décrivant la formation continuelle de l’argent dans les mines, Frézier fait surgir une image saisissante qui rappelle la réalité de l’exploitation des mines, transformant cet espace en tombeau des indigènes et en symbole de leur domination et extermination :}
              \begin{teiCit}[xmlid={cit07-4}]

                \begin{teiQuote}
L’expérience a prouvé cette opinion dans la montagne du Potosi, où l’on a tant creusé en différents endroits, que plusieurs mines ont abîmé \& enseveli les Indiens qui y travailloient avec leurs outils \& étançons. Dans la suite des temps on est venu refouiller les mêmes mines, \& l’on a trouvé dans le bois, dans les crânes, \& dans les os, des filets d’argent qui les penetroient comme la veine même\teiNote[xmlid={ftn59-7},n={30},place={foot}]{\teiP[xmlid={p-439}]{\teiHi[rend={italic}]{Ibid.}, t. I, p. 281.}}.
\end{teiQuote}
              
\end{teiCit}
              \teiP[xmlid={p23-5}]{À travers cette image de la mine, associant l’observation d’un phénomène naturel à un spectacle effroyable, Frézier produit une lecture résolument politique du monde américain. Sans en faire un défenseur des indigènes opprimés ou un critique de la colonisation en soi, il faut retenir que Frézier ne se limite pas à l’observation des fortifications. L’espionnage dépasse sa fonction purement militaire pour faire de l’espion un observateur capable de montrer tous les aspects des lieux traversés. Il en ressort que la faiblesse des Espagnols n’est pas seulement matérielle, elle est surtout morale : le voyageur dessine en creux l’image d’un monde \teiHi[rend={italic}]{miné} de l’intérieur par l’asservissement des indigènes.}
              \teiP[xmlid={p24-5}]{L’ouvrage de Frézier présente un rapport plus ambigu qu’il n’y paraît à la fonction visuelle du récit de voyage. La description et la transmission, textuelle ou iconique, des choses vues lors du voyage dit finalement peu de nouveautés sur les lieux visités et leurs secrets. La vision est un préalable nécessaire, mais elle ne saurait être le tout du récit de voyage. Le voyageur, espion ou naturaliste, doit apprendre à voir et restituer le savoir acquis au cours du voyage, seule véritable voie d’accès aux secrets de l’Amérique, qu’on ne peut déceler en se bornant aux apparences. À ce titre, ce voyage permet de concevoir le secret de l’Amérique du Sud non seulement comme relevant du domaine du visible, mais surtout de l’intelligible et du sensible. Le savoir que rapporte Frézier mériterait d’être interrogé, dans la mesure où, dans le cas des mines, il relève aussi d’un lieu commun. Le regard du voyageur est complexe, à la fois aguerri par l’expérience, formé par la pratique scientifique et la familiarité de ces nouveaux lieux ; il demeure cependant informé par des représentations qui le précèdent et peuvent être partagées par des individus n’ayant jamais voyagé. Le secret de l’asservissement des indigènes en est-il vraiment un, de même que la médiocrité des fortifications espagnoles ? À la lecture de ce voyage, il serait tentant de dire qu’il n’existe pas de secret \teiHi[rend={italic}]{en soi}, mais des secrets construits par le texte : le récit de voyage joue sur la tension entre la découverte, le nouveau, le déjà connu et le lieu commun. C’est dans l’écriture elle-même qu’émerge l’idée ou l’image du secret, dans la mesure où elle cherche à révéler ou dévoiler les objets décrits pour les rendre visibles et sensibles au public, et cela s’effectue avec davantage de force en suggérant que cette révélation émerge d’un secret préalable.}
            
\end{teiDiv}
          
\end{teiElement}
        
\end{teiElement}
        \begin{teiElement}[name={group},type={chapter},data-page-title={L’œil du pouvoir},data-page-subtitle={Les lieux de Richelieu selon Georges de Scudéry},data-page-authors={Maxime Cartron@@FNRS, UCLouvain, GEMCA, Ottignies-Louvain-la-Neuve, Belgique},data-include-href={../XML/chap\_11\_L\_oeil\_du\_pouvoir\_Locus\_politicus.xml},xmlid={chapter-010}]

          \begin{teiElement}[name={front}]

            \begin{teiDiv}[type={titlePage},xmlid={titlepage-011}]

              \teiP[rend={title-main},xmlid={p-440}]{L’œil du pouvoir}
              \teiP[rend={title-sub},xmlid={p-441}]{Les lieux de Richelieu selon Georges de Scudéry}
              \teiP[rend={author-aut},xmlid={p-442}]{Maxime \teiHi[rend={small-caps}]{Cartron}}
              \teiP[rend={authority\_affiliation},xmlid={p-443}]{FNRS, UCLouvain, GEMCA, Ottignies-Louvain-la-Neuve, Belgique}
            
\end{teiDiv}
          
\end{teiElement}
          \begin{teiElement}[name={body},xmlid={body-011}]

            \teiP[xmlid={p1-11}]{Créature autoproclamée de Richelieu\teiNote[xmlid={ftn1-11},n={1},place={foot}]{\teiP[xmlid={p-444}]{Voir, sur cette terminologie, Orest Ranum, \teiHi[rend={italic}]{Les Créatures de Richelieu. Secrétaires d’État et Surintendants des Finances (1635-1642)}, Simonne Guenée (trad.), Roland Mousnier (préf.), Paris, A. Pedone, « Bibliothèque de la revue d’histoire diplomatique », 1966. Il me semble que l’on peut étendre cette appellation à Scudéry, qui a tâché de placer sa carrière sous l’égide du cardinal jusqu’après sa mort. La démarche de Hugh Gaston Hall est similaire dans son \teiHi[rend={italic}]{Richelieu’s Desmarets and the Century of Louis XIV} (Oxford, Clarendon Press, 1990), à la réserve que l’auteur de \teiHi[rend={italic}]{Clovis ou la France chrétienne} était véritablement proche de Richelieu, tandis que Scudéry fait de sa supposée familiarité avec le ministre d’État une fiction d’accréditation.}}, le poète Georges de Scudéry (1601-1667) est notamment connu pour ses éloges clientélistes de la figure du ministre d’État, en particulier après la mort de ce dernier. Le dénominateur commun des textes dans lesquels il dispense ses \teiHi[rend={italic}]{laudationes }est l’insistance sur le regard, motif structurant investi à des fins propagandistes et théocratiques : en (d)écrivant le regard de Richelieu et ses effets, Scudéry se pose en défenseur et en exégète privilégié de son legs politique, opération par laquelle il prétend dévoiler ses arcanes. Or cette stratégie d’écriture, qui se présente comme un geste de divulgation du secret attaché à l’action de Richelieu, est étroitement liée à la question des lieux, puisque le poète s’ingénie systématiquement à inscrire le regard de Richelieu dans une géographie tantôt mythique et topique, voire universelle (le monde est le terrain d’action du Cardinal, qui rayonne tel un soleil), tantôt biographique (le lieu natal). On ne saurait réduire ce schème d’écriture à ce que Richard Sayce désigne comme « la fâcheuse habitude qu’a Scudéry de se répéter continuellement », laquelle, il est vrai, « présente au moins l’avantage de montrer ce qui lui tient véritablement à cœur\teiNote[xmlid={ftn51-7},n={2},place={foot}]{\teiP[xmlid={p-445}]{Richard Sayce, « L’architecture dans l’\teiHi[rend={italic}]{Alaric} de Scudéry », \teiHi[rend={italic}]{Mélanges d’histoire littéraire (}\teiHi[rend={ small-caps italic}]{xvi}\teiHi[rend={italic sup}]{e}\teiHi[rend={italic}]{-}\teiHi[rend={ small-caps italic}]{xvii}\teiHi[rend={italic sup}]{e}\teiHi[rend={italic}]{ siècle) offerts à Raymond Lebègue}, Paris, Nizet, 1969, p. 190.}} ». Loin de constituer une maladresse d’expression, cette insistance permanente relève plutôt de ce que Françoise Graziani appelle le « mythe de l’Œil », puisque selon elle, dans la littérature baroque, « tout s’articule autour de l’œil, en tant que figure analogique céleste, et du regard\teiNote[xmlid={ftn52-7},n={3},place={foot}]{\teiP[xmlid={p-446}]{Françoise Graziani, \teiHi[rend={italic}]{L’Œil et le simulacre. Conscience et formes de la représentation dans « L’Adone » de G. B. Marino et dans la poésie baroque en Espagne, en Italie et en France}, thèse de troisième cycle, Université de Provence, 1980, p. 446.}} ». Il s’agit donc ici d’éclairer cette accommodation de l’œil du pouvoir à son environnement immédiat, en mettant en évidence les déterminations herméneutiques de l’association de l’espace et du regard politiques. Une telle démarche permettra d’expliciter les enjeux de ce qui constitue, par-delà la dimension mythographique, une véritable phénoménologie\teiNote[xmlid={ftn53-8},n={4},place={foot}]{\teiP[xmlid={p-447}]{J’entends le terme dans son sens classique, de recherche des phénomènes perceptifs originaux à partir d’une expérience subjectivante. On verra que cette démarche, pour anachronique qu’elle paraisse, jette une lumière inédite sur l’écriture de la propagande politique pratiquée par Scudéry.}} du politique dévoilant son essence secrète. On rappellera à cet égard avec Louis Marin que le « geste primitif du secret » relève de « \teiHi[rend={italic}]{ce qui} est mis à part\teiNote[xmlid={ftn54-8},n={5},place={foot}]{\teiP[xmlid={p-448}]{Louis Marin, « Logiques du secret », \teiHi[rend={italic}]{Traverses}, n\teiHi[rend={sup}]{o} 30-31, 1984, p. 60.}} ». En d’autres termes, par son insistance sur les lieux de Richelieu, Scudéry prétend révéler la singularité totale du ministre d’État et instruire par là son lecteur de ce qu’il ne peut apercevoir, ébloui qu’il est de sa splendeur.}
            \teiP[xmlid={p2-11}]{Cette question des espaces du Cardinal a déjà suscité des travaux, notamment sur Desmarets de Saint-Sorlin et ses \teiHi[rend={italic}]{Promenades de Richelieu }de 1653\teiNote[xmlid={ftn55-8},n={6},place={foot}]{\teiP[xmlid={p-449}]{Jean-Vincent Blanchard, « \teiHi[rend={italic}]{Les Promenades de Richelieu} de Jean Desmarets et la modernité du théâtre à machines », \teiHi[rend={italic}]{Dix-septième siècle}, n\teiHi[rend={sup}]{o} 280, 2018, p. 461-474.}}. Mais là où l’exploration des lieux du politique s’impose pour ainsi dire d’elle-même chez Desmarets au vu du sujet et de l’organisation de son œuvre, chez Scudéry ces idéologèmes sont disséminés dans l’ensemble de l’œuvre. Il importe par conséquent, en premier lieu, de les identifier et de prendre en compte leur inscription éditoriale, dont les variations influent sur le sens à donner au geste d’écriture. Il est en effet significatif que Scudéry ait réuni dans son volume de \teiHi[rend={italic}]{Poésies diverses }plusieurs pièces initialement publiées en plaquette\teiNote[xmlid={ftn56-8},n={7},place={foot}]{\teiP[xmlid={p-450}]{Seront abordés ici plus particulièrement « Le Temple » (« Autres œuvres », \teiHi[rend={italic}]{La Mort de César, tragédie}, Paris, Augustin Courbé, 1636, mais publication initiale en plaquette : \teiHi[rend={italic}]{Le Temple. Poème à la gloire du Roi et de Monseigneur le Cardinal duc de Richelieu}, Paris, François Targa, 1633. Texte également présent dans le \teiHi[rend={italic}]{Sacrifice des Muses au grand Cardinal de Richelieu}, Paris, Cramoisy, 1635) ; « L’Ombre du grand Armand » (\teiHi[rend={italic}]{Poésies diverses}, 1649, mais publication initiale en plaquette : \teiHi[rend={italic}]{L’Ombre du grand Armand}, Paris, Pierre Anguerrant, 1643) ; « Le Portrait du Grand Cardinal. Fait par Champagne » (\teiHi[rend={italic}]{Le Cabinet de Monsieur de Scudéry}, Paris, Augustin Courbé, 1646) ; « Le Grand Exemple » (\teiHi[rend={italic}]{Poésies diverses}, Paris, Augustin Courbé, 1649, mais publication initiale en plaquette : \teiHi[rend={italic}]{Le Grand exemple, à Mgr le duc de Richelieu}, s. l., 1647) ; « Les Muses. À Monsieur l’Abbé de Richelieu » (\teiHi[rend={italic}]{Poésies diverses dédiées à Monseigneur le Duc de Richelieu}, Paris, Augustin Courbé, 1649).}} : il rassemble de la sorte tous les lieux marqués par la présence de Richelieu, afin de lui ériger un monument imaginaire voué à acquérir une pérennité d’ordre dynastique, puisque le petit-neveu du Cardinal est invité par la dédicace et par plusieurs pièces du volume, que cet agencement éditorial fait résonner entre elles, à imiter son illustre oncle.}
            \teiP[xmlid={p3-11}]{On constate également que plusieurs de ces pièces poétiques ont été initialement publiées à la suite de pièces de théâtre ; c’est le cas en particulier des textes parus dans \teiHi[rend={italic}]{La Mort de César}, qui constitue, avant les \teiHi[rend={italic}]{Poésies diverses}, un premier massif à la gloire de Richelieu, comme l’explique Dominique Moncond’huy :}
            \begin{teiCit}[xmlid={cit01-11}]

              \begin{teiQuote}
En marge de la pièce, certes, les \teiHi[rend={italic}]{Autres Œuvres} font plus que la compléter : elles participent d’une stratégie encomiastique. Sur les arceaux qu’elles composent avec les pièces liminaires s’arqueboute le sens du volume, vaste édifice érigé en l’honneur de Richelieu\teiNote[xmlid={ftn57-8},n={8},place={foot}]{\teiP[xmlid={p-451}]{Georges de Scudéry, \teiHi[rend={italic}]{La Mort de César}, Dominique Moncond’huy (éd.), Paris, STFM, 1992, p. 274-275. Dominique Moncond’huy note de plus qu’« aucun des précédents recueils publiés à la suite des pièces de Scudéry ne consacrait une telle place à l’éloge de caractère politique » (\teiHi[rend={italic}]{ibid.}, p. 273).}}.
\end{teiQuote}
            
\end{teiCit}
            \teiP[xmlid={p4-11}]{Dès lors, s’interroger sur la question des lieux du politique revient à explorer le rapport consubstantiel entre l’écriture et son cadrage éditorial, soit la cohérence travaillée d’un système de représentation investissant l’espace imaginaire et l’espace du livre afin de les charger de toute l’intensité de la réflexion propagandiste.}
            \begin{teiDiv}[type={section1},xmlid={div01-10}]

              \teiHead[xmlid={ladresse-a-richelieu-et-a-ses-descendants-cadrage-incitatif-des-lieux-oculaires-001}]{L’adresse à Richelieu et à ses descendants : cadrage incitatif des lieux oculaires}
              \teiP[xmlid={p5-11}]{L’écriture du regard scudéryenne cherche à forcer l’adhésion politique\teiNote[xmlid={ftn58-8},n={9},place={foot}]{\teiP[xmlid={p-452}]{« Si faire autorité c’est affecter à grande échelle, alors il n’est pas d’enjeu plus central que de faire passer les idées pures impuissantes à l’état d’\teiHi[rend={italic}]{adhæsiones}, c’est-à-dire d’idées-affects, par-là dotées d’une force mobilisatrice – littéralement parlant : une force de mise en mouvement des corps » (Frédéric Lordon, \teiHi[rend={italic}]{La Société des affects. Pour un structuralisme des passions}, Paris, Éditions du Seuil, « L’Ordre philosophique », 2013, p. 201).}} en faisant de l’image des lieux de Richelieu une « \teiHi[rend={italic}]{instance d’intervention}\teiNote[xmlid={ftn59-8},n={10},place={foot}]{\teiP[xmlid={p-453}]{Karl Sierek, \teiHi[rend={italic}]{Images oiseaux. Aby Warburg et la théorie des médias}, Pierre Rusch (trad.), Paris, Klincksieck, 2009, p. 15.}} », de sorte que « l’image intervient désormais dans l’espace culturel comme une force agissante\teiNote[xmlid={ftn60-7},n={11},place={foot}]{\teiP[xmlid={p-454}]{\teiHi[rend={italic}]{Ibid.}, p. 43.}} ». Tel est le sens de l’appel savamment orchestré au Cardinal, depuis la demeure éternelle qu’il occupe :}
              \begin{teiLg}

                \teiL{O trois fois Grand Armand, qui regnes dans les Cieux,}
                \teiL{Des choses de là haut, détourne un peu les yeux ;}
                \teiL{Et sans troubler ta paix, considere la guerre :}
                \teiL{Vois ce jeune Heros, paroistre au premier rang,}
                \teiL{Et digne de ton Nom ; et digne de ton Sang ;}
                \teiL{Faisant voller sa gloire, aux deux bouts de la Terre\teiNote[xmlid={ftn61-7},n={12},place={foot}]{\teiP[xmlid={p-455}]{Georges de Scudéry, \teiHi[rend={italic}]{Poésies diverses I}, \teiHi[rend={italic}]{op. cit.}, « Sonnet. Au mesme (NDLR : « M\teiHi[rend={sup}]{r} le Duc de Richelieu »), sur son Combat naval », p. 104, v. 9-14.}}.}
              
\end{teiLg}
              \teiP[xmlid={p6-11}]{Le regard de Richelieu est appelé à \teiHi[rend={italic}]{reconnaître} sa descendance et à authentifier sa valeur afin de lui passer le témoin politique et de marquer la persistance de son élection divine : terre et ciel sont en parfaite adéquation. Il n’est en rien innocent que ce dispositif oculaire ait son exact pendant dans « Le Grand Exemple », la pièce la plus didactique du volume des \teiHi[rend={italic}]{Poésies diverses}, qui s’ouvre comme suit :}
              \begin{teiLg}

                \teiL{Illustre et jeune Aiglon, leve, leve les yeux ;}
                \teiL{Regarde fixement un bel Astre des Cieux ;}
                \teiL{Fais toy voir à la Terre en ta force premiere,}
                \teiL{Digne de ses rayons, digne de sa lumiere :}
                \teiL{Vois le donc sans ciller, et conçois aujourd’huy,}
                \teiL{Le genereux dessein d’estre semblable à luy.}
                \teiL{(S’il est permis pourtant, à la vertu mortelle,}
                \teiL{D’arriver jusqu’au point où sa gloire t’apelle ;)}
                \teiL{Car marcher sur les pas du divin RICHELIEU,}
                \teiL{C’est suivre pour le moins Hercule Demy-Dieu.}
                \teiL{Ton Pere sur la Mer, ton Ayeul sur la Terre,}
                \teiL{Ont aquis tout l’honneur qu’on aquiert à la guerre ;}
                \teiL{Mais devant le Soleil aucun Astre ne luit,}
                \teiL{Et c’est au plus parfait que mon Art te conduit.}
                \teiL{Observe ce Soleil, que l’Univers contemple ;}
                \teiL{Propose à ton esprit ce rare et GRAND EXEMPLE ;}
                \teiL{Tu prens le plus Grand Nom, que l’on porte icy bas,}
                \teiL{Mais il le faut remplir, où ne le porter pas\teiNote[xmlid={ftn62-7},n={13},place={foot}]{\teiP[xmlid={p-456}]{Georges de Scudéry, \teiHi[rend={italic}]{Poésies diverses II}, \teiHi[rend={italic}]{op. cit.}, « Le Grand Exemple. À Monseigneur le Duc de Richelieu », p. 96, v. 1-18.}}.}
              
\end{teiLg}
              \teiP[xmlid={p7-11}]{Le double mouvement descendant-ascendant est assorti d’une description que l’on peut pleinement qualifier de phénoménologique : c’est l’origine pure de la valeur politique qui est exposée ici, renforcée par l’assimilation de Richelieu au Soleil comme métaphore de l’originaire et, de ce fait, synecdoque de l’univers. Scudéry procède au dévoilement éclatant de ce qui était jusque-là caché aux yeux de tous, cette exposition solaire faisant surgir la charge d’éblouissement qui se trouvait auparavant dans un état de latence caractéristique du secret : le lieu révélé dessine en creux une dialectique du clair-obscur, ce qui était dans l’obscurité devenant visible de manière aveuglante par le biais d’une exposition ostentatoire.}
            
\end{teiDiv}
            \begin{teiDiv}[type={section1},xmlid={div02-10}]

              \teiHead[xmlid={nekuia-et-convocation-le-portrait-du-grand-cardinal-et-lombre-du-grand-armand-001}]{\teiHi[rend={italic}]{Nekuia} et convocation : \teiHi[rend={italic}]{Le Portrait du Grand Cardinal} et \teiHi[rend={italic}]{L’Ombre du grand Armand}}
              \teiP[xmlid={p8-11}]{Ceci nous conduit tout naturellement à interroger l’effet de présence recherché par Scudéry, qui évoque constamment les lieux de mémoire du politique absolu, en un processus relevant de la \teiHi[rend={italic}]{nekuia}, c’est-à-dire de la convocation des morts\teiNote[xmlid={ftn63-7},n={14},place={foot}]{\teiP[xmlid={p-457}]{Sur ce phénomène sous l’Ancien Régime, on se reportera à Neil Kenny, \teiHi[rend={italic}]{Death and tenses: posthumous presence in early modern France}, Oxford, Oxford University Press, 2015.}} : l’écriture du regard incarne la force originaire de Richelieu dans des lieux symboliques clés de l’exercice de son pouvoir\teiNote[xmlid={ftn64-6},n={15},place={foot}]{\teiP[xmlid={p-458}]{On se souvient à ce titre de cette analyse capitale de Louis Marin : « représenter \lbrack{}…\rbrack{}, c’est faire revenir le mort comme s’il était présent et vivant, et c’est aussi redoubler le présent et intensifier la présence dans l’institution d’un sujet de représentation. Comment donc la représentation est-elle accomplissement du désir d’absolu qui anime l’essence de tout pouvoir, sinon en étant le substitut imaginaire de cet accomplissement, sinon en étant son image ? Le portrait du roi que le roi contemple lui offre l’icône du monarque absolu qu’il désire être au point de se reconnaitre et de s’identifier par lui et en lui au moment même où le référent du portrait s’en absente. Le roi n’est vraiment roi, c’est-à-dire monarque, que dans des images. Elles sont sa \teiHi[rend={italic}]{présence réelle} : une croyance dans l’efficacité et l’opérativité de \teiHi[rend={italic}]{ses }signes iconiques est obligatoire, sinon le monarque se vide de toute sa substance par défaut de transsubstantiation et il n’en reste plus que le simulacre ; mais, à l’inverse, parce que \teiHi[rend={italic}]{ses} signes \teiHi[rend={italic}]{sont }la \teiHi[rend={italic}]{réalité} royale, l’être et la substance du prince, cette croyance est nécessairement exigée par les signes eux-mêmes ; son défaut est à la fois hérésie et sacrilège, erreur et crime » (Louis Marin, \teiHi[rend={italic}]{Le Portrait du roi}, Paris, Éditions de Minuit, « Le sens commun », 1981, p. 12-13).}}. Tel est le sens de l’ouverture du \teiHi[rend={italic}]{Portrait du Grand Cardinal. Fait par Champagne }:}
              \begin{teiLg}

                \teiL{Prends tes pinceaux, docte Champagne,}
                \teiL{Choisis tes plus belles couleurs,}
                \teiL{Et pour soulager mes douleurs,}
                \teiL{Viens faire encor trembler l’Espagne.}
                \teiL{Rappelle dans ton souvenir,}
                \teiL{D’un héros qu’on a vu finir}
                \teiL{La grande et merveilleuse Idée :}
                \teiL{Revois cet objet glorieux}
                \teiL{Dont mon âme était possédée,}
                \teiL{Et le fais revoir à mes yeux\teiNote[xmlid={ftn65-6},n={16},place={foot}]{\teiP[xmlid={p-459}]{Georges de Scudéry, \teiHi[rend={italic}]{Le Cabinet de Monsieur de Scudéry} \lbrack{}1646\rbrack{}, Christian Biet et Dominique Moncond’huy (éd.), Paris, Klincksieck, « Théorie et critique à l’âge classique », 1991, p. 130, v. 1-10.}}.}
              
\end{teiLg}
              \teiP[xmlid={p9-11}]{Cette invocation, fondée sur l’itératif, charge l’espace, avant même sa mise en place, d’une présence rémanente – « le secret est l’effet présent dans le présent d’un état passé », selon les mots de Marin\teiNote[xmlid={ftn66-6},n={17},place={foot}]{\teiP[xmlid={p-460}]{Louis Marin, « Logiques du secret », art. cité, p. 60.}} – qui va l’irradier et lui conférer une force de représentation visant à frapper l’œil du spectateur-lecteur :}
              \begin{teiLg}

                \teiL{D’une ordonnance ingénieuse,}
                \teiL{Imagine bien ce tableau :}
                \teiL{Fais-y paraître un grand rideau,}
                \teiL{Dont l’étoffe soit précieuse\teiNote[xmlid={ftn67-6},n={18},place={foot}]{\teiP[xmlid={p-461}]{\teiHi[rend={italic}]{Ibid.}, p. 130, v. 11-24.}}.}
              
\end{teiLg}
              \teiP[xmlid={p10-11}]{Le rideau matérialise l’incarnation progressive de Richelieu au cœur du lieu. Annonçant sa venue, il marque la frontière entre la dissimulation et la révélation du secret. Lieu qui n’en est pas un à proprement parler, il constitue un accessoire privilégié de la représentation du secret, puisqu’il métaphorise l’entre-deux de la révélation, en apportant à l’image cette temporalité unique, à la fois présente et passée. Simultanément espace et moment, il peut par conséquent être appréhendé comme une métonymie de l’apparition, laquelle est entée sur l’endroit de l’un des plus grands triomphes du Cardinal\teiNote[xmlid={ftn68-3},n={19},place={foot}]{\teiP[xmlid={p-462}]{« Il ne s’agit pas seulement de décrire physiquement l’homme d’État, mais de dire comment il doit être mis en scène dans l’espace de la représentation, quelle doit y être sa place, quelle doit être sa posture (une main appuyée sur une balustrade, une autre désignant « un rempart foudroyé » ; ses ennemis, symboliquement représentés par un aigle, des lions et des léopards, à ses pieds ; ses yeux « d’une bonté lumineuse », son air noble…), ce que l’on doit voir à l’arrière-plan (La Rochelle)… ; tout doit faire sens pour susciter une interprétation politique et Scudéry n’oublie rien, imaginant un tableau politiquement idéal, c’est-à-dire reprenant tout le dispositif de la représentation politique mis en œuvre dans différents tableaux réellement réalisés par divers artistes, notamment par le même Champaigne » (Dominique Moncond’huy, « Le poète commande au peintre : enjeux et effets d’un modèle poétique (de Ronsard à Scudéry) », Pascaline Mourier\teiHi[rend={small-caps}]{-}Casile et Dominique Moncond’huy (dir.), \teiHi[rend={italic}]{Lisible/visible : problématiques}, Poitiers, La Licorne, 1992, p. 23).}} :}
              \begin{teiLg}

                \teiL{Que sous ce rideau retroussé,}
                \teiL{Dont le velours sera pressé,}
                \teiL{Paraisse une large fenêtre,}
                \teiL{D’où le jour tombe doucement,}
                \teiL{Et qui serve à faire paraître}
                \teiL{La Rochelle en éloignement.}
                \teiL{Sur le devant de ta peinture,}
                \teiL{Par les traits hardis de ta main,}
                \teiL{Pour achever ce beau dessein,}
                \teiL{Mets la principale figure.}
                \teiL{Peintre fameux, c’est en ce lieu}
                \teiL{Qu’il faut que le grand RICHELIEU}
                \teiL{Comme un astre brillant éclate :}
                \teiL{Et je prétends à cette fois}
                \teiL{Que le lustre de l’écarlate}
                \teiL{Efface la pourpre des rois\teiNote[xmlid={ftn69-3},n={20},place={foot}]{\teiP[xmlid={p-463}]{Georges de Scudéry, Le Portrait du Grand Cardinal, \teiHi[rend={italic}]{op. cit.}, p. 130-131, v. 25-50.}}.}
              
\end{teiLg}
              \teiP[xmlid={p11-11}]{La peinture de Champaigne est appréhendée, en termes narratifs, comme la construction progressive du lieu par l’inscription du corps de Richelieu ; la commande exige de ménager un dévoilement progressif, à la force de suspens considérable. Le poète peut effectivement s’exclamer à ce moment, feignant de croire à la venue effective de Richelieu pour asseoir l’efficacité de son effet de présence\teiNote[xmlid={ftn70-3},n={21},place={foot}]{\teiP[xmlid={p-464}]{« Ce que Scudéry demande à Champaigne, c’est bien un portrait politique, qui ait une efficacité politique » (Dominique Moncond’huy, « Poésie, peinture et politique : la place de Richelieu dans \teiHi[rend={italic}]{Le Cabinet de M. de Scudéry} », \teiHi[rend={italic}]{Dix-septième siècle}, n\teiHi[rend={sup}]{o} 165, 1989, p. 426).}} :}
              \begin{teiLg}

                \teiL{Ah ! je le vois, c’est lui sans doute,}
                \teiL{Ton Art a suivi mon désir ;}
                \teiL{Et je crois, tant j’ai de plaisir,}
                \teiL{Qu’il me regarde et qu’il m’écoute\teiNote[xmlid={ftn71-3},n={22},place={foot}]{\teiP[xmlid={p-465}]{Georges de Scudéry, « Le Portrait du Grand Cardinal », \teiHi[rend={italic}]{op. cit.}, p. 131-132, v. 61-70.}}.}
              
\end{teiLg}
              \teiP[xmlid={p12-11}]{Amenée subtilement par la conjonction d’un jeu du regard et de l’ouïe, l’intervention omnisciente du Cardinal, explicitement construite en objet de « désir », peut enfin avoir lieu, le cadre de son exercice ayant été posé ; l’œil de Richelieu emplit les lieux dans lesquels il inscrit la puissance de son discernement\teiNote[xmlid={ftn72-3},n={23},place={foot}]{\teiP[xmlid={p-466}]{« Ce visage qui dit tout du personnage : c’est là que réside sa puissance, fruit de sa nature même, de son essence et non de ce qu’il a pu acquérir par le mérite. C’est le regard qui compte plus que tout » (Dominique Moncond’huy, « Poésie, peinture et politique… », art. cité, p. 427).}} :}
              \begin{teiLg}

                \teiL{En vain l’ennemi de la France}
                \teiL{Avait des ministres discrets,}
                \teiL{Car il n’était point de secrets}
                \teiL{À l’épreuve de sa prudence.}
                \teiL{Cet esprit toujours agissant}
                \teiL{Pénétrait d’un regard perçant}
                \teiL{Jusques au fond de leur pensée :}
                \teiL{Il présidait à leur conseil ;}
                \teiL{Et leur trame était renversée,}
                \teiL{Comme ils en formaient l’appareil.}
                \teiL{Oui, toutes les choses futures}
                \teiL{Étaient présentes à ses yeux ;}
                \teiL{Et l’on eût dit que dans les cieux}
                \teiL{Il ordonnait nos aventures ;}
                \teiL{Que les vents respectaient sa voix ;}
                \teiL{Que les flots recevaient ses lois ;}
                \teiL{Que tout avait soin de lui plaire ;}
                \teiL{Que tout redoutait son pouvoir ;}
                \teiL{Et que l’astre qui nous éclaire}
                \teiL{Ne luisait plus que pour le voir\teiNote[xmlid={ftn73-3},n={24},place={foot}]{\teiP[xmlid={p-467}]{Georges de Scudéry, « Le Portrait du Grand Cardinal », \teiHi[rend={italic}]{op. cit.}, p. 133-134, v. 131-150.}}.}
              
\end{teiLg}
              \teiP[xmlid={p13-11}]{Le retournement de Richelieu-Soleil en Richelieu objet de regard du Soleil parachève la dialectique d’incarnation du regard au cœur des lieux du pouvoir : le Cardinal emplit le monde de sa présence omnisciente et surpuissante. Il est notable que celle-ci soit assise sur la « prudence », qui le rend apte à déjouer toute dissimulation par l’intermédiaire de son regard. La \teiHi[rend={italic}]{phronesis} de Richelieu est présentée comme un acte de vision avant tout, en tant qu’elle exprime la connaissance absolue du politique, connaissance qui est par elle-même un \teiHi[rend={italic}]{pouvoir agir}. L’ouverture de \teiHi[rend={italic}]{L’Ombre du Grand Armand}, pièce reprise dans le volume de 1649 mais initialement parue en 1643, est également orientée dans le but de faire ressentir tout l’impact politique de l’apparition du fantôme de Richelieu :}
              \begin{teiLg}

                \teiL{Lors que le plus grand Astre, achevant sa carriere,}
                \teiL{Mesloit la nuit au jour, et l’ombre à la lumiere ;}
                \teiL{Au Temple de Sorbonne, une vive clarté,}
                \teiL{Vint faire un nouveau jour, dans cette obscurité.}
                \teiL{Du creux du grand Tombeau, la clarté jaillissante,}
                \teiL{Imitant du Soleil la lumiere naissante}
                \teiL{Paroist confusément ; s’acroist avec ardeur ;}
                \teiL{S’esleve ; et remplit tout d’éclat et de splendeur.}
                \teiL{Mille rayons dorez, dissipant les ténèbres,}
                \teiL{Y servent d’ornement, aux Ornemens funebres ;}
                \teiL{Et parmi cet objet lugubre, mais charmant ;}
                \teiL{Aparoist à nos yeux, l’OMBRE DU GRAND ARMAND.}
                \teiL{De la Pourpre Romaine il conserve l’usage ;}
                \teiL{La Majesté des Rois esclate en son visage ;}
                \teiL{Et bien qu’il semble triste, il plaist encor aux yeux,}
                \teiL{Plus que le plus beau jour qui nous tombe des Cieux\teiNote[xmlid={ftn74-2},n={25},place={foot}]{\teiP[xmlid={p-468}]{Georges de Scudéry, \teiHi[rend={italic}]{Poésies diverses II}, \teiHi[rend={italic}]{op. cit.}, p. 89, v. 1-16.}}.}
              
\end{teiLg}
              \teiP[xmlid={p14-10}]{L’importance du Temple de Sorbonne réside dans son alternance avec des lieux qui s’opposent dialectiquement à la précision de sa localisation ; le regard de Richelieu « esclate » au-delà du noir secret\teiNote[xmlid={ftn75-2},n={26},place={foot}]{\teiP[xmlid={p-469}]{Entendu ici au sens de « qui est ou doit être caché aux autres ».}} de sa résidence éternelle dans laquelle ses adversaires souhaiteraient le laisser disparaître peu à peu des mémoires, pour rayonner sur l’ensemble de la terre :}
              \begin{teiLg}

                \teiL{J’ay fait trembler l’Europe, et l’Afrique, et l’Asie ;}
                \teiL{J’ay vaincu les mutins ; j’ay dompté l’Heresie ;}
                \teiL{Sous le plus grand des Rois, par mes conseils prudents,}
                \teiL{J’ay surmonté le Sort, et la Mer, et les Vents.}
                \teiL{J’ay fait voir à ses pieds, l’orgueil de la Rochelle ;}
                \teiL{J’ay fait voir à ses pieds, tout un Peuple rebelle ;}
                \teiL{Qui depuis si long temps par un lasche attentat,}
                \teiL{Paroissoit Estranger au milieu de l’Estat\teiNote[xmlid={ftn76-2},n={27},place={foot}]{\teiP[xmlid={p-470}]{Georges de Scudéry, « L’Ombre du grand Armand », \teiHi[rend={italic}]{op. cit.}, p. 91-92, v. 89-96.}}.}
              
\end{teiLg}
              \teiP[xmlid={p15-9}]{La géographie de Richelieu procède, même après sa mort, par expansion, ce mouvement étant métaphorique de la propagation idéologique que Scudéry entend réaliser. La mémoire est vouée à se transformer en présent, le lieu passé à s’actualiser dans la présence d’un regard encore actif, comme le montre la revue des hauts faits et exploits guerriers du ministre d’État, dont la scansion des lieux réactive le choc politique :}
              \begin{teiLg}

                \teiL{Par les mesmes conseils, et dans la mesme guerre,}
                \teiL{J’ay chassé de nos bords, les Armes d’Angleterre ;}
                \teiL{Repoussé leur puissance, et leur ambition,}
                \teiL{Jusqu’aux flots recullez de la Grande Albion :}
                \teiL{Et fait voir à ce Peuple, apres son entreprise,}
                \teiL{Que la Seine aujourd’huy ne craint point la Tamise\teiNote[xmlid={ftn77-2},n={28},place={foot}]{\teiP[xmlid={p-471}]{\teiHi[rend={italic}]{Ibid.}, p. 92, v. 97-110.}}.}
              
\end{teiLg}
              \teiP[xmlid={p16-9}]{Si ce catalogue se poursuit encore longuement\teiNote[xmlid={ftn78-2},n={29},place={foot}]{\teiP[xmlid={p-472}]{« Casal trois fois sauvé par mon prudent conseil, / Esleva mon renom, plus haut que le Soleil : / Et comme j’eus tousjours la vertu pour compagne, / La vertu triompha de l’orgueil de l’Espagne. / Pignerol par mes soings, et par mes grands Exploits, / Ouvre encor l’Italie aux armes des François : / Et ne m’arrestant point dans cette illustre voye, / Des Rochers de Piémont, aux rochers de Savoye, / Je fais passer un Prince aussi Grand que cheri ; / Il attaque ; il emporte ; et puis rend Chamberi. / Des hauts murs de Paris, et des bords de la Seine, / Son Tonnerre va fondre aux terres de Lorraine : / Nanci par mes conseils, cede, et reçoit sa loy ; / Et le Vassal rebelle est aux pieds de son Roy. / Mais si mes grands conseils font tomber des Murailles, / En suite mes conseils font gagner des Batailles : / Et les Plaines d’Avein, feront voir dans cent ans, / Ce que par mes conseils firent nos Combatans » (\teiHi[rend={italic}]{ibid.}, p. 92, v. 111-128).}}, on s’avise surtout que c’est toujours le regard omniscient de Richelieu qui rend possibles ces victoires et cette gouvernance sur l’ensemble de la terre, laquelle constitue par essence le champ d’exercice et d’application de ses compétences hors normes. La conséquence logique de cette scénographie des lieux du pouvoir est l’adhésion sans réserve du public, que Scudéry anticipe avec une certitude absolue :}
              \begin{teiLg}

                \teiL{Mille et mille flambeaux, à lumiere esclatante,}
                \teiL{Environnent le Temple, et la Chapelle ardante :}
                \teiL{Et parmi tant de feux, saintement allumez,}
                \teiL{Par une voix de feu les Peuples sont charmez\teiNote[xmlid={ftn79-2},n={30},place={foot}]{\teiP[xmlid={p-473}]{\teiHi[rend={italic}]{Ibid.}, p. 94, v. 173-188.}}.}
              
\end{teiLg}
              \teiP[xmlid={p17-9}]{L’alliance du catalogue des lieux et de la prosopopée est un stratagème littéraire et politique particulièrement efficace, puisqu’il permet à Scudéry de sortir Richelieu du tombeau dans lequel ses supposés ennemis politiques voudraient le maintenir afin de rappeler et de renforcer sa place éminente dans l’ensemble des conquêtes françaises : pour gouverner le monde, il faut se souvenir de l’apport considérable du Cardinal dans ce grand dessein.}
            
\end{teiDiv}
            \begin{teiDiv}[type={section1},xmlid={div03-9}]

              \teiHead[xmlid={epiphanies-du-visible-les-muse-s-et-le-grand-exemple-001}]{Épiphanies du visible : \teiHi[rend={italic}]{Les Muse}s et \teiHi[rend={italic}]{Le Grand Exemple}}
              \teiP[xmlid={p18-8}]{Avec \teiHi[rend={italic}]{Les Muses}, paru dans les \teiHi[rend={italic}]{Poésies diverses }de 1649, c’est le pays natal du Cardinal, le domaine de Richelieu\teiNote[xmlid={ftn80-2},n={31},place={foot}]{\teiP[xmlid={p-474}]{Sur ce lieu, on se reportera à Christian Jouhaud, « Le territoire richelais du cardinal duc », \teiHi[rend={italic}]{Richelieu à Richelieu. Architecture et décors d’un château disparu}, Milan, Silvana Editoriale, 2011, p. 23-29 ; Christine Toulier, « La capitale d’un duché : jardins, parc, ville et territoire », \teiHi[rend={italic}]{ibid.}, p. 55-61 ; John Schloder, « Les héritiers de Richelieu », \teiHi[rend={italic}]{ibid.}, p. 449-455 ; Marie-Pierre Terrien, \teiHi[rend={italic}]{Richelieu. Histoire d’une cité idéale (1631-2011)}, Rennes, PUR, « Art et Société », 2011.}}, qui est mis en adéquation avec l’hypostase du visible que constitue son regard. On retrouve ici le Richelieu céleste du \teiHi[rend={italic}]{Portrait du grand Cardinal}\teiNote[xmlid={ftn81-2},n={32},place={foot}]{\teiP[xmlid={p-475}]{« Armand le plus Grand des hommes ; / Armand mis entre les Dieux ; / Qui de la Terre où nous sommes, / Porta son Nom dans les Cieux » (Georges de Scudéry, \teiHi[rend={italic}]{Poésies diverses II}, \teiHi[rend={italic}]{op. cit.}, p. 50, v. 11-14).}}, \teiHi[rend={italic}]{via }l’incarnation d’un regard politique absolu au cœur du monde et du sensible, le lieu natal apportant une épaisseur biographique supplémentaire :}
              \begin{teiLg}

                \teiL{Assez près de cette ville,}
                \teiL{Qu’admire tout l’Univers,}
                \teiL{L’on voit un Païs fertile,}
                \teiL{Où les Champs sont tousjours vers.}
                \teiL{L’Art joint avec la Nature,}
                \teiL{D’une eternelle peinture,}
                \teiL{Orne ce lieu si vanté :}
                \teiL{Et jamais le froid Borée,}
                \teiL{N’a pû se donner entrée,}
                \teiL{Dans ce sejour enchanté.}
                \teiL{La Perspective sçavante,}
                \teiL{Y trompe l’oeil curieux :}
                \teiL{Et Rome estant Triomphante,}
                \teiL{N’eut point d’Arc plus glorieux :}
                \teiL{De cent superbes Statuës,}
                \teiL{Toutes richement vestuës,}
                \teiL{Un Palais est decoré :}
                \teiL{Et les Phrises ; les Corniches ;}
                \teiL{Et les Festons ; et les Niches ;}
                \teiL{Vont jusqu’au Lambris doré.}
                \teiL{Les Arbres hauts et superbes,}
                \teiL{Y couvrent de leurs rameaux,}
                \teiL{L’humide fraischeur des herbes,}
                \teiL{Qui les rend, et grands, et beaux.}
                \teiL{Sous l’espace de leurs ombres,}
                \teiL{Des Promenoirs vers et sombres,}
                \teiL{Font un Objet sans pareil :}
                \teiL{Mais un Objet solitaire,}
                \teiL{Qui dans l’heure la plus claire,}
                \teiL{N’est jamais veû du Soleil\teiNote[xmlid={ftn82-2},n={33},place={foot}]{\teiP[xmlid={p-476}]{\teiHi[rend={italic}]{Ibid.}, p. 51, v. 31-34.}}.}
              
\end{teiLg}
              \teiP[xmlid={p19-8}]{L’écriture du \teiHi[rend={italic}]{locus amoenus} et son alliance du naturel et de l’artificiel est parfaitement topique pour l’époque, et elle vise bien entendu à susciter l’émerveillement du potentiel lecteur. Mais cette fantaisie descriptive prend un tout autre sens si on l’observe du point de vue de l’inscription du corps de Richelieu en son sein, ce dernier investissant progressivement l’ensemble du paysage sensible :}
              \begin{teiLg}

                \teiL{Dans ce beau sejour des Graces,}
                \teiL{Est un Promenoir charmant,}
                \teiL{Qui garde encore les traces,}
                \teiL{Des pas du divin Armand.}
                \teiL{Ce fut dans la sombre Allée,}
                \teiL{Du jour la plus recullée,}
                \teiL{Et sous ces feüillages vers ;}
                \teiL{Que cet homme incomparable,}
                \teiL{Prit le dessein memorable,}
                \teiL{D’affranchir tout l’Univers\teiNote[xmlid={ftn83-2},n={34},place={foot}]{\teiP[xmlid={p-477}]{\teiHi[rend={italic}]{Ibid.}, p. 54, v. 141-150.}}.}
              
\end{teiLg}
              \teiP[xmlid={p20-8}]{La mémoire du corps de Richelieu, imprimée profondément dans le lieu, justifie \teiHi[rend={italic}]{a posteriori} son caractère mirifique. Plus encore, elle engage l’instauration d’une dynastie politique dans la mesure où elle incite son descendant à l’imitation, tout en saturant ses affects par la performativité que le poète prête à la prophétie révélatrice suivante :}
              \begin{teiLg}

                \teiL{Ce fut à la mesme place,}
                \teiL{Apres son triste trespas,}
                \teiL{Que pour pleurer sa disgrace,}
                \teiL{Son Neveu porta ses pas.}
                \teiL{Mais à peine d’Erotime,}
                \teiL{La douleur si legitime,}
                \teiL{Poussa ses premiers accens ;}
                \teiL{Qu’un object rempli de gloire,}
                \teiL{Vint dissiper l’humeur noire,}
                \teiL{Qui s’emparoit de ses sens.}
                \teiL{Cent rayons d’or parmi l’ombre,}
                \teiL{Brillerent de toutes parts ;}
                \teiL{Qui dans ce lieu frais et sombre,}
                \teiL{Esbloüirent ses regards.}
                \teiL{Dans cette vive lumiere,}
                \teiL{Pallas fit voir la premiere,}
                \teiL{Les charmes de sa beauté :}
                \teiL{Et dès que’on la vit paroistre,}
                \teiL{L’AEgide la fit connoistre,}
                \teiL{Et plus encore sa fierté.}
                \teiL{Le Dieu que Dellos admire,}
                \teiL{Marche à costé de Pallas ;}
                \teiL{L’Arc, le Carquois, et la Lire,}
                \teiL{Font que l’on n’en doute pas.}
                \teiL{Dans cette rare avanture,}
                \teiL{L’on vit encore Mercure,}
                \teiL{Suivre Apollon en ce lieu :}
                \teiL{Et quand la peur fut passée,}
                \teiL{Ses Aisles, son Caducée,}
                \teiL{Firent connoistre ce Dieu\teiNote[xmlid={ftn84-2},n={35},place={foot}]{\teiP[xmlid={p-478}]{\teiHi[rend={italic}]{Ibid.}, p. 54-55, v. 151-180.}}.}
              
\end{teiLg}
              \teiP[xmlid={p21-8}]{Le clair-obscur qui baigne ce passage métaphorise les états émotionnels extrêmes se succédant, le deuil laissant place à l’enthousiasme devant la révélation de l’élection divine de l’ensemble de la famille de Richelieu\teiNote[xmlid={ftn85-2},n={36},place={foot}]{\teiP[xmlid={p-479}]{C’est ainsi en ces termes qu’Uranie s’adresse plus loin à Erotime : « Si tu veux suivre ma regle, / Je te feray comme une Aigle, / Voir fixement le Soleil : / J’illumineray ton ame, / De son esclat sans pareil » (\teiHi[rend={italic}]{ibid.}, p. 62, v. 385-390).}}. Or ceci n’est rendu possible que par la persistance à l’identique du lieu autobiographique, qui devient celui d’un dévoilement théocratique et d’une initiation aux arcanes du politique que Scudéry, en phénoménologue écartant le voile des apparences et annulant de la sorte la prégnance du secret, rend pour ainsi dire aveuglément visible : Erotime (le nom du petit-neveu dans ce poème), affirment les Dieux dont le poète manipule les prosopopées, « Connoistra tout l’Univers\teiNote[xmlid={ftn86-2},n={37},place={foot}]{\teiP[xmlid={p-480}]{\teiHi[rend={italic}]{Ibid.}, p. 64, v. 470.}} ». Exactement comme son illustre ancêtre, il pourra ainsi « fai(re) trembler l’Europe, et l’Afrique, et l’Asie\teiNote[xmlid={ftn87-2},n={38},place={foot}]{\teiP[xmlid={p-481}]{\teiHi[rend={italic}]{Ibid.}}} ». On s’avise par conséquent de la cohérence du système de représentation scudéryen, qui relie étroitement les textes entre eux, par l’intermédiaire des lieux, lesquels engagent une véritable poétique du discernement, et donc de révélation du secret :}
              \begin{teiLg}

                \teiL{Tel parut aux yeux du Monde,}
                \teiL{Armand l’objet de nos Vers ;}
                \teiL{Sa connoissance profonde,}
                \teiL{Embrassoit tout l’Univers.}
                \teiL{Rien n’échapoit à sa veuë ;}
                \teiL{Et toute cause inconnuë,}
                \teiL{Pour luy seul ne l’estoit pas :}
                \teiL{Bref, par ses intelligences,}
                \teiL{Il fut l’Atlas des Sciences ;}
                \teiL{Aussi bien que des Estats\teiNote[xmlid={ftn88-2},n={39},place={foot}]{\teiP[xmlid={p-482}]{\teiHi[rend={italic}]{Ibid.}, p. 65, v. 481-490.}}.}
              
\end{teiLg}
              \teiP[xmlid={p22-7}]{On retrouve le regard omniscient du ministre d’État et son pouvoir de connaissance et d’action, une dialectique supplémentaire s’ajoutant ici : Richelieu voit tout et est vu de tous, en une forme de chiasme oculaire faisant de son corps, en quelque manière, un lieu de mémoire du politique à part entière. Ceci permet à Scudéry d’amorcer l’horizon final de son dispositif de représentation, qui se trouve être la perpétuation dans le regard de son descendant, laquelle se réalise à travers la fictionnalisation de son adhésion enthousiaste :}
              \begin{teiLg}

                \teiL{Vois, vois ces superbes Temples,}
                \teiL{Ces Villes, et ces Palais,}
                \teiL{Et dis, si tu les contemples,}
                \teiL{C’est Armand qui les a faits :}
                \teiL{Toutes ces marques publiques,}
                \teiL{De ses pensers magnifiques,}
                \teiL{La grandeur font esclater :}
                \teiL{Ces effets font voir leurs Causes,}
                \teiL{Il aimoit les belles choses,}
                \teiL{Et je le veux imiter.}
                \teiL{Erotime à ce langage,}
                \teiL{Tout transporté de plaisir,}
                \teiL{Fait voir sur son beau visage,}
                \teiL{Des marques de son désir.}
                \teiL{Il demande à ces Deesses,}
                \teiL{Les effets de leurs promesses ;}
                \teiL{Il promet d’en bien user :}
                \teiL{Et que son ame hautaine,}
                \teiL{Ne craindra jamais la peine,}
                \teiL{Voulant s’immortaliser\teiNote[xmlid={ftn89-2},n={40},place={foot}]{\teiP[xmlid={p-483}]{\teiHi[rend={italic}]{Ibid.}, p. 73, v. 731-750 (c’est Thalie qui parle).}}.}
              
\end{teiLg}
              \teiP[xmlid={p23-6}]{Un autre regard s’insinue ici : celui du poète, dont l’injonction à l’endroit du successeur de Richelieu lui assure une place de choix au cœur de ce dispositif, celle d’interface ayant à charge de passer le témoin d’un Richelieu à l’autre. On peut ainsi observer un lien étroit entre ces passages oculaires et les dispositifs spatiaux d’adhésion, en particulier dans \teiHi[rend={italic}]{Le Grand Exemple} :}
              \begin{teiLg}

                \teiL{Chacun dans son Palais voulut le retenir :}
                \teiL{Et LA VERTU VISIBLE, admirable en ses charmes,}
                \teiL{Pour retenir ses pas, n’espargna pas ses larmes.}
                \teiL{Il sortit toutefois ; et sans plus l’escouter,}
                \teiL{Ce Lion, genereux ne pût rien redouter.}
                \teiL{Il sortit (je le vy) presques sans nulle suite :}
                \teiL{Et remettant au Ciel le soing de sa conduite,}
                \teiL{Il traversa Paris, d’un visage assuré ;}
                \teiL{Il traversa les flots d’un Peuple conjuré ;}
                \teiL{Et d’un œil tout puissant, et d’un œil tout de flame,}
                \teiL{Il luy remit l’amour ; et le respect en l’ame\teiNote[xmlid={ftn90-2},n={41},place={foot}]{\teiP[xmlid={p-484}]{Georges de Scudéry, \teiHi[rend={italic}]{Poésies diverses II}, \teiHi[rend={italic}]{op. cit.}, p. 101, v. 162-178.}} ;}
              
\end{teiLg}
              \teiP[xmlid={p24-6}]{À la stase dans un lieu protecteur, Scudéry oppose le mouvement de Richelieu à travers un lieu hostile et périlleux, l’action de son regard ayant pour effet immédiat de désarmer la conjuration : la dynamique des lieux et la dynamique oculaire sont en symbiose parfaite, de sorte que le poème s’achève sur des espaces génériques présentés comme signes de l’excellence de la politique du Cardinal :}
              \begin{teiLg}

                \teiL{Vois de cette Vertu mille Marques publiques ;}
                \teiL{Des Villes, des Palais, des Temples magnifiques ;}
                \teiL{Monumens eternels d’un Esprit genereux,}
                \teiL{Qui fit la France illustre, et nostre Siecle heureux\teiNote[xmlid={ftn91-2},n={42},place={foot}]{\teiP[xmlid={p-485}]{\teiHi[rend={italic}]{Ibid.}, p. 104-105, v. 299-302.}}.}
              
\end{teiLg}
              \teiP[xmlid={p25-5}]{La mention de ces lieux produit un effet d’expansion synecdochique – tous les lieux pouvant être désignés par de tels termes sont rassemblés – qui atteste de la grandeur et de la prospérité de la France gouvernée par le regard du ministre d’État.}
            
\end{teiDiv}
            \begin{teiDiv}[type={section1},xmlid={div04-3}]

              \teiHead[xmlid={inscrire-et-eterniser-le-temple-001}]{Inscrire et éterniser : \teiHi[rend={italic}]{Le Temple}}
              \teiP[xmlid={p26-4}]{Cette mention de lieux stables et apaisés conduit à une dialectique de l’inscription et de l’éternisation du regard de Richelieu en leur cœur. C’est \teiHi[rend={italic}]{Le Temple} qui est sollicité ici, en ce qu’il procure en quelque sorte par anticipation, avant les autres textes déjà analysés (on rappellera qu’il est paru initialement en 1633), la synecdoque de toutes les victoires du ministre. Ce poème narre l’arrivée, après un naufrage, d’une troupe française sur « L’Isle de la Valeur, et de la Renommée\teiNote[xmlid={ftn92-2},n={43},place={foot}]{\teiP[xmlid={p-486}]{Georges de Scudéry, \teiHi[rend={italic}]{Poésies diverses II}, \teiHi[rend={italic}]{op. cit.}, p. 227, v. 166.}} » où se trouve le fameux Temple « que bastit Apollon pour un grand demy-Dieu\teiNote[xmlid={ftn93-2},n={44},place={foot}]{\teiP[xmlid={p-487}]{\teiHi[rend={italic}]{Ibid.}, p. 227, v. 172.}} ». Sa description surpasse celles analysées précédemment, en ce qu’elle se manifeste comme l’expression de l’originaire le plus pur, dont le poète va patiemment percer les arcanes politiques :}
              \begin{teiLg}

                \teiL{Jamais forest ne fut si verte, ny si sombre ;}
                \teiL{Mil ans ont travaillé pour former sa belle ombre,}
                \teiL{Et jamais le Soleil, d’un rayon curieux,}
                \teiL{Ne sçauroit penetrer au secret de ces lieux.}
                \teiL{Là tous les animaux vivent francs de querelle,}
                \teiL{Et suivent la douceur qui leur est naturelle\teiNote[xmlid={ftn94-2},n={45},place={foot}]{\teiP[xmlid={p-488}]{\teiHi[rend={italic}]{Ibid.}, p. 226, v. 133-158.}}.}
              
\end{teiLg}
              \teiP[xmlid={p27-4}]{Le regard du poète, lui, est capable de « penetrer au secret de ces lieux » et de rendre compte des enseignements politiques qu’ils comportent. L’inscription du Temple dans ces bois qui semblent détachés de l’histoire humaine relève là encore d’une fictionnalisation théocratique ; la venue de Richelieu était prophétisée bien avant son apparition, la description constituant en fait une promenade parmi les éminentes victoires de Richelieu : La Rochelle\teiNote[xmlid={ftn95-2},n={46},place={foot}]{\teiP[xmlid={p-489}]{« C’est là qu’on aperçoit ceste fiere Rebelle, / Cet objet de terreur, la superbe Rochelle, / Perdre le nom de Ville, aussi bien que le cœur, / Pour donner à mon ROY, celuy de son Vainqueur. / Là, tous ces grands travaux se font paroistre insignes, / L’on y voit tous nos Forts, attachez par des lignes ; / Et l’Ocean reçoit des mains de RICHELIEU, / Un frain, qu’il n’eut jamais, que de celles de Dieu. / Le Peintre industrieux, pour troupes ennemies, / Sur le haut de Tadon, fait marcher des Momies, / Et ces hardis Pinceaux figurent enfermez, / Des Phantosmes en garde, et des Spectres armez. / On voit que leur orgueil est tout près de s’abatre. / Ils ont un ennemy qu’ils ne sçauroient combatre ; / Haves, tous descharnez, et demy-morts de faim, / Cents mets extravagants paroissent en leur main » (\teiHi[rend={italic}]{ibid.}, p. 230, v. 277-292).}} et Cazal\teiNote[xmlid={ftn96-2},n={47},place={foot}]{\teiP[xmlid={p-490}]{« Un autre me fait voir les Espagnols rusez, / Le Siege de Cazal, et deux camps opposez, / L’on y voit la poussiere, on y voit la fumée, / Des chevaux, et des feux de ces grands Corps d’Armée ; / Ny Guidons, ny Drapeaux n’y sont pas oubliez ; / Les blancs sont ondoyants, les rouges sont pliez : / Icy des Esquadrons ; là de l’Infanterie ; / Là se voit le Bagage ; icy l’Artillerie ; / Un Champ reste au milieu, vuide pour le Vainqueur ; / Mais trente mille bras ne treuvent pas un Cœur : / Et ces fiers Conquerans, que la frayeur tenaille, / De peur de la donner, perdent une Bataille » (\teiHi[rend={italic}]{ibid.}, p. 232-233, v. 367-378).}} couronnent l’écriture par les lieux des hauts faits du ministre, leur enchaînement produisant un effet synesthétique puissant : « O merveilleux effect de ce rare Pinceau, / Qui rend le bruit visible, en ce divin Tableau\teiNote[xmlid={ftn97-2},n={48},place={foot}]{\teiP[xmlid={p-491}]{\teiHi[rend={italic}]{Ibid.}, p. 232, v. 347-348.}}. » Ce dispositif suscite tous ces sens pour mieux forcer l’adhésion politique et culmine dans les propos d’« Apollon au Roy », qui déclare au sujet de Richelieu :}
              \begin{teiLg}

                \teiL{Sa prevoyance est sans seconde ;}
                \teiL{Comme son zele est sans pareil ;}
                \teiL{Et crois ce qu’en dit le Soleil,}
                \teiL{Qui chaque jour voit tout le Monde :}
                \teiL{Je vay chez tous les Potentats ;}
                \teiL{Je visite tous les Estats ;}
                \teiL{Mais par tout sa gloire est premiere :}
                \teiL{Il passe tout ce qu’on escrit ;}
                \teiL{Et mon Char a moins de lumiere,}
                \teiL{Que ce rare et divin Esprit\teiNote[xmlid={ftn98-2},n={49},place={foot}]{\teiP[xmlid={p-492}]{\teiHi[rend={italic}]{Ibid.}, p. 235, v. 453-462.}}.}
              
\end{teiLg}
              \teiP[xmlid={p28-4}]{On assiste ici au transfert de la phénoménologie des lieux originaires sur l’être humain, dont les qualités sont également perçues comme « premières », c’est-à-dire, sur le plan hyperbolique de l’éloge, superlatives et supérieures même à l’humanité de leur détenteur. De ce fait, elles apparaissent, sur le plan de la théocratie politique de Scudéry, comme l’expression de la lumière élective originelle du politique, et donc comme son incarnation la plus achevée.}
              \teiP[xmlid={p29-3}]{Il apparaît en somme que Scudéry a intensément réfléchi à la manière la plus efficace de ramener ses contemporains au culte de Richelieu, en organisant un dispositif de représentation fondé sur l’incarnation du regard du ministre au cœur de lieux systématiquement chargés d’histoire, ou plus exactement de son histoire politique. En procédant ainsi, le poète réactive – ou tente de le faire – la force émotionnelle de ces événements : les lieux sont partie prenante d’une stratégie de re-médiatisation des exploits de Richelieu. Mais la carte maîtresse du propagandiste réside dans le fait qu’il a d’emblée verrouillé ce dispositif de représentation sur la réussite prétendue d’une tentative d’instauration dynastique. En effet, avant même les textes l’appelant à imiter son illustre ancêtre, le petit-neveu semble avoir d’ores et déjà suivi les conseils du poète, comme la dédicace des \teiHi[rend={italic}]{Poésies diverses} le donne à penser :}
              \begin{teiCit}[xmlid={cit02-10}]

                \begin{teiQuote}
La Mer a esté le vaste Theatre de vostre gloire : la fameuse Ville de Naples en a esté le tesmoin : et l’espoisseur obscure de la fumée des Canons, n’a fait qu’en augmenter l’esclat, et que la rendre plus brillante\teiNote[xmlid={ftn99-2},n={50},place={foot}]{\teiP[xmlid={p-493}]{Georges de Scudéry, \teiHi[rend={italic}]{Poésies diverses I}, \teiHi[rend={italic}]{op. cit.}, p. 40.}}.
\end{teiQuote}
              
\end{teiCit}
              \teiP[xmlid={p30-3}]{L’expansion territoriale se poursuit ; le nom de Richelieu demeurera, dans l’esprit du poète, éternellement associé aux lieux du politique triomphant. Par la grâce d’une écriture se rêvant performative, Scudéry met en scène l’étape suivant la révélation du secret des lieux politiques : l’accomplissement des vérités théocratiques qu’ils contenaient et que le secret empêchait justement de voir clairement. Pour autant, ceci ne signifie pas que le mystère attaché à ces représentations disparaisse : l’enjeu, pour le poète, est de maintenir la dimension sacrale du politique, tout en la mettant dans le même temps à nu. De ce fait, on peut considérer que de tels écrits procèdent d’une mythographie du regard, active tout spécialement dans les lieux qui balisent et scandent la mise en œuvre et en récit du secret politique.}
            
\end{teiDiv}
          
\end{teiElement}
        
\end{teiElement}
        \begin{teiElement}[name={group},type={chapter},data-page-title={Secrets de théâtre et héros « en déshabillé », de la comédie des Lumières au drame romantique},data-page-authors={Thibaut Julian@@Université Lumière Lyon 2, IHRIM, UMR 5317, Lyon, France},data-include-href={../XML/chap\_12\_Secrets\_de\_theatre\_Locus\_politicus.xml},xmlid={chapter-011}]

          \begin{teiElement}[name={front}]

            \begin{teiDiv}[type={titlePage},xmlid={titlepage-012}]

              \teiP[rend={title-main},xmlid={p-494}]{Secrets de théâtre et héros « en déshabillé », de la comédie des Lumières au drame romantique}
              \teiP[rend={author-aut},xmlid={p-495}]{Thibaut \teiHi[rend={small-caps}]{Julian}}
              \teiP[rend={authority\_affiliation},xmlid={p-496}]{Université Lumière Lyon 2, IHRIM, UMR 5317, Lyon, France}
            
\end{teiDiv}
          
\end{teiElement}
          \begin{teiElement}[name={body},xmlid={body-012}]

            \teiP[xmlid={p1-12}]{Contrairement au tour de magie, le théâtre place le public dans la confidence : il n’est \teiHi[rend={italic}]{a priori} point de secret pour lui. En position de voyeur, il perçoit et déchiffre un ensemble de signes, activité qui contribue à son plaisir\teiNote[xmlid={ftn1-12},n={1},place={foot}]{\teiP[xmlid={p-497}]{Nous renvoyons aux trois ouvrages fondamentaux d’Anne Ubersfeld réunis depuis 1977 sous le titre \teiHi[rend={italic}]{Lire le théâtre} \lbrack{}rééd. Paris, Belin, 1996\rbrack{}, ainsi qu’à Florence Naugrette, \teiHi[rend={italic}]{Le Plaisir du spectateur de théâtre}, Rosny-sous-Bois, Bréal, 2002.}}. Les personnages, à l’inverse, ignorent tout ou partie des intentions, des « fourberies » et des secrets de leurs semblables : Tartuffe ne sait pas qu’Elmire a caché Orgon sous la table, Géronte que Scapin le bat enveloppé dans son sac, ni Britannicus que son entretien avec Junie est épié par Néron dans ce « lieu \teiHi[rend={italic}]{stupéfié}\teiNote[xmlid={ftn61-8},n={2},place={foot}]{\teiP[xmlid={p-498}]{Roland Barthes, \teiHi[rend={italic}]{Sur Racine}, Paris, Éditions du Seuil, 1963, p. 19-20.}} » que constitue la scène tragique racinienne. Dans l’univers socialement hiérarchisé que le théâtre réfléchit, aussi bien parmi le monde domestique de la comédie que dans l’univers curial qui organise traditionnellement la tragédie, l’action est mue par des conflits et des relations de pouvoir qui jouent sur des secrets plus ou moins bien gardés.}
            \teiP[xmlid={p2-12}]{Ainsi, le travestissement est un \teiHi[rend={italic}]{leitmotiv} pour les personnages de grands dissimulés dès les premières comédies de Marivaux, comme \teiHi[rend={italic}]{Le Prince travesti} et \teiHi[rend={italic}]{La Fausse Suivante ou le Fourbe puni} (1724) : il s’agit alors de personnages strictement fictifs qui se camouflent pour mener une enquête, sans être détenteurs du pouvoir suprême dans la plupart des cas. Une exception : le Prince « tendre » de \teiHi[rend={italic}]{La Double Inconstance} (1723), d’abord spectateur de l’intrigue conduite par Flaminia, puis acteur déguisé en officier courtisant Silvia dès la scène 2 de l’acte II, qui lève le masque à l’acte III auprès de son rival Arlequin puis devant la proie qu’il a séduite : « Oui Silvia, je vous ai jusqu’ici caché mon rang, pour essayer de ne devoir votre tendresse qu’à la mienne : je ne voulais rien perdre du plaisir qu’elle pouvait me faire\teiNote[xmlid={ftn62-8},n={3},place={foot}]{\teiP[xmlid={p-499}]{Marivaux, \teiHi[rend={italic}]{La Double Inconstance} \lbrack{}1723\rbrack{}, Frédéric Deloffre (éd.), \teiHi[rend={italic}]{Théâtre complet}, Paris, Garnier, 1968, 2 vol., ici t. I, p. 314-315. Voir Jean Goldzink, \teiHi[rend={italic}]{Comique et comédie au siècle des Lumières}, Paris, L’Harmattan, 2000, « L’archipel Marivaux », p. 169-369.}}. »}
            \teiP[xmlid={p3-12}]{Prolongeant ce processus de révélation, la comédie historique qui se développe dans la seconde moitié du \teiHi[rend={small-caps}]{xviii}\teiHi[rend={sup}]{e} siècle vient accentuer la transgression des normes politiques et des frontières entre le haut et le bas, l’histoire et la fiction : il n’est pas rare qu’un « grand », généralement de statut royal, se transporte dans un milieu bourgeois ou populaire, en décalage avec son code de référence. L’intrigue elle-même est en partie construite sur un secret assumé par cette autorité masquée, avant qu’elle ne se dévoile et soit « restaurée » en majesté. En outre, la cour apparaît comme un lieu tendu de pièges et hanté par les conspirateurs et les espions avec lesquels tout prince ou ministre doit ruser, en tirant profit des espaces et des accessoires que la scène lui fournit. Nous proposons ici de saisir le fil reliant la comédie historique des Lumières au drame romantique, à partir des espaces et de la dramaturgie du secret qui les sous-tend ; il s’agit d’en faire ressortir la permanence et d’en questionner les implications politiques entre 1760 et 1830, au cours d’une période qui, malgré les césures des régimes entre monarchie(s), république et Empire, est marquée par le mouvement conjoint et continu de la désacralisation de la royauté et du débordement de l’esthétique « classique ».}
            \begin{teiDiv}[type={section1},xmlid={div01-11}]

              \teiHead[xmlid={deux-modeles-colle-et-beaumarchais-001}]{Deux modèles : Collé et Beaumarchais}
              \teiP[xmlid={p4-12}]{Dans le discours\teiHi[rend={italic}]{ De la poésie dramatique} (1758), Diderot postule que « La comédie veut être jouée en déshabillé. Il ne faut être sur la scène ni plus apprêté ni plus négligé que chez soi\teiNote[xmlid={ftn63-8},n={4},place={foot}]{\teiP[xmlid={p-500}]{Denis Diderot, \teiHi[rend={italic}]{De la poésie dramatique} \lbrack{}1758\rbrack{}, Laurent Versini (éd.), \teiHi[rend={italic}]{Œuvres}, t. IV, \teiHi[rend={italic}]{Esthétique-Théâtre}, Paris, Robert Laffont, « Bouquins », 1996, p. 1334.}} ». Le « déshabillé » est défini dans l’\teiHi[rend={italic}]{Encyclopédie }comme un « terme fort en usage en France, et que les Anglais ont adopté depuis peu. Il signifie proprement une \teiHi[rend={italic}]{robe de chambre}, et les autres choses dont on se couvre quand on est chez soi en négligé\teiNote[xmlid={ftn64-7},n={5},place={foot}]{\teiP[xmlid={p-501}]{« Déshabillé », \teiHi[rend={italic}]{Encyclopédie ou Dictionnaire raisonné des sciences, des arts et des métiers}, 1754, t. IV, p. 882, en ligne : \teiRef[target={http://enccre.academie-sciences.fr/encyclopedie/article/v4-2231-0/}]{\teiHi[rend={underline}]{http://enccre.academie-sciences.fr/encyclopedie/article/v4-2231-0/}}.}} ». Cela annonce à un siècle de distance la série des \teiHi[rend={italic}]{Grands hommes en robe de chambre}, Alexandre Dumas préférant à l’histoire des « héros drapés dans des habits de cérémonie », le « déshabillé » tiré des « notes empruntées aux valets de chambre\teiNote[xmlid={ftn65-7},n={6},place={foot}]{\teiP[xmlid={p-502}]{Alexandre Dumas, \teiHi[rend={italic}]{Les Grands Hommes en robe de chambre}, Paris, Michel Lévy, 1866, 2 vol., ici t. II, p. ii.}} ». Mais Dumas n’aurait probablement pas repris les anecdotes sur Henri \teiHi[rend={small-caps}]{iv}, Louis \teiHi[rend={small-caps}]{xiii} et Richelieu si le théâtre n’avait abondamment popularisé leur vie privée à la lisière de la sphère publique. C’est à Charles Collé, secrétaire et lecteur du duc d’Orléans, que revient la primeur de la représentation du « bon roi Henri » au théâtre à partir de 1762 dans \teiHi[rend={italic}]{La Partie de chasse de Henri }\teiHi[rend={ small-caps italic}]{IV} ; c’est à lui aussi qu’est due la popularisation de la formule :}
              \begin{teiCit}[xmlid={cit01-12}]

                \begin{teiQuote}
Le titre seul de la pièce annonce assez que je n’ai point eu la prétention de montrer dans une comédie le grand roi, le premier capitaine de son siècle, le politique équitable, le conquérant légitime \lbrack{}…\rbrack{}. Ce sont seulement quelques instants de sa vie privée que j’ai saisis. C’est (si l’on veut me passer cette expression) le \teiHi[rend={italic}]{héros en déshabillé} que j’ai essayé de peindre\teiNote[xmlid={ftn66-7},n={7},place={foot}]{\teiP[xmlid={p-503}]{Charles Collé, « Avertissement » de \teiHi[rend={italic}]{La Partie de chasse de Henri }\teiHi[rend={ small-caps italic}]{IV}, Jacques Truchet (éd.), \teiHi[rend={italic}]{Théâtre du }\teiHi[rend={ small-caps italic}]{xviii}\teiHi[rend={italic sup}]{e}\teiHi[rend={italic}]{ siècle}, Paris, Gallimard, « Bibliothèque de la Pléiade », 1974, t. II, p. 599.}}.
\end{teiQuote}
              
\end{teiCit}
              \teiP[xmlid={p5-12}]{L’intrigue de cette « comédie » imitée d’une pièce anglaise de Dodsley est très simple\teiNote[xmlid={ftn67-7},n={8},place={foot}]{\teiP[xmlid={p-504}]{Voir Marie-Emmanuelle Plagnol\teiHi[rend={small-caps}]{-}Diéval et Dominique Quéro (dir.), \teiHi[rend={italic}]{Charles Collé (1709-1783), au cœur de la République des Lettres}, Rennes, PUR, « Interférences », 2013.}}. La structure en trois actes oppose sphère publique et sphère privée grâce à la variation des espaces\teiNote[xmlid={ftn68-4},n={9},place={foot}]{\teiP[xmlid={p-505}]{Pour les distinctions entre les notions de « lieu » et d’« espace » (théâtral, scénique, ludique, dramatique, etc.), voir Georges Zaragoza, \teiHi[rend={italic}]{Faire jouer l’espace dans le théâtre romantique européen. Essai de dramaturgie comparée}, Paris, Honoré Champion, 1999, p. 20-22.}} : le premier acte – ajout de Collé – se situe en journée au château de Fontainebleau, au milieu des intrigues de cour culminant avec la réconciliation entre Henri et son ministre Sully, dont la réputation était ternie par des ragots, tandis que les deux suivants retracent l’égarement du roi à la chasse dans la forêt de Sénart et sa rencontre, à la nuit venue, avec le meunier Michau, qui le prend d’abord pour un braconnier puis offre l’hospitalité à celui qui se présente comme un officier de la Cour. Le roi, spectateur interne, jouit \teiHi[rend={italic}]{incognito} des éloges dont il fait l’objet, notamment lors du repas rustique du troisième acte où il déclare en aparté : « Quel plaisir ! Je vais donc avoir encore une fois la satisfaction d’être traité comme un homme ordinaire… de voir la nature humaine sans déguisement ; cela est charmant\teiNote[xmlid={ftn69-4},n={10},place={foot}]{\teiP[xmlid={p-506}]{Charles Collé, \teiHi[rend={italic}]{La Partie de chasse de Henri }\teiHi[rend={ small-caps italic}]{IV}, \teiHi[rend={italic}]{op. cit.}, III, 2, p. 637.}} ! » La jonction des intrigues curiale et populaire est opérée par une amourette de paysans que le roi unit en libérant la fille des rets du libertin Conchiny. L’arrivée de Sully, Bellegarde et Conchiny révèle son identité et renforce l’admiration des personnages populaires. Comédie la plus jouée à la Comédie-Française derrière \teiHi[rend={italic}]{Le Mariage de Figaro} entre sa création publique en 1774 et 1793, elle connaît plus de cent représentations parisiennes durant la Révolution et des reprises dès 1814-1815, ainsi que de nombreuses imitations\teiNote[xmlid={ftn70-4},n={11},place={foot}]{\teiP[xmlid={p-507}]{Par exemple \teiHi[rend={italic}]{Une Journée de Henri }\teiHi[rend={ small-caps italic}]{IV}, comédie de Villemain d’Abancourt (Théâtre Molière, 12 octobre 1791), ou \teiHi[rend={italic}]{La Forêt de Sénart ou la Partie de chasse de Henri }\teiHi[rend={ small-caps italic}]{IV}, opéra-comique de Castil-Blaze (Odéon, 14 janvier 1826).}} ; cette popularité atteste du succès de l’humanisation du monarque, marquant selon Sophie Mentzel « l’acceptation d’une forme de désacralisation\teiNote[xmlid={ftn71-4},n={12},place={foot}]{\teiP[xmlid={p-508}]{Sophie Mentzel, \teiHi[rend={italic}]{Trônes vacillants. La représentation de la royauté sur la scène romantique, 1820-1840}, thèse de Littérature française, Université de Nantes, 2016, p. 84. Voir aussi la section « Les Bourbons en déshabillé », p. 529-537.}} », accentuée par la Révolution. Collé est donc la première référence convoquée par les créateurs de la « comédie historique », dont les recueils d’œuvres dramatiques paraissent sous la Restauration : Alexandre Duval\teiNote[xmlid={ftn72-4},n={13},place={foot}]{\teiP[xmlid={p-509}]{Alexandre Duval, \teiHi[rend={italic}]{Œuvres complètes d’Alexandre Duval}, Paris, Barba, 1822, t. I, p. xx.}} (\teiHi[rend={italic}]{La Jeunesse de Richelieu ou le Lovelace français}, coécrite avec Monvel, date de 1796) et Népomucène Lemercier (\teiHi[rend={italic}]{Pinto}, créée en 1800, sera reprise à la Porte Saint-Martin en 1834\teiNote[xmlid={ftn73-4},n={14},place={foot}]{\teiP[xmlid={p-510}]{Louis Népomucène Lemercier, \teiHi[rend={italic}]{Comédies historiques}, Paris, Ambroise Dupont et C\teiHi[rend={sup}]{ie}, 1828.}}). Deux auteurs alors hostiles aux romantiques, mais que les expérimentations dramatiques rapprochent pourtant d’eux.}
              \teiP[xmlid={p6-12}]{L’autre grand « précurseur\teiNote[xmlid={ftn74-3},n={15},place={foot}]{\teiP[xmlid={p-511}]{Georges Zaragoza attribue à Beaumarchais le rôle de « précurseur » ou maillon entre l’esthétique classique et l’esthétique romantique en fait d’utilisation de l’espace. Voir Georges Zaragoza, \teiHi[rend={italic}]{Faire jouer l’espace dans le théâtre romantique européen},\teiHi[rend={italic}]{ op. cit.}, p. 34-36\teiHi[rend={italic}]{.}}} » est Beaumarchais qui, sans avoir puisé dans l’histoire, a popularisé avec sa trilogie de Figaro un ensemble de recettes à succès pour cacher des secrets dans une atmosphère libertine, en particulier grâce à ce que Jacques Scherer a appelé le « troisième lieu\teiNote[xmlid={ftn75-3},n={16},place={foot}]{\teiP[xmlid={p-512}]{Jacques Scherer, \teiHi[rend={italic}]{La Dramaturgie de Beaumarchais}, Paris, Nizet, 1954, chap. 5, « Les instruments de tension ».}} ». Cet espace en scène distinct des coulisses sert de cachette à la fois visible et dérobée aux regards – tels sont les pavillons du parc à l’acte V du \teiHi[rend={italic}]{Mariage de Figaro} ainsi que la scène du cabinet de toilette accolé à la chambre de la comtesse dans l’acte II, où Suzanne cache Chérubin déguisé en jeune fille, tandis que le comte Almaviva furieux, prêt à enfoncer la porte, s’attend à y trouver un amant – le page ayant entre-temps sauté par la fenêtre\teiNote[xmlid={ftn76-3},n={17},place={foot}]{\teiP[xmlid={p-513}]{Beaumarchais, \teiHi[rend={italic}]{Le Mariage de Figaro} (1778 et 1784), Pierre Larthomas (éd.), \teiHi[rend={italic}]{Œuvres}, Paris, Gallimard, « Bibliothèque de la Pléiade », 1988, II, 10-20, p. 411-420.}}. Dumas prétendra que la scène provient d’une infidélité faite à Henri \teiHi[rend={small-caps}]{IV} par Gabrielle d’Estrées chez elle avec le duc de Bellegarde, sorti par la fenêtre d’un cellier fermé à clé, rempli de pots de confiture\teiNote[xmlid={ftn77-3},n={18},place={foot}]{\teiP[xmlid={p-514}]{Dumas, \teiHi[rend={italic}]{Les Grands Hommes en robe de chambre},\teiHi[rend={italic}]{ op. cit}., t. II, p. 78.}}. Il s’en est certainement souvenu pour le dénouement d’\teiHi[rend={italic}]{Henri }\teiHi[rend={ small-caps italic}]{III}\teiHi[rend={italic}]{ et sa cour }en 1829 : Arthur, le « petit page » de la duchesse de Guise rappelle Chérubin, dont la charge érotique est transférée sur Saint-Mégrin, amant de la duchesse s’échappant lui aussi par la fenêtre avant d’être tué par les hommes de Guise\teiNote[xmlid={ftn78-3},n={19},place={foot}]{\teiP[xmlid={p-515}]{Dumas, \teiHi[rend={italic}]{Henri III et sa cour} \lbrack{}1829\rbrack{}, Sylvain Ledda (éd.), Paris, Flammarion, « GF », 2016, V, 2-3, p. 162-164.}}. Mais dès 1800, \teiHi[rend={italic}]{Pinto ou la Journée d’une conspiration} déplaçait le procédé : le décor du troisième acte rappelle la scène du cabinet de la comtesse, mais par inversion, car la fenêtre de la chambre de la duchesse de Bragance sert à l’entrée par effraction nocturne de l’amiral espagnol Lopez Ozorio, croyant avoir reçu d’elle des marques de faveur… Détrompé et penaud, il lui révèle que son mari doit être arrêté le lendemain. Du reste, la comédie a souvent été comparée à \teiHi[rend={italic}]{La Folle Journée}, et Lemercier fait de Beaumarchais plus un modèle d’émulation pour renouveler le comique qu’une source, alors même que le style et les monologues de Pinto ont plusieurs similitudes avec ceux de Figaro\teiNote[xmlid={ftn79-3},n={20},place={foot}]{\teiP[xmlid={p-516}]{Voir Vincenzo De Santis, \teiHi[rend={italic}]{Le Théâtre de Louis Lemercier, entre Lumières et romantisme}, Paris, Classiques Garnier, « L’Europe des Lumières », 2015, p. 191-221.}}. Parallèlement, l’ouverture de la pièce en extérieur durant la chasse fait signe vers la pièce de Collé, dont la formule est glosée par Lemercier pour préciser ses objectifs poétique et politique :}
              \begin{teiCit}[xmlid={cit02-11}]

                \begin{teiQuote}
\lbrack{}…\rbrack{} j’aperçus qu’en dépouillant ces éminents personnages du faux appareil qui les conserve, et qu’en appliquant à leurs vices et à leurs actions perverses la force du ridicule, il en résulterait un genre vrai, moral, instructif, qui apprendrait au peuple à démasquer la basse politique, et lui montrerait les grands en déshabillé, et, pour ainsi dire, mis à nu sous le fouet de la satire\teiNote[xmlid={ftn80-3},n={21},place={foot}]{\teiP[xmlid={p-517}]{Louis Népomucène Lemercier, \teiHi[rend={italic}]{Comédies historiques},\teiHi[rend={italic}]{ op. cit.}, p. iii.}}.
\end{teiQuote}
              
\end{teiCit}
            
\end{teiDiv}
            \begin{teiDiv}[type={section1},xmlid={div02-11}]

              \teiHead[xmlid={espaces-ambivalents-001}]{Espaces ambivalents}
              \teiP[xmlid={p7-12}]{Dans la comédie historique des Lumières, le dispositif du « déshabillé » permet une forme d’effraction douce ; cette expérience liminale symbolique se déroule souvent dans des cadres champêtres et festifs en deux étapes : le roi d’abord camouflé revient en majesté pour « dénouer » la pièce et ses secrets. C’est le cas dans deux comédies au tournant du siècle situées en extérieur, où François \teiHi[rend={small-caps}]{I}\teiHi[rend={sup}]{er} est souverain. Dans \teiHi[rend={italic}]{Le Chevalier sans peur et sans reproche ou les Amours de Bayard} (« comédie héroïque » de Monvel créée en 1786, reprise en 1808), dès la scène d’exposition, le roi venu chez la belle veuve M\teiHi[rend={sup}]{me} de Rendan, courtisée par de nombreux prétendants de haut rang, demande la discrétion à un seigneur de sa cour : « Ah ! point de Majesté, je vous prie, Imbercourt. Souvenez-vous que je ne suis ici qu’un très petit particulier, un pauvre amant rebuté : ce n’est pas en matière de galanterie, et surtout quand on éprouve l’humiliation d’un refus, qu’il convient de faire le roi. Gardons l’incognito, mon ami\teiNote[xmlid={ftn81-3},n={22},place={foot}]{\teiP[xmlid={p-518}]{Jacques-Marie Boutet de Monvel, \teiHi[rend={italic}]{Le Chevalier sans peur et sans reproche ou les Amours de Bayard}, dans Roselyne Laplace (éd.), \teiHi[rend={italic}]{Théâtre, discours politiques et réflexions diverses}, Paris, Honoré Champion, 2001, I, 1, p. 209.}}. » De fait, il s’éclipse très vite avant d’avoir pu rencontrer la dame et ne revient qu’au dernier acte pour présider très officiellement le duel du chevalier Bayard en lice contre son rival espagnol Sotomayor, tué en scène, avant que le roi débonnaire n’unisse les amants. Dans la comédie en deux actes et en vers mêlée d’ariettes de Sewrin et Chazet, \teiHi[rend={italic}]{François }\teiHi[rend={ small-caps italic}]{I}\teiHi[rend={italic sup}]{er}\teiHi[rend={italic}]{ ou la Fête mystérieuse} (1807), le secret est plus durablement tu : le roi – substitué à Henri \teiHi[rend={small-caps}]{IV}, interdit par la censure impériale\teiNote[xmlid={ftn82-3},n={23},place={foot}]{\teiP[xmlid={p-519}]{Voir Odile Krakovitch, « Les mythes du bon et du mauvais roi : Henri \teiHi[rend={small-caps}]{IV} et François\teiHi[rend={small-caps}]{ I}\teiHi[rend={sup}]{er} dans le théâtre de la première moitié du \teiHi[rend={small-caps}]{xix}\teiHi[rend={sup}]{e} siècle », \teiHi[rend={italic}]{La Légende d’Henri IV}, Pau, Société Henri IV, 1995, p. 215-241.}} – n’apparaît en personne qu’au début du second acte, dans le parc d’un château à la tombée de la nuit, couvert d’un long manteau, avant la fête qu’il a secrètement ordonnée pour D’Aubigny (\teiHi[rend={italic}]{alias} D’Aubigné), ne dévoilant son statut de commanditaire qu’au seul intendant du château dans une lettre. Se pensant évincé par un rival, D’Aubigny « ne doute pas que la fête qu’on prépare ne soit toute entière pour son heureux compétiteur. Le Roi s’amuse quelque temps des alarmes du courtisan car il ne l’a exilé qu’en apparence, et lui conserve toujours une tendre affection. Il révoque ensuite l’exil, et reçoit son ami au milieu des fêtes, en lui annonçant que c’est pour lui seul qu’elles ont été préparées\teiNote[xmlid={ftn83-3},n={24},place={foot}]{\teiP[xmlid={p-520}]{\teiHi[rend={italic}]{Courrier des spectacles}, 15 mars 1807, p. 2.}}. » L’entrée du roi est célébrée en fanfare dans un chœur vantant le paternalisme bienveillant de son règne, et la pièce se clôt par la célébration des fiançailles préparées à l’insu des amants\teiNote[xmlid={ftn84-3},n={25},place={foot}]{\teiP[xmlid={p-521}]{Sewrin et René de Chazet, \teiHi[rend={italic}]{François I}\teiHi[rend={italic sup}]{er}\teiHi[rend={italic}]{, ou la Fête mystérieuse}, Paris, Barba, 1807.}}. Dans ces deux cas, ce sont moins les lieux qui sont évocateurs, que les deux temps de l’intrigue avec le retour du personnage sous deux facettes, passant du statut privé dans l’ombre à la lumière de la majesté acclamée.}
              \teiP[xmlid={p8-12}]{Pour autant, on l’a vu avec Collé, la comédie historique des Lumières s’affranchit souvent déjà de l’unité de lieu, proposant une alternance signifiante que l’on retrouve dans le drame romantique, qui joue avec la contiguïté et la succession des espaces sociaux (publics) et privés (intimes)\teiNote[xmlid={ftn85-3},n={26},place={foot}]{\teiP[xmlid={p-522}]{Voir Anne Ubersfeld, « L’espace du drame romantique », Isabelle Moindrot (dir.), \teiHi[rend={italic}]{Le Spectaculaire dans les arts de la scène, du romantisme à la Belle-Époque}, Paris, CNRS Éditions, « Arts du spectacle », 2006, p. 16-23.}}. Prenons \teiHi[rend={italic}]{Le Roi s’amuse} de Hugo : les actes I et III sont au Louvre avec les courtisans, alors que les actes II, IV et V se situent de nuit dans des lieux parisiens interlopes ou tenus secrets : la maison où le bouffon Triboulet cache sa fille Blanche dans le « recoin le plus désert du cul-de-sac Bussy\teiNote[xmlid={ftn86-3},n={27},place={foot}]{\teiP[xmlid={p-523}]{Victor Hugo, \teiHi[rend={italic}]{Le Roi s’amuse} \lbrack{}1832\rbrack{}, Jean-Jacques Thierry et Josette Mélèze (éd.), \teiHi[rend={italic}]{Théâtre complet}, Paris, Gallimard, « Bibliothèque de la Pléiade », 1963, t. I, p. 1372.}} », puis le bouge de Saltabadil en bord de Seine, « lieu-Cour des Miracles » et « espace de la dérision » anti-tragique\teiNote[xmlid={ftn87-3},n={28},place={foot}]{\teiP[xmlid={p-524}]{Anne Ubersfeld, \teiHi[rend={italic}]{Le Roi et le Bouffon : étude sur le théâtre de Hugo de 1830 à 1839 }\lbrack{}1974\rbrack{}, Paris, J. Corti, 2001, p. 636.}}, tout à la fois lupanar royal et lieu du crime commandité par son bouffon. De même, \teiHi[rend={italic}]{Henri }\teiHi[rend={ small-caps italic}]{III}\teiHi[rend={italic}]{ et sa cour} de Dumas offre « trois types de décorations pour cinq actes : deux espaces politiques (II, IV), deux lieux intimes (III, V) et un décor fantastique aux rémanences gothiques (I)\teiNote[xmlid={ftn88-3},n={29},place={foot}]{\teiP[xmlid={p-525}]{Sylvain Ledda, « Notice » de \teiHi[rend={italic}]{Henri }\teiHi[rend={ small-caps italic}]{III}\teiHi[rend={italic}]{ et sa cour}, \teiHi[rend={italic}]{op. cit.}, p. 45.}}. »}
              \teiP[xmlid={p9-12}]{Ce va-et-vient de tableaux contrastés emprunte aussi à la tradition de la comédie satirique visant des figures célèbres, qui s’est développée à la faveur de la Révolution en même temps que le mélodrame\teiNote[xmlid={ftn89-3},n={30},place={foot}]{\teiP[xmlid={p-526}]{Sur la veine pamphlétaire, pornographique ou carnavalesque du théâtre de la Révolution, voir mon ouvrage \teiHi[rend={italic}]{Un théâtre pour la Nation. L’histoire en scène (1765-1806)}, Lyon, PUL, « Théâtre et société », 2022, p. 163-169.}}. L’action du \teiHi[rend={italic}]{Lovelace français}, comédie créée au Théâtre de la République en décembre 1796, se situe sous la Régence et les cinq actes font alterner trois lieux : « La scène se passe à Paris, chez M. Michelin \lbrack{}I, III et V\rbrack{}, à l’hôtel de Richelieu \lbrack{}II\rbrack{}, et dans la petite maison de ce dernier \lbrack{}IV\rbrack{}\teiNote[xmlid={ftn90-3},n={31},place={foot}]{\teiP[xmlid={p-527}]{Alexandre Duval et Jacques Marie Boutet dit Monvel, \teiHi[rend={italic}]{La Jeunesse du duc de Richelieu, ou le Lovelace français}, Paris, Barba, \lbrack{}an V\rbrack{}, n. p. (didascalie liminaire). J’ajoute entre crochets les actes correspondants.}}. » Le duc de Richelieu, « nouveau César \lbrack{}qui fait\rbrack{} marcher ensemble les affaires et les plaisirs\teiNote[xmlid={ftn91-3},n={32},place={foot}]{\teiP[xmlid={p-528}]{\teiHi[rend={italic}]{Ibid}., II, 5, p. 34.}} », entreprend de séduire M\teiHi[rend={sup}]{me} Michelin qui lui résiste malgré son inclination ; il se fait passer pour son « valet de chambre » La Fosse lorsqu’il rend visite aux Michelin, dont la maison devient un terrain de chasse à l’insu du mari. Un dualisme fonctionnel distingue les domestiques du duc : si La Fosse s’occupe de ses menus plaisirs, Armand, épris de justice et ami de la famille, est son « secrétaire » politique depuis un an. On voit en effet l’homme à bonne fortune exercer son pouvoir sur les cœurs de maintes conquêtes éplorées et s’en servir pour obtenir des charges politiques. Après avoir fait conduire M\teiHi[rend={sup}]{me} Michelin contre son gré dans sa petite maison, où elle rencontre l’une de ses rivales (acte IV), le séducteur se repent et M. Michelin pardonne à sa femme : la « comédie » partage le code du mélodrame en soutenant l’idéologie du Directoire qui promeut la famille comme socle de la société bourgeoise, à l’opposé des valeurs aristocratiques libertines épinglées ici dans la débauche de la Régence et ses lieux de pouvoir, ternis par la cabale ou la dissimulation.}
              \teiP[xmlid={p10-12}]{Outre ce principe d’alternance des milieux sociaux, les espaces du secret sont à proprement parler ambivalents, marqués par des tensions fonctionnelles (ouverture/fermeture ; refuge/prison) et des frontières poreuses, que symbolisent des interfaces tels que portes, fenêtres, balcons et murs fissurés. Les palais romantiques regorgent de passages secrets, comme celui qui relie dans \teiHi[rend={italic}]{Henri }\teiHi[rend={ small-caps italic}]{III }\teiHi[rend={italic}]{et sa cour} l’hôtel de Catherine de Médicis au cabinet de l’astrologue Côme Ruggieri, qui comporte une « alcôve cachée dans la boiserie » permettant à la reine de « voir et entendre au besoin, sans être vue\teiNote[xmlid={ftn92-3},n={33},place={foot}]{\teiP[xmlid={p-529}]{Dumas, \teiHi[rend={italic}]{Henri }\teiHi[rend={ small-caps italic}]{III}\teiHi[rend={italic}]{ et sa cour}, \teiHi[rend={italic}]{op. cit.}, I, 1, p. 66 et 72-73.}} » ; ou bien chez Hugo le palais du podesta Angelo, où l’espion Homodei s’introduit dans la chambre de Catarina, femme du « tyran de Padoue », en surprenant une de ses domestiques au seuil de la « deuxième journée ». Les décors (donc la scène) matérialisent ainsi le secret et son renversement, tandis que les personnages sont sous surveillance, comme le traduit la reprise ironique par le ténébreux agent du Conseil des Dix :}
              \begin{teiCit}[xmlid={cit03-10}]

                \begin{teiQuote}
\teiHi[rend={small-caps}]{Reginella\teiLb{}}\lbrack{}…\rbrack{} Mais, vois-tu, Dafne, je te recommande le silence dans ce maudit palais. Il n’y a que cette chambre où l’on soit en sûreté. Ah ! ici, du moins, on est tranquille. On peut dire tout ce qu’on veut. C’est le seul endroit où, quand on parle, on soit sûr de ne pas être écouté.\teiLb{}\teiHi[rend={italic}]{Pendant qu’elle prononce ces derniers mots, un dressoir adossé au mur à droite tourne sur lui-même, donne passage à Homodei sans qu’elle s’en aperçoive, et se referme}.\teiHi[rend={small-caps}]{\teiLb{}Homodei\teiLb{}}C’est le seul endroit où quand on parle on soit sûr de ne pas être écouté\teiNote[xmlid={ftn93-3},n={34},place={foot}]{\teiP[xmlid={p-530}]{Victor Hugo, \teiHi[rend={italic}]{Angelo, tyran de Padoue} \lbrack{}1835\rbrack{}, Olivier Bara (éd.), Paris, Gallimard, « Folio théâtre », 2022, p. 80. Dafne est déjà sortie de scène.}}.
\end{teiQuote}
              
\end{teiCit}
              \teiP[xmlid={p11-12}]{Ce que l’on croit refuge est en fait piège et laboratoire, où la parole et les actions des personnages sont épiées. Ce dispositif voyeuriste voire panoptique était déjà central dans des pièces de Marivaux (dont \teiHi[rend={italic}]{La Dispute} fournit la forme de mise en abyme la plus explicite), ainsi que dans les comédies de conspiration de Lemercier, qui prétend avoir mis « les mémoires en action » pour exhiber « la vicieuse conduite des grandes affaires d’État\teiNote[xmlid={ftn94-3},n={35},place={foot}]{\teiP[xmlid={p-531}]{Lemercier, « Avant-propos » des \teiHi[rend={italic}]{Comédies historiques}, \teiHi[rend={italic}]{op. cit}., p. viii.}} ». C’est alors la ville entière, lieu de l’action politique excédant la scène comme Florence dans \teiHi[rend={italic}]{Lorenzaccio}\teiNote[xmlid={ftn95-3},n={36},place={foot}]{\teiP[xmlid={p-532}]{Voir Didier Alexandre, « Florence extra muros : le traitement de l’espace dans \teiHi[rend={italic}]{Lorenzaccio} », \teiHi[rend={italic}]{Littératures}, n\teiHi[rend={sup}]{o} 23, 1990, p. 117-134.}}, qui se trouve sous surveillance. Dans \teiHi[rend={italic}]{Richelieu ou la Journée des dupes }(écrite en 1804 mais créée en 1835 seulement\teiNote[xmlid={ftn96-3},n={37},place={foot}]{\teiP[xmlid={p-533}]{Voir Thibaut Julian, « Louis \teiHi[rend={small-caps}]{XIII }et Richelieu au théâtre sous Napoléon, ou Napoléon au théâtre sous Louis XIII et Richelieu », Laurent Angard, Guillaume Cousin et Blandine Poirier (dir.), \teiHi[rend={italic}]{Le Lys recomposé. La représentation des pouvoirs sous l’Ancien Régime dans la littérature fictionnelle du }\teiHi[rend={ small-caps italic}]{xix}\teiHi[rend={italic sup}]{e}\teiHi[rend={italic}]{ siècle (1800-1850)}, Publications numériques du CÉRÉDI, 2019, \teiRef[target={http://publis-shs.univ-rouen.fr/ceredi/index.php?id=680}]{\teiHi[rend={underline}]{http://publis-shs.univ-rouen.fr/ceredi/index.php?id=680}}.}}), la conspiration des partisans de la reine-mère échoue, ce que traduisent les quatre derniers actes situés dans le foyer du pouvoir, les palais du Louvre et du Luxembourg ; seul le premier acte peint « une rue voisine du couvent du Val de Grâce », lieu de réunion nocturne des comploteurs. Or le capucin Jacques, mouche du cardinal, suit leurs traces\teiNote[xmlid={ftn97-3},n={38},place={foot}]{\teiP[xmlid={p-534}]{Lemercier, \teiHi[rend={italic}]{Richelieu ou la Journée des dupes},\teiHi[rend={italic}]{ op}. \teiHi[rend={italic}]{cit}., I, 7, p. 196 : « Autour du Val-de-Grâce, où se glissent, le soir, / Des agents que, dit-on, la reine daigne voir. / Je n’espionne pas, Dieu m’en garde ! j’épie. / Je sers le cardinal, partant c’est œuvre pie. »}}, et le « grand orage » déclenché par Marie de Médicis au 4\teiHi[rend={sup}]{e} acte confirmera \teiHi[rend={italic}]{in fine} le pouvoir du cardinal et la disgrâce de la reine-mère. À l’inverse, dans \teiHi[rend={italic}]{Pinto}, la conjuration contre le pouvoir espagnol de Philippe \teiHi[rend={small-caps}]{II} aboutit à la révolution portugaise. Au terme du conciliabule matinal dans la forêt (acte I), le secrétaire du duc de Bragance, héros plébéien éponyme, ordonne aux seigneurs la plus grande dissimulation avant leur dispersion : « Vous, par le bois. Vous, le long du village. Vous, un fusil en main, chassant l’oiseau. Passez le Tage séparément ; surtout l’air désoccupé, le sourire à la bouche, le front libre, et point de ces rides de conspirateur\teiNote[xmlid={ftn98-3},n={39},place={foot}]{\teiP[xmlid={p-535}]{Lemercier, \teiHi[rend={italic}]{Pinto, ou la Journée d’une conspiration}, Vincenzo De Santis (éd.), dans \teiHi[rend={italic}]{Le Théâtre de Louis Lemercier, entre Lumières et romantisme}, \teiHi[rend={italic}]{op. cit}., p. 511 (fin de l’acte I).}}. » L’acte II déplace la scène dans « un salon du palais de la vice-reine à Lisbonne », au cœur du pouvoir et sous le regard soupçonneux du secrétaire d’État Vasconcellos. Tel Figaro, c’est par la ruse et le masque que Pinto détrompe les bruits de conjuration parvenus aux oreilles de la vice-reine : « Il me faut tout voir, sans avoir l’air de regarder, caresser ceux que j’abhorre, perdre un temps qui me presse, causer et folâtrer quand un serpent me ronge. \lbrack{}…\rbrack{} On vient. Reprenons le masque\teiNote[xmlid={ftn99-3},n={40},place={foot}]{\teiP[xmlid={p-536}]{\teiHi[rend={italic}]{Ibid}., II, 5, p. 521.}}. » Le ralliement secret des conjurés jurant la mort des Castillans a lieu à l’aube chez Pinto (acte IV), tandis que le dernier acte nous ramène au palais où les Espagnols et leurs partisans sont finalement assaillis. Le motif du secret, moteur de l’intrigue politique, investit un espace qui déborde la scène et franchit le quatrième mur de la fiction, pour renvoyer à la cité (lieu du réel) dans laquelle le théâtre forme une sorte d’îlot.}
            
\end{teiDiv}
            \begin{teiDiv}[type={section1},xmlid={div03-10}]

              \teiHead[xmlid={accessoires-reversibles-001}]{Accessoires réversibles}
              \teiP[xmlid={p12-12}]{L’ambivalence des espaces dans les pièces évoquées est accentuée par des accessoires qui étoffent le décor ou servent de support au partage du secret, c’est-à-dire à sa divulgation par interception ou effraction. Trois motifs sont particulièrement saillants : la lettre, l’armoire et le lit.}
              \teiP[xmlid={p13-12}]{La lettre est ce par quoi le secret réservé à un destinataire restreint est trahi quand elle est interceptée. Toute l’intrigue de \teiHi[rend={italic}]{La Mère coupable} repose sur le poids du passé de la Comtesse, conservé dans son écrin à double-fond que Bégearss révèle habilement au comte :}
              \begin{teiCit}[xmlid={cit04-8}]

                \begin{teiQuote}
\teiHi[rend={small-caps}]{Bégearss} \teiHi[rend={italic}]{prend l’écrin} : Monsieur, au nom du véritable honneur…\teiLb{}\teiHi[rend={small-caps}]{Le Comte} \teiHi[rend={italic}]{a enlevé le bracelet de l’écrin} : Ah ! mon cher portrait, je te tiens ! J’aurai du moins la joie d’en orner le bras de ma fille, cent fois plus digne de le porter…\teiLb{}\teiHi[rend={italic}]{Il y substitue l’autre}.\teiLb{}\teiHi[rend={small-caps}]{Bégearss} \teiHi[rend={italic}]{feint de s’y opposer. Ils tirent chacun l’écrin de leur côté ; Bégearss fait ouvrir adroitement le double fond et dit avec colère} : Ah ! voilà la boîte brisée !\teiLb{}\teiHi[rend={small-caps}]{Le Comte} \teiHi[rend={italic}]{regarde} : Non ; ce n’est qu’un secret\teiNote[xmlid={ftn100-2},n={41},place={foot}]{\teiP[xmlid={p-537}]{Au sens mécanique ici, défini dans le \teiHi[rend={italic}]{Dictionnaire de l’Académie française} (3\teiHi[rend={sup}]{e} édition, 1740) : « On appelle, dans quelques Arts mécaniques, \teiHi[rend={italic}]{Secrets}, Certains ressorts particuliers qui servent à divers usages. \teiHi[rend={italic}]{On ne peut ouvrir ce coffre-fort, si l’on n’en sait le secret. Il y a un secret qui fait qu’une arquebuse tire deux, trois coups. }/ On appelle aussi, \teiHi[rend={italic}]{Secret}, Une cache qui est pratiquée dans un coffre-fort, dans un cabinet » (https://www.dictionnaire-academie.fr/article/A3S0340).}} que le débat a fait ouvrir. Ce double fond renferme des papiers\teiNote[xmlid={ftn101-2},n={42},place={foot}]{\teiP[xmlid={p-538}]{Beaumarchais, \teiHi[rend={italic}]{La Mère coupable} \lbrack{}1792\rbrack{}, dans \teiHi[rend={italic}]{Œuvres}, \teiHi[rend={italic}]{op. cit.}, I, 8, p. 615-616.}} !
\end{teiQuote}
              
\end{teiCit}
              \teiP[xmlid={p14-11}]{Surgit ainsi le fantôme de Chérubin, à travers la lecture frénétique par le comte de sa lettre d’adieu ensanglantée, que la comtesse a fétichisée durant vingt ans et que Bégearss, pour camoufler sa divulgation, prétend faire disparaître en lui imposant de brûler la liasse (III, 6). Mais ayant conservé la lettre qui révèle la culpabilité de sa femme, le Comte revient dévoiler le secret dont il a connaissance dans un tableau-comble d’une rare violence, qui inverse la « surprise » nocturne du page vingt ans plus tôt en métaphorisant un viol jusqu’au délire et à l’évanouissement de la Comtesse, à mesure que le mari fait entendre la voix de l’amant défunt (IV, 13). Cet artifice a probablement inspiré Alexandre Duval pour sa « comédie en cinq actes et en prose\teiNote[xmlid={ftn102-2},n={43},place={foot}]{\teiP[xmlid={p-539}]{Tel est le genre indiqué sur la page de garde. Mais Duval, dans sa « Notice » introductive, explique avoir conçu une tragédie sur le modèle étranger : « cet ouvrage, \teiHi[rend={italic}]{drame} ou \teiHi[rend={italic}]{tragédie}, comme on voudra le nommer, fut longtemps l’objet de mes rêves \lbrack{}…\rbrack{} » (Alexandre Duval, \teiHi[rend={italic}]{Struensé ou le ministre d’État}, dans \teiHi[rend={italic}]{Œuvres complètes d’Alexandre Duval}, \teiHi[rend={italic}]{op. cit.}, t. IV, p. 106).}} » \teiHi[rend={italic}]{Struensé ou le Ministre d’État} (\teiHi[rend={italic}]{ca}. 1802, non représentée), inspirée de l’histoire de Johann Friedrich Struensee, médecin nommé ministre du roi de Danemark après avoir guéri la reine dont il est secrètement amoureux, qui finit décapité à Copenhague en avril 1772. Dans la pièce de Duval\teiNote[xmlid={ftn103-3},n={44},place={foot}]{\teiP[xmlid={p-540}]{Le même sujet sera repris par Gaillardet dans \teiHi[rend={italic}]{Struensée ou le médecin de la reine}, Paris, Barba, 1833.}}, comme plus tard dans \teiHi[rend={italic}]{Ruy Blas}, le ministre parvenu est en butte à l’hostilité méprisante des grands et du ministre de la police Hocborn, qui parviennent à le disgracier en interceptant une lettre de la reine Caroline. La reine-mère Julianne force sa bru à lui lire son message\teiNote[xmlid={ftn104-2},n={45},place={foot}]{\teiP[xmlid={p-541}]{Alexandre Duval, \teiHi[rend={italic}]{Struensé ou le ministre d’État}, \teiHi[rend={italic}]{op. cit.}, p. 173 : « Enfin, madame, comme mère de votre époux, j’ai le droit de savoir des secrets qui intéressent son honneur. C’est à ce titre sacré que je vous demande la lecture de cette lettre » (III, 9).}}, avant de le communiquer par ruse à son fils. Alors que tous les actes ont pour cadre le palais, le roi n’est qu’un actant de l’ombre, majoritairement absent sauf par intermittence dans le hors-scène à l’acte IV, où il s’entretient avec son ministre dans son cabinet, d’où émanent ses ordres. L’objet théâtral qu’est la lettre – et sa variante, le billet – constituent généralement une réduction métonymique de l’action et participent d’une économie du secret, avec une valeur d’usage et une valeur d’échange\teiNote[xmlid={ftn105-2},n={46},place={foot}]{\teiP[xmlid={p-542}]{Voir Jean-Marc Hovasse (dir.), \teiHi[rend={italic}]{Correspondance et théâtre}, Rennes, PUR, « Interférences », 2012, notamment l’étude d’Éric Francalanza, « La lettre dans le théâtre de Marivaux », p. 275-291.}} : ils circulent, cachent et dévoilent, ils servent et perdent les personnages, comme Ruy Blas manipulé dès l’acte I par Don Salluste qui lui fait rédiger deux billets visant, l’un à compromettre la reine en la convoquant à un entretien secret (acte V), l’autre à révéler sa propre identité de laquais en livrée cachée sous son manteau de velours écarlate avec la Toison d’or. Salluste est la « main » qui maîtrise l’espace comme le destin de son valet :}
              \begin{teiLg}

                \teiL{Mais lui ! par quelle porte, ô Dieu, par quelle trappe,}
                \teiL{Par où va-t-il venir, l’homme de trahison ?}
                \teiL{Dans ma vie et dans moi, comme en cette maison,}
                \teiL{Il est maître. Il en peut arracher les dorures.}
                \teiL{Il a toutes les clefs de toutes les serrures.}
                \teiL{Il peut entrer, sortir, dans l’ombre s’approcher,}
                \teiL{Et marcher sur mon cœur comme sur ce plancher\teiNote[xmlid={ftn106-2},n={47},place={foot}]{\teiP[xmlid={p-543}]{Hugo, \teiHi[rend={italic}]{Ruy Blas} \lbrack{}1838\rbrack{}, dans \teiHi[rend={italic}]{Théâtre complet}, \teiHi[rend={italic}]{op. cit.}, IV, 1, p. 1600.}}.}
              
\end{teiLg}
              \teiP[xmlid={p15-10}]{L’armoire, ou toute cachette comparable (alcôve, statue…), constitue un autre artifice significatif dans les pièces historiques. Les drames de Hugo tirent profit des « jeux de cache-cache propres au mélodrame et au vaudeville\teiNote[xmlid={ftn107-2},n={48},place={foot}]{\teiP[xmlid={p-544}]{Florence Naugrette, « Pantomime et tableau », Florence Naugrette et Guy Rosa (dir.), \teiHi[rend={italic}]{Victor Hugo et la langue}, Rosny-sous-Bois, Bréal, 2005, p. 429-450 (435).}} », à l’instar de l’entrée par ruse de François \teiHi[rend={small-caps}]{I}\teiHi[rend={sup}]{er} dans le jardin de Triboulet (\teiHi[rend={italic}]{Le Roi s’amuse}, acte II, où le roi se fait passer pour un certain Gaucher Mahiet, écolier très pauvre de son état), ou du refuge de Don Carlos dans l’armoire à l’exposition d’\teiHi[rend={italic}]{Hernani}, motif qui est ensuite renversé lorsque don Ruy cache le héros derrière son propre portrait (III, 5) : « Le spectateur ne peut percevoir que confusément, partiellement, et en grande partie rétrospectivement le symbolisme de cet espace machiné », analyse Florence Naugrette ; « derrière la généreuse protection de son rival au nom des lois de l’hospitalité, se dessine ici le pacte faustien par lequel Hernani prêtera sa jeunesse au vieillard \lbrack{}…\rbrack{} mais aussi la force mortifère de la piété filiale\teiNote[xmlid={ftn108-2},n={49},place={foot}]{\teiP[xmlid={p-545}]{\teiHi[rend={italic}]{Ibid}., p. 442.}} ». Trente ans plus tôt, Lemercier recourait déjà au signe \teiHi[rend={italic}]{duplex} de l’armoire, simultanément cachette et prison, au dénouement de \teiHi[rend={italic}]{Pinto }; comme Vasconcellos songe à fuir et détruire par le feu ses papiers compromettants, il opte finalement pour les cacher dans l’armoire où s’est réfugié le valet muet de Pinto : « \teiHi[rend={italic}]{Il ouvre l’armoire, le muet l’y pousse et l’y enferme à sa place} », puis les Portugais le découvrent par surprise et l’arrêtent : « \teiHi[rend={italic}]{Vasconcellos sort tout armé ; tout le monde recule avec effroi}\teiNote[xmlid={ftn109-2},n={50},place={foot}]{\teiP[xmlid={p-546}]{Lemercier, \teiHi[rend={italic}]{Pinto}, \teiHi[rend={italic}]{op. cit.}, V, 10-13, p. 614 et 619.}}. » Ouverture-fermeture-ouverture ludiques, corrélat objectif du mouvement politique de fermeture-ouverture au terme de la pièce, puisqu’à la mort des ennemis est préférée l’option de la prison et du procès, avant qu’une amnistie générale ne soit accordée sur demande de Pinto : « Non, non, sire ! grâce entière, oubli généreux du passé ; unissez-vous tous à ma voix qui répond au cœur d’un digne souverain, et criez amnistie\teiNote[xmlid={ftn110-2},n={51},place={foot}]{\teiP[xmlid={p-547}]{\teiHi[rend={italic}]{Ibid}., p. 622.}}. » C’est l’héritage problématique de la Révolution (française) qui est métaphorisé – question que les dramaturges romantiques posent eux aussi, en donnant à voir après le retour des rois, des rois historiques dégradés ou avilis sur les planches.}
              \teiP[xmlid={p16-10}]{En effet, le lit du souverain, ou la paillasse sur laquelle il se vautre, sont également porteurs de cette charge critique. Disconvenant dans la tragédie, le lit appartient à une tradition comique qui l’associe à la mort ou maladie simulée, comme dans \teiHi[rend={italic}]{Le Légataire universel} de Regnard (1708) où le valet Crispin se fait passer pour un Géronte moribond alité devant les notaires, convoqués pour dresser son testament ; le vieil oncle d’Éraste ne revient qu’au 5\teiHi[rend={sup}]{e} acte après le forfait du valet. Lemercier reprend le procédé à l’acte III de \teiHi[rend={italic}]{Pinto}, où le héros charge son amante M\teiHi[rend={sup}]{me} Dolmar de se substituer à la duchesse de Bragance, qui a déjà rejoint les conjurés, en feignant d’être malade et alitée pour éviter de se rendre aux ordres de la vice-reine : l’expédient leurre les commis et sauve ainsi le secret des révolutionnaires. La dégradation atteint son paroxysme dans \teiHi[rend={italic}]{Le Roi s’amuse}, où le lit, invisible mais obsédant, est ici placé sous le signe d’Éros\teiNote[xmlid={ftn111-2},n={52},place={foot}]{\teiP[xmlid={p-548}]{Inversement, dans la « troisième journée » d’\teiHi[rend={italic}]{Angelo, tyran de Padoue}, le lit devient le signe métonymique de la mort de Catarina commanditée par son mari, qui déclare : « Mon lit souillé se change en tombe » (\teiHi[rend={italic}]{op. cit.}, p. 144). Sa femme découvre peu après le funèbre appareil : « \teiHi[rend={italic}]{Elle tire le rideau et recule avec horreur. À la place du lit il y a un billot couvert d’un drap noir et une hache} » (p. 152). Mais Tisbe, sa rivale, la sauve en substituant un narcotique au poison, et la recueille dans son propre lit où elle se réveille au dénouement, alors que Tisbe s’est empoisonnée.}}. Le troisième acte représente :}
              \begin{teiCit}[xmlid={cit05-8}]

                \begin{teiQuote}
\teiHi[rend={italic}]{L’antichambre du roi, au Louvre. Dorures, ciselures, meubles, tapisseries, dans le goût de la Renaissance. – Sur le devant, une table, un fauteuil, un pliant. – Au fond, une grande porte dorée. – À gauche, la porte de la chambre à coucher du roi, revêtue d’une portière en tapisserie. À droite, un dressoir chargé de vaisselles d’or et d’émaux. – La porte du fond s’ouvre sur un mail}\teiNote[xmlid={ftn112-2},n={53},place={foot}]{\teiP[xmlid={p-549}]{Hugo, \teiHi[rend={italic}]{Le Roi s’amuse}, \teiHi[rend={italic}]{op. cit.}, p. 1410.}}.
\end{teiQuote}
              
\end{teiCit}
              \teiP[xmlid={p17-10}]{C’est l’acte dans lequel le roi, « \teiHi[rend={italic}]{vêtu d’un magnifique négligé du matin} », va littéralement se déshabiller pour violer Blanche. La dégradation lubrique, inversant le cérémonial du « petit lever » du roi, est soulignée par Marot qui recourt au bestiaire : « Le lion a traîné la brebis dans son antre\teiNote[xmlid={ftn113-2},n={54},place={foot}]{\teiP[xmlid={p-550}]{\teiHi[rend={italic}]{Ibid}., p. 1419 (III, 3).}}. » Le bouffon révèle alors le secret de sa paternité, mais les courtisans l’empêchent d’accéder à la chambre du roi. Tableau pathétique et véhément du père meurtri suppliant à genoux les « gentilshommes » complices, puis recevant sa fille souillée dans cette « cour prostituée au mal » qui n’est « Que vice, effronterie, impudeur, adultère, / Infamie et débauche\teiNote[xmlid={ftn114-2},n={55},place={foot}]{\teiP[xmlid={p-551}]{\teiHi[rend={italic}]{Ibid}., III, 4, p. 1435.}} »… De fait, l’acte IV expose aux yeux de la pauvre « brebis » le donjuanisme de son séducteur, déguisé en simple officier dans le bouge de Saltabadil :}
              \begin{teiCit}[xmlid={cit06-7}]

  \begin{teiQuote}

    \begin{teiElement}[name={sp}]

      \begin{teiElement}[name={speaker}]
Maguelonne
\end{teiElement}
      \begin{teiElement}[name={stage},type={delivery}]
riant
\end{teiElement}
      \teiL{Monsieur, vous m’avez l’air d’un libertin parfait !}
    
\end{teiElement}

    \begin{teiElement}[name={sp}]

      \begin{teiElement}[name={speaker}]
Le Roi
\end{teiElement}
      \begin{teiElement}[name={stage},type={delivery}]
riant aussi
\end{teiElement}
      \teiL{Oui, j’ai fait le malheur de plus d’une, en effet.}
      \teiL{
        C’est vrai, je suis un monstre \lbrack{}…\rbrack{}
        \teiNote[xmlid={ftn115-2},n={56},place={foot}]{
          \teiP[xmlid={p-552}]{\teiHi[rend={italic}]{Ibid}., IV, 2, p. 1444-1445.}
        }
      }
    
\end{teiElement}
  
\end{teiQuote}

\end{teiCit}
              \teiP[xmlid={p18-9}]{Si, en 1832, le roi de Hugo « s’amuse », il est bien seul : le public et les censeurs n’ont guère de quoi rire. « L’inversion du Roi noble en roi de carnaval et sa détronisation rêvée forment le canevas de la pièce », conclut Anne Ubersfeld, pointant l’absence quasi-totale du comique dans ce drame qu’elle qualifie de « tragédie \teiHi[rend={italic}]{grotesque}\teiNote[xmlid={ftn116-2},n={57},place={foot}]{\teiP[xmlid={p-553}]{Anne Ubersfeld, \teiHi[rend={italic}]{Le Roi et le Bouffon},\teiHi[rend={italic}]{ op. cit}., p. 637-640.}} ». Nous sommes loin, en apparence, de la comédie historique des Lumières, où l’attendrissement admiratif du peuple permettait de communier avec ses souverains par le truchement des grands « en déshabillé ». Dans les drames romantiques, l’heure est à la satire grinçante qui flétrit le pouvoir\teiNote[xmlid={ftn117-2},n={58},place={foot}]{\teiP[xmlid={p-554}]{Voir sur ce point la thèse de Sophie Mentzel, \teiHi[rend={italic}]{Trônes vacillants},\teiHi[rend={italic}]{ op. cit}., \teiHi[rend={italic}]{passim}.}}, détruisant le héros. Pour autant, au-delà de ce déplacement tonal éminemment politique, nous retrouvons des usages semblables des lieux et du secret, suivant un processus d’hybridation sous l’influence de l’esthétique du drame et du mélodrame bien relevé par les historiens du théâtre\teiNote[xmlid={ftn118-2},n={59},place={foot}]{\teiP[xmlid={p-555}]{Voir notamment Florence Naugrette, \teiHi[rend={italic}]{Le Théâtre romantique. Histoire, écriture, mise en scène}, Paris, Seuil, 2001, chap. 1, « Les origines du drame romantique ».}}, qui ont cependant moins repéré l’héritage de la comédie historique conçue au \teiHi[rend={small-caps}]{xviii}\teiHi[rend={sup}]{e} siècle. Espaces et accessoires \teiHi[rend={italic}]{signifient} autant que les personnages et leurs discours, en ce qu’ils sont les véhicules matériels des secrets qu’ils rendent sensibles pour les lecteurs-spectateurs. La « comédie qui peint l’humanité » pour « les penseurs\teiNote[xmlid={ftn119-2},n={60},place={foot}]{\teiP[xmlid={p-556}]{D’après la partition générique établie par Hugo au début de la préface de \teiHi[rend={italic}]{Ruy Blas},\teiHi[rend={italic}]{ op. cit}., p. 1489.}} » recourt ainsi à des procédés et use de motifs récurrents, selon une plasticité d’usages : leur analyse révèle le fruit de l’Histoire, la trace de la Révolution et des positions politiques que médiatise le théâtre dans un contexte socialement situé.}
            
\end{teiDiv}
          
\end{teiElement}
        
\end{teiElement}
      
\end{teiElement}
    
\end{teiElement}
    \begin{teiElement}[name={back}]

\end{teiElement}
  
\end{teiElement}

\end{teiElement}
\end{lateiDocument}

\cleardoublepage
\pagestyle{plain}
\titleformat{\chapter}[display]{\PURHTitleFont\bfseries\fontsize{16pt}{19pt}\selectfont\centering}{}{10pt}{\MakeUppercase}
\tableofcontents

\end{document}
